\documentclass[11pt,oneside,openany]{book}
\usepackage[T1]{fontenc}
\usepackage[utf8]{inputenc}
\usepackage{lmodern}
\usepackage{amsmath,amssymb,amsthm,mathtools}
\usepackage{tikz-cd}
\usepackage{needspace}
\usepackage{xcolor}
\usepackage{hyperref}
\hypersetup{
  colorlinks=true,
  linkcolor=blue!50!black,
  citecolor=blue!50!black,
  urlcolor=blue!50!black,
  hypertexnames=false,
  pdftitle={Lecture Notes on Algebra: Groups, Rings, Modules, Vector Spaces, and Fields},
  pdfsubject={Introductory graduate algebra},
  pdfkeywords={algebra, groups, rings, modules, vector spaces, fields, Galois theory}
}
\raggedbottom
\makeatletter
\renewcommand{\@makechapterhead}[1]{%
  \vspace*{30\p@}%
  {\parindent \z@ \raggedright \normalfont
    \ifnum \c@secnumdepth >\m@ne
      \huge\bfseries \@chapapp\space \thechapter\par\nobreak
      \vskip 14\p@
    \fi
    \Huge \bfseries #1\par\nobreak
    \vskip 30\p@
  }}
\renewcommand{\@makeschapterhead}[1]{%
  \vspace*{30\p@}%
  {\parindent \z@ \raggedright \normalfont
    \Huge \bfseries #1\par\nobreak
    \vskip 30\p@
  }}
\renewcommand\section{\@startsection{section}{1}{\z@}%
  {-3.5ex \@plus -1ex \@minus -.2ex}%
  {2.3ex \@plus .2ex}%
  {\normalfont\large\bfseries}}
\makeatother

\theoremstyle{definition}
\newtheorem{theorem}{Theorem}[chapter]
\newtheorem{proposition}[theorem]{Proposition}
\newtheorem{lemma}[theorem]{Lemma}
\newtheorem{corollary}[theorem]{Corollary}
\newtheorem{definition}[theorem]{Definition}
\newtheorem{example}[theorem]{Example}
\newtheorem*{applicationtheorem}{Theorem}
\newtheorem*{undefinition}{Definition}
% Exercises have their own counter within each chapter; they do not share the
% theorem/proposition sequence.
\newtheorem{exercise}{Exercise}[chapter]
\theoremstyle{remark}
\newtheorem{remark}[theorem]{Remark}
\newtheorem*{unremark}{Remark}
\newcommand{\highlight}[1]{\textbf{#1}}
\renewcommand{\emph}[1]{\textbf{#1}}
\newenvironment{proofidea}{\par\medskip\noindent\textbf{Proof Idea.}\ }{\par\medskip}

\begin{document}
\frontmatter
\hypersetup{pageanchor=false}
\pdfbookmark[0]{Cover}{cover}
\begin{titlepage}
\centering
\vspace*{0.18\textheight}
{\Huge\bfseries Lecture Notes on Algebra\par}
\vspace{0.8em}
{\Large Groups, Rings, Modules, Vector Spaces, and Fields\par}
\vspace{2.25em}
{\Large Tommy Chen\par}
\vfill
% Version 1.12 is a small maintenance update refining Chapter 6 foundations.
{\large Version 1.12\par}
\vspace{0.45em}
{\small Licensed under CC BY 4.0.\par}
\end{titlepage}
\chapter*{Preface}
\addcontentsline{toc}{chapter}{Preface}
These notes offer a self-contained introduction to the central ideas of modern algebra, beginning with groups and proceeding through rings, modules, linear algebra, fields, and Galois theory.  Their guiding ambition is an elementary entry point with a graduate destination: definitions are motivated by familiar examples, while the later chapters develop the structural language needed for further study in algebra, number theory, geometry, and representation theory.

The exposition is cumulative.  Results are proved before they are used, and the hypotheses that make a construction meaningful are stated explicitly.  In particular, the notes promote important structural results that are often assigned as exercises in the primary reference into the main narrative when later arguments depend on them.  The exercises are intentionally selective: they are designed to consolidate the main ideas rather than to conceal essential theory.

The chapters may be read in order as a first course in abstract algebra with substantial linear-algebraic and homological preparation.  Readers already familiar with an earlier portion may also use the later parts as a reference, subject to the stated dependencies.  The author hopes that the recurring themes---symmetry, quotient constructions, universal properties, and decomposition---will make the subject feel like a connected whole rather than a collection of techniques.

\vspace{1.5em}
\noindent
Tommy Chen
\pdfbookmark[0]{Contents}{contents}
\tableofcontents
\clearpage
\hypersetup{pageanchor=true}
\chapter*{Editorial convention}
\addcontentsline{toc}{chapter}{Editorial convention}
Completed portions of these notes are intended to be logically complete within their chosen scope.  Important
results in the primary reference are incorporated and proved according to
mathematical dependency, whether the source presents them in its text or its
exercises.  The accompanying coverage audit records each verified promotion
with an exact source reference.
\mainmatter
\chapter{Preliminaries and Guide to the Notes}\label{chap:preliminaries}

\section{Purpose and point of view}

Graduate algebra is not chiefly a collection of classifications.  It is a
language for recognizing when apparently different problems have the same
form. A quotient group, a quotient ring, and a quotient module all arise by
collapsing a chosen subobject to the trivial one: a normal subgroup becomes
the identity in a group, while an ideal or submodule becomes zero. An action describes symmetry; a
polynomial in an operator turns linear algebra into module theory; and a
Galois group records the symmetries of a field extension.  These notes develop
that common language with an eye toward the places it will be used next.

The intended reader has strong high-school or early undergraduate mathematical
maturity: comfort with sets, functions, elementary proof, and familiar number
systems such as \(\mathbb Z\), \(\mathbb Q\), and \(\mathbb R\).  No prior
course in abstract algebra is assumed.  Definitions begin from that elementary
entry point, examples arrive before abstraction, and proofs introduce each new
move before relying on it.  The destination is nevertheless graduate algebra:
the resulting core is designed for later commutative algebra, homological
algebra, representation theory, algebraic geometry, and algebraic number
theory.  The last chapters give a complete first treatment of classical Galois
theory.

This cumulative policy is deliberate.  Early chapters do not appeal to an
unseen ``standard fact from algebra''; later chapters may use a result only
after it has been established here.  Readers who have already studied a first
course can move more quickly through the examples, but they are not the
prerequisite for understanding the logical development.

Readers who would like additional preparation in mathematical language,
proof-writing, foundational constructions, and the number systems may consult
Chapter~2 of \emph{Lecture Notes on Mathematical Analysis} \cite{MA}.
Its discussions of sets, functions, proof conventions, induction, and the
construction and basic properties of \(\mathbb N,\mathbb Z,\mathbb Q,\mathbb R\)
can help bridge to the proof-based style used here.  This is useful
preparation, not a prerequisite: all algebraic material required in these
notes is developed within the manuscript.

Two habits are central.  First, definitions are treated as constructions with
universal features, rather than as names for collections of examples.  Second,
proofs are organized around reusable moves: pass to a quotient, use a universal
property, compare kernels and images, act on a set, reduce a statement to a
cyclic object, or use induction on a degree or a composition length.  The same
moves will recur throughout the book.

\section{The mathematical route}

Part~I establishes the group-theoretic language of normal subgroups, quotients,
actions, and composition factors.  Sylow theory is included for its structural
consequences, while lengthy classifications of groups of individual small
orders are not.  Part~II develops division rings, ideals, quotients, group
rings, and localization, then the chain
of implications
\[
  \text{Euclidean domain}\quad\Longrightarrow\quad\text{PID}
  \quad\Longrightarrow\quad\text{UFD}.
\]
It also develops Noetherian finiteness, Hilbert's basis theorem, formal power
series, one- and several-variable polynomial rings, finite-field applications,
a short quadratic-reciprocity perspective, and an introduction to Gr\"obner
bases.  Localization is included because it is the first algebraic operation
that makes a global ring look local.

Part~III develops four connected chapters. Chapter~5 introduces modules
together with the finite-dimensional linear algebra needed later: bases,
dimension, matrices, determinants, eigenvalues, diagonalization, duality,
and a modest introduction to bilinear and quadratic forms.
Chapter~6 proves a strengthening theorem for submodules of finite free
modules over a PID, producing compatible ambient coordinates one
invariant-factor direction at a time. It then derives the PID structure
theorem. Uniqueness in elementary-divisor form is recovered from intrinsic
prime-primary layers, while uniqueness in invariant-factor form is also
proved directly by annihilator, cyclic cancellation, and induction. Smith
normal form gives the computational
matrix manifestation before the chapter regards an operator as an
\(F[x]\)-module, leading to minimal polynomials and rational and Jordan
canonical forms. Chapter~7 adds inner-product geometry, adjoints, spectral
theory, and singular value decomposition. Chapter~8 introduces tensor
products, exterior powers, and exterior algebra: the module constructions
that linearize bilinear and alternating multilinear maps and prepare the way
for later commutative and homological algebra. Thus vector spaces are
modules over fields, abelian groups are modules over \(\mathbb Z\), and a
linear operator on an \(F\)-vector space makes that space into an
\(F[x]\)-module. The structure theorem for finitely generated modules over a
PID therefore simultaneously explains finitely generated abelian groups and
the rational canonical form of a linear map. Jordan form appears as the
split-polynomial refinement of the same idea.

Part~IV studies field extensions, their embeddings, separability, normality,
and the correspondence between intermediate fields and subgroups of a Galois
group.  Finite fields and cyclotomic fields supply recurring examples, while
direct and inverse limits explain how finite Galois symmetries assemble into
profinite and absolute Galois groups.  Chapter~11 includes a brief local
bridge through point-set topology and topological groups, solely to make the
profinite and infinite-Galois language self-contained; Tychonoff's theorem is
used there as an external input.  It closes with a non-prerequisite survey of
the later directions opened by all four parts. The later discussion of
radicals illustrates how questions about polynomial formulas are controlled by
group-theoretic properties.

\section{How to read a proof}

Every result used later is stated with its hypotheses and proved before its
first essential use.  A proof normally begins by identifying its idea: for
example, the first isomorphism theorem compares an object with the quotient by
the kernel of a map, whereas the orbit--stabilizer theorem compares a group
action with a quotient by a stabilizer.  Readers should pause at these ideas,
not merely at the final calculation.

Some results are deliberately stated in a general form.  This is done only
when the extra generality clarifies later work.  For instance, the module
isomorphism theorems subsume their group and ring analogues, but the earlier
versions are still proved because quotient constructions must be understood in
their native settings.  Conversely, the classification of finitely generated
abelian groups is presented as a corollary of the PID module theorem rather
than as an unrelated classification.

\section{Exercises and completeness}

Exercises have two roles: they give essential practice, and they sometimes
contain theorems that deserve a place in the main narrative.  We do not follow
the source's typography as a measure of importance.  When an exercise supplies
a standard lemma, a useful converse, a structural construction, or a result
needed later, the result is promoted and proved in the text.  The companion
coverage audit records the exact source exercise, its location here, and the
reason for promotion.

At the end of each chapter, exercises are grouped as \emph{Basic},
\emph{Core}, \emph{Further}, and \emph{Looking Ahead}.  The first group tests
definitions; the second develops results expected of every reader; the third
extends the chapter; and the last points toward later subjects without being a
logical prerequisite.  A routine exercise remains an exercise even when its
solution is long.

\section{Conventions and notation}

Rings in these notes are associative and have a multiplicative identity; the
zero ring is allowed.  Ring homomorphisms preserve identity unless explicitly
stated otherwise, and a subring contains the same identity as its ambient
ring.  Terminology concerning rings is not completely uniform: some authors
allow rings without a multiplicative identity, and some do not require a
subring to contain the ambient identity.  The conventions used here will
always be stated explicitly when a nonunital object is needed.  Unless
explicitly qualified otherwise, an ideal means a two-sided ideal.  When
one-sided ideals are relevant, we explicitly say left ideal or right ideal;
in a commutative ring these notions coincide.
Modules are unital left modules unless the contrary is specified. For a general
group \(G\), the identity is \(e\), and for an additive group it is \(0\).
We retain \(1\) for concrete multiplicative groups and \(I_n\) for matrix
groups; \(\operatorname{id}_X\) denotes an identity map. The
notation \(H\leq G\) permits equality, while \(H\trianglelefteq G\) denotes a
normal subgroup.  If \(R\) is a ring and \(M\) an \(R\)-module, then
\(\operatorname{Ann}_R(M)\) is its annihilator.

For a regular \(n\)-gon, \(D_n\) denotes its full dihedral symmetry group, so
\(|D_n|=2n\).  Many authors write \(D_{2n}\) for this same group; the notation
in these notes records the number of polygon vertices rather than the group
order.  Thus \(D_4\) is the symmetry group of a square.

\begin{remark}[Looking Ahead]
The repeated pattern ``kernel, image, quotient'' is the beginning of the
language of exact sequences.  We will use exact sequences of modules in an
elementary form, but derived functors and category theory are not prerequisites
for these notes.
\end{remark}

\section{Elementary complex arithmetic}

The complex numbers extend the real numbers by an element \(i\) satisfying
\(i^2=-1\). Every complex number has a unique form
\[
z=a+bi,\qquad a,b\in\mathbb R.
\]
The numbers \(a=\operatorname{Re}(z)\) and \(b=\operatorname{Im}(z)\) are,
respectively, its real and imaginary parts. We add and multiply by collecting
real and imaginary parts and using \(i^2=-1\):
\[
(a+bi)+(c+di)=(a+c)+(b+d)i,
\quad
(a+bi)(c+di)=(ac-bd)+(ad+bc)i.
\]
For instance,
\[
(2+3i)+(1-4i)=3-i,\qquad
(2+3i)(1-4i)=14-5i.
\]

These definitions inherit the familiar algebraic laws of real-number
arithmetic. They are closed under addition and multiplication; addition and
multiplication are associative and commutative; \(0=0+0i\) is the additive
identity, \(1=1+0i\) is the multiplicative identity, and
\(-(a+bi)=-a-bi\) is the additive inverse. For example, coordinatewise
addition gives
\begin{align*}
((a+bi)+(c+di))+(p+qi)
&=(a+c+p)+(b+d+q)i\\
&=(a+bi)+((c+di)+(p+qi)).
\end{align*}
The same calculation gives commutativity and the additive identity and inverse
laws. Expanding both sides of \(z(w+u)=zw+zu\) and collecting real and
imaginary parts proves distributivity; associativity and commutativity of
multiplication follow similarly from the corresponding real-number laws.
The usual cancellation arguments prove uniqueness of the additive and
multiplicative identities and of additive inverses.

Conjugation is a useful device for converting a complex number into a real
number. Define
\[
\overline z=a-bi,\qquad |z|=\sqrt{a^2+b^2}.
\]
Thus for \(z=3+4i\), \(\overline z=3-4i\) and \(|z|=5\). Direct coordinate
calculations give
\[
\overline{z+w}=\overline z+\overline w,\qquad
\overline{zw}=\overline z\,\overline w,\qquad
\overline{\overline z}=z,
\]
and
\[
z\overline z=(a+bi)(a-bi)=a^2+b^2=|z|^2.
\]
Consequently
\[
|zw|^2=(zw)\overline{zw}=z\overline z\,w\overline w=|z|^2|w|^2,
\]
so \(|zw|=|z||w|\), since moduli are nonnegative.

If \(z=a+bi\ne0\), then \(a\) and \(b\) are not both zero, so
\(a^2+b^2>0\). Hence
\[
z^{-1}=\frac{a-bi}{a^2+b^2}=\frac{\overline z}{|z|^2}
\]
is defined, and \(zz^{-1}=z\overline z/|z|^2=1\). For example,
\[
(3+4i)^{-1}=\frac{3-4i}{25}.
\]
This inverse is unique: if \(u\) and \(v\) both invert \(z\), then
\[
u=u\cdot1=u(zv)=(uz)v=1\cdot v=v.
\]
These properties---together with the existence of multiplicative inverses for
every nonzero element---are exactly the familiar algebraic laws that will
later be summarized by saying that \(\mathbb C\) is a field.

Geometrically, \(a+bi\) is the point \((a,b)\) in the plane, and \(|z|\) is
its Euclidean distance from the origin. Thus the complex \(n\)-th roots of
unity
\[
\mu_n=\{z\in\mathbb C:z^n=1\}
=\left\{\cos\frac{2\pi k}{n}+i\sin\frac{2\pi k}{n}:0\leq k<n\right\}
\]
are the \(n\) equally spaced points on the unit circle. For example,
\[
\mu_2=\{1,-1\},\qquad \mu_4=\{1,i,-1,-i\}.
\]
They will later provide important group examples.

\section{Sources}

The primary mathematical reference is \cite{DF}.  The separate coverage audit
records editorial provenance and verified exercise references; it is not part
of the reader-facing source apparatus.

\part{Groups}
\chapter{Groups, Quotients, and Homomorphisms}\label{chap:groups-foundations}

\section{Groups, Subgroups, and Generated Subgroups}

Groups record a collection of reversible symmetries or operations.  The
axioms are deliberately sparse: associativity makes an unparenthesized product
meaningful, the identity records ``doing nothing,'' and inverses record the
possibility of undoing an operation.  The familiar examples \(\mathbb Z\),
\(\mathbb Z/n\mathbb Z\), the symmetric groups \(S_n\), and the dihedral groups
\(D_n\) will recur throughout Part~I.

\begin{definition}
A \emph{group} is a set \(G\) equipped with a binary operation
\(G\times G\to G\), \((x,y)\mapsto xy\), satisfying the following axioms.
\begin{enumerate}
\setlength{\itemsep}{7pt}
\item \textbf{Associativity:} for all \(x,y,z\in G\),
      \[(xy)z=x(yz).\]
\item \textbf{Identity:} there is an element \(1\in G\) such that, for every
      \(x\in G\),
      \[1x=x=x1.\]
\item \textbf{Inverses:} for every \(x\in G\), there is an element
      \(x^{-1}\in G\) such that
      \[xx^{-1}=1=x^{-1}x.\]
\end{enumerate}
A nonempty subset \(H\subseteq G\) is a \emph{subgroup} when the same
operation makes \(H\) into a group; equivalently, it is closed under products
and inverses.
\end{definition}

Associativity permits \(g_1g_2\cdots g_n\) without specifying parentheses.

\begin{example}[Additive Number Systems]
The integers \((\mathbb Z,+)\) form a group: associativity is inherited from
ordinary addition, the identity is \(0\), and the inverse of \(m\) is \(-m\).
The same verification gives additive groups \((\mathbb Q,+)\) and
\((\mathbb R,+)\).  These examples explain the word ``inverse'' in an additive
setting: it is an additive inverse, not a reciprocal.  The natural numbers
\((\mathbb N,+)\) are not a group, because a positive integer has no additive
inverse in \(\mathbb N\).
\end{example}

\begin{example}[Multiplicative Number Systems]
The nonzero rational numbers \(\mathbb Q^\times=\mathbb Q\setminus\{0\}\)
form a group under multiplication.  The identity is \(1\), and the inverse of
\(a/b\) is \(b/a\).  Likewise \(\mathbb R^\times\) is a group.  The zero
rational or real number must be removed: it has no multiplicative inverse.
This small change illustrates a common algebraic move---restrict to the
elements for which an operation is reversible.
\end{example}

\begin{example}[Matrices]
For the familiar number systems \(F=\mathbb Q\) or \(\mathbb R\), the
invertible \(n\times n\) matrices form the group \(\operatorname{GL}_n(F)\).
Matrix multiplication is associative, the identity is \(I_n\), and membership
means precisely that a matrix inverse exists.  (The general notion of a field
will be developed later.)  For \(n\geq2\) this group is generally nonabelian:
two reversible linear transformations need not commute.  This is the first
indication that groups capture more than arithmetic.
\end{example}

\begin{proposition}[Elementary Group Identities]
The identity and inverses are unique.  Cancellation holds, and
\((xy)^{-1}=y^{-1}x^{-1}\).
\end{proposition}
\begin{proof}
If \(e,e'\) are identities, use first the right-identity property of \(e'\)
and then the left-identity property of \(e\): \(e=ee'=e'\).  If \(y,z\) both
invert \(x\), then associativity lets us insert parentheses as follows:
\[y=y1=y(xz)=(yx)z=1z=z.\]
For cancellation, if \(ax=ay\), multiply the equality on the left by
\(a^{-1}\):
\[x=1x=(a^{-1}a)x=a^{-1}(ax)=a^{-1}(ay)=(a^{-1}a)y=1y=y.\]
Right cancellation is similar.  Finally,
\[(xy)(y^{-1}x^{-1})=x(yy^{-1})x^{-1}=x1x^{-1}=1,\]
and the same calculation in the opposite order shows that \(y^{-1}x^{-1}\)
is the inverse of \(xy\).
\end{proof}

Thus a subgroup is not merely a subset: it is a collection of operations that
is itself closed under composition and reversal.  For example,
\(n\mathbb Z\leq\mathbb Z\), while the transpositions in \(S_3\) do not form a
subgroup because their product need not be a transposition.

We will use one kind of map before studying maps systematically.  A function
\(f:G\to H\) between groups is a \emph{homomorphism} if
\(f(xy)=f(x)f(y)\) for all \(x,y\in G\).  Its \emph{kernel} is the set of
elements sent to the identity of \(H\).  The structural consequences of these
definitions---images, normality, and quotient groups---are developed later in
this chapter.

\begin{proposition}[Subgroup Test]
A nonempty subset \(H\) of \(G\) is a subgroup if and only if
\(xy^{-1}\in H\) for all \(x,y\in H\).
\end{proposition}
\begin{proof}
If \(H\) is already a subgroup, then \(y^{-1}\in H\) and closure under
products gives \(xy^{-1}\in H\).  Conversely, assume the displayed condition.
Choose \(h\in H\), possible because \(H\) is nonempty.  Taking \(x=y=h\)
gives \(hh^{-1}=1\in H\).  Now take \(x=1\) and any \(y\in H\); the condition
gives \(y^{-1}\in H\).  Finally, if \(x,y\in H\), then \(y^{-1}\in H\), so
the condition applied to \(x\) and \(y^{-1}\) gives
\(x(y^{-1})^{-1}=xy\in H\).  Thus \(H\) contains the identity and is closed
under products and inverses.
\end{proof}

The subgroup test is useful because it replaces three separate closure checks
by a single expression.  It also anticipates a recurring algebraic pattern:
subobjects are often recognized by closure under a difference-like operation.

\begin{definition}
For a subset \(S\subseteq G\), the \emph{subgroup generated by \(S\)} is
\[
\langle S\rangle=\bigcap_{\substack{H\leq G\\S\subseteq H}}H.
\]
It is the smallest subgroup of \(G\) containing \(S\).  A group is
\emph{finitely generated} if \(G=\langle S\rangle\) for some finite \(S\).
\end{definition}

The intersection is a subgroup, and it is contained in every subgroup that
contains \(S\), which proves the asserted minimality.  Equivalently,
\(\langle S\rangle\) consists exactly of finite products of elements of \(S\)
and their inverses (including the empty product \(1\)).  This formulation is
often the practical one; the intersection formulation is the structural one.
The same ``smallest subobject containing a set'' construction will reappear for
ideals, modules, algebras, and field extensions.

\begin{example}
In the additive group \(\mathbb Z\), \(\langle6,10\rangle=2\mathbb Z\), since
\(2=2\cdot6-10\).  In \(S_3\), \((12)\) and \((23)\) generate \(S_3\).
\end{example}

\begin{proposition}[Finite Products from a Generating Set]
For \(S\subseteq G\), the subgroup \(\langle S\rangle\) consists exactly of
the finite products \(s_1^{\varepsilon_1}\cdots s_k^{\varepsilon_k}\), where
\(s_i\in S\) and \(\varepsilon_i\in\{1,-1\}\), together with the empty
product \(1\).
\end{proposition}
\begin{proof}
Let \(W\) be the set of these finite products in the already given group
\(G\).  Concatenation gives closure under products, and reversal with inverse
exponents gives closure under inverses.
Thus \(W\) is a subgroup containing \(S\), so minimality gives
\(\langle S\rangle\subseteq W\).  Every subgroup containing \(S\) contains
all these finite products, so its intersection contains \(W\).
\end{proof}

This two-containment argument will recur whenever we describe an object
generated by a specified set of elements.

\section{Cyclic Groups}

For \(g\in G\), \(\langle g\rangle=\{g^n:n\in\mathbb Z\}\) is the special
case \(S=\{g\}\), the \emph{cyclic subgroup generated by \(g\)}.  The order
\(|g|\) is the least positive \(n\) with \(g^n=1\), if it exists.

The cyclic group \(\langle g\rangle\) is the portion of \(G\) generated by one
operation.  Additively, \(\langle 1\rangle=\mathbb Z\); multiplicatively, the
rotations of a regular \(n\)-gon form \(C_n\).  The order of \(g\) measures the
first time repeated application returns to the identity.

\begin{proposition}
Every subgroup of a cyclic group is cyclic.  If \(G=\langle g\rangle\) has
order \(n\), then for each divisor \(d\mid n\), \(G\) has a unique subgroup of
order \(d\), namely \(\langle g^{n/d}\rangle\).
\end{proposition}
\begin{proof}
For a nontrivial subgroup \(H\), let \(r\) be the least positive integer with
\(g^r\in H\).  Given \(g^k\in H\), divide \(k\) by \(r\): write
\(k=qr+s\), \(0\leq s<r\).  Then \(g^s=g^k(g^r)^{-q}\in H\).  Minimality of
\(r\) forces \(s=0\), so every element of \(H\) is a power of \(g^r\), and
\(H=\langle g^r\rangle\).  In a finite cyclic group, \(g^a\) has order
\(n/\gcd(n,a)\).  Thus \(g^{n/d}\) has order \(d\).  If a subgroup has order
\(d\), write it as \(H=\langle g^r\rangle\).  Then
\(\gcd(n,r)=n/d\), so \(r=(n/d)u\) with \(\gcd(u,d)=1\).  Since
\(g^{n/d}\) has order \(d\), there are integers \(v,w\) with
\(uv+dw=1\).  Consequently
\[
 g^{n/d}=(g^r)^v(g^{n/d})^{dw}=(g^r)^v,
\]
because \((g^{n/d})^d=1\).  Hence
\(\langle g^{n/d}\rangle\leq\langle g^r\rangle\); the reverse inclusion is
immediate from \(g^r=(g^{n/d})^u\).  Thus
\(H=\langle g^{n/d}\rangle\), which proves uniqueness precisely.
\end{proof}

\begin{example}[The Lattice of \(C_{12}\)]
\leavevmode\\
If \(C_{12}=\langle g\rangle\), its subgroups are
\(\langle g\rangle,\langle g^2\rangle,\langle g^3\rangle,
\langle g^4\rangle,\langle g^6\rangle\), and \(\{1\}\), of orders
\(12,6,4,3,2,1\), respectively.  The divisors of \(12\) predict exactly this
subgroup lattice.  In contrast, the analogous subgroup problem for a general
finite group is substantially harder, which is why cyclic groups are a useful
first laboratory.
\end{example}

\begin{proposition}
The homomorphism \(\varphi_g:\mathbb Z\to G\), \(n\mapsto g^n\), has image
\(\langle g\rangle\), and
\[
 \ker\varphi_g=
 \begin{cases}
 \{0\},&\text{if \(g\) has infinite order},\\[2pt]
 n\mathbb Z,&\text{if \(|g|=n\)}.
 \end{cases}
\]
Thus an infinite cyclic group is isomorphic to \(\mathbb Z\), whereas a
cyclic group of order \(n\) is isomorphic to \(\mathbb Z/n\mathbb Z\).
\end{proposition}
\begin{proof}
For integers \(a,b\),
\(\varphi_g(a+b)=g^{a+b}=g^ag^b=\varphi_g(a)\varphi_g(b)\), so
\(\varphi_g\) is indeed a homomorphism.  Its image is precisely the set of
all integer powers of \(g\), namely \(\langle g\rangle\).
The kernel consists of precisely those integers \(a\)
for which \(g^a=1\).  If \(g\) has infinite order, this happens only for
\(a=0\).  If \(|g|=n\), division with remainder shows that \(g^a=1\) if and
only if \(n\mid a\), so the kernel is \(n\mathbb Z\).  In the infinite-order
case, \(\varphi_g\) is one-to-one: \(g^a=g^b\) implies
\(g^{a-b}=1\), hence \(a=b\).  It is onto \(\langle g\rangle\) by its
definition, and thus gives \(\mathbb Z\cong\langle g\rangle\).
If \(|g|=n\), the map
\[a+n\mathbb Z\longmapsto g^a\]
is well-defined because \(g^n=1\); it is a bijective homomorphism onto
\(\langle g\rangle\).  The First Isomorphism Theorem will later combine these
two direct arguments into the single statement
\(\langle g\rangle\cong\mathbb Z/\ker\varphi_g\).
\end{proof}

\begin{proposition}
If \(G=\langle g\rangle\) has order \(n\), then \(g^k\) is a generator if and
only if \(\gcd(k,n)=1\).
\end{proposition}
\begin{proof}
The order of \(g^k\) is \(n/\gcd(n,k)\).  Thus \(C_{12}\) has generators
\(g,g^5,g^7,g^{11}\).
\end{proof}

\section{Dihedral Groups}

\begin{definition}
The \emph{dihedral group} \(D_n\) is the full symmetry group of a regular
\(n\)-gon.  If \(r\) is rotation through one vertex and \(s\) a reflection,
then
\[D_n=\langle r,s\mid r^n=s^2=1,\ srs^{-1}=r^{-1}\rangle.\]
Many authors call this order-\(2n\) group \(D_{2n}\); these notes use \(D_n\).
\end{definition}

The relation \(sr=r^{-1}s\) lets us move every occurrence of \(s\) to the
left in a product.  Since \(s^2=1\) and \(r^n=1\), every product of rotations
and reflections therefore has the form \(r^i\) or \(sr^i\), with
\(0\leq i<n\).  These \(2n\) symmetries are distinct geometrically: the first
family preserves orientation and the second reverses it.  Hence
\(|D_n|=2n\) and \(\langle r\rangle\cong C_n\).  Notice that multiplying
\(srs^{-1}=r^{-1}\) on the right by \(s\) gives the useful rewriting rule
\(sr=r^{-1}s\): reflecting after rotating has the same effect as rotating in
the opposite direction after reflecting.

For the square, choose \(r\) to be the quarter-turn and \(s\) a reflection
in an axis through opposite vertices.  The complete list is
\[
 D_4=\{1,r,r^2,r^3,s,sr,sr^2,sr^3\}.
\]
The first four are rotations, and the last four are reflections.  In
particular, \(sr=r^{-1}s=r^3s\), whereas \(rs\) is a different reflection:
\(sr\ne rs\).  This is a concrete calculation of noncommutativity, not just
an abstract consequence of the presentation.

\begin{example}[Dihedral Groups as Permutation Groups]
Label the vertices of a regular \(n\)-gon by \(1,\ldots,n\).  Every symmetry
permutes these vertices, so it determines an element of \(S_n\).  The action
is injective: a geometric symmetry that fixes every vertex is the identity.
Thus \(D_n\) is concretely a subgroup of \(S_n\).  The later Cayley theorem
will generalize this phenomenon from dihedral groups to every finite group.
\end{example}

\section{Symmetric and Alternating Groups}

\begin{definition}
The \emph{symmetric group} \(S_n\) is the group of all bijections of the set
\(\{1,2,\ldots,n\}\), with composition as the operation.  A \(k\)-\emph{cycle}
\((a_1\ a_2\ \cdots\ a_k)\) sends \(a_i\) to \(a_{i+1}\) for
\(1\leq i<k\), sends \(a_k\) to \(a_1\), and fixes every other letter.
\end{definition}

Throughout this section, products of permutations are evaluated from right to
left: \(\sigma\tau\) means ``first apply \(\tau\), then apply \(\sigma\).''
Cycle notation records only the letters that move.  For example,
\((1\ 3\ 2)\in S_4\) sends \(1\mapsto3\), \(3\mapsto2\),
\(2\mapsto1\), and fixes \(4\).  A cycle may be cyclically rewritten without
changing the permutation: \((1\ 3\ 2)=(3\ 2\ 1)\).  Cycles with disjoint
supports commute because they act independently on every letter.  For a
worked noncommuting product,
\[
 (1\ 2\ 3)(1\ 2)=(1\ 3).
\]
Indeed, \(1\mapsto2\mapsto3\), \(3\mapsto3\mapsto1\), and \(2\) is fixed
by the composite.  Reading this calculation from the right is essential.

\begin{proposition}[Cycle Decomposition]
Every permutation in \(S_n\) is a product of disjoint cycles.  This
decomposition is unique up to reordering the cycles and cyclically rewriting
each individual cycle.  The order of a permutation is the least common
multiple of its disjoint cycle lengths.
\end{proposition}
\begin{proof}
Begin with a letter \(a\) and follow its successive images under \(\sigma\):
\[a,\ \sigma(a),\ \sigma^2(a),\ \ldots.\]
Because the set is finite and \(\sigma\)
is injective, the first repetition must be \(a\), giving one cycle.  Repeat
with a letter not yet used.  The resulting cycles have disjoint supports and
their product agrees with \(\sigma\) on every letter.  The orbit of each letter
under repeated application of \(\sigma\) determines its cycle, proving
uniqueness.  A power fixes every cycle exactly when its exponent is divisible
by that cycle length, so the least positive such exponent is the stated lcm.
\end{proof}

\begin{definition}
A \emph{transposition} is a two-cycle \((a\ b)\).  An \emph{adjacent
transposition} is one of the form \((i\ i+1)\).
\end{definition}

\begin{proposition}[Transpositions Generate \(S_n\)]
Every cycle is a product of transpositions, every permutation is a product of
transpositions, and every transposition is a product of adjacent transpositions.
\end{proposition}
\begin{proof}
Directly checking the images of the letters gives
\[(a_1\ a_2\ \cdots\ a_k)=(a_1\ a_k)(a_1\ a_{k-1})\cdots(a_1\ a_2).\]
The cycle-decomposition theorem therefore gives a transposition expression for
every permutation.  If \(i<j\), then
\[(i\ j)=(i\ i+1)\cdots(j-2\ j-1)(j-1\ j)(j-2\ j-1)\cdots(i\ i+1),\]
as is verified by following the images of \(i\), \(j\), and every other
letter.  Thus adjacent transpositions also generate \(S_n\).
\end{proof}

The preliminary definition of homomorphism introduced above is all that is
needed for the next result; its kernel and quotient theory are still deferred
to the later systematic treatment.

\begin{theorem}[Parity and the Sign Homomorphism]
Every expression of a fixed permutation as a product of transpositions has the
same parity of length.  Consequently
\[\operatorname{sgn}:S_n\longrightarrow\{1,-1\},\qquad
\operatorname{sgn}(\sigma)=(-1)^r\]
when \(\sigma\) is written as a product of \(r\) transpositions, is a
well-defined surjective homomorphism.
\end{theorem}
\begin{proof}
For \(\sigma\in S_n\), define its inversion number to be the number of pairs
\((i,j)\) with \(i<j\) and \(\sigma(i)>\sigma(j)\).  Multiplying by an
adjacent transposition on the right swaps two adjacent entries in the list
\(\sigma(1),\ldots,\sigma(n)\).  All inversions involving another entry are
unchanged in pairs, while the inversion between these two entries is reversed;
therefore the inversion number changes by exactly one.  The displayed
adjacent-transposition expression for \((i\ j)\) has
\(2(j-i)-1\) factors.  Thus every transposition changes inversion parity,
and a product of \(r\) transpositions has inversion number congruent to \(r\)
modulo two.  Two transposition expressions for the same permutation therefore
have the same parity.  The definition of \(\operatorname{sgn}\) is
well-defined.  Concatenating a transposition expression for \(\sigma\) with
one for \(\tau\) gives an expression for \(\sigma\tau\), so
\(\operatorname{sgn}(\sigma\tau)=\operatorname{sgn}(\sigma)
\operatorname{sgn}(\tau)\).  A transposition has sign \(-1\), so the map is
onto.
\end{proof}

\begin{definition}
The \emph{alternating group} \(A_n\) is the set of even permutations, namely
\(\{\sigma\in S_n:\operatorname{sgn}(\sigma)=1\}\).
\end{definition}

The sign rule shows directly that the product and inverse of even permutations
are even, so \(A_n\) is a subgroup.  The displayed sign rule says that
\(\operatorname{sgn}\) is a homomorphism.  The general theory of such maps
follows after cosets.  There we will prove that \(A_n\), as the kernel of the
sign map, is normal of index two.
For \(S_3\), \(A_3=\langle(123)\rangle\), and the three transpositions are
the odd permutations.  Later quotient theory will identify the two parity
classes with a copy of \(C_2\).

\section{Cosets, Index, and Lagrange's Theorem}

Every subgroup divides a group into pieces of equal size, whether or not its
left and right translates agree.  A coset is a translate of the subgroup: it remembers
how far an element is from belonging to \(H\), while forgetting the particular
element of \(H\) used to describe that displacement.

For \(H\leq G\), a \emph{left coset} is \(gH=\{gh:h\in H\}\); its number is
\([G:H]\), the \emph{index} of \(H\), when finite.

There are also right cosets \(Hg\).  They need not equal left cosets in a
nonabelian group.  A later condition, called normality, will make the two
notions coincide and permit a consistent multiplication of cosets.

\begin{lemma}
The left cosets of \(H\) partition \(G\), and each has the cardinality of
\(H\).
\end{lemma}
\begin{proof}
If \(g_1h_1=g_2h_2\), then \(g_2^{-1}g_1=h_2h_1^{-1}\in H\), which implies
\(g_1H=g_2H\).  Thus intersecting cosets are equal.  Every element lies in a
coset, and \(h\mapsto gh\) is a bijection \(H\to gH\).
\end{proof}

\begin{proposition}[Coset Equality Criterion]
For \(g,h\in G\), \(gH=hH\) if and only if \(h^{-1}g\in H\).
\end{proposition}
\begin{proof}
If \(gH=hH\), then \(g=ha\) for some \(a\in H\), so \(h^{-1}g\in H\).
Conversely, \(h^{-1}g=a\in H\) gives \(gH=haH=hH\).
\end{proof}

\begin{example}
The cosets of \(n\mathbb Z\) in \(\mathbb Z\) are the residue classes
\(r+n\mathbb Z\).  In \(S_3\), the left cosets of \(\langle(12)\rangle\) each
have two elements, so there are three of them.  This latter partition exists,
but its left and right cosets differ; later we will identify the precise
condition that permits their consistent multiplication.
\end{example}

\begin{example}
In \(D_4\), \(r\langle s\rangle=\{r,rs\}\), whereas
\(\langle s\rangle r=\{r,sr\}=\{r,r^{-1}s\}\).  The two sets differ, just as
a reflection and a rotation of a square do not commute.
\end{example}

\begin{theorem}[Lagrange]
If \(G\) is finite and \(H\leq G\), then \(|G|=[G:H]|H|\).  Consequently,
\(|g|\) divides \(|G|\) for every \(g\in G\).
\end{theorem}
\begin{proofidea}
The preceding lemma has already done the conceptual work.  Once \(G\) is a
disjoint union of translates of \(H\), counting the pieces gives the formula.
\end{proofidea}
\begin{proof}
Choose one representative from each left coset.  The preceding lemma says that
these cosets are pairwise disjoint, cover \(G\), and each has exactly \(|H|\)
elements.  If there are \([G:H]\) of them, counting the elements in this
disjoint union gives \(|G|=[G:H]|H|\).  Taking \(H=\langle g\rangle\) says
that the order of \(g\) divides \(|G|\).
\end{proof}

\begin{corollary}
If \(G\) is finite, then \(g^{|G|}=1\) for every \(g\in G\).
\end{corollary}
\begin{proof}
The order of \(g\) divides \(|G|\).
\end{proof}

\begin{corollary}
A finite group of prime order is cyclic.
\end{corollary}
\begin{proof}
Any element other than the identity has order dividing the prime group order,
so its order is the whole prime and it generates the group.
\end{proof}

The converse to Lagrange's theorem is false: divisibility of \(|H|\) by
\(|G|\) does not guarantee a subgroup of order \(|H|\).  For example, \(A_4\)
has no subgroup of order six, although six divides twelve.  This limitation is
one reason later existence theorems such as Cauchy's and Sylow's are valuable.

\section{Homomorphisms, Kernels, and Images}\label{sec:homomorphisms}

A homomorphism is the appropriate map between groups because it preserves the
operation.  Its kernel measures exactly what information the map discards,
and its image records exactly what it reaches.  These two subgroups are the
bridge between maps and quotient constructions.

\begin{definition}
\leavevmode\\
A \emph{homomorphism} \(\varphi:G\to K\) satisfies
\(\varphi(gh)=\varphi(g)\varphi(h)\).

Its \emph{kernel} is
\(\ker\varphi=\{g:\varphi(g)=1\}\), and its \emph{image} is
\(\operatorname{im}\varphi=\{\varphi(g):g\in G\}\).
\end{definition}

A homomorphism preserves the basic group operations automatically:
\(\varphi(1)=1\), \(\varphi(g^{-1})=\varphi(g)^{-1}\), and
\(\varphi(g^n)=\varphi(g)^n\).  Its kernel records exactly which elements the
map identifies, since \(\varphi(x)=\varphi(y)\) if and only if
\(x^{-1}y\in\ker\varphi\).  Images, by contrast, are subgroups but need not
be normal: \(\langle(12)\rangle\) is the image of \(C_2\to S_3\) but is not
invariant under conjugation.

For example, reduction modulo \(n\), \(\mathbb Z\to\mathbb Z/n\mathbb Z\),
has kernel \(n\mathbb Z\) and is onto.  The sign map
\(S_n\to\{\pm1\}\) has kernel \(A_n\).  The latter example is a useful source
of normal subgroups that is not obtained by inspection.  Likewise, the map
\(D_n\to\{\pm1\}\) that sends rotations to \(1\) and reflections to \(-1\)
has kernel \(\langle r\rangle\).

\begin{proposition}
The kernel and image of a homomorphism are subgroups; moreover,
\(\varphi\) is injective if and only if \(\ker\varphi=\{1\}\).
\end{proposition}
\begin{proof}
The subgroup test gives both subgroup assertions.  Also,
\(\varphi(x)=\varphi(y)\) exactly when
\(\varphi(xy^{-1})=1\), equivalently when \(xy^{-1}\in\ker\varphi\).
\end{proof}

\begin{definition}
A subgroup \(N\leq G\) is \emph{normal}, written \(N\trianglelefteq G\), if
\(gNg^{-1}=N\) for every \(g\in G\).
\end{definition}

\begin{proposition}[Equivalent Forms of Normality]
For \(N\leq G\), the conditions \(gNg^{-1}=N\) and \(gN=Ng\) for every
\(g\in G\) are equivalent.  Thus the left and right cosets of \(N\) coincide
exactly when \(N\) is normal.
\end{proposition}
\begin{proof}
Multiplying \(gNg^{-1}=N\) on the right by \(g\) gives \(gN=Ng\), and the
reverse implication follows by multiplying on the right by \(g^{-1}\).  The
coset statement is merely this equality written for each element \(g\).
\end{proof}

Normality says that left and right translates of \(N\) agree:
\(gN=Ng\) for every \(g\in G\).  Equivalently, membership in \(N\) is
unaffected by conjugation.  This condition is natural rather than cosmetic:
it is exactly what allows the product of two cosets to be independent of the
chosen representatives.

\begin{proposition}
\(\ker\varphi\trianglelefteq G\) for every homomorphism \(\varphi:G\to K\).
\end{proposition}
\begin{proof}
For \(n\in\ker\varphi\), \(\varphi(gng^{-1})=\varphi(g)1\varphi(g)^{-1}=1\).
The reverse inclusion follows on replacing \(g\) by \(g^{-1}\).
\end{proof}

\begin{corollary}
The alternating group \(A_n\) is a normal subgroup of \(S_n\) of index two.
\end{corollary}
\begin{proof}
The sign homomorphism has kernel \(A_n\), so the preceding proposition gives
normality.  It is surjective because a transposition has sign \(-1\).  Its
two fibers are the even and odd permutations; multiplication by a fixed
transposition pairs these two sets bijectively.  Thus each has \(n!/2\)
elements, and \([S_n:A_n]=2\).
\end{proof}

\section{Quotient Groups and Isomorphism Theorems}

The progression \(N\trianglelefteq G\Rightarrow G/N\) is fundamental.  We
identify two elements when they differ by an element of \(N\), then ask whether
the group operation descends to the resulting equivalence classes.  The
quotient criterion answers this question precisely.

\begin{theorem}[Quotient Criterion]
For \(N\leq G\), the rule \((gN)(hN)=(gh)N\) defines a group operation on
the left cosets if and only if \(N\trianglelefteq G\).  The group is \(G/N\).
\end{theorem}
\begin{proofidea}
Changing a representative of \(gN\) multiplies it by an element of \(N\).
Normality is exactly what lets that extra factor move past the representative
of the other coset and be absorbed back into \(N\).
\end{proofidea}
\begin{proof}
If \(N\) is normal, \(g'=gn_1\) and \(h'=hn_2\) imply
\[
g'h'=gh(h^{-1}n_1h)n_2\in ghN,
\]
so the rule is well-defined; associativity, identity \(N\), and inverse
\(g^{-1}N\) follow from \(G\).  Conversely, well-definedness of
\((gN)(nN)=(gN)N\) shows \(gng^{-1}\in N\) for every \(n\in N\), which is
normality.
\end{proof}

\begin{example}
The quotient \(\mathbb Z/n\mathbb Z\) identifies integers whose difference is
a multiple of \(n\).  The sign homomorphism has kernel \(A_n\), so
\(S_n/A_n\cong C_2\).  For the dihedral group
\(D_n=\langle r,s\mid r^n=s^2=1,\ srs^{-1}=r^{-1}\rangle\), the rotation
subgroup \(\langle r\rangle\) is normal and the quotient has two elements:
it records whether a symmetry preserves or reverses orientation.
\end{example}

\begin{example}[Why Normality Is Necessary]
In \(S_3\), let \(H=\langle(12)\rangle\).  It is not normal because
\((123)(12)(123)^{-1}=(23)\notin H\).  Indeed, the proposed coset product is
ambiguous: \(H=(12)H\), but
\[
H(13)H=(13)H\qquad\text{whereas}\qquad(12)H(13)H=(132)H,
\]
and these are different cosets.  The coset set exists, but it is not a
quotient group.
\end{example}

\begin{theorem}[Universal Property of a Quotient]
Let \(N\trianglelefteq G\), let \(\pi:G\to G/N\) be the quotient map, and
let \(f:G\to H\) be a homomorphism with \(N\subseteq\ker f\).  There is a
unique homomorphism \(\overline f:G/N\to H\) satisfying
\(f=\overline f\circ\pi\).
\end{theorem}
The following diagram gives a compact visual record of this factorization:
\[
\begin{tikzcd}[column sep=large, row sep=large]
G \arrow[r, "\pi"] \arrow[dr, "f"'] & G/N \arrow[d, "\overline f"] \\
& H
\end{tikzcd}
\]
A diagram \emph{commutes} when any two directed paths with the same source and
target define the same map.  Here commutativity says exactly
\(f=\overline f\circ\pi\).  This elementary language will be used whenever a
canonical map and a factorization clarify a later construction.
\begin{proofidea}
The quotient identifies \(g\) and \(gn\).  The condition \(f(n)=1\) says
that \(f\) also identifies them, so its value should depend only on the coset.
\end{proofidea}
\begin{proof}
Define \(\overline f(gN)=f(g)\).  If \(gN=g'N\), then \(g'=gn\) for some
\(n\in N\), and \(f(g')=f(g)f(n)=f(g)\); thus the definition is
well-defined.  It is a homomorphism because
\(\overline f((gN)(hN))=f(gh)=f(g)f(h)\).  The equation
\(f=\overline f\circ\pi\) holds by construction.  Finally every coset is
\(\pi(g)\), so this equation determines \(\overline f\) on every element of
\(G/N\), proving uniqueness.
\end{proof}

\begin{example}[A Factorization Computed in Full]
Let \(f:\mathbb Z\to C_6\) send \(1\) to a chosen generator \(u\).  Since
\(f(6k)=u^{6k}=1\), the subgroup \(6\mathbb Z\) is contained in \(\ker f\).
The universal property produces a unique map
\[\overline f:\mathbb Z/6\mathbb Z\longrightarrow C_6,
\qquad \overline f(\overline k)=u^k.\]
Why is this formula safe?  If \(\overline k=\overline\ell\), then
\(k-\ell=6q\), and \(u^k=u^{\ell+6q}=u^\ell(u^6)^q=u^\ell\).  The quotient
is not merely a convenient notation for residues: it is the unique group
through which every map that kills multiples of six must pass.
\end{example}

\begin{remark}[A Reusable Descent Test]
Whenever a proposed formula is defined on cosets, the first question is not
whether it is a homomorphism but whether it is well-defined.  The standard
test is: replace a representative by an equivalent one and show that the
output does not change.  In quotient groups this is usually equivalent to
showing that the subgroup being collapsed lies in a kernel.  The same sequence
of checks will be used for quotient rings and quotient modules.
\end{remark}

\begin{example}[Revisiting Cyclic Groups]
For \(\varphi_g:\mathbb Z\to G\), the image is \(\langle g\rangle\).  The
kernel calculation in the cyclic-group section therefore becomes the uniform
formula \(\langle g\rangle\cong\mathbb Z/\ker\varphi_g\).  When \(g\) has
infinite order this is \(\mathbb Z/\{0\}\cong\mathbb Z\), and when
\(|g|=n\) it is \(\mathbb Z/n\mathbb Z\).  This is precisely the later
theorem promised in the earlier direct proof.
\end{example}

\begin{theorem}[First Isomorphism Theorem]
For a homomorphism \(\varphi:G\to K\),
\[
G/\ker\varphi\cong\operatorname{im}\varphi.
\]
\end{theorem}
\begin{proofidea}
Two elements have the same image under \(\varphi\) precisely when their
quotient lies in \(\ker\varphi\).  Thus the fibers of \(\varphi\) are the
cosets of its kernel, so collapsing those cosets should leave exactly the
image.
\end{proofidea}
\begin{proof}
The map \(g\ker\varphi\mapsto\varphi(g)\) is well-defined because equal cosets
differ by an element of the kernel.  It is a surjective homomorphism with
trivial kernel, hence an isomorphism.
\end{proof}

\begin{remark}[How to Use the First Isomorphism Theorem]
The theorem is not a formula to apply after the fact.  It gives a workflow:
construct a homomorphism that records the feature of interest, calculate what
the homomorphism forgets, and then identify its image.  For the sign map the
feature is parity, the forgotten permutations are the even ones, and the image
is the two-element group.  For reduction modulo \(n\), the feature is a
residue class, the forgotten integers are the multiples of \(n\), and the
image is \(\mathbb Z/n\mathbb Z\).
\end{remark}

\begin{example}
The sign map \(\operatorname{sgn}:S_n\to\{\pm1\}\) is surjective and has
kernel \(A_n\).  The First Isomorphism Theorem therefore gives
\(S_n/A_n\cong\{\pm1\}\cong C_2\).  This is a typical use: calculate a
kernel, then recognize a quotient without multiplying cosets by hand.
\end{example}

\begin{theorem}[Second Isomorphism Theorem]
Let \(N\trianglelefteq G\) and \(H\leq G\).  Then \(HN\leq G\),
\(H\cap N\trianglelefteq H\), and
\[HN/N\cong H/(H\cap N).\]
\end{theorem}
\begin{proof}
The product \(HN\) is a subgroup because normality of \(N\) gives
\[
  (h_1n_1)(h_2n_2)=h_1h_2(h_2^{-1}n_1h_2)n_2\in HN.
\]
The subgroup test gives the remaining closure.  Conjugating an element
of \(H\cap N\) by an element of \(H\) keeps it in both \(H\) and \(N\), so it
is normal in \(H\).  Restrict the quotient map \(G\to G/N\) to \(H\).  Its
kernel is \(H\cap N\), and its image is \(HN/N\).  The First Isomorphism
Theorem gives the claimed isomorphism.
\end{proof}

The theorem says that when \(H\) enters the quotient by \(N\), precisely the
elements already in \(N\), namely \(H\cap N\), are killed.

\begin{theorem}[Third Isomorphism Theorem]
If \(N\trianglelefteq G\) and \(N\leq H\trianglelefteq G\), then
\(H/N\trianglelefteq G/N\) and \((G/N)/(H/N)\cong G/H\).
\end{theorem}
\begin{proof}
Define \(\psi:G/N\to G/H\) by \(\psi(gN)=gH\).  If \(gN=g'N\), then
\(g^{-1}g'\in N\subseteq H\), so \(gH=g'H\); hence \(\psi\) is well-defined.
It is a surjective homomorphism, and its kernel consists of the cosets \(gN\)
with \(g\in H\), which is \(H/N\).  The First Isomorphism Theorem proves the
claim, including normality of the kernel.
\end{proof}

\begin{example}[Quotienting in Stages]
The chain \(12\mathbb Z\subseteq4\mathbb Z\subseteq\mathbb Z\) yields
\[(\mathbb Z/12\mathbb Z)/(4\mathbb Z/12\mathbb Z)
\cong\mathbb Z/4\mathbb Z.\]
The left side first remembers residues modulo twelve and then identifies
residues differing by four; the right side simply remembers residues modulo
four.  This is precisely the content of quotienting in stages.
\end{example}

This is the slogan \highlight{quotienting in stages is the same as quotienting
once}.  It is a basic organizational principle in all later quotient theories.

\begin{theorem}[Correspondence Theorem]
Let \(N\trianglelefteq G\), and let \(\pi:G\to G/N\) be the quotient map.
Then \(H\mapsto H/N\) is an inclusion-preserving bijection between subgroups
of \(G\) containing \(N\) and subgroups of \(G/N\).  Its inverse is
\(K\mapsto\pi^{-1}(K)\).  Under this bijection, normal subgroups correspond,
and \(H\trianglelefteq G\) implies \((G/N)/(H/N)\cong G/H\).
\end{theorem}
\begin{proof}
If \(N\leq H\), then \(H/N\leq G/N\).  If \(K\leq G/N\), then
\(\pi^{-1}(K)\leq G\) and contains \(N\).  The two operations are inverse
because \(\pi\) is surjective: \(\pi^{-1}(H/N)=H\) and
\(\pi(\pi^{-1}(K))=K\).  Conjugation commutes with \(\pi\), proving the
normality assertion.  The final statement is the Third Isomorphism Theorem.
\end{proof}

For instance, subgroups of \(\mathbb Z/n\mathbb Z\) correspond to subgroups
of \(\mathbb Z\) that contain \(n\mathbb Z\).  The containment condition is
necessary: only such subgroups are unions of cosets of \(n\mathbb Z\), so only
they have a meaningful image as a subgroup of the quotient.

\subsection*{Worked Quotient Examples}

The following two examples gather the ideas of normality, quotient groups,
universal properties, and isomorphism theorems in concrete settings.

\begin{example}[Reduction Modulo \(12\)]
Consider \(\pi:\mathbb Z\to\mathbb Z/12\mathbb Z\).  Two integers have the
same image exactly when their difference is a multiple of twelve:
\[
\pi(a)=\pi(b)\quad\Longleftrightarrow\quad a-b\in12\mathbb Z.
\]
Thus the quotient identifies exactly the elements that the map cannot
distinguish.  The subgroup correspondence says that subgroups of
\(\mathbb Z/12\mathbb Z\) correspond to the subgroups of \(\mathbb Z\)
containing \(12\mathbb Z\).  For instance, the subgroup generated by
\(\overline3\) has four elements and corresponds to \(3\mathbb Z\), while the
subgroup generated by \(\overline4\) has three elements and corresponds to
\(4\mathbb Z\).  Inclusion reverses neither upstairs nor downstairs: the
containment \(12\mathbb Z\subseteq4\mathbb Z\subseteq\mathbb Z\) becomes
\(\{\overline0\}\subseteq\langle\overline4\rangle\subseteq
\mathbb Z/12\mathbb Z\).
\end{example}

\begin{example}[The Sign Quotient of \(S_3\)]
The sign map separates the six permutations of \(S_3\) into the three even
permutations \(A_3=\{1,(123),(132)\}\) and the three odd permutations.  The
two cosets of \(A_3\) are therefore \(A_3\) itself and \((12)A_3\).  In the
quotient, every even permutation becomes the identity coset and every odd
permutation becomes the same nontrivial coset.  Its square is the identity,
because the product of two odd permutations is even.  Hence the quotient is
the two-element group.  This makes the abstract statement
\(S_3/A_3\cong C_2\) concrete: the quotient retains parity and forgets every
other feature of a permutation.
\end{example}

\begin{proposition}[A Direct Product of Cyclic Groups]
For positive integers \(m,n\), the group \(C_m\times C_n\) is cyclic if and
only if \(\gcd(m,n)=1\).  In that case \((g,h)\), where \(g\) and \(h\) are
generators of the two factors, generates the product.
\end{proposition}
\begin{proof}
The order of \((g,h)\) is the least positive integer \(k\) for which both
\(g^k=1\) and \(h^k=1\).  Thus its order is \(\operatorname{lcm}(m,n)\).
If \(m\) and \(n\) are coprime, this least common multiple is \(mn\), the
order of the product, so \((g,h)\) is a generator.  Conversely, suppose a
single element generated \(C_m\times C_n\).  Its order would be \(mn\), but
the order of every element divides \(\operatorname{lcm}(m,n)\).  Therefore
\(mn\mid\operatorname{lcm}(m,n)\), which forces \(\gcd(m,n)=1\).
\end{proof}

This example contrasts direct products with quotients.  A quotient imposes
identifications and thereby discards distinctions; a direct product places two
independent systems side by side.  Their universal properties will recur
throughout algebra, but their concrete group-theoretic forms already show two
basic operations of algebraic construction.

\subsection*{A Guided Tour of the Quotient Construction}

It is useful to see the quotient construction operate several times before
leaving Chapter~2.  The same four questions organize each instance:

\begin{enumerate}
\item What elements should be regarded as indistinguishable?
\item Do these equivalence classes form cosets of a normal subgroup?
\item Does the proposed operation depend only on the classes, rather than on
      chosen representatives?
\item What information is retained by the resulting quotient?
\end{enumerate}

For \(\mathbb Z/n\mathbb Z\), the answer to the first question is that two
integers are indistinguishable when their difference is a multiple of \(n\).
The retained information is a remainder.  For \(S_3/A_3\), two permutations are
indistinguishable when they have the same parity; the retained information is
whether the permutation is even or odd.  For \(D_n/\langle r\rangle\), two
symmetries are indistinguishable when they differ by a rotation; the quotient
retains whether the symmetry preserves or reverses orientation.

\begin{example}[The Dihedral Quotient in Detail]
Let \(N=\langle r\rangle\leq D_n\).  The relation \(srs^{-1}=r^{-1}\)
implies \(sr^is^{-1}=r^{-i}\in N\), so conjugation by the reflection keeps
the rotation subgroup inside itself.  Conjugation by a rotation plainly does
the same; therefore \(N\) is normal.  Its two cosets are
\[N=\{1,r,\ldots,r^{n-1}\}\quad\text{and}\quad sN=
\{s,sr,\ldots,sr^{n-1}\}.\]
The product of two elements of the second coset lies in \(N\), because
\[(sr^i)(sr^j)=s(r^is)r^j=r^{j-i}.\]
Thus the nontrivial coset has order two
in the quotient.  This gives a direct, representative-level verification that
\(D_n/N\cong C_2\).
\end{example}

\begin{example}[A Nonnormal Subgroup Revisited]
Let \(H=\langle(12)\rangle\leq S_3\).  It has three left cosets, each with two
elements.  The set of cosets is therefore perfectly meaningful, but the
attempted product of cosets depends on representatives.  Indeed, \(H=(12)H\),
yet multiplying each representative by the coset \((13)H\) produces different
cosets.  The obstruction is not a flaw in the formula; it is the fact that
conjugating \((12)\) by \((13)\) produces a transposition outside \(H\).
The normality condition is exactly what removes this obstruction in the
quotient criterion.
\end{example}

\begin{remark}[Proof Technique: Pass to the Right Quotient]
When a problem contains a homomorphism, a normal subgroup, or a statement that
is unchanged after ``forgetting'' some information, try forming a quotient.
The quotient may turn a complicated relation into the identity relation.  The
First Isomorphism Theorem then identifies the simplified object with an image,
and the Correspondence Theorem translates subgroup questions upstairs and
downstairs.  This is one of the central reusable moves of graduate algebra.
\end{remark}

\subsection*{Why These Constructions Matter Later}

The material of this chapter forms a prototype for several later branches of
algebra.  A quotient group is the first example of imposing an equivalence
relation compatible with an operation.  In ring theory, ideals play the role
of normal subgroups, and the same well-definedness question produces quotient
rings.  In module theory, submodules replace normal subgroups; the isomorphism
theorems and correspondence theorem reappear with almost identical proofs.

Group actions supply a second durable pattern.  When a group acts on a set of
geometric objects, on roots of a polynomial, or on a vector space, the orbit
records what can be moved into what, and the stabilizer records the residual
symmetry at a chosen point.  This is why the elementary formula
\(|G|=|Gx|\,|G_x|\) later becomes a tool in Galois theory, representation
theory, algebraic geometry, and number theory.

Finally, the examples \(S_3\) and \(D_4\) should be regarded as small models
of two different sources of noncommutativity.  In \(S_3\), noncommutativity
comes from rearranging different objects; in \(D_4\), it comes from composing
rotations with reflections.  Both groups have meaningful quotients and
semidirect-product decompositions.  Revisiting them when later constructions
become abstract is an efficient way to keep the structural ideas grounded.

\section{Direct Products and Extensions}

The external direct product \(G\times H\) combines independent symmetries:
one may act in the first coordinate without affecting the second.  For example,
the symmetries of a rectangle with unequal side lengths form \(C_2\times C_2\).
The internal criterion below recognizes this construction inside a larger group.

\begin{proposition}
Suppose \(N,H\trianglelefteq G\), \(G=NH\), and \(N\cap H=\{1\}\).  Then
\(G\cong N\times H\).
\end{proposition}
\begin{proof}
Every element is \(nh\).  Such an expression is unique because
\(n_1h_1=n_2h_2\) implies \(n_2^{-1}n_1=h_2h_1^{-1}\in N\cap H\).
Normality makes the commutator of an element of \(N\) with one of \(H\) lie in
both subgroups, hence be trivial.  Therefore \((n,h)\mapsto nh\) is a
bijective homomorphism.
\end{proof}

\begin{remark}[Looking Ahead]
The quotient map \(G\to G/N\) is the prototype for quotient rings and quotient
modules.  Kernels, images, and quotients will later be packaged into exact
sequences.
\end{remark}

\section*{Exercises}
\begin{exercise}[Basic]
In the multiplicative group \(\mathbb Q^\times\), prove that
\[\langle-1,2\rangle=\{\pm2^k:k\in\mathbb Z\}.\]
Decide whether this subgroup is cyclic.
\end{exercise}
\begin{exercise}[Core]
Prove that a subgroup of index two is normal.
\end{exercise}
\begin{exercise}[Core]
Let \(C_{12}=\langle g\rangle\).  Determine
\(\langle g^3\rangle\cap\langle g^4\rangle\) and the subgroup generated by
\(\langle g^3\rangle\cup\langle g^4\rangle\).  State the corresponding
gcd/lcm pattern for two subgroups of a finite cyclic group.
\end{exercise}
\begin{exercise}[Core]
Let \(H\leq G\) be the unique subgroup of \(G\) having its order.  Prove
that \(H\trianglelefteq G\).  Apply this criterion to the subgroup of order
three in a cyclic group of order twelve.
\end{exercise}
\begin{exercise}[Core]
Let \(N\trianglelefteq G\), and let \(gN\in G/N\).  Prove that the order of
\(gN\) is the least positive integer \(r\) for which \(g^r\in N\), when
such an integer exists.  Apply this to compute the order of \(r\langle r^d\rangle\)
in \(D_n/\langle r^d\rangle\) for \(d\mid n\).
\end{exercise}
\begin{exercise}[Further]
If \(G\) is cyclic of order \(n\), prove
\(\operatorname{Aut}(G)\cong(\mathbb Z/n\mathbb Z)^\times\).
\end{exercise}
\begin{exercise}[Looking Ahead]
Let \(N\trianglelefteq G\).  Use the Correspondence Theorem to show that
maximal subgroups of \(G/N\) correspond to maximal subgroups of \(G\) that
contain \(N\).
\end{exercise}

\chapter[Group Actions and Structure]{Group Actions and Structural Group Theory}

\section{Actions, Orbits, Stabilizers, and Permutation Representations}

An action studies a group by observing the objects it moves.  This makes an
abstract group concrete without choosing generators or a multiplication table:
an action is a systematic way to realize its elements as permutations.

\begin{definition}
An \emph{action} of \(G\) on \(X\) is a map \(G\times X\to X\), written
\((g,x)\mapsto g\cdot x\), satisfying the following two requirements:
\begin{enumerate}
\setlength{\itemsep}{5pt}
\item the group identity does nothing: \(1\cdot x=x\) for all \(x\in X\);
\item successive motions combine according to multiplication:
      \((gh)\cdot x=g\cdot(h\cdot x)\).
\end{enumerate}
The \emph{orbit} of \(x\) is \(Gx=\{g\cdot x:g\in G\}\), and its
\emph{stabilizer} is \(G_x=\{g\in G:g\cdot x=x\}\).
\end{definition}

Equivalently, an action is a homomorphism \(G\to\operatorname{Sym}(X)\).
This equivalence explains the two action axioms: they say exactly that the
identity acts as the identity permutation and a product acts by composition.
The natural action of \(S_n\) on \(\{1,\ldots,n\}\), the action of \(D_n\)
on the vertices of a regular polygon, and the left-translation action of \(G\)
on itself are the basic examples.

It is useful to pause over these examples.  The natural action of \(S_n\) is
transitive, and the stabilizer of \(1\) consists of permutations of the other
\(n-1\) letters.  The geometric action of \(D_n\) is also transitive; the
stabilizer of a vertex has two elements, the identity and the reflection through
that vertex.  The left regular action of \(G\) is transitive with trivial
stabilizer.  Finally, the coset action on \(G/H\) is transitive and has
stabilizer \(H\) at the coset \(H\).  These are the examples to test whenever
an abstract action appears later.

\begin{definition}
Let \(\rho:G\to\operatorname{Sym}(X)\) be the homomorphism associated to an
action.  Its \emph{kernel} is
\[\ker\rho=\{g\in G:g\cdot x=x\text{ for every }x\in X\}.\]
The action is \emph{faithful} if this kernel is trivial.
\end{definition}

Thus an action is faithful exactly when no nonidentity group element is
invisible on \(X\), equivalently when \(\rho\) is injective.  The natural
action of \(S_n\) and the left regular action of any group are faithful.  By
contrast, the action of \(D_4\) on the set of diagonals of a square is not
faithful: distinct geometric symmetries can induce the same diagonal action.

\begin{proposition}
The stabilizer is a subgroup, the orbits partition \(X\), and an action gives
a homomorphism \(G\to\operatorname{Sym}(X)\).
\end{proposition}
\begin{proof}
For the stabilizer, \(1\cdot x=x\), so \(1\in G_x\).  If \(g,h\in G_x\),
then \((gh^{-1})\cdot x=g\cdot(h^{-1}\cdot x)=g\cdot x=x\), so the subgroup
test applies.  If two orbits meet, say \(g\cdot x=h\cdot y\), then applying
\(g^{-1}\) yields \(x=g^{-1}h\cdot y\).  Hence \(x\in Gy\), and applying
group elements shows \(Gx\subseteq Gy\); symmetry gives equality.  Finally,
the assignment \(g\mapsto\rho(g)\), where \(\rho(g)(x)=g\cdot x\), satisfies
\(\rho(gh)=\rho(g)\rho(h)\) by the second action axiom.
\end{proof}

\begin{theorem}[Cayley]
Every group \(G\) is isomorphic to a subgroup of \(\operatorname{Sym}(G)\).
\end{theorem}
\begin{proof}
Let \(G\) act on itself by left translation, \(g\cdot x=gx\).  The associated
homomorphism sends \(g\) to the permutation \(x\mapsto gx\).  If this
permutation is the identity, then evaluating it at \(1\) gives \(g=1\).
Its kernel is therefore trivial, so the homomorphism is injective.
\end{proof}

\begin{example}
For the left action of \(G\) on itself, the orbit of \(1\) is all of \(G\)
and its stabilizer is trivial.  For the conjugation action, the orbit and
stabilizer become a conjugacy class and a centralizer.  These examples show
that one definition organizes both the regular action and the internal
symmetry of a group.
\end{example}

\begin{example}[The Coset Action]
For \(H\leq G\), left multiplication defines \(g\cdot(aH)=(ga)H\) on
\(G/H\).  The action is transitive and the stabilizer of \(H\) is \(H\).
Its kernel is
\[\bigcap_{g\in G}gHg^{-1}.\]
Indeed, an element fixes every coset exactly when it belongs to every conjugate
of \(H\).  This intersection, called the \emph{core} of \(H\), is the largest
normal subgroup of \(G\) contained in \(H\).  The example is an important way
to construct permutation representations from subgroups.
\end{example}

\begin{theorem}[Orbit--Stabilizer]
For a finite group action, \(|Gx|=[G:G_x]\), and therefore
\(|G|=|Gx|\,|G_x|\).
\end{theorem}
\begin{proofidea}
The cosets of \(G_x\) are exactly the different ways to move \(x\).  Two
elements \(g,h\) send \(x\) to the same point precisely when
\(h^{-1}g\) fixes \(x\), which is the equivalence relation defining the cosets.
\end{proofidea}
\begin{proof}
Define \(\Phi:G/G_x\to Gx\) by \(\Phi(gG_x)=g\cdot x\).  To check
well-definedness, suppose \(gG_x=hG_x\).  Then \(h^{-1}g\in G_x\), whence
\(g\cdot x=h\cdot((h^{-1}g)\cdot x)=h\cdot x\).  Conversely, if
\(g\cdot x=h\cdot x\), applying \(h^{-1}\) gives
\((h^{-1}g)\cdot x=x\); thus \(h^{-1}g\in G_x\) and \(gG_x=hG_x\).  This
proves injectivity at the same time.  Surjectivity holds because every point
of the orbit has the form \(g\cdot x\).  Therefore \(|Gx|=[G:G_x]\), and
Lagrange's theorem yields \(|G|=|Gx|\,|G_x|\).
\end{proof}

\begin{example}
Let \(D_4\) act on the four vertices of a square.  The orbit of any vertex is
the whole vertex set, so it has size four.  Its stabilizer contains the identity
and the reflection across the diagonal through that vertex, so it has size two.
Thus \(|D_4|=4\cdot2\), as orbit--stabilizer predicts.
\end{example}

\begin{corollary}[Burnside's Counting Lemma]
The number of orbits of a finite action on \(X\) is
\[\frac{1}{|G|}\sum_{g\in G}|\{x\in X:g\cdot x=x\}|.\]
\end{corollary}
\begin{proof}
Count pairs \((g,x)\) with \(g\cdot x=x\).  Counting by \(g\) gives the
numerator.  Counting by orbits gives \(|G|\) for each orbit, by
orbit--stabilizer.
\end{proof}

\begin{example}[Counting Colourings up to Symmetry]
Let the rotation group \(C_4\) act on the two-colourings of the vertices of a
square.  The identity fixes all \(2^4\) colourings, a quarter-turn fixes two,
the half-turn fixes four, and the other quarter-turn fixes two.  Burnside's
lemma gives \((16+2+4+2)/4=6\) rotational symmetry types.  The point is not
the answer itself: the action replaces a potentially awkward case distinction
by a fixed-point count.
\end{example}

\section{Conjugation, Centralizers, and the Class Equation}

Conjugation answers the question ``which elements look the same from inside
the group?''  It preserves order and many other structural properties.

\begin{definition}
For a group \(G\) and an element \(x\in G\), the \emph{centralizer} of \(x\)
is
\[C_G(x)=\{g\in G:gx=xg\}.\]
The \emph{center} of \(G\) is
\[Z(G)=\{g\in G:gx=xg\text{ for every }x\in G\};\]
thus \(Z(G)\) is the collection of elements that commute with all of \(G\).
For a subgroup \(H\leq G\), the \emph{normalizer} of \(H\) is
\[N_G(H)=\{g\in G:gHg^{-1}=H\}.\]
It consists of the elements whose conjugation preserves \(H\) as a set.
\end{definition}

Each of these sets is a subgroup of \(G\).  For example, the subgroup test
for \(C_G(x)\) follows from the fact that if \(a,b\) commute with \(x\), then
\((ab^{-1})x=a(b^{-1}x)=a(xb^{-1})=(ax)b^{-1}=x(ab^{-1})\).
The same calculation with every \(x\) gives the assertion for \(Z(G)\), and
the identity \((ab^{-1})H(ab^{-1})^{-1}=H\) gives it for \(N_G(H)\).

Conjugation, \(g\cdot x=gxg^{-1}\), has conjugacy classes as its orbits and
centralizer \(C_G(x)\) as the stabilizer.  The elements fixed by all of \(G\)
are precisely the center \(Z(G)\).  The normalizer similarly records the
largest subgroup of \(G\) in which \(H\) is normal.

\begin{lemma}[\(p\)-Group Action Congruence]
If a finite \(p\)-group \(P\) acts on a finite set \(X\), then
\[|X|\equiv|X^P|\pmod p,\]
where \(X^P=\{x\in X:g\cdot x=x\text{ for every }g\in P\}\).
\end{lemma}
\begin{proof}
Partition \(X\) into \(P\)-orbits.  Orbit--stabilizer makes every orbit size
a power of \(p\).  The orbits of size one are exactly the points fixed by all
of \(P\); every other orbit has size divisible by \(p\).  Adding orbit sizes
gives the stated congruence.
\end{proof}

This elementary congruence is a reusable engine: it converts a group action
into the existence of fixed points whenever \(|X|\) is not divisible by \(p\).

\begin{example}[A Fixed Point Forced Modulo \(p\)]
Let a group of order \(p\) act on a set with \(p+1\) elements.  The congruence
gives \(|X^P|\equiv p+1\equiv1\pmod p\), so there is a global fixed point.
The same orbit-counting idea powers the center theorem and Sylow congruence.
\end{example}

\begin{remark}[How to Design an Action Argument]
An action proof begins by choosing a set whose points encode the objects one
wants to study.  To prove a divisibility statement, choose a set with an
accessible cardinality and compare orbit sizes.  To prove a containment
statement, let a subgroup act on a coset space, so that a fixed coset translates
into a conjugate containment.  To study elements inside one group, use
conjugation: its orbits are conjugacy classes and its stabilizers are
centralizers.  In each case the formal calculation is short only after the
right action and the right set have been chosen.
\end{remark}

\begin{theorem}[Class Equation]
If \(x_1,\ldots,x_r\) represent the noncentral conjugacy classes of finite
\(G\), then
\[
|G|=|Z(G)|+\sum_i[G:C_G(x_i)].
\]
\end{theorem}
\begin{proof}
Partition \(G\) into conjugacy classes.  Central elements yield singleton
classes, and orbit--stabilizer gives the size of every other class.
\end{proof}

\begin{example}
In \(S_3\), the three transpositions form one conjugacy class, the two
three-cycles form another, and the identity is central.  Thus
\(6=1+3+2\).  This small example already illustrates the class equation as a
partition into symmetry types rather than a numerical trick.
\end{example}

\begin{example}[Centralizers in \(S_3\)]
The preceding equality can be recovered directly from centralizers.  The
centralizer of \((12)\) is \(\{1,(12)\}\), so orbit--stabilizer gives a
conjugacy class of size \([S_3:C_{S_3}((12))]=3\).  The centralizer of
\((123)\) is \(\langle(123)\rangle\), so its class has size two.  This is a
useful practical lesson: to count conjugates of an element, compute the
subgroup of elements that commute with it.
\end{example}

\begin{corollary}
Every finite \(p\)-group has nontrivial center.
\end{corollary}
\begin{proofidea}
For a \(p\)-group, every noncentral orbit under conjugation has size divisible
by \(p\).  The class equation therefore forces the fixed-point part, the
center, to have size divisible by \(p\) as well.
\end{proofidea}
\begin{proof}
For a noncentral \(x\), the centralizer \(C_G(x)\) is proper.  Its order is a
proper divisor of the order of the \(p\)-group \(G\), hence
\([G:C_G(x)]\) is divisible by \(p\).  Therefore the class equation gives
\[|G|=|Z(G)|+\sum_i[G:C_G(x_i)]\equiv|Z(G)|\pmod p.\]
Since \(p\mid|G|\), also \(p\mid|Z(G)|\).  In particular the center has more
than its identity element.
\end{proof}

\section{The Sylow Theorems and Structural Applications}

Sylow theory turns the divisibility of \(|G|\) into actual subgroups.  Its
proofs are a model use of actions: choose a finite set on which a \(p\)-group
acts, observe that nontrivial orbit sizes are divisible by \(p\), and extract
a fixed point from a count not divisible by \(p\).

\begin{theorem}[Cauchy]
If a prime \(p\) divides the order of a finite group \(G\), then \(G\) has an
element of order \(p\).
\end{theorem}
The unusual-looking tuple set is chosen so that it has a readily computed size
and a natural cyclic symmetry.  Once the first \(p-1\) entries are selected,
the last must be \((g_1\cdots g_{p-1})^{-1}\), so there are exactly
\(|G|^{p-1}\) tuples.  Rotation preserves the condition that the product is
one because a cyclic rotation conjugates the total product by its first entry.
\begin{proofidea}
The set of tuples has cardinality divisible by \(p\), while cyclic rotation
organizes it into orbits of size \(1\) or \(p\).  The fixed points are exactly
the repeated tuples, which encode elements satisfying \(g^p=1\).
\end{proofidea}
\begin{proof}
Let \(X\) be the set of \(p\)-tuples \((g_1,\ldots,g_p)\) with
\(g_1\cdots g_p=1\).  Its first \(p-1\) entries may be chosen freely, so
\(|X|=|G|^{p-1}\), which is divisible by \(p\).  Cyclically rotate these
tuples.  Every orbit has size \(1\) or \(p\), and the fixed tuples are exactly
\((g,\ldots,g)\) with \(g^p=1\).  Thus the number of such \(g\) is divisible
by \(p\).  It includes \(1\), so there is a nonidentity one; its order divides
the prime \(p\), hence equals \(p\).
\end{proof}

This proof illustrates a recurring discovery principle: construct a finite set
whose cardinality is divisible by \(p\), then let a group of order \(p\) act so
that its fixed points encode the desired elements.

\begin{theorem}[Sylow I: existence]
Let \(|G|=p^am\), where \(p\) is prime and \(p\nmid m\).  Then \(G\) has a
subgroup of order \(p^a\).  Such a subgroup is called a \emph{Sylow
\(p\)-subgroup}.
\end{theorem}
\begin{proofidea}
Induct on the group order.  A central element of order \(p\) lets us reduce to
a quotient.  If the center does not supply one, the class equation locates a
proper centralizer that still contains the full \(p\)-part of \(|G|\).
\end{proofidea}
\begin{proof}
We use strong induction on \(|G|\); the assertion is immediate when \(a=0\).
First suppose that \(p\mid |Z(G)|\).  By Cauchy's theorem, \(Z(G)\) contains
a subgroup \(C\) of order \(p\).  It is normal in \(G\), and
\(|G/C|=p^{a-1}m<|G|\).  Let \(\pi:G\to G/C\) be the quotient map.  The
induction hypothesis applied to \(G/C\) gives a subgroup
\(\overline P\leq G/C\) of order \(p^{a-1}\).  Define
\(P=\pi^{-1}(\overline P)\).  Then \(P/C\cong\overline P\), so
\(|P|=|C|\,|\overline P|=p\cdot p^{a-1}=p^a\).

Now suppose \(p\nmid |Z(G)|\).  In the class equation, the noncentral terms
are conjugacy-class sizes \([G:C_G(x_i)]\).  Since \(|G|\equiv0\pmod p\) but
\(|Z(G)|\not\equiv0\pmod p\), at least one of these indices is not divisible
by \(p\).  For that noncentral \(x_i\), the centralizer \(C_G(x_i)\) is a
proper subgroup of \(G\), while
\[
 |C_G(x_i)|=\frac{|G|}{[G:C_G(x_i)]}
\]
still has \(p^a\) as its full \(p\)-part.  The induction hypothesis inside
this strictly smaller group produces a subgroup of order \(p^a\), which is
also a subgroup of \(G\).
\end{proof}

\begin{theorem}[Sylow II: containment and conjugacy]
Let \(Q\) be a Sylow \(p\)-subgroup of a finite group \(G\).  Every
\(p\)-subgroup \(P\leq G\) has a conjugate contained in \(Q\).  In particular,
all Sylow \(p\)-subgroups are conjugate.
\end{theorem}
\begin{proof}
Let \(P\) act by left multiplication on the coset set \(G/Q\).  Its size is
\([G:Q]=m\), which is not divisible by \(p\).  Every \(P\)-orbit has size a
power of \(p\), by orbit--stabilizer.  If there were no fixed coset, the sum
of the orbit sizes would be divisible by \(p\), contrary to \(p\nmid m\).
Thus some coset \(gQ\) is fixed by every element of \(P\).  The equality
\(pgQ=gQ\) for all \(p\in P\) is equivalent to
\(g^{-1}pg\in Q\) for all \(p\in P\), hence to \(g^{-1}Pg\leq Q\).
If \(P\) is itself Sylow, the two groups in this containment have the same
finite order, so they are equal.  This proves conjugacy.
\end{proof}

\begin{theorem}[Sylow III: number of Sylow subgroups]
If \(|G|=p^am\) with \(p\nmid m\), the number \(n_p\) of Sylow
\(p\)-subgroups divides \(m\) and satisfies \(n_p\equiv1\pmod p\).
\end{theorem}
\begin{proof}
Fix a Sylow subgroup \(Q\).  Sylow II makes conjugation by \(G\) transitive
on the set \(\mathcal S\) of Sylow \(p\)-subgroups, so orbit--stabilizer gives
\[n_p=|\mathcal S|=[G:N_G(Q)].\]
Because \(Q\leq N_G(Q)\), the index divides \(|G|/|Q|=m\).

Now let \(Q\) act by conjugation on \(\mathcal S\).  Suppose that \(R\in
\mathcal S\) is fixed.  Then \(Q\leq N_G(R)\); since \(R\trianglelefteq
N_G(R)\), the product \(QR\) is a subgroup.  It is a \(p\)-subgroup
containing the Sylow subgroup \(R\), hence \(QR=R\), and therefore \(Q\leq
R\).  Equal orders give \(Q=R\).  Thus \(Q\) is the unique fixed point of
this action.  Every remaining orbit has size divisible by \(p\), so
\(|\mathcal S|\equiv1\pmod p\).
\end{proof}

\begin{example}
If \(|G|=21\), then \(n_7\) divides \(3\) and is congruent to \(1\) modulo
\(7\); hence \(n_7=1\), so the Sylow \(7\)-subgroup is normal.  This is a
structural conclusion, not a classification of all groups of order \(21\).
\end{example}

\begin{corollary}
A Sylow \(p\)-subgroup is normal if and only if it is the unique Sylow
\(p\)-subgroup.  In particular, every group of order \(45\) has a normal
Sylow \(5\)-subgroup.
\end{corollary}
\begin{proof}
Conjugation carries a Sylow \(p\)-subgroup to another one.  Thus a normal one
is fixed by all conjugations and, by Sylow II, is the only one.  Conversely,
a unique Sylow subgroup is fixed under conjugation, so it is normal.  For
\(|G|=45=5\cdot9\), Sylow III says \(n_5\mid9\) and \(n_5\equiv1\pmod5\);
among \(1,3,9\), only \(1\) has this congruence.
\end{proof}

\section{Semidirect Products and Split Extensions}

A direct product describes two components that commute and have no interaction.
A semidirect product records the next possibility: the second component acts on
the first by automorphisms.  The homomorphism \(\alpha:H\to\operatorname{Aut}(N)\)
is therefore not auxiliary data; it is the rule describing that interaction.

The multiplication formula is forced by the internal situation.  Suppose
\(N\trianglelefteq G\), \(H\leq G\), and every element of \(G\) has a unique
form \(nh\).  Then
\[
 (nh)(n'h')=n\bigl(hn'h^{-1}\bigr)hh'.
\]
Normality puts \(hn'h^{-1}\) back in \(N\), while the element \(h\) acts on
\(N\) by conjugation.  The external definition simply preserves this rule
when the two groups have not yet been placed inside a common ambient group.

\begin{definition}
For \(\alpha:H\to\operatorname{Aut}(N)\), the \emph{semidirect product}
\(N\rtimes_\alpha H\) is \(N\times H\) with
\((n,h)(n',h')=(n\alpha(h)(n'),hh')\).
\end{definition}

\begin{proposition}
This operation makes \(N\rtimes_\alpha H\) a group; \(N\times\{1\}\) is
normal, \(\{1\}\times H\) is a subgroup, and they have trivial intersection.
\end{proposition}
\begin{proof}
The identity is \((1,1)\).  Since every \(\alpha(h)\) is an automorphism,
\[(n,h)^{-1}=(\alpha(h^{-1})(n^{-1}),h^{-1}).\]
For associativity, the first coordinates of the two possible products are
\begin{align*}
 ((n,h)(n',h'))(n'',h'')&:
 n\alpha(h)(n')\alpha(hh')(n''),\\
 (n,h)((n',h')(n'',h''))&:
 n\alpha(h)\bigl(n'\alpha(h')(n'')\bigr).
\end{align*}
They agree because \(\alpha(hh')=\alpha(h)\alpha(h')\).  The two displayed subsets have the
claimed subgroup properties.  Finally,
\[
 (n,h)(n',1)(n,h)^{-1}=
 \bigl(n\alpha(h)(n')n^{-1},1\bigr)\in N\times\{1\},
\]
so \(N\times\{1\}\) is normal.  The two subgroups meet only in \((1,1)\).
\end{proof}

An internal semidirect product is equivalently a split short exact sequence
\[1\longrightarrow N\longrightarrow G\longrightarrow H\longrightarrow1.\]
Here ``split'' means that the quotient map \(q:G\to H\) has a homomorphic
section \(s:H\to G\). Then \(q\) restricted to \(s(H)\) is inverse to \(s\),
so \(N\cap s(H)=1\), where \(N=\ker q\). Moreover, if \(g\in G\) and
\(h=q(g)\), then
\[
q\bigl(gs(h)^{-1}\bigr)=hh^{-1}=1,
\]
so \(gs(h)^{-1}\in N\) and \(g\in Ns(H)\). Hence \(G=Ns(H)\), and
Proposition~\ref{prop:internal-semiproduct} applies. Conversely, an internal
semidirect product gives \(G/N\cong H\), and the inclusion of the complement
gives a homomorphic section. This is only a first encounter with extensions,
but it explains why semidirect products are ubiquitous.

\begin{example}
Let \(C_2=\langle s\rangle\) act on \(C_3=\langle r\rangle\) by inversion,
\(srs^{-1}=r^{-1}\).  Then \(C_3\rtimes C_2\) has presentation
\(\langle r,s\mid r^3=s^2=1,\ srs^{-1}=r^{-1}\rangle\), and the map sending
\(r\) to \((123)\) and \(s\) to \((12)\) identifies it with \(S_3\).  It is
not a direct product because its two factors do not commute.
\end{example}

\begin{example}
More generally, \(D_n\cong C_n\rtimes C_2\), where the nontrivial element of
\(C_2\) acts on \(C_n\) by inversion.  The relation \(srs^{-1}=r^{-1}\) is
exactly the twisting encoded by this action.  Thus the dihedral examples
introduced in Chapter~2 provide a geometric model for every semidirect product
in which a reflection reverses a cyclic rotation.
\end{example}

\begin{proposition}[Internal Semidirect Product]\label{prop:internal-semiproduct}
If \(N\trianglelefteq G\), \(H\leq G\), \(G=NH\), and \(N\cap H=1\), then
\(G\cong N\rtimes H\), where \(H\) acts on \(N\) by conjugation.
\end{proposition}
\begin{proof}
Every element has an expression \(nh\).  It is unique: if \(nh=n'h'\), then
\(n^{-1}n'=hh'^{-1}\) belongs to both \(N\) and \(H\), hence is \(1\).
Sending \((n,h)\) to \(nh\) is therefore bijective.  Its compatibility with multiplication is
\(nhn'h'=n(hn'h^{-1})hh'\), which is exactly the external semidirect-product
law for conjugation.
\end{proof}

\section{Composition Series, Solvability, and Nilpotence}

Composition series make precise the idea of breaking a group into indivisible
layers.  A simple group has no nontrivial proper normal subgroup, so it cannot
be decomposed further by quotients.  Jordan--H\"older says that, although the
intermediate subgroups may differ, the resulting simple layers are intrinsic.

\begin{definition}
A \emph{composition series} is a finite normal series with simple factors.  A
group is \emph{solvable} when it has a finite normal series with abelian factors.
\end{definition}

\begin{example}
The chain \(S_3\trianglerighteq A_3\trianglerighteq1\) is a composition
series, with factors \(C_2\) and \(C_3\).  The order in which the factors are
listed has no intrinsic significance; their isomorphism types do.
\end{example}

\begin{proposition}
Every finite group has a composition series.
\end{proposition}
\begin{proof}
Induct on the group order.  If \(G\ne1\), choose a maximal proper normal
subgroup \(N\), which exists because the group is finite.  Then \(G/N\) is
simple, and induction gives a composition series for \(N\).  Prepending \(G\)
to that series gives one for \(G\).
\end{proof}

\begin{theorem}[Jordan--H\"older]
Any two composition series of a finite group have the same length and the same
simple factors, up to ordering and isomorphism.
\end{theorem}
\begin{proof}
We argue by induction on \(|G|\).  There is nothing to prove for \(G=1\).
Write the two series as
\[
 G=G_0\trianglerighteq G_1=N\trianglerighteq\cdots\trianglerighteq1,
 \qquad
 G=H_0\trianglerighteq H_1=M\trianglerighteq\cdots\trianglerighteq1.
\]
The quotients \(G/N\) and \(G/M\) are simple, so \(N\) and \(M\) are maximal
proper normal subgroups of \(G\).

If \(N=M\), the tails are composition series of the smaller group \(N\).
The induction hypothesis says that these tails have the same length and the
same factors, and adding their common first factor \(G/N\) proves the claim.

Suppose \(N\ne M\).  The product \(NM\) is normal in \(G\): for \(g\in G\),
\(g(NM)g^{-1}=(gNg^{-1})(gMg^{-1})=NM\).  It strictly contains \(N\), since
otherwise \(M\leq N\), and then the maximality of \(M\) would force
\(N=M\).  Maximality of \(N\) consequently gives \(NM=G\).  Put
\(I=N\cap M\).  Since both \(N\) and \(M\) are normal in \(G\), so is \(I\),
and in particular \(I\) is normal in both \(N\) and \(M\).  Applying the Second
Isomorphism Theorem in the forms
\[
 M/(M\cap N)\cong MN/N=G/N,\qquad
 N/(N\cap M)\cong NM/M=G/M
\]
shows that \(M/I\) and \(N/I\) are simple.

Choose a composition series from \(I\) down to \(1\).  Because \(N/I\) is
simple, adjoining \(N\) as a new initial term produces a composition series
of \(N\), whose new factor is \(N/I\); similarly, adjoining \(M\) produces a
composition series of \(M\) with new factor \(M/I\).  By induction, the first constructed series has the same factors as
the tail of the original \(G\)-series beginning at \(N\), and the second has
the same factors as the tail beginning at \(M\).  The complete first factor
list is therefore the factors below \(I\), followed by \(N/I\) and \(G/N\);
the second is the same factors below \(I\), followed by \(M/I\) and \(G/M\).
The displayed isomorphisms identify \(N/I\cong G/M\) and
\(M/I\cong G/N\), so the two lists agree up to interchanging these final two
factors.
\end{proof}

For Galois theory it is useful to recognize solvability without guessing a
normal series.  The commutator \([x,y]=xyx^{-1}y^{-1}\) measures the failure
of \(x\) and \(y\) to commute: \([x,y]=1\) if and only if \(xy=yx\).  Indeed,
\(xyx^{-1}y^{-1}=1\) is equivalent, after multiplying on the right by \(yx\),
to \(xy=yx\).  The subgroup generated by all commutators is the
\emph{derived subgroup} \([G,G]\); define \(G^{(0)}=G\) and
\(G^{(i+1)}=[G^{(i)},G^{(i)}]\).

For later quotient constructions, it matters that \([G,G]\) is normal.
Conjugation sends a commutator to a commutator:
\[
 g[x,y]g^{-1}=[gxg^{-1},gyg^{-1}].
\]
Thus conjugation preserves the subgroup generated by all commutators, and
\([G,G]\trianglelefteq G\).

\begin{proposition}
The quotient \(G/[G,G]\) is abelian and is the largest abelian quotient of
\(G\): every homomorphism from \(G\) to an abelian group factors uniquely
through it.
\end{proposition}
\begin{proof}
Every commutator becomes trivial in the quotient, so its cosets commute.  If
\(f:G\to A\) and \(A\) is abelian, then \(f([x,y])=1\) for all \(x,y\), hence
\([G,G]\leq\ker f\).  The universal property of the quotient map now gives
the unique factorization.
\end{proof}

\begin{example}[The Derived Subgroup of \(S_3\)]
The quotient \(S_3/A_3\cong C_2\) is abelian, so the universal property above
gives \([S_3,S_3]\leq A_3\).  On the other hand, direct calculation gives
\[
 [(12),(23)]=(12)(23)(12)(23)=(123)^2=(132).
\]
This nonidentity element generates \(A_3\), so \(A_3\leq[S_3,S_3]\).  Hence
\([S_3,S_3]=A_3\), and \(S_3/[S_3,S_3]\cong C_2\).
\end{example}

\begin{theorem}
A group \(G\) is solvable if and only if \(G^{(r)}=1\) for some \(r\geq0\).
\end{theorem}
\begin{proof}
If the derived series reaches \(1\), its successive quotients are abelian by
the preceding proposition.  Conversely, let
\(G=G_0\trianglerighteq\cdots\trianglerighteq G_r=1\) have abelian factors.
Inductively, \(G^{(i)}\leq G_i\): if \(G^{(i)}\leq G_i\), then the image of
\(G^{(i)}\) in the abelian group \(G_i/G_{i+1}\) has trivial derived subgroup,
so \(G^{(i+1)}\leq G_{i+1}\).  Hence \(G^{(r)}=1\).
\end{proof}

Solvability asks for a group built by successive abelian quotients.  Nilpotence
is stronger: it asks that the construction proceed from the center outward.
This distinction matters already for \(S_3\), which is solvable but not
nilpotent.

\begin{definition}
Set \(Z_0(G)=1\) and define the \emph{upper central series} inductively by
\[Z_{i+1}(G)/Z_i(G)=Z\bigl(G/Z_i(G)\bigr).\]
The group \(G\) is \emph{nilpotent} if \(Z_c(G)=G\) for some \(c\).  The least
such \(c\) is its nilpotence class (unless \(G=1\)).
\end{definition}

Equivalently, a nilpotent group admits a finite chain
\[
1=N_0\leq N_1\leq\cdots\leq N_c=G
\]
with \(N_{i+1}/N_i\leq Z(G/N_i)\).  Indeed, inductively
\(N_i\leq Z_i(G)\): this is clear for \(i=0\), and centrality of
\(N_{i+1}/N_i\) gives \(N_{i+1}\leq Z_{i+1}(G)\).  Thus
\(G=N_c\leq Z_c(G)\), so \(G\) is nilpotent.  Conversely, the upper central
series itself is such a chain.  The upper central series is therefore the
canonical maximal central series, obtained by repeatedly pulling back the
center of the current quotient.  Thus abelian groups are nilpotent of class
at most one.  A nontrivial nilpotent group has nontrivial center; therefore
\(S_3\), whose center is trivial, is not nilpotent.

\begin{theorem}
Every finite \(p\)-group is nilpotent.
\end{theorem}
\begin{proof}
Induct on \(|G|\).  The center of a nontrivial finite \(p\)-group is
nontrivial, and Cauchy's theorem gives a central subgroup \(C\) of order
\(p\).  By induction, \(G/C\) has an upper central series
\[
 1=\overline Z_0\leq\overline Z_1\leq\cdots\leq\overline Z_c=G/C.
\]
Let \(Z_i\) be the inverse image of \(\overline Z_i\) in \(G\).  Then
\(Z_0=C\), and the chain
\[
 1\leq C=Z_0\leq Z_1\leq\cdots\leq Z_c=G
\]
is central: if \(z\in Z_{i+1}\) and \(g\in G\), then the image of \(z\)
is central modulo \(\overline Z_i\), so \([z,g]\in Z_i\).  The first step
is central because \(C\leq Z(G)\).  Thus this is a central series from
\(1\) to \(G\), proving that \(G\) is nilpotent.
\end{proof}

\begin{proposition}
Every nilpotent group is solvable.
\end{proposition}
\begin{proof}
In the upper central series, each factor
\(Z_{i+1}(G)/Z_i(G)\) lies in the center of \(G/Z_i(G)\), so it is abelian.
Thus the upper central series is a solvable series.
\end{proof}

\begin{proposition}[Subgroups and quotients]
Subgroups and quotients of solvable groups are solvable.
\end{proposition}
\begin{proof}
Let \(G=G_0\trianglerighteq G_1\trianglerighteq\cdots\trianglerighteq
G_r=1\) be a solvable series, and let \(H\leq G\).  Intersecting gives
\[
 H=H\cap G_0\trianglerighteq H\cap G_1\trianglerighteq\cdots
 \trianglerighteq H\cap G_r=1.
\]
Each term is normal in the preceding one.  The map
\(H\cap G_i\to G_i/G_{i+1}\) has kernel \(H\cap G_{i+1}\), so its quotient
\((H\cap G_i)/(H\cap G_{i+1})\) is isomorphic to a subgroup of the abelian
group \(G_i/G_{i+1}\).  Removing repeated terms leaves a solvable series for
\(H\).

If \(q:G\to Q\) is onto, the chain
\(Q=q(G_0)\trianglerighteq q(G_1)\trianglerighteq\cdots\trianglerighteq
q(G_r)=1\) has normal terms.  Each nontrivial factor is an image of a
subgroup of some abelian factor \(G_i/G_{i+1}\), and is therefore abelian.
After repetitions are removed, this is a solvable series for \(Q\).
\end{proof}

\begin{proposition}[Extensions]
If \(N\trianglelefteq G\), and both \(N\) and \(G/N\) are solvable, then
\(G\) is solvable.
\end{proposition}
\begin{proof}
Take a solvable series for \(N\).  Take also a solvable series for \(G/N\)
and pull its terms back under the quotient map \(G\to G/N\).  Splicing the
first series to these inverse images gives a normal series for \(G\); its
factors are, respectively, factors from the series for \(N\) and factors
isomorphic to those from the series for \(G/N\).  They are all abelian.
\end{proof}

\begin{remark}[Looking Ahead]
The action of a Galois group on roots and intermediate fields will make these
methods central again.  Solvability of groups eventually characterizes
solvability of polynomial equations by radicals.
\end{remark}

\section{Two Recurring Structural Laboratories}

The abstract constructions of this chapter become much easier to retain when
they are repeatedly tested on a few familiar groups.  The following two
examples summarize how the same group can illuminate several different ideas.

\subsection*{The Group \(S_3\)}

The six-element symmetric group is the smallest nonabelian group.  It acts
faithfully and transitively on three letters.  Orbit--stabilizer at the letter
\(1\) gives
\[|S_3|=|S_3\cdot1|\,|(S_3)_1|=3\cdot2.\]
The stabilizer is \(\{1,(23)\}\), a copy of \(C_2\).  Thus the action turns a
concrete observation about permutations into a coset calculation.

Under conjugation, its elements separate into the classes
\[\{1\},\qquad\{(12),(13),(23)\},\qquad\{(123),(132)\}.
\]
The class equation is \(6=1+3+2\).  This is more than bookkeeping: it shows
that the three transpositions are structurally indistinguishable within
\(S_3\), whereas a transposition subgroup is not normal because conjugation
moves it to another such subgroup.

The subgroup \(A_3=\langle(123)\rangle\) is normal, and the quotient
\(S_3/A_3\cong C_2\) records parity.  The subgroup \(A_3\) is complemented by
\(\langle(12)\rangle\), and conjugation by \((12)\) inverts \((123)\).  Hence
\[S_3\cong C_3\rtimes C_2.\]
Finally the chain \(S_3\trianglerighteq A_3\trianglerighteq1\) shows that
\(S_3\) is solvable.  Its derived subgroup is \(A_3\), so its abelianization
is \(C_2\).  One group has therefore illustrated actions, cosets, normality,
quotients, conjugacy, semidirect products, composition factors, and
solvability.

\subsection*{The Group \(D_4\)}

The square group \(D_4=\langle r,s\mid r^4=s^2=1,\ srs^{-1}=r^{-1}\rangle\)
has eight elements.  Its geometric action on the four vertices is transitive,
and the stabilizer of a vertex has two elements.  This gives \(8=4\cdot2\),
the orbit--stabilizer formula in a visibly geometric form.

The rotation subgroup \(\langle r\rangle\cong C_4\) is normal because
\(s r^i s^{-1}=r^{-i}\) remains a rotation.  The quotient by rotations is
\(C_2\): it distinguishes orientation-preserving symmetries from reflections.
The relation \(srs^{-1}=r^{-1}\) also displays \(D_4\) as
\(C_4\rtimes C_2\).  This is an especially useful contrast with
\(C_4\times C_2\), whose two factors commute; the dihedral relation gives the
nontrivial twisting.

The center of \(D_4\) is \(\{1,r^2\}\).  Since the quotient by this subgroup
is abelian, the chain \(1\leq\langle r^2\rangle\leq D_4\) is a central series.
Thus \(D_4\) is nilpotent.  Comparing it with \(S_3\), which is solvable but
has trivial center and is not nilpotent, makes the hierarchy
\[\text{abelian}\ \Longrightarrow\ \text{nilpotent}\ \Longrightarrow\
\text{solvable}\]
concrete rather than merely formal.

\begin{remark}[A Study Method]
When a new group-theoretic construction is introduced, revisit \(S_3\) and
\(D_4\).  Ask which subgroups are normal, what the natural actions are, what a
quotient remembers, and whether the group decomposes as a direct or semidirect
product.  This practice develops structural intuition more effectively than a
long list of unrelated calculations.
\end{remark}

\section*{Exercises}
\begin{exercise}[Basic]
Show that the kernel of a permutation representation is the intersection of
all point stabilizers.
\end{exercise}
\begin{exercise}[Core]
Let \(G\) act transitively on a set \(X\), and fix \(x\in X\).  Prove that
the stabilizers of the points of \(X\) are conjugate.  Interpret this result
for the left action of \(G\) on the cosets of a subgroup \(H\).
\end{exercise}
\begin{exercise}[Core]
Compute the conjugacy classes, centralizers, and class equation of \(S_3\).
Explain which features of the computation reflect the geometry of a triangle.
\end{exercise}
\begin{exercise}[Core]
Let \(P\) be a Sylow \(p\)-subgroup of a finite group \(G\).  Prove that
\(P\trianglelefteq G\) if and only if it is the unique Sylow \(p\)-subgroup.
Deduce that any Sylow subgroup whose number is forced to be one by the Sylow
divisibility and congruence conditions is normal.
\end{exercise}
\begin{exercise}[Core]
Let \(G\) be a finite \(p\)-group.  Prove that every nontrivial normal
subgroup of \(G\) meets \(Z(G)\) nontrivially.
\end{exercise}
\begin{exercise}[Further]
Show that \(S_3\cong C_3\rtimes C_2\), but is not a direct product of them.
\end{exercise}
\begin{exercise}[Looking Ahead]
Let \(|G|=p^2\), where \(p\) is prime.  Use the center theorem and the
quotient by a central subgroup of order \(p\) to prove that \(G\) is abelian.
Explain why this does not classify such groups.
\end{exercise}

\part{Rings and Polynomial Rings}
\chapter{Rings, Ideals, and Localization}

The integers have addition and multiplication, tied together by the distributive
laws. A ring retains exactly this structure. The familiar examples---integers,
residue classes, polynomials, and matrices---will give the definitions their
meaning. Ideals permit us to force selected elements to be zero; localization
instead forces selected elements to be invertible.

\section{Rings, Subrings, and Homomorphisms}
\begin{definition}
A \emph{ring} is a set \(R\) with an abelian-group structure under \(+\),
an associative multiplication, and a multiplicative identity \(1\), such
that, for all \(a,b,c\in R\),
\[
\begin{array}{rclcrcl}
(a+b)+c&=&a+(b+c),&&a+b&=&b+a,\\
a+0&=&a,&&a+(-a)&=&0,\\
(ab)c&=&a(bc),&&1a&=&a1=a,\\
a(b+c)&=&ab+ac,&&(a+b)c&=&ac+bc.
\end{array}
\]
A ring is \emph{commutative} if \(ab=ba\) for all \(a,b\). A subring
contains the ambient identity \(1_R\) and is closed under subtraction and
multiplication.
\end{definition}

If \(1=0\), then \(R\) has only one element: for every \(a\in R\),
\[
a=1a=0a=0.
\]
Thus the zero ring is included in the definition.

\begin{unremark}[Ring conventions]
Terminology varies among authors.  Some allow rings without an identity,
and some call a subset a subring when it is closed under subtraction and
multiplication without requiring it to contain the ambient identity.  Here
rings have identities, ring homomorphisms preserve them, and subrings contain
the same identity as the ambient ring.  For example, \(2\mathbb Z\) is
closed under subtraction and multiplication, so it is a subring of
\(\mathbb Z\) under the looser convention, but it is not a subring in these
notes because \(1\notin2\mathbb Z\).
\end{unremark}

The rings \(\mathbb Z,\mathbb Q,\mathbb R\), and a field \(F\), are
commutative. So are \(\mathbb Z/n\mathbb Z\), \(F[x]\), and \(R\times S\),
where operations are coordinatewise. Matrix multiplication makes \(M_n(F)\)
a ring, usually not a commutative one. Thus commutativity must be an
additional hypothesis, not a consequence of the axioms.

\begin{proposition}[Basic identities]
For \(a,b\in R\), one has
\[0a=a0=0,\qquad (-a)b=a(-b)=-(ab),\qquad(-a)(-b)=ab.\]
\end{proposition}
\begin{proof}
Since \(0a=(0+0)a=0a+0a\), additive cancellation gives \(0a=0\);
the proof of \(a0=0\) is analogous. Now
\(ab+(-a)b=(a+(-a))b=0\), giving the next identities. Applying this twice
gives \((-a)(-b)=-\bigl(a(-b)\bigr)=ab\).
\end{proof}

\begin{definition}
A \emph{unit} is an element \(u\) with \(uv=vu=1\) for some \(v\).
A nonzero \(a\) is a \emph{zero divisor} if \(ab=0\) or \(ba=0\) for some
nonzero \(b\). A \emph{division ring} is a nonzero ring in which every
nonzero element is a unit; commutativity is not assumed.  A \emph{field} is
a commutative division ring.  An \emph{integral domain} (or simply a
\emph{domain}) is a commutative ring with \(1\ne0\) and no zero divisors.
\end{definition}

Every field is a domain: if \(ab=0\) and \(a\ne0\), multiply by \(a^{-1}\).
The converse fails, since \(\mathbb Z\) is a domain but \(2\) has no inverse.
In \(\mathbb Z/12\mathbb Z\), \(5\) is a unit, while \(2\cdot6=0\) shows
that \(2\) and \(6\) are zero divisors. The class of \(a\) modulo \(n\) is
a unit exactly when \(\gcd(a,n)=1\), by B\'ezout's identity.

\begin{example}[Hamilton's quaternions]
The set
\[
\mathbb H=\{a+bi+cj+dk:a,b,c,d\in\mathbb R\}
\]
is a division ring under distributive multiplication determined by
\[
i^2=j^2=k^2=ijk=-1.
\]
For example, \(ij=k\), whereas \(ji=-k\); thus \(\mathbb H\) is not
commutative.  Define
\[
\overline{a+bi+cj+dk}=a-bi-cj-dk,
\qquad N(a+bi+cj+dk)=a^2+b^2+c^2+d^2.
\]
Expanding with the defining relations gives
\(q\overline q=\overline q q=N(q)\) and
\(N(qr)=N(q)N(r)\).  Hence, if \(q\ne0\), then \(N(q)>0\) and
\[
q^{-1}=\frac{\overline q}{N(q)}.
\]
The quaternions are therefore the standard example of a noncommutative
division ring.
\end{example}

\begin{theorem}
Every finite integral domain is a field.
\end{theorem}
\begin{proof}
Let \(0\ne a\in R\).  The map \(R\to R\), \(x\mapsto ax\), is injective:
\(ax=ay\) implies \(a(x-y)=0\), so cancellation gives \(x=y\).  Since
\(R\) is finite, the map is surjective.  Thus \(ax=1\) for some \(x\in R\),
and commutativity gives \(xa=1\).  Every nonzero element is invertible.
\end{proof}
\begin{applicationtheorem}[Wedderburn's little theorem]
Every finite division ring is a field.
\end{applicationtheorem}
\begin{remark}
This assertion is considerably deeper than the finite-domain theorem: the
latter uses cancellation and a finite-set argument, whereas the former must
first force commutativity.  Its proof is beyond the needs of the present
discussion.
\end{remark}

\begin{example}[A useful warning]
Over a domain, every unit of \(R[x]\) is constant. This is false in a ring
with nilpotents: in \((\mathbb Z/4\mathbb Z)[x]\),
\[(1+2x)^2=1+4x+4x^2=1.\]
Thus \(1+2x\) is a nonconstant unit. Domain hypotheses in polynomial
arguments are therefore essential.
\end{example}

\begin{definition}
A ring homomorphism \(f:R\to S\) preserves addition, multiplication, and
identity. Its kernel is \(\ker f=\{a:f(a)=0\}\), and its image is
\(\operatorname{im}f=\{f(a):a\in R\}\).
\end{definition}
For example, reduction \(\mathbb Z\to\mathbb Z/n\mathbb Z\) is onto and has
kernel \(n\mathbb Z\). As for group homomorphisms, \(f(0)=0\) and
\(f(-a)=-f(a)\). The determinant is not a ring homomorphism on matrices,
because it does not in general preserve addition.

\section{Ideals and Quotient Rings}
An additive subgroup is not automatically suitable for quotienting: it must
remain stable when multiplied by arbitrary ring elements. This is precisely
the ring analogue of the normal-subgroup condition.

\begin{definition}
A \emph{left ideal} of \(R\) is an additive subgroup \(I\) such that
\(ra\in I\) for every \(r\in R\), \(a\in I\).  A \emph{right ideal} is
defined by the condition \(ar\in I\).  A \emph{two-sided ideal} is both a
left and a right ideal.  Throughout these notes, unless explicitly
qualified, an \emph{ideal} means a two-sided ideal, and we write
\(I\trianglelefteq R\).

For \(A\subseteq R\), \((A)\) is the smallest two-sided ideal containing
\(A\).  If \(A=\{a_1,\ldots,a_m\}\), then
\[(a_1,\ldots,a_m)=
\left\{\sum_{\ell=1}^N r_\ell a_{i_\ell}s_\ell:
N\ge0,\ 1\le i_\ell\le m,\ r_\ell,s_\ell\in R\right\}.\]
When \(R\) is commutative, left, right, and two-sided ideals coincide, and
this simplifies to the familiar sums
\[
\sum_k r_ka_k.
\]
\end{definition}

The ideals \((0)\) and \((1)=R\) are trivial. In \(\mathbb Z\),
\((n)=n\mathbb Z\), while in \(F[x]\), \((x)\) is the set of polynomials
with zero constant term. In \(F[x,y]\), \((x^2,xy)\) consists of sums
\(Ax^2+Bxy\). The odd integers are not an ideal of \(\mathbb Z\), since
\(2\cdot1\) is not odd.  In \(M_n(F)\), the matrices whose last row is zero
form a right ideal but generally not a left ideal.  Quotient rings therefore
require two-sided, not merely one-sided, ideals.

The notation \((A)\) is deliberately parallel to the subgroup generated by
a subset: it records the smallest ideal forced by the elements of \(A\).
In a commutative ring, the principal ideal generated by \(a\) is simply
\((a)=\{ra:r\in R\}\).  Principal ideals reverse divisibility:
\[(a)\subseteq(b)\quad\Longleftrightarrow\quad b\mid a.\]
Indeed, \(a\in(b)\) is exactly the assertion \(a=br\).  Thus, in
\(\mathbb Z\), \((12)\subseteq(6)\subseteq(2)\).  This reversal will
underlie the use of principal ideals, gcds, ascending chains, and
factorization in Chapter~4.

\begin{proposition}
Every ideal of \(\mathbb Z\) is \(n\mathbb Z\) for a unique \(n\ge0\).
\end{proposition}
\begin{proof}
For a nonzero ideal \(I\), choose its least positive member \(n\). For
\(a\in I\), write \(a=qn+r\), \(0\le r<n\). Since \(r=a-qn\in I\),
minimality gives \(r=0\); hence \(I=n\mathbb Z\). The reverse containment
is ideal closure, and \(n\ge0\) makes the generator unique.
\end{proof}

For ideals \(I,J\), define their sum and product by
\[
I+J=\{a+b:a\in I,\ b\in J\},\qquad
IJ=\left\{\sum_{k=1}^n a_kb_k:a_k\in I,\ b_k\in J\right\}.
\]
The sums are finite; the empty sum is \(0\).  Together with \(I\cap J\),
these are ideals.  The sum is the smallest ideal containing both: any ideal
containing \(I\) and \(J\) contains every \(a+b\).  The intersection is their
largest common subideal. Always \(IJ\subseteq I\cap J\).  Finite sums and
products \(I_1+\cdots+I_n\) and \(I_1\cdots I_n\) are defined similarly.
For example, in \(\mathbb Z\),
\[(6)+(10)=(2),\qquad (6)\cap(10)=(30),\qquad (6)(10)=(60).\]
Thus intersection and product need not agree when the ideals are not
comaximal.

\begin{theorem}[Quotient Ring Criterion]
For an additive subgroup \(I\le R\), the rule
\[(a+I)(b+I)=ab+I\]
is well-defined exactly when \(I\) is a two-sided ideal. In that case \(R/I\)
is a ring and \(\pi(a)=a+I\) is a surjective homomorphism.
\end{theorem}
\begin{proof}
If \(a'=a+i\), \(b'=b+j\) with \(i,j\in I\), then
\[(a+i)(b+j)-ab=aj+ib+ij\in I,\]
using closure under multiplication on both sides.  Thus the product is
independent of representatives.  All ring axioms descend from \(R\).
Conversely, \(i+I=0+I\) implies
\((r+I)(i+I)=(r+I)(0+I)\), hence \(ri\in I\), while
\((i+I)(r+I)=(0+I)(r+I)\) gives \(ir\in I\). Thus \(I\) is two-sided.
\end{proof}

\begin{theorem}[Universal Property of Quotients]
If \(f:R\to T\) is a homomorphism and \(I\subseteq\ker f\), then there is a
unique homomorphism \(\bar f:R/I\to T\) such that \(f=\bar f\circ\pi\).
\end{theorem}
\[
\begin{tikzcd}[column sep=large, row sep=large]
R \arrow[r, "\pi"] \arrow[dr, "f"'] & R/I \arrow[d, "\bar f"] \\
& T
\end{tikzcd}
\]
\begin{proof}
Put \(\bar f(a+I)=f(a)\). If \(a+I=b+I\), then \(a-b\in I\), so
\(f(a)=f(b)\); this proves well-definedness. The homomorphism laws follow
from those for \(f\), and every coset is \(\pi(a)\), proving uniqueness.
\end{proof}

Thus \(R/I\) is universal among maps that force elements of \(I\) to zero,
just as quotient groups were in Part~I.  The structural parallel is
\[
\text{normal subgroup }N\trianglelefteq G\ \longleftrightarrow\ G/N,
\qquad
\text{two-sided ideal }I\trianglelefteq R\ \longleftrightarrow\ R/I.
\]
In particular, taking \(I=R\) gives \(R/R\), the zero ring.

\begin{theorem}[Isomorphism and correspondence]
For \(f:R\to S\), \(\ker f\) is an ideal and
\[R/\ker f\cong\operatorname{im}f.\]
Ideals of \(R/I\) correspond to ideals of \(R\) containing \(I\), by
\(J\mapsto\pi^{-1}(J)\) and \(K\mapsto K/I\); in particular,
\((R/I)/(K/I)\cong R/K\) if \(I\subseteq K\).
\end{theorem}
\begin{proof}
The kernel is additive and \(f(ra)=f(r)f(a)=0=f(ar)\) for \(a\in\ker f\).
The map \(a+\ker f\mapsto f(a)\) is well-defined, onto the image, and has
zero kernel. Directly checking the two displayed assignments shows they
are inverse ideal-preserving maps. Finally,
\((r+I)+(K/I)\mapsto r+K\) is well-defined and has zero kernel.
\end{proof}

\section{Characteristic and the Prime Subring}
Let \(R\) be a ring with identity. The canonical homomorphism
\[
\iota:\mathbb Z\longrightarrow R,\qquad n\longmapsto n1_R
\]
records how the integers act inside \(R\). Its kernel is an ideal of
\(\mathbb Z\), so it is either \(\{0\}\) or \(n\mathbb Z\) for a unique
integer \(n>0\).

\begin{definition}\label{def:ring-characteristic}
If \(\ker\iota=\{0\}\), the \emph{characteristic} of \(R\) is \(0\). If
\(\ker\iota=n\mathbb Z\) with \(n>0\), its characteristic is \(n\). The image
of \(\iota\) is the \emph{prime subring} of \(R\).
\end{definition}

Under this definition the zero ring has characteristic \(1\), since
\(1_R=0\) and therefore \(\ker\iota=\mathbb Z\).

The First Isomorphism Theorem gives
\[
\operatorname{im}\iota\cong\mathbb Z/\ker\iota.
\]
Thus in characteristic \(0\) the prime subring is isomorphic to
\(\mathbb Z\), while in characteristic \(p>0\) it is
\(\mathbb Z/p\mathbb Z\), and hence isomorphic to \(\mathbb F_p\) when
\(p\) is prime.

\begin{proposition}[Characteristic of a domain]
\label{prop:domain-characteristic}
If \(R\) is a nonzero integral domain, then
\(\operatorname{char}R\) is either \(0\) or a prime.
\end{proposition}
\begin{proof}
Suppose \(\operatorname{char}R=n>0\). If \(n=ab\) with
\(1<a,b<n\), then minimality of \(n\) gives
\[
a1_R\ne0,\qquad b1_R\ne0,
\]
whereas
\[
(a1_R)(b1_R)=n1_R=0.
\]
This contradicts the absence of zero divisors. Hence \(n\) is prime.
\end{proof}

For example,
\[
\operatorname{char}\mathbb Z=\operatorname{char}\mathbb Q=0,\qquad
\operatorname{char}\mathbb F_p=p,\qquad
\operatorname{char}(\mathbb Z/n\mathbb Z)=n\quad(n\geq2).
\]

\begin{example}[Group rings]
Let \(R\) be a commutative ring and \(G\) a group.  The \emph{group ring}
\(R[G]\) consists of finite formal sums
\[
\sum_{g\in G}a_g g,
\]
with coefficientwise addition and multiplication determined by distributivity
and the multiplication in \(G\):
\[
\left(\sum_g a_g g\right)\left(\sum_h b_h h\right)
=\sum_{g,h}a_gb_h(gh).
\]
The finiteness condition makes each coefficient in the product a finite sum.
For the cyclic group \(C_2=\{1,s\}\), \(s^2=1\), every element is uniquely
\(a+bs\), and the map sending \(s\) to the class of \(x\) gives
\[
R[C_2]\cong R[x]/(x^2-1).
\]
This is an elementary way in which a group relation becomes a polynomial
relation.  The polynomial-ring construction is developed formally in the next
chapter; this is an intentional forward preview.  No representation theory is
needed here.
\end{example}

\section{Prime Ideals, Maximal Ideals, Domains, and Fields}
From this point through localization, rings are assumed to be commutative
and unital; the zero ring remains allowed.  Thus the unqualified word
\emph{ideal} continues to mean a two-sided ideal, but sidedness is automatic.
\begin{definition}
A proper ideal \(\mathfrak p\) is \emph{prime} if
\(ab\in\mathfrak p\) implies \(a\in\mathfrak p\) or \(b\in\mathfrak p\).
A proper ideal \(\mathfrak m\) is \emph{maximal} if no ideal is strictly
between \(\mathfrak m\) and \(R\).
\end{definition}
\begin{theorem}\label{thm:prime-maximal-quotient}
\(\mathfrak p\) is prime iff \(R/\mathfrak p\) is a domain, and
\(\mathfrak m\) is maximal iff \(R/\mathfrak m\) is a field. Hence every
maximal ideal is prime.
\end{theorem}
\begin{proof}
\((a+\mathfrak p)(b+\mathfrak p)=0\) iff \(ab\in\mathfrak p\), giving the
first claim. If \(\mathfrak m\) is maximal and \(a\notin\mathfrak m\), then
\((a)+\mathfrak m=R\), so \(ra+m=1\) for suitable \(r,m\), and
\(r+\mathfrak m\) inverts \(a+\mathfrak m\). Conversely, a field has only
the ideals \(0\) and itself; apply the correspondence theorem.
\end{proof}
In \(\mathbb Z\), \((p)\) is prime and maximal for prime \(p\), while
\((6)\) is neither: \(2\cdot3\in(6)\), but \(2,3\notin(6)\). In
\(\mathbb Z/12\mathbb Z\), the ideal \((2)\) is
maximal because its quotient is \(\mathbb Z/2\mathbb Z\). Later, the
division algorithm will show that \((f)\subset F[x]\) is maximal exactly
when \(f\) is irreducible.

\begin{applicationtheorem}[Maximal ideals above proper ideals]
\label{thm:maximal-ideal-existence}
Let \(R\) be a commutative ring with identity. Every proper ideal \(I\) of
\(R\) is contained in a maximal ideal.
\end{applicationtheorem}
\begin{proof}
Let \(\mathcal P\) be the set of proper ideals of \(R\) that contain \(I\),
ordered by inclusion. It is nonempty because \(I\in\mathcal P\). If
\(\mathcal C\) is a chain in \(\mathcal P\), then
\[
J=\bigcup_{K\in\mathcal C}K
\]
is an ideal containing \(I\): the chain condition puts any two elements of
\(J\) in a common member of \(\mathcal C\), where the ideal operations may
be performed. Moreover \(J\) is proper. Indeed, if \(1\in J\), then
\(1\in K\) for some \(K\in\mathcal C\), contradicting the propriety of
\(K\). Thus every chain has an upper bound in \(\mathcal P\). Zorn's Lemma
now gives a maximal member \(\mathfrak m\) of \(\mathcal P\), which is
maximal among all proper ideals of \(R\), hence is a maximal ideal.
\end{proof}
\begin{unremark}[A set-theoretic qualification]
The preceding proof invokes Zorn's Lemma. For Noetherian rings, the
corresponding maximality conclusion follows from the usual ascending-chain
maximality argument, so no appeal to Zorn's Lemma is needed in the standard
treatment. We do not pursue the precise foundational strength of the general
maximal-ideal theorem here.
\end{unremark}

\section{The Chinese Remainder Theorem}
\begin{definition}
Ideals \(I,J\) of a commutative ring are \emph{comaximal} if \(I+J=R\).
Equivalently,
\[
I+J=R\quad\Longleftrightarrow\quad 1=a+b
\quad\text{for some }a\in I,\ b\in J.
\]
Indeed, \(1\in I+J\) is exactly the displayed assertion, and an ideal
containing \(1\) is all of \(R\).
\end{definition}
\begin{lemma}
If \(I+J=R\), then \(I\cap J=IJ\).
\end{lemma}
\begin{proof}
Certainly \(IJ\subseteq I\cap J\). Choose \(u\in I,v\in J\) with
\(u+v=1\). If \(x\in I\cap J\), then \(x=xu+xv\), and both summands lie in
\(IJ\).
\end{proof}
\begin{theorem}[Chinese Remainder Theorem]
If \(I+J=R\), then \(R/(I\cap J)\cong R/I\times R/J\).
\end{theorem}
\begin{proof}
The homomorphism \(\Phi(r)=(r+I,r+J)\) has kernel \(I\cap J\). With
\(u+v=1\), \(u\in I,v\in J\), the element \(av+bu\) maps to
\((a+I,b+J)\); hence \(\Phi\) is onto. Apply the First Isomorphism Theorem.
\end{proof}
\begin{example}
To solve \(x\equiv2\pmod3\), \(x\equiv3\pmod5\), use
\(6+(-5)=1\), with \(6\in(3)\), \(-5\in(5)\). The proof gives
\(x=2(-5)+3(6)=8\). All solutions are \(8\pmod {15}\).
\end{example}
\begin{theorem}[Finite Chinese Remainder Theorem]
If \(I_1,\ldots,I_n\) are pairwise comaximal ideals of a commutative ring
\(R\), then
\[
R/\bigcap_{i=1}^n I_i\cong\prod_{i=1}^n R/I_i.
\]
Moreover,
\[
\bigcap_{i=1}^n I_i=\prod_{i=1}^n I_i.
\]
\end{theorem}
\begin{proof}
For \(n=2\), both assertions were just proved.  Suppose they hold for
\(I_1,\ldots,I_{n-1}\), and put \(J=I_1\cap\cdots\cap I_{n-1}\).  For each
\(i<n\), choose \(a_i\in I_i\), \(b_i\in I_n\) with \(a_i+b_i=1\).
Then \(\prod_{i<n}a_i\in J\), while
\(1-\prod_{i<n}a_i\in I_n\), by expanding the product.  Thus \(J+I_n=R\).
The two-ideal theorem and the induction hypothesis now give
\[
R/(J\cap I_n)\cong R/J\times R/I_n
\cong\prod_{i=1}^n R/I_i.
\]
The two-ideal product formula gives
\(J\cap I_n=JI_n\), and induction identifies this with
\(I_1\cdots I_n\).
\end{proof}

\section{Localization and Local Rings}
The inclusion \(\mathbb Z\subseteq\mathbb Q\) permits division by every
nonzero integer.  Often we want only selected denominators.  For instance,
\[\mathbb Z[1/2]=\{a/2^n:a\in\mathbb Z,\ n\ge0\}\]
allows division by powers of \(2\), without requiring arbitrary rational
denominators.  Localization is the controlled construction that introduces
chosen denominators: it makes selected elements units while otherwise
changing the ring as little as the required universal property permits.
The canonical map need not be injective when zero divisors are present.

If \(s\in R\), the basic case is
\[R[1/s]=\{1,s,s^2,\ldots\}^{-1}R.\]
Its elements are represented by \(a/s^n\), and \(s\) becomes a unit with
inverse \(1/s\).  Equivalently, \(R[1/s]\) is universal among rings receiving
a homomorphism from \(R\) in which the image of \(s\) is invertible.
\(\mathbb Z[1/2]\) is the first example.  Inverting \(6\) also makes both
\(2\) and \(3\) invertible, so one should not expect the only new units to
be powers of the chosen element.

\begin{definition}
A set \(S\subseteq R\) is \emph{multiplicatively closed} if \(1\in S\) and
\(s,t\in S\) imply \(st\in S\). Define \((a,s)\sim(b,t)\) in \(R\times S\)
when \(u(at-bs)=0\) for some \(u\in S\). The class of \((a,s)\) is \(a/s\),
and \(S^{-1}R\) has operations
\[\frac as+\frac bt=\frac{at+bs}{st},\qquad
\frac as\frac bt=\frac{ab}{st}.\]
\end{definition}
Ordinary fraction arithmetic suggests
\[
\frac as=\frac bt\quad\Longleftrightarrow\quad at=bs.
\]
In a general ring, however, a denominator may be a zero divisor, so the
cancellation implicit in cross multiplication can fail after denominators
are inverted.  The condition \(u(at-bs)=0\) for some \(u\in S\) says exactly
that the cross-difference becomes zero after multiplying by an allowed
denominator; it is therefore the correct replacement.  If every element of
\(S\) is a non-zero-divisor, this condition simplifies to \(at=bs\), since
one may cancel \(u\).  The conditions in the definition are also forced by
fraction arithmetic: if \(s,t\) are allowed denominators, then
\((a/s)(b/t)=ab/(st)\) requires \(st\) to be allowed as well; \(1\) is needed
to represent each original element as \(a/1\).
\begin{proposition}
The relation is an equivalence relation, and these operations are
well-defined.
\end{proposition}
\begin{proof}
Reflexivity uses \(1(as-as)=0\), and symmetry changes the sign. If
\(u(at-bs)=0\), \(v(bq-ct)=0\), then multiplying the first equality by
\(vq\), the second by \(us\), and adding gives \(uv(atq-cst)=0\).
This proves transitivity. If representatives are changed, multiplying the
two witnessing equalities by the other numerator and denominator proves
\(ab/(st)=a'b'/(s't')\). For sums the needed cross-difference is
\(tt'(as'-a's)+ss'(bt'-b't)\), after multiplication by the two witnesses;
it is zero. The remaining ring laws follow from common denominators.
\end{proof}
\begin{proposition}[Fraction fields as localizations]
If \(R\) is an integral domain and \(S=R\setminus\{0\}\), then
\[\operatorname{Frac}(R)=S^{-1}R.\]
Under this identification, \(a/b=c/d\) if and only if \(ad=bc\).
In particular, \(\mathbb Q=\operatorname{Frac}(\mathbb Z)\).
\end{proposition}
\begin{proof}
The set \(S\) is multiplicatively closed, and every one of its elements is
regular.  The defining equivalence relation is consequently ordinary cross
multiplication.  Thus the usual fractions and the localization have the same
representatives, equality relation, and operations.  Taking \(R=\mathbb Z\)
gives the final assertion.
\end{proof}
\begin{example}
In \(R=\mathbb Z/6\mathbb Z\), take \(S=\{1,3\}\). Then
\(0/1=2/1\), since \(3(0-2)=0\), although \(0\ne2\). The witness \(u\) is
therefore essential in rings with zero divisors.  It also shows directly that
the canonical map \(R\to S^{-1}R\) need not be injective.
\end{example}
\begin{proposition}
The canonical map \(\iota:R\to S^{-1}R\), \(a\mapsto a/1\), has kernel
\[\{a\in R:sa=0\text{ for some }s\in S\}.\]
It is injective exactly when every \(s\in S\) is \emph{regular}, meaning
\(sa=0\Rightarrow a=0\).  In the present commutative setting, this says
precisely that \(S\) contains no zero divisors.  Moreover, \(s/1\) is a unit
with inverse \(1/s\).
\end{proposition}
\begin{proof}
\(a/1=0/1\) is exactly the displayed condition. The fraction formulas
prove the homomorphism statement and \((s/1)(1/s)=1\).
\end{proof}
\begin{theorem}[Universal Property of Localization]
If \(f:R\to T\) maps each \(s\in S\) to a unit, there is a unique
\(\widetilde f:S^{-1}R\to T\) such that \(f=\widetilde f\circ\iota\).
\end{theorem}
\[
\begin{tikzcd}[column sep=large, row sep=large]
R \arrow[r, "\iota"] \arrow[dr, "f"'] & S^{-1}R \arrow[d, "\widetilde f"] \\
& T
\end{tikzcd}
\]
\begin{proof}
Set \(\widetilde f(a/s)=f(a)f(s)^{-1}\). A witnessing equality
\(u(at-bs)=0\), after applying \(f\) and cancelling the unit \(f(u)\),
proves this is well-defined. The fraction formulas prove the homomorphism
laws. Finally \(a/s=(a/1)(s/1)^{-1}\), which forces uniqueness.
\end{proof}
Quotients force elements of \(I\) to be zero; localizations force elements
of \(S\) to be units. These are complementary universal constructions.  In
particular, \(R/(s)\) forces \(s\) to be zero, whereas \(R[1/s]\) forces
\(s\) to be invertible.  For example,
\(\mathbb Z/(2)=\mathbb F_2\), while \(\mathbb Z[1/2]\) still has
characteristic zero.

\begin{definition}
A commutative ring is \emph{local} if it has a unique maximal ideal.
\end{definition}
\begin{proposition}
A commutative ring with identity is local if and only if its nonunits form an
ideal.
\end{proposition}
\begin{proof}
If \(\mathfrak m\) is the unique maximal ideal, every element outside it is a
unit: otherwise its principal ideal is proper and, by the
\hyperref[thm:maximal-ideal-existence]{maximal-ideal theorem}, lies in a
maximal ideal, which must be \(\mathfrak m\). Hence the nonunits are exactly
\(\mathfrak m\). Conversely, suppose that the nonunits form an ideal \(N\).
It is proper because \(1\) is a unit. Every proper ideal is contained in
\(N\), since an ideal containing a unit is all of \(R\). Thus \(N\) itself
is maximal, and every maximal ideal, being proper, equals \(N\). This proves
both existence and uniqueness.
\end{proof}
\begin{unremark}
Under this convention the zero ring is not local: it has no maximal ideals.
\end{unremark}

For a prime ideal \(\mathfrak p\), \(1\notin\mathfrak p\), and its complement
\(S=R\setminus\mathfrak p\) is multiplicatively closed.  Indeed, if
\(a,b\notin\mathfrak p\) but \(ab\in\mathfrak p\), primality would put
\(a\) or \(b\) in \(\mathfrak p\), a contradiction.  Hence
\[R_{\mathfrak p}=(R\setminus\mathfrak p)^{-1}R\]
is defined, and its elements are fractions \(a/s\) with \(s\notin\mathfrak p\).
It makes every element outside \(\mathfrak p\) invertible, leaving the
behavior attached to \(\mathfrak p\) visible; informally, it focuses attention
at \(\mathfrak p\).
\begin{theorem}
For a prime ideal \(\mathfrak p\), \(R_{\mathfrak p}\) is local with unique
maximal ideal
\[\mathfrak pR_{\mathfrak p}=\{a/s:a\in\mathfrak p,\ s\notin\mathfrak p\}.\]
\end{theorem}
\begin{proof}
\(R\setminus\mathfrak p\) is multiplicatively closed by primality. The
displayed set is an ideal: numerators in \(\mathfrak p\) remain there under
addition and multiplication. It is proper, since an equality \(1=a/s\)
with \(a\in\mathfrak p\) would say \(u(s-a)=0\) for some
\(u\notin\mathfrak p\), contradicting primality. If \(a\notin\mathfrak p\),
then \(a/s\) has inverse \(s/a\). If \(a\in\mathfrak p\), it cannot be a
unit by the same argument applied to \((a/s)(b/t)=1\). Hence this ideal is
exactly the nonunits; every proper ideal is contained in it.
\end{proof}
For \(\mathbb Z_{(p)}\), the unique maximal ideal is
\(p\mathbb Z_{(p)}\).  Concretely,
\[\mathbb Z_{(p)}=\{a/b\in\mathbb Q:p\nmid b\};\]
every integer not divisible by \(p\) is a unit, and its unique maximal ideal
is formed by the fractions whose numerator is divisible by \(p\).

\[
\begin{array}{c|c|c}
S&\text{localization}&\text{elements made invertible}\\ \hline
\{1,s,s^2,\ldots\}&R[1/s]&\text{one chosen element }s\\
R\setminus\mathfrak p&R_{\mathfrak p}&\text{everything outside }\mathfrak p\\
R\setminus\{0\}\ (R\text{ a domain})&\operatorname{Frac}(R)&
\text{every nonzero element}
\end{array}
\]

\section*{Exercises}
\begin{exercise}[Basic] Determine the units and zero divisors of
\(\mathbb Z/10\mathbb Z\). \end{exercise}
\begin{exercise}[Core] Prove that the preimage of an ideal under a ring
homomorphism is an ideal. Use correspondence to describe ideals of
\(\mathbb Z/18\mathbb Z\). \end{exercise}
\begin{exercise}[Core] Show \((x)\) and \((x-1)\) are comaximal in \(F[x]\),
and describe \(F[x]/(x(x-1))\). \end{exercise}
\begin{exercise}[Core] Determine which of \((2),(3),(4),(6)\) are prime or
maximal in \(\mathbb Z/12\mathbb Z\). \end{exercise}
\begin{exercise}[Further] Construct \(S^{-1}\mathbb Z\cong\mathbb Z[1/2]\)
for \(S=\{1,2,2^2,\ldots\}\), and identify its units. \end{exercise}
\begin{exercise}[Further] Prove
\(\mathbb Z_{(p)}/p\mathbb Z_{(p)}\cong\mathbb F_p\). \end{exercise}

\chapter{Polynomial Rings and Factorization}

Polynomials turn arithmetic in a ring into a tool for studying roots, maps,
and later linear operators.  Over a field, division of polynomials recovers
the Euclidean algorithm; this leads from fields to PIDs and UFDs.

\section{Polynomial Rings and Evaluation}
For a commutative ring \(R\), \(R[x]\) consists of finite sums
\(a_0+a_1x+\cdots+a_nx^n\), with coefficientwise addition and convolution
multiplication.  If \(R\) is a domain and \(f,g\ne0\), then
\(\deg(fg)=\deg f+\deg g\), because the product of leading coefficients is
nonzero.
\begin{proposition}[Degree Rule]
If \(R\) is a domain and \(f,g\in R[x]\) are nonzero, then
\(\deg(fg)=\deg f+\deg g\).  The units of \(R[x]\) are exactly the constant
polynomials whose value is a unit of \(R\).
\end{proposition}
\begin{proof}
Writing \(f=a_mx^m+\cdots\), \(g=b_nx^n+\cdots\), the coefficient of
\(x^{m+n}\) in \(fg\) is \(a_mb_n\), which is nonzero precisely because a
domain has no zero divisors.  Thus a product can have degree zero only when
both factors have degree zero, proving the assertion about units.
\end{proof}
\begin{proposition}[Evaluation] For \(a\in R\), evaluation
\(\operatorname{ev}_a:R[x]\to R\), \(f\mapsto f(a)\), is a ring
homomorphism.  If \(R\) is a field, its kernel is \((x-a)\). \end{proposition}
\begin{proof} The first assertion follows by distributing finite sums.  The
division algorithm below writes \(f=(x-a)q+r\), where \(r\) is constant;
then \(r=f(a)\), proving the kernel statement. \end{proof}

\section{Division Algorithms and Roots}
\begin{theorem}[Division Algorithm] If \(F\) is a field and \(f,g\in F[x]\),
\(g\ne0\), then there are unique \(q,r\in F[x]\) with
\[f=qg+r,\qquad r=0\ \text{or}\ \deg r<\deg g.\]
\end{theorem}
\begin{proof} Write \(g=b_nx^n+\cdots\), with \(b_n\ne0\).  If the current
dividend is \(h=c_mx^m+\cdots\) with \(m\ge n\), subtract
\((c_mb_n^{-1})x^{m-n}g\).  This is possible because \(F\) is a field, and
it cancels the \(x^m\)-term, leaving a polynomial of smaller degree.  Iterating
must stop, since nonnegative degrees cannot decrease forever; collecting the
subtracted terms gives \(f=qg+r\).  If two such expressions exist, then
\((q-q')g=r'-r\).  Unless both sides are zero, the left side has degree at
least \(\deg g\), while the right side has smaller degree, a contradiction.
\end{proof}
\begin{example}[Long Division]
Dividing \(x^3-2x+4\) by \(x-1\), first subtract \(x^2(x-1)\), then
\(x(x-1)\), then \(-1(x-1)\).  Thus
\[x^3-2x+4=(x^2+x-1)(x-1)+3.\]
\end{example}
\begin{corollary}[Factor Theorem] For \(f\in F[x]\), \(f(a)=0\) iff
\(x-a\) divides \(f\). \end{corollary}
\begin{proof} Apply the division algorithm to \(f\) by \(x-a\) and evaluate.
\end{proof}
\begin{corollary} A nonzero polynomial of degree \(n\) over a field has at
most \(n\) roots. \end{corollary}
\begin{proof} Each root supplies a distinct linear factor; induction on
degree proves the assertion. \end{proof}

\section{Euclidean Domains and Principal Ideal Domains}
\begin{definition} A domain \(R\) is \emph{Euclidean} if it has a function
\(\delta:R\setminus\{0\}\to\mathbb N\) such that division with remainder
reduces \(\delta\).  A \emph{principal ideal domain} (PID) is a domain in
which every ideal is principal. \end{definition}
\begin{definition}
For \(a,b\) in a domain, a \emph{greatest common divisor} is an element \(d\)
such that \(d\mid a\), \(d\mid b\), and every common divisor of \(a,b\)
divides \(d\).  A gcd, when it exists, is unique only up to associates.  In
\(\mathbb Z\) one chooses the positive representative, but a general domain
need not have a preferred associate.
\end{definition}
\begin{theorem} Every Euclidean domain is a PID.  In particular \(F[x]\) is a
PID for every field \(F\). \end{theorem}
\begin{proof} Let \(I\ne0\) be an ideal and choose \(d\in I\) with minimal
\(\delta(d)\).  Divide any \(a\in I\) by \(d\): \(a=qd+r\).  Then
\(r\in I\), so minimality gives \(r=0\), hence \(I=(d)\).  Polynomial degree
is a Euclidean function by the division algorithm. \end{proof}
\begin{proposition}[B\'ezout and gcd]
In a PID, \((a,b)=(d)\) if and only if \(d\) is a gcd of \(a,b\), up to
associates.  Consequently every pair has a gcd and
\[d=ra+sb\]
for suitable \(r,s\).
\end{proposition}
\begin{proof}
If \((a,b)=(d)\), then \(a,b\in(d)\), so \(d\mid a,b\).  A common divisor
\(c\) divides every linear combination of \(a,b\), and hence divides
\(d\in(a,b)\).  Thus \(d\) is a gcd.  Conversely, write
\((a,b)=(g)\), which is possible because \(R\) is a PID.  The forward
direction shows that \(g\) is a gcd.  If \(d\) is another gcd, then
\(d\mid g\), since \(d\) is a common divisor, and \(g\mid d\), since \(g\)
is a common divisor.  Thus \(d\) and \(g\) are associates, so
\((d)=(g)=(a,b)\).  Consequently \(d\in(a,b)\), giving the displayed
B\'ezout representation.  This connects gcds, generated ideals, and
B\'ezout identities.
\end{proof}

\section{Unique Factorization Domains and Gauss's Lemma}
The factorization of an integer into primes is a model for a property enjoyed
by many domains.  The central question here is whether the property survives
the passage from a UFD \(R\) to \(R[x]\).  The answer is yes, and Gauss's
lemma explains why coefficients do not introduce hidden factorizations.

\begin{definition}
Let \(R\) be a domain.  Nonzero \(a,b\in R\) are \emph{associates} if
\(a=ub\) for a unit \(u\).  A nonzero nonunit \(p\) is \emph{irreducible} if
\(p=ab\) implies that \(a\) or \(b\) is a unit.  It is \emph{prime} if
\(p\mid ab\) implies \(p\mid a\) or \(p\mid b\).  A \emph{UFD} is a domain
where every nonzero nonunit has a factorization into irreducibles, unique up
to ordering and replacement by associates.
\end{definition}

Every prime is irreducible: if \(p=ab\), then \(p\mid a\) or \(p\mid b\);
cancellation makes the other factor a unit.  The converse need not hold in
an arbitrary domain, but it does hold in a PID.

\begin{proposition}[Irreducibles are prime in a PID]
Every irreducible element of a principal ideal domain is prime.
\end{proposition}
\begin{proof}
Let \(p\) be irreducible, suppose \(p\mid ab\), and assume \(p\nmid a\).
The gcd \(d\) of \(p,a\) cannot be associated to \(p\), so it is a unit.
Thus \((p,a)=R\), and B\'ezout gives \(rp+sa=1\).  After multiplication by
\(b\), both terms on the left of \(rpb+sab=b\) are divisible by \(p\).
Hence \(p\mid b\).
\end{proof}

\phantomsection\label{thm:irreducible-maximal-field}
\begin{applicationtheorem}[Irreducible polynomials and field quotients]
Let \(F\) be a field and let \(f\in F[x]\) be nonconstant.  The following are
equivalent:
\[
f\text{ is irreducible},
\qquad
(f)\text{ is a maximal ideal of }F[x],
\qquad
F[x]/(f)\text{ is a field}.
\]
\end{applicationtheorem}
\begin{proof}
Since \(F[x]\) is a PID, every ideal containing \((f)\) has the form
\((g)\), where \(g\mid f\).  If \(f\) is irreducible, then \(g\) is either a
unit or an associate of \(f\).  Thus there is no ideal strictly between
\((f)\) and \(F[x]\), so \((f)\) is maximal.  Conversely, if
\(f=gh\) with both \(g\) and \(h\) nonunits, then
\[
(f)\subsetneq(g)\subsetneq F[x],
\]
contradicting maximality.  The equivalence between maximal ideals and field
quotients is the quotient characterization of maximal ideals from Chapter~4.
\end{proof}

Irreducibility therefore does more than forbid a factorization: it is exactly
what makes the quotient obtained by adjoining a formal root into a field.

\begin{theorem}[Every PID is a UFD]
Every principal ideal domain is a unique factorization domain.
\end{theorem}
\begin{proofidea}
Existence comes from the ascending-chain condition on ideals: an element that
could not be factored would force a strictly increasing chain of principal
ideals.  Uniqueness follows by cancelling prime factors.
\end{proofidea}
\begin{proof}
An ascending chain of ideals in a PID stabilizes, since its union is an ideal
\((d)\), and \(d\) belongs to one member of the chain.  Suppose some nonzero
nonunit does not factor into irreducibles, and choose \(a\) with \((a)\)
maximal among ideals generated by such elements.  Then \(a=bc\), with both
\(b,c\) nonunits.  Therefore \((a)\subsetneq(b)\) and
\((a)\subsetneq(c)\), so maximality gives factorizations of \(b,c\), a
contradiction.  Thus factorizations exist.

If \(p_1\cdots p_m=q_1\cdots q_n\) are products of irreducibles, \(p_1\) is
prime by the proposition, hence divides some \(q_j\).  Since \(q_j\) is
irreducible, they are associates.  Reorder, absorb the unit, and cancel.
Induction gives uniqueness.
\end{proof}

Let \(R\) now be a UFD.  A gcd of coefficients is only determined up to a
unit, precisely the ambiguity that factorization theory permits.

Gcds exist in every UFD, although they need not have B\'ezout
representations.  After choosing representatives of irreducible associate
classes, write
\[a=u\prod p_i^{\alpha_i},\qquad b=v\prod p_i^{\beta_i}.\]
Then, up to associates,
\[\gcd(a,b)\sim\prod p_i^{\min(\alpha_i,\beta_i)}.\]
The exponents show at once that this element divides both \(a\) and \(b\),
and that every common divisor divides it.  In a PID, gcds are B\'ezout
combinations; in a general UFD this is false.  For example,
\(\gcd(x,y)=1\) in \(F[x,y]\), but \((x,y)\ne F[x,y]\): an identity
\(fx+gy=1\) would become \(0=1\) after evaluation at \(x=y=0\).
Thus factorization-coprime and comaximal agree in a PID but not in every UFD.

\begin{definition}[Content]
For \(0\ne f=a_0+\cdots+a_nx^n\in R[x]\), its \emph{content} is
\[\operatorname{cont}(f)=\gcd(a_0,\ldots,a_n),\]
where the gcd is understood up to multiplication by a unit.  The polynomial
is \emph{primitive} exactly when \(\operatorname{cont}(f)\) is a unit,
equivalently when no nonunit divides all of its coefficients.
\end{definition}

The content is therefore simply the gcd of the coefficients, whose existence
is supplied by the UFD property just discussed.

Every nonzero \(f\) has the form
\[f=\operatorname{cont}(f)f_0\]
up to multiplication by a unit, with \(f_0\) primitive.
For example,
\[6x^2-9x+3=3(2x^2-3x+1),\]
and the parenthesized factor is primitive in \(\mathbb Z[x]\).

\begin{lemma}[Primitive Product Lemma]
The product of two primitive polynomials in \(R[x]\) is primitive.
\end{lemma}
\begin{proof}
If a nonunit divided every coefficient of \(fg\), some irreducible factor
\(p\) would do so.  In a UFD, irreducibles are prime, so \((p)\) is a prime
ideal: \(ab\in(p)\) means \(p\mid ab\).  Hence \(R/(p)\) is a domain by
Chapter~4.  The reductions of \(f\) and \(g\) modulo \(p\) are nonzero,
because \(f,g\) are primitive, but their product is zero.  This contradicts
the degree rule in the polynomial ring over a domain.
\end{proof}

\begin{corollary}[Content formula]
For nonzero \(f,g\in R[x]\),
\[\operatorname{cont}(fg)\sim\operatorname{cont}(f)\operatorname{cont}(g),\]
where \(\sim\) denotes association.  Equality is naturally only up to
associates in a general UFD.  For \(R=\mathbb Z\), normalizing content to
be positive gives the literal equality
\(\operatorname{cont}(fg)=\operatorname{cont}(f)\operatorname{cont}(g)\).
\end{corollary}
\begin{proof}
Write \(f=\operatorname{cont}(f)f_0\),
\(g=\operatorname{cont}(g)g_0\), where \(f_0,g_0\) are primitive.
Then \(fg=\operatorname{cont}(f)\operatorname{cont}(g)f_0g_0\), and the
lemma says the final factor has unit
content.
\end{proof}

Let \(K=\operatorname{Frac}(R)\).  The next statement is the precise form
of clearing denominators.

\begin{lemma}[Gauss's factorization lemma]
If \(f\in R[x]\) is primitive and \(f=gh\) in \(K[x]\), then there are
primitive \(g_0,h_0\in R[x]\) and a unit \(\varepsilon\in R\) such that
\[f=(\varepsilon g_0)h_0.\]
Thus a positive-degree factorization over \(K\) yields one over \(R\).
\end{lemma}
\begin{proof}
Choose nonzero \(a,b\in R\) such that \(G=ag\) and \(H=bh\) lie in \(R[x]\).
Write \(G=\operatorname{cont}(G)g_0\),
\(H=\operatorname{cont}(H)h_0\) with primitive parts.  Since
\(GH=abf\), the content formula gives
\[\operatorname{cont}(G)\operatorname{cont}(H)
\sim\operatorname{cont}(GH)=\operatorname{cont}(abf)\sim ab.\]
Thus \(\operatorname{cont}(G)\operatorname{cont}(H)=\varepsilon ab\) for a
unit \(\varepsilon\), and
\[abf=GH=\varepsilon abg_0h_0.\]
Cancellation in the domain \(R[x]\) gives the result.  Scaling does not
change degree.
\end{proof}

\begin{corollary}[Irreducibility transfer]
For primitive \(f\in R[x]\),
\[f\text{ is irreducible in }R[x]\quad\Longleftrightarrow\quad
f\text{ is irreducible in }K[x].\]
\end{corollary}
\begin{proof}
A factorization in \(R[x]\) is one in \(K[x]\).  Conversely, a nontrivial
factorization in \(K[x]\) has two positive-degree factors; the lemma turns it
into one in \(R[x]\).
\end{proof}

\begin{lemma}
If primitive \(g,h\in R[x]\) are associates in \(K[x]\), then they are
associates in \(R[x]\).
\end{lemma}
\begin{proof}
Write \(h=(a/b)g\) with \(a,b\in R\setminus\{0\}\).  Then \(bh=ag\).
Taking contents gives \(b\sim a\), so \(a/b\) is a unit of \(R\).
\end{proof}

\begin{theorem}[Polynomial UFD Theorem]
If \(R\) is a UFD, then \(R[x]\) is a UFD.  Hence
\(R[x_1,\ldots,x_n]\) is a UFD for every \(n\ge1\).
\end{theorem}
\begin{proof}
The field \(K\) has \(K[x]\) a PID, hence a UFD.  For \(0\ne f\in R[x]\),
factor its content in \(R\), then factor its primitive part in \(K[x]\).
The factorization lemma clears the denominators of the latter factors; the
irreducibility-transfer corollary makes the resulting primitive factors
irreducible in \(R[x]\).  Irreducible constants of \(R\) remain irreducible
in \(R[x]\), by the degree rule.  This proves existence.

For uniqueness, compare two factorizations in the UFD \(K[x]\).  Its
uniqueness pairs the primitive factors up to association, and the preceding
lemma upgrades these associations to \(R[x]\).  Unique factorization of the
remaining contents occurs in \(R\).  Repeated application of the one-variable
result gives the multivariable assertion.
\end{proof}

\begin{corollary}
\(\mathbb Z[x]\) and \(\mathbb Q[x]\) are UFDs.  A primitive polynomial in
\(\mathbb Z[x]\) is irreducible over \(\mathbb Z\) exactly when it is
irreducible over \(\mathbb Q\).
\end{corollary}
\begin{proof}
Apply the theorem to the UFD \(\mathbb Z\), whose fraction field is
\(\mathbb Q\), and use irreducibility transfer.
\end{proof}

\section{Irreducibility Criteria and Factorization over Fields}
\begin{proposition}
A polynomial of degree two or three over a field is reducible if and only if
it has a root in that field.
\end{proposition}
\begin{proof} A root yields a linear factor by the Factor Theorem.  Conversely,
a proper factorization of degree two or three has a factor of degree one,
which has a root. \end{proof}
\begin{theorem}[Eisenstein's Criterion]
Let \(f=a_nx^n+\cdots+a_0\in\mathbb Z[x]\).  If a prime \(p\) divides every
\(a_i\) for \(i<n\), does not divide \(a_n\), and \(p^2\nmid a_0\), then
\(f\) is irreducible over \(\mathbb Q\).
\end{theorem}
\begin{proof} By Gauss's Lemma it is enough to rule out a factorization in
\(\mathbb Z[x]\); the factors may be chosen primitive.  Hence neither has
all coefficients divisible by \(p\), so their reductions modulo \(p\) are
nonzero.  Reducing a hypothetical \(f=gh\) modulo \(p\) gives
\(\overline g\,\overline h=\overline {a_n}x^n\).  Since the reduced factors
are nonzero and \(\mathbb F_p[x]\) is a domain, each is a nonzero scalar
times a power of \(x\).  Thus every nonleading coefficient of \(g,h\) is
divisible by \(p\).  Their constant terms are therefore both divisible by
\(p\), contradicting \(p^2\nmid a_0\).
\end{proof}
\begin{example} Eisenstein at \(p=2\) proves \(x^n-2\) irreducible over
\(\mathbb Q\) for every \(n\ge1\). \end{example}

\section*{Exercises}
\begin{exercise}[Basic] Divide \(x^3-1\) by \(x-1\) and factor the result over \(\mathbb Q\). \end{exercise}
\begin{exercise}[Core] Use the Euclidean algorithm in \(\mathbb Q[x]\) to
find a gcd of \(x^3-1\) and \(x^2-1\), and express it as a B\'ezout combination.
\end{exercise}
\begin{exercise}[Core] Determine whether \(x^3-2\) and \(x^4+1\) are irreducible over \(\mathbb Q\). \end{exercise}
\begin{exercise}[Core] Find the content and primitive part of each of
\(12x^3-18x^2+6x\) and \(15x^2+10x+5\).  Verify the content formula for
their product. \end{exercise}
\begin{exercise}[Core] Prove directly from the definition that a primitive
polynomial in \(\mathbb Z[x]\) cannot have all of its coefficients divisible
by the same prime.  Use reduction modulo \(2\) to show that
\(x^3+x+1\) is irreducible over \(\mathbb Q\).
\emph{Hint:} use Gauss's Lemma; a nontrivial factorization of this primitive
polynomial would reduce to one in \(\mathbb F_2[x]\), where the cubic has no
root. \end{exercise}
\begin{exercise}[Further] Use B\'ezout to show that coprime polynomials in \(F[x]\) have no common root in any extension field. \end{exercise}
\begin{exercise}[Further, Challenging] Let \(R\) be a UFD and let
\(f\in R[x]\) be primitive.  Show that if \(f\) divides \(g\) in
\(\operatorname{Frac}(R)[x]\) and \(g\in R[x]\), then there is a
primitive polynomial \(h\in R[x]\) and \(a\in R\) such that
\(g=ahf\) up to multiplication by a unit. \end{exercise}

\part{Modules and Linear Algebra}
\chapter{Modules and Linear Algebra}\label{chap:modules-linear-algebra}

Vector spaces combine addition with scalar multiplication by numbers from a
field. Modules keep the same formal operations while allowing scalars from a
ring. This modest-looking change includes abelian groups, ideals, quotient
rings, and vector spaces in one language. It also explains why linear algebra
is the first important special case of module theory: over a field one can
divide by every nonzero scalar, whereas a general module can have torsion and
need not possess coordinates.

Throughout \(R\) is a ring with identity, and a module means a unital left
\(R\)-module unless a different convention is stated. For a noncommutative
ring this is important: a right module has its scalar multiplication on the
other side. From \S6.8 onward, \(F\) denotes a field and \(V,W\) denote
finite-dimensional \(F\)-vector spaces.

\section{Modules and Submodules}
The vector-space axioms make sense before one assumes that nonzero scalars
are invertible. That observation gives the correct setting for relations such
as \(n\bar a=0\) in \(\mathbb Z/n\mathbb Z\), which cannot occur in a
nonzero vector space.

\begin{definition}
An \emph{\(R\)-module} is an abelian group \((M,+)\) equipped with a map
\(R\times M\to M\), \((r,m)\mapsto rm\), such that, for \(r,s\in R\) and
\(m,n\in M\),
\[
\begin{aligned}
(r+s)m&=rm+sm, & r(m+n)&=rm+rn,\\
(rs)m&=r(sm), & 1m&=m.
\end{aligned}
\]
A subset \(N\subseteq M\) is a \emph{submodule}, written \(N\leq M\), if
\(m-n\in N\) and \(rm\in N\) whenever \(m,n\in N\) and \(r\in R\).
An \emph{\(R\)-module homomorphism} \(f:M\to P\) preserves addition and
scalar multiplication. Its \emph{kernel} and \emph{image} are
\(\ker f=\{m:f(m)=0\}\) and \(\operatorname{im}f=\{f(m):m\in M\}\).
An isomorphism is a bijective module homomorphism.
\end{definition}

The principal examples should be kept in view. A vector space is a module
over its scalar field; abelian groups are exactly \(\mathbb Z\)-modules;
\(R\) and \(R^n\) are left \(R\)-modules; and every left ideal is a
submodule of the regular module \(R\). If \(R\) is commutative, every ideal
has this role. The module \(\mathbb Z/n\mathbb Z\) is not free over
\(\mathbb Z\), because its nonzero elements are killed by a nonzero scalar.

There is a useful one-line test hidden in the definition. A nonempty subset
\(N\) of a module is a submodule precisely when it is closed under
subtraction and under multiplication by every scalar. The subtraction test
contains both additive closure and additive inverses, just as it did for
subgroups. For instance, \(2\mathbb Z\leq\mathbb Z\) as a
\(\mathbb Z\)-module, while the set of odd integers is not a submodule:
the difference of two odd integers is even. Such examples are elementary,
but they train the reader to distinguish closure conditions from merely
familiar-looking subsets.

\begin{proposition}
For a module homomorphism \(f:M\to P\), both \(\ker f\) and
\(\operatorname{im}f\) are submodules.
\end{proposition}
\begin{proof}
If \(m,n\in\ker f\), then \(f(m-n)=0\), and \(f(rm)=rf(m)=0\). Thus the
kernel is a submodule. If \(f(m),f(n)\in\operatorname{im}f\), then
\(f(m)-f(n)=f(m-n)\) and \(rf(m)=f(rm)\), so the image is one as well.
\end{proof}

\section{Quotient Modules and the Isomorphism Theorems}
Submodules are precisely the pieces that can be collapsed to zero without
destroying scalar multiplication. Quotient modules therefore generalize the
quotient groups and quotient rings developed earlier.

\begin{definition}
If \(N\leq M\), the \emph{quotient module} \(M/N\) is the additive quotient
group with scalar multiplication \(r(m+N)=rm+N\).
\end{definition}
\begin{proof}
If \(m+N=m'+N\), then \(m-m'\in N\), whence \(r(m-m')\in N\). Therefore
\(rm+N=rm'+N\), so the operation is well-defined. Each module axiom follows
from the corresponding axiom in \(M\).
\end{proof}

\begin{theorem}[Module isomorphism theorems]
Let \(f:M\to P\) be linear.
\begin{enumerate}
\item \(M/\ker f\cong\operatorname{im}f\).
\item If \(A,B\leq M\), then \((A+B)/B\cong A/(A\cap B)\).
\item If \(N\leq L\leq M\), then \((M/N)/(L/N)\cong M/L\).
\item Submodules of \(M/N\) correspond bijectively to submodules of \(M\)
containing \(N\), under inverse image and quotient by \(N\).
\end{enumerate}
\end{theorem}
\begin{proof}
For (1), \(m+\ker f\mapsto f(m)\) is well-defined, onto the image, and has
zero kernel. For (2), restrict \(M\to M/B\) to \(A\): its kernel is
\(A\cap B\), and its image is \((A+B)/B\). For (3), map
\((m+N)+(L/N)\) to \(m+L\), checking directly that the map is well-defined
and has zero kernel. Finally, if \(\pi:M\to M/N\) is the quotient map,
then \(Q\mapsto\pi^{-1}(Q)\) and \(K\mapsto K/N\) are inverse operations.
\end{proof}

\begin{theorem}[Universal property of a quotient]
If \(f:M\to P\) is linear and \(N\subseteq\ker f\), there is a unique
linear map \(\bar f:M/N\to P\) with \(f=\bar f\pi\), where \(\pi\) is the
quotient map.
\end{theorem}
\[
\begin{tikzcd}[column sep=large, row sep=large]
M \arrow[r, "\pi"] \arrow[dr, "f"'] & M/N \arrow[d, "\bar f"] \\
& P
\end{tikzcd}
\]
\begin{proof}
Set \(\bar f(m+N)=f(m)\). If \(m-m'\in N\), then \(f(m-m')=0\), so this
does not depend on the representative. It is linear, and surjectivity of
\(\pi\) forces uniqueness.
\end{proof}

Thus the group, ring, and module isomorphism theorems are not unrelated
coincidences: they express one recurring quotient pattern.

\begin{example}[A quotient module]
The quotient \(\mathbb Z/6\mathbb Z\) is a \(\mathbb Z\)-module, and its
submodule generated by \(\bar2\) is \(\{\bar0,\bar2,\bar4\}\). Collapsing
this submodule gives a quotient with two elements,
\[(\mathbb Z/6\mathbb Z)/\langle\bar2\rangle\cong\mathbb Z/2\mathbb Z.
\]
The example is the module version of reducing a group by a normal subgroup;
the point of the module language is that multiplication by every integer also
descends automatically.
\end{example}

\section{Direct Sums and Products}
Products collect arbitrary coordinate choices; direct sums collect finite
linear combinations of contributions from separate modules. The latter is the
algebraic construction needed for bases and for decompositions.

\begin{definition}
For a family \((M_i)_{i\in I}\), the product
\[
\prod_{i\in I}M_i
\]
consists of all tuples \((m_i)\), with coordinatewise operations. The direct
sum
\[
\bigoplus_{i\in I}M_i
\]
is its submodule consisting of
tuples with \emph{finite support}: only finitely many \(m_i\) are nonzero.
For a finite family, product and direct sum have the same elements and are
both denoted
\[
M_1\oplus\cdots\oplus M_n.
\]
\end{definition}

There are canonical injections and projections
\[
\iota_i:M_i\longrightarrow\bigoplus_{i\in I}M_i,
\qquad
\pi_i:\prod_{i\in I}M_i\longrightarrow M_i.
\]
A map
\[
X\longrightarrow\prod_{i\in I}M_i
\]
is
equivalent to a family of maps \(X\to M_i\), obtained by composing with the
projections. Conversely, a family \(M_i\to Y\) determines a unique map
\[
\bigoplus_{i\in I}M_i\longrightarrow Y,
\]
because every element is a finite sum of injected
components. For example,
\(\mathbb Z/2\mathbb Z\oplus\mathbb Z/3\mathbb Z\cong\mathbb Z/6\mathbb Z\).

For an infinite family the finite-support requirement matters. The element
\[
(1,1,1,\ldots)\in\prod_{n\geq1}\mathbb Z
\]
does not belong to
\[
\bigoplus_{n\geq1}\mathbb Z.
\]
A map out of the direct sum can be defined
by adding only finitely many contributions, whereas there is no reason an
infinite sum should make sense in an arbitrary module. This is why direct
sums, rather than products, are the natural home for bases.

\section{Cyclic Modules, Annihilators, Torsion, and Exact Sequences}
The simplest module is generated by one element. Understanding such modules
is useful because a generator has exactly one kind of obstruction: some
scalars may kill it. The annihilator records those scalar relations and turns
cyclic modules into quotients of the ring itself.

\begin{definition}
An \(R\)-module is \emph{cyclic} if \(M=Rm\) for some \(m\in M\). The
\emph{annihilator} of \(m\) is
\[\operatorname{Ann}_R(m)=\{r\in R:rm=0\};\]
the annihilator of \(M\) consists of scalars that kill every element of
\(M\). If \(R\) is a domain, an element \(m\) is \emph{torsion} if
\(rm=0\) for some nonzero \(r\in R\). A module is torsion-free if its only
torsion element is zero.
\end{definition}

For a left module, an annihilator is a left ideal; in the commutative rings
used later it is an ideal in the unqualified sense. Over \(\mathbb Z\),
torsion means finite additive order. Thus \(\mathbb Z^n\) is torsion-free,
whereas every element of \(\mathbb Z/n\mathbb Z\) is torsion.

When \(R\) is a domain, the set \(T(M)\) of torsion elements is a submodule,
called the \emph{torsion submodule}. Indeed, if \(am=0\) and \(bn=0\) with
\(a,b\ne0\), then \(ab(m-n)=0\); and \(a(rm)=0\) for every \(r\in R\).
The domain hypothesis ensures \(ab\ne0\). This construction will make the
free part of a finitely generated PID module intrinsic in Chapter~7.

\begin{theorem}[Cyclic-module classification]
For \(m\in M\), there is an isomorphism
\[Rm\cong R/\operatorname{Ann}_R(m).\]
Consequently every cyclic module over a commutative ring is isomorphic to
\(R/I\) for an ideal \(I\).
\end{theorem}
\begin{proof}
The map \(R\to Rm\), \(r\mapsto rm\), is a surjective module
homomorphism, and its kernel is exactly \(\operatorname{Ann}_R(m)\). The
first isomorphism theorem supplies the stated isomorphism. Conversely,
\(R/I\) is generated by \(1+I\).
\end{proof}

Exact sequences provide a compact language for saying that a module is built
from a submodule and a quotient. They will later organize module
decompositions and further constructions.
\begin{definition}
A sequence \(A\xrightarrow{u}B\xrightarrow{v}C\) is \emph{exact at
\(B\)} if \(\operatorname{im}u=\ker v\). A \emph{short exact sequence}
\[0\longrightarrow A\xrightarrow{i}B\xrightarrow{q}C\longrightarrow0\]
is exact everywhere; equivalently, \(i\) is injective, \(q\) is surjective,
and \(\operatorname{im}i=\ker q\).
\end{definition}

\begin{undefinition}[Cokernel]
For a module homomorphism \(f:M\to N\), its \emph{cokernel} is
\[
\operatorname{coker}(f)=N/\operatorname{im}f.
\]
The canonical map \(N\to\operatorname{coker}(f)\) is the quotient map.
\end{undefinition}

Kernels measure what a map kills; cokernels measure what remains in its target
after the image is collapsed.  Concretely, if \(g:N\to P\) satisfies
\(gf=0\), then \(g\) vanishes on \(\operatorname{im}f\), so there is a
unique map \(\widetilde g:\operatorname{coker}(f)\to P\) with
\[
\begin{tikzcd}[column sep=large,row sep=large]
N\arrow[r]\arrow[dr,"g"']&\operatorname{coker}(f)\arrow[d,"\widetilde g"]\\
&P.
\end{tikzcd}
\]

The quotient sequence \(0\to N\to M\to M/N\to0\) is the basic example.
It is useful to distinguish equality of short exact sequences from the
weaker fact that their three terms happen to be isomorphic.

\begin{definition}
An \emph{isomorphism of short exact sequences} is a commutative diagram
\[
\begin{tikzcd}[column sep=large]
0\arrow[r]&A\arrow[r,"i"]\arrow[d,"\alpha", "\sim"']&
B\arrow[r,"q"]\arrow[d,"\beta", "\sim"']&
C\arrow[r]\arrow[d,"\gamma", "\sim"']&0\\
0\arrow[r]&A'\arrow[r,"i'"]&B'\arrow[r,"q'"]&C'\arrow[r]&0
\end{tikzcd}
\]
in which the three vertical maps are module isomorphisms.
\end{definition}

\begin{proposition}[The canonical split sequence]\label{prop:canonical-split-sequence}
For modules \(A\) and \(C\), the sequence
\[
0\longrightarrow A\xrightarrow{a\mapsto(a,0)}A\oplus C
\xrightarrow{(a,c)\mapsto c}C\longrightarrow0
\]
is short exact.
\end{proposition}
\begin{proof}
The first map is injective and the second is surjective.  Its kernel consists
precisely of the pairs \((a,0)\), which are exactly the image of the first
map.  Thus exactness says, quite literally, that \(A\oplus C\) contains
\(A\) as its first coordinate and has quotient \(C\).
\end{proof}

\section{Splitting and Diagram Chases}
The preceding observation asks when a short exact sequence contains no
extension information beyond its two ends.  A map \(s:C\to B\) with
\(qs=1_C\) is a \emph{section}; a map \(r:B\to A\) with \(ri=1_A\) is a
\emph{retraction}.

\begin{theorem}[Splitting Lemma]\label{thm:splitting-lemma}
For a short exact sequence
\[
0\longrightarrow A\xrightarrow{i}B\xrightarrow{q}C\longrightarrow0,
\]
the following are equivalent: \(q\) has a section; \(i\) has a retraction;
and the given short exact sequence is isomorphic, as a short exact sequence,
to the canonical split sequence of
Proposition~\ref{prop:canonical-split-sequence}.  Concretely, the end maps
are \(1_A\) and \(1_C\), and the middle isomorphism
\(\Phi:A\oplus C\to B\) satisfies
\[
\Phi(a,0)=i(a),
\qquad
q\Phi(a,c)=c.
\]
\end{theorem}
\begin{proof}
A section gives the map
\[
\Phi:A\oplus C\longrightarrow B,\qquad
\Phi(a,c)=i(a)+s(c).
\]
It is onto: for \(b\in B\), the element \(b-sq(b)\) lies in \(\ker q\), hence
is \(i(a)\) for some \(a\).  If \(\Phi(a,c)=0\), applying \(q\) gives
\(c=0\), and then injectivity of \(i\) gives \(a=0\).  Thus \(\Phi\) is an
isomorphism, and it gives an isomorphism to the canonical sequence in
Proposition~\ref{prop:canonical-split-sequence}.  The map
\(r=\operatorname{pr}_A\Phi^{-1}\) is a retraction of \(i\).

Conversely, suppose that \(ri=1_A\).  For \(c\in C\), choose \(b\in B\) with
\(q(b)=c\), and put
\[
s(c)=b-i(r(b)).
\]
If \(b'=b+i(a)\) is another lift, then
\[
b'-i(r(b'))=b+i(a)-i(r(b)+a)=b-i(r(b)),
\]
so the formula is independent of the lift.  It is linear, and
\(q(s(c))=c\).  Hence it is a section.
\end{proof}

Commutative diagrams organize a \emph{diagram chase}: start with an element,
use exactness to replace a kernel element by a preimage, and use commutativity
to transport the resulting equality vertically.  The Five Lemma is the basic
instance.

\begin{theorem}[Five Lemma]\label{thm:five-lemma}
In a commutative diagram with exact rows
\[
\begin{tikzcd}[column sep=small]
A_1\arrow[r]\arrow[d,"f_1"]&A_2\arrow[r]\arrow[d,"f_2"]&A_3\arrow[r]\arrow[d,"f_3"]&A_4\arrow[r]\arrow[d,"f_4"]&A_5\arrow[d,"f_5"]\\
B_1\arrow[r]&B_2\arrow[r]&B_3\arrow[r]&B_4\arrow[r]&B_5
\end{tikzcd}
\]
if \(f_1\) is surjective, \(f_2,f_4\) are isomorphisms, and \(f_5\) is
injective, then \(f_3\) is an isomorphism.
\end{theorem}
\begin{proof}
Write the horizontal maps as \(d_i:A_i\to A_{i+1}\) and
\(e_i:B_i\to B_{i+1}\).  Suppose first that \(f_3(x)=0\).  Then
\(f_4d_3(x)=e_3f_3(x)=0\).  Since \(f_4\) is injective,
\(d_3(x)=0\), so \(x=d_2(y)\).  As
\[
0=f_3d_2(y)=e_2f_2(y),
\]
there is \(z\in B_1\) with \(e_1(z)=f_2(y)\).  Choose \(w\in A_1\) with
\(f_1(w)=z\).  Then
\[
f_2(y-d_1(w))=0.
\]
Injectivity of \(f_2\) gives \(y=d_1(w)\), hence \(x=d_2d_1(w)=0\).
This proves the \emph{injectivity} of \(f_3\).

\emph{Surjectivity.} Now take \(b\in B_3\).  Since \(e_3(b)\in B_4\), choose
\(a_4\in A_4\) with \(f_4(a_4)=e_3(b)\).  Exactness gives
\(d_4(a_4)=0\), because \(f_5d_4(a_4)=e_4f_4(a_4)=0\) and \(f_5\) is
injective.  Choose \(a_3\in A_3\) with \(d_3(a_3)=a_4\).  Then
\[
e_3(b-f_3(a_3))=0,
\]
so \(b-f_3(a_3)=e_2(b_2)\) for some \(b_2\in B_2\).  Choose
\(a_2\in A_2\) with \(f_2(a_2)=b_2\).  Replacing \(a_3\) by
\(a_3+d_2(a_2)\) gives \(f_3(a_3)=b\).  This proves the \emph{surjectivity}
of \(f_3\).

\emph{Isomorphism.} A module homomorphism is an isomorphism precisely when it
is both injective and surjective.  The preceding two parts prove the claim.
\end{proof}

\begin{theorem}[Snake Lemma]\label{thm:snake-lemma}
In a commutative diagram with exact rows
\[
\begin{tikzcd}[column sep=large]
0\arrow[r]&A'\arrow[r,"i'"]\arrow[d,"f_A"]&
B'\arrow[r,"q'"]\arrow[d,"f_B"]&
C'\arrow[r]\arrow[d,"f_C"]&0\\
0\arrow[r]&A\arrow[r,"i"]&
B\arrow[r,"q"]&
C\arrow[r]&0,
\end{tikzcd}
\]
there is an exact sequence
\[
\begin{gathered}
0\longrightarrow\ker f_A\longrightarrow\ker f_B
\longrightarrow\ker f_C\xrightarrow{\delta}\operatorname{coker}f_A\\
\longrightarrow\operatorname{coker}f_B\longrightarrow
\operatorname{coker}f_C\longrightarrow0.
\end{gathered}
\]
\end{theorem}
\begin{proof}
For \(c'\in\ker f_C\), choose \(b'\in B'\) with \(q'(b')=c'\).  Then
\[
q(f_B(b'))=f_C(q'(b'))=0,
\]
so \(f_B(b')=i(a)\) for a unique \(a\in A\).  Define
\[
\delta(c')=a+\operatorname{im}f_A.
\]
If \(b'\) is replaced by \(b'+i'(a')\), then \(a\) is replaced by
\(a+f_A(a')\); hence \(\delta\) is well-defined.  Choosing lifts of two
elements and adding them also proves that \(\delta\) is linear.

The maps before and after \(\delta\) are the maps induced by the vertical
arrows.  The first one is injective because \(i'\) is injective.  Exactness
at the first two kernels follows immediately by chasing the two squares and
using injectivity of \(i\).  If \(c'=q'(b')\) comes from
\(\ker f_B\), then \(f_B(b')=0\), so \(\delta(c')=0\).  Conversely,
\(\delta(c')=0\) means that \(a=f_A(a')\); then
\(b'-i'(a')\) lies in \(\ker f_B\) and still maps to \(c'\).  This proves
exactness at \(\ker f_C\).

If \(a+\operatorname{im}f_A\) maps to zero in \(\operatorname{coker}f_B\),
then \(i(a)=f_B(b')\) for some \(b'\).  Hence
\(f_C(q'(b'))=0\), and the construction gives
\(\delta(q'(b'))=a+\operatorname{im}f_A\).  This is exactness at
\(\operatorname{coker}f_A\).  Finally, exactness at
\(\operatorname{coker}f_B\) and \(\operatorname{coker}f_C\) follows by the
same two-square chase: a class \(b+\operatorname{im}f_B\) maps to zero
precisely when \(q(b)=f_C(c')\), and a lift of \(c'\) corrects \(b\) to an
element of \(i(A)\).  The last induced map is surjective because \(q\) is
surjective.  This also proves the asserted endpoint exactness.
\end{proof}

The construction of \(\delta\) is the model diagram chase: choose a lift,
measure the obstruction in the next kernel, and record it modulo the
ambiguity forced by the first vertical map.  It is natural with respect to
maps of such diagrams, a fact used later in homological algebra.

\section{Free Modules and Module Bases}
Coordinates are useful only when expansions are unique. Free modules record
that condition over a general ring; vector spaces are the especially well
behaved free modules over fields.

\begin{definition}
A subset \(S\subseteq M\) \emph{spans} \(M\) if every element is a finite
linear combination of members of \(S\). It is \emph{linearly independent}
if every relation
\[
\sum_i r_is_i=0,
\]
with distinct \(s_i\in S\), has all
coefficients zero. A \emph{basis} is an independent spanning set, and a
module with a basis is \emph{free}. The free module on a set \(X\) is the
module \(R^{(X)}\) of finitely supported functions \(X\to R\), with basis
\(\{e_x:x\in X\}\).
\end{definition}

\begin{theorem}[Universal property of a free module]
For every function \(u:X\to M\), there is a unique linear map
\(R^{(X)}\to M\) taking \(e_x\) to \(u(x)\).
\end{theorem}
\begin{proof}
Every element of \(R^{(X)}\) has a unique finite expression
\[
\sum_{x\in X}r_xe_x.
\]
Send it to
\[
\sum_{x\in X}r_xu(x).
\]
This is linear, and any
linear map with the required values has this formula.
\end{proof}

Not every module is free: \(\mathbb Z/n\mathbb Z\) has no \(\mathbb Z\)-basis
for \(n>1\). Indeed, if \(\bar a\) belonged to a \(\mathbb Z\)-basis, then
\(n\bar a=0\) would be a nontrivial relation. This is the first warning that
the familiar basis theory of vector spaces cannot simply be imported into
module theory. In particular, a submodule of a free module need not be free
over an arbitrary ring, and the apparent number of generators of \(R^n\)
requires justification before it is called a rank.

\section{Projective and Injective Modules}
Free modules make lifting possible because a map out of a free module is
controlled by its values on a basis.  Projective modules are precisely the
modules that retain this lifting property, even when they need not themselves
have a basis.

\begin{definition}
A module \(P\) is \emph{projective} if, for every surjection
\(\pi:M\twoheadrightarrow N\) and every map \(f:P\to N\), there is a map
\(\widetilde f:P\to M\) such that \(\pi\widetilde f=f\).  A module \(E\) is
\emph{injective} if, for every injection \(j:A\hookrightarrow B\) and every
map \(f:A\to E\), there is a map \(\widetilde f:B\to E\) such that
\(\widetilde f j=f\).
\end{definition}
\[
\begin{tikzcd}[column sep=large,row sep=large]
&P\arrow[dl,dashed,"\widetilde f"']\arrow[d,"f"]\\
M\arrow[r,two heads,"\pi"']&N\arrow[r]&0
\end{tikzcd}
\qquad
\begin{tikzcd}[column sep=large,row sep=large]
&E\\
0\arrow[r]&A\arrow[u,"f"]\arrow[r,hook,"j"']&B\arrow[ul,dashed,"\widetilde f"']
\end{tikzcd}
\]
These matching triangles display the dual geometry of the two definitions:
projectivity is a lifting problem, while injectivity is an extension problem.

\begin{definition}
For left \(R\)-modules \(M,N\), write
\[
\operatorname{Hom}_R(M,N)
\]
for the set of \(R\)-linear maps \(M\to N\).  It is an abelian group under
pointwise addition:
\[
(f+g)(m)=f(m)+g(m).
\]
If \(R\) is commutative, it is also an \(R\)-module by
\[
(rf)(m)=r f(m).
\]
\end{definition}
For a noncommutative ring the displayed scalar formula need not define a
left-module map, since \(r f(sm)\) need not equal \(s r f(m)\).  Accordingly,
the exact sequences below are always sequences of abelian groups, and only
under commutativity are they automatically sequences of \(R\)-modules.

\begin{proposition}[Hom detects exactness]\label{prop:hom-exactness}
Suppose
\[
0\longrightarrow A\xrightarrow{i}B\xrightarrow{q}C\longrightarrow0
\]
satisfies \(qi=0\), without yet assuming exactness.  Define
\[
\begin{aligned}
i_*:\operatorname{Hom}_R(M,A)&\longrightarrow\operatorname{Hom}_R(M,B),
&f&\longmapsto i\circ f,\\
q_*:\operatorname{Hom}_R(M,B)&\longrightarrow\operatorname{Hom}_R(M,C),
&g&\longmapsto q\circ g,
\end{aligned}
\]
\[
\begin{aligned}
q^*:\operatorname{Hom}_R(C,N)&\longrightarrow\operatorname{Hom}_R(B,N),
&h&\longmapsto h\circ q,\\
i^*:\operatorname{Hom}_R(B,N)&\longrightarrow\operatorname{Hom}_R(A,N),
&k&\longmapsto k\circ i.
\end{aligned}
\]
Then
\[
0\longrightarrow\operatorname{Hom}_R(M,A)\xrightarrow{i_*}
\operatorname{Hom}_R(M,B)\xrightarrow{q_*}\operatorname{Hom}_R(M,C)
\]
is exact for every \(R\)-module \(M\) if and only if
\[
0\longrightarrow A\xrightarrow{i}B\xrightarrow{q}C
\]
is exact.  Dually,
\[
0\longrightarrow\operatorname{Hom}_R(C,N)\xrightarrow{q^*}
\operatorname{Hom}_R(B,N)\xrightarrow{i^*}\operatorname{Hom}_R(A,N)
\]
is exact for every \(R\)-module \(N\) if and only if
\[
A\xrightarrow{i}B\xrightarrow{q}C\longrightarrow0
\]
is exact.  Thus the original sequence is short exact precisely when both
displayed Hom tests are exact for all \(M\) and \(N\).
\end{proposition}
\begin{proof}
Assume first that \(0\to A\to B\to C\) is exact.  A map \(M\to B\) has
zero composite with \(q\) exactly when its image lies in
\(\ker q=\operatorname{im}i\), and it then factors uniquely through \(A\);
this proves exactness of the first Hom sequence.  Conversely, suppose this
Hom sequence is exact for every \(M\).  With \(K=\ker i\), the inclusion
\(K\to A\) has zero image under \(i_*\), hence is zero; so \(i\) is
injective.  With \(K=\ker q\), the inclusion \(K\to B\) has zero image
under \(q_*\), hence factors through \(i\).  Thus
\(\ker q\subseteq\operatorname{im}i\), and the reverse inclusion follows
from \(qi=0\).

Now assume that \(A\to B\to C\to0\) is exact.  A map \(B\to N\) vanishes on
\(\operatorname{im}i=\ker q\) exactly when it factors uniquely through
\(C\), proving exactness of the second Hom sequence.  Conversely, suppose
that sequence is exact for every \(N\).  Taking
\[
N=C/\operatorname{im}q
\]
and the quotient map \(C\to N\), injectivity of \(q^*\) forces that quotient
map to be zero; hence \(q\) is surjective.  Taking
\[
N=B/\operatorname{im}i
\]
and the quotient map \(B\to N\), exactness gives a factorization through
\(q\).  It vanishes on \(\ker q\), so
\(\ker q\subseteq\operatorname{im}i\); the opposite inclusion again follows
from \(qi=0\).
\end{proof}

For a short exact sequence, both Hom sequences are exact at their initial
terms.  This familiar property is often called left exactness of Hom.
Exactness through the final Hom term is strictly stronger: with \(M=P\) it
is the lifting property defining projectivity, and with \(N=E\) it is the
extension property defining injectivity.

\begin{theorem}[Projective-module characterizations]\label{thm:projective-characterizations}
For an \(R\)-module \(P\), the following are equivalent:
\begin{enumerate}
\item \(P\) is projective;
\item every short exact sequence ending in \(P\) splits;
\item \(P\) is isomorphic to a direct summand of a free module;
\item for every short exact sequence, the first sequence in
Proposition~\ref{prop:hom-exactness}, with \(M=P\), is exact at its final
term.
\end{enumerate}
In particular, every free module is projective.
\end{theorem}
\begin{proof}
If \(P\) is projective and
\[
0\longrightarrow A\longrightarrow B\xrightarrow{q}P\longrightarrow0
\]
is exact, apply the lifting property to \(1_P:P\to P\).  The resulting
section of \(q\) splits the sequence by the Splitting Lemma.

Suppose every such sequence splits.  Choose a generating set \(X\) of \(P\).
The universal property of \(R^{(X)}\) gives a surjection
\(R^{(X)}\twoheadrightarrow P\).  Its kernel produces a short exact
sequence ending in \(P\), so it splits.  Thus
\[
R^{(X)}\cong\ker(R^{(X)}\to P)\oplus P,
\]
and \(P\) is a direct summand of a free module.

Finally, let \(P\) be a direct summand of a free module \(F\), with
inclusion \(u:P\to F\) and retraction \(r:F\to P\).  Given
\(\pi:M\twoheadrightarrow N\) and \(f:P\to N\), choose a basis of \(F\).
For each basis vector \(e_x\), choose a lift in \(M\) of
\(f(r(e_x))\).  The universal property of \(F\) gives
\(g:F\to M\) with \(\pi g=fr\).  Then \(\widetilde f=gu\) satisfies
\(\pi\widetilde f=f\).  Thus \(P\) is projective.  Taking \(P=F\) proves
the final assertion.  Finally, Proposition~\ref{prop:hom-exactness} says
exactly that the final Hom map is surjective precisely when every map
\(P\to C\) lifts through \(B\).  This proves the equivalence of (1) and
(4), and distinguishes full exactness here from the left exactness enjoyed
by every test module.
\end{proof}

\noindent\textbf{Example.} For \(n>1\), the \(\mathbb Z\)-module
\(\mathbb Z/n\mathbb Z\) is not projective.  If
\[
0\longrightarrow\mathbb Z\xrightarrow{\times n}\mathbb Z
\longrightarrow\mathbb Z/n\mathbb Z\longrightarrow0
\]
split, then \(\mathbb Z\) would contain a subgroup isomorphic to
\(\mathbb Z/n\mathbb Z\).  This is impossible because \(\mathbb Z\) has no
nonzero torsion.

\begin{theorem}[Injective-module characterization]\label{thm:injective-characterization}
An \(R\)-module \(E\) is injective if and only if every short exact sequence
\[
0\longrightarrow E\longrightarrow B\longrightarrow C\longrightarrow0
\]
splits.
\; Equivalently, for every short exact sequence, the second sequence in
Proposition~\ref{prop:hom-exactness}, with \(N=E\), is exact at its final
term.
\end{theorem}
\begin{proof}
If \(E\) is injective, apply its extension property to the identity map of
\(E\) and the first injection in the sequence.  The resulting retraction
splits the sequence.  Conversely, let \(j:A\hookrightarrow B\) and
\(f:A\to E\) be given.  Form
\[
S=\{(f(a),-j(a)):a\in A\}\subseteq E\oplus B,
\qquad
Q=(E\oplus B)/S.
\]
The set \(S\) is a submodule: it is the image of the linear map
\(A\to E\oplus B\), \(a\mapsto(f(a),-j(a))\).  The map
\(E\to Q\), \(e\mapsto(e,0)+S\), is injective: if
\[
(e,0)=(f(a),-j(a))\in S,
\]
then \(j(a)=0\), so \(a=0\) and \(e=0\).  The map
\[
Q\longrightarrow B/j(A),
\qquad
(e,b)+S\longmapsto b+j(A)
\]
is well-defined, since adding \((f(a),-j(a))\) changes \(b\) by an element
of \(j(A)\).  Its kernel is the image of \(E\): if \(b=j(a)\), then
\[
(e,b)+S=(e+f(a),0)+S,
\]
and the reverse inclusion is immediate.  It is onto, so the First
Isomorphism Theorem identifies the quotient of \(Q\) by the image of \(E\)
with \(B/j(A)\).  By hypothesis the resulting short exact sequence splits,
giving a retraction \(Q\to E\).  Composing
\(B\to Q\), \(b\mapsto(0,b)+S\), with that retraction gives an extension of
\(f\).  Hence \(E\) is injective.  The final Hom characterization follows
from Proposition~\ref{prop:hom-exactness}: surjectivity of its final map is
exactly the assertion that maps \(A\to E\) extend to \(B\).
\end{proof}

\begin{proposition}[Enough projectives]\label{prop:enough-projectives}
Every \(R\)-module is a quotient of a free, hence projective, \(R\)-module.
\end{proposition}
\begin{proof}
Let \(X\) be the underlying set of \(M\).  The free module
\[
R^{(X)}
\]
has basis \(\{e_x:x\in X\}\).  The rule \(e_x\mapsto x\) extends uniquely to
a surjective map \(R^{(X)}\to M\).  The final assertion follows from
Theorem~\ref{thm:projective-characterizations}.
\end{proof}

\begin{theorem}[Baer's criterion]\label{thm:baer-criterion}
An \(R\)-module \(E\) is injective if and only if every \(R\)-linear map
\[
f:J\longrightarrow E
\]
from a left ideal \(J\) of \(R\) extends to a map \(R\to E\).
\end{theorem}
\begin{proof}
The forward implication applies injectivity to \(J\hookrightarrow R\).
Conversely, given \(A\subseteq B\) and \(f:A\to E\), order extensions
\[
(C,g),\qquad A\subseteq C\subseteq B,\qquad g|_A=f,
\]
by extension.  Unions of chains give extensions, so Zorn's Lemma gives a
maximal pair \((C,g)\).  If \(b\notin C\), then
\[
J=\{r\in R:rb\in C\}
\]
is a left ideal.  The map \(r\mapsto g(rb)\) extends to \(r\mapsto re\) on
\(R\).  Hence
\[
g'(c+rb)=g(c)+re
\]
is well-defined: if \(c+rb=c'+r'b\), then
\[
(r-r')b=c'-c\in C,
\]
so \(r-r'\in J\).  The chosen extension agrees with
\(s\mapsto g(sb)\) on \(J\), hence
\[
(r-r')e=g((r-r')b)=g(c'-c).
\]
Rearranging gives \(g(c)+re=g(c')+r'e\).  Thus \(g'\) extends \(g\) to
\(C+Rb\), contradicting maximality.  Therefore \(C=B\).
\end{proof}

\begin{remark}
The standard converse proof of Baer's criterion uses Zorn's Lemma; its
forward implication does not.  Over Noetherian rings, left ideals are
finitely generated, so many applications reduce ideal-level arguments to
finite-generation arguments.  The precise foundational strength of the full
theorem is not pursued here.  Looking ahead, every module embeds in an
injective module; standard proofs use additional constructions, often
injective abelian groups such as \(\mathbb Q/\mathbb Z\) and suitable Hom
modules.
\end{remark}

\begin{proposition}[Injective abelian groups]\label{prop:injective-divisible}
An abelian group is injective as a \(\mathbb Z\)-module if and only if it is
\emph{divisible}: multiplication by every nonzero integer is surjective.
\end{proposition}
\begin{proof}
If \(E\) is injective, extend
\[
n\mathbb Z\longrightarrow E,\qquad nk\longmapsto ka
\]
along \(n\mathbb Z\hookrightarrow\mathbb Z\); the image of \(1\) is an
\(n\)-th root of \(a\).  Conversely, every ideal of \(\mathbb Z\) is
\(n\mathbb Z\), and divisibility extends any map from it by sending \(1\) to
an \(n\)-th root of the image of \(n\).  Apply Baer's criterion.
\end{proof}

\begin{example}
\(\mathbb Q\), \(\mathbb Q/\mathbb Z\), and
\[
\mathbb Q^r
\]
are divisible, hence injective \(\mathbb Z\)-modules; \(\mathbb Z\) is not.
\end{example}

\begin{proposition}[Enough injectives over \(\mathbb Z\)]
Every abelian group embeds in a divisible abelian group.  Equivalently, every
\(\mathbb Z\)-module embeds in an injective \(\mathbb Z\)-module.
\end{proposition}
\begin{proof}
The group \(\mathbb Q/\mathbb Z\) is divisible, hence injective by
Proposition~\ref{prop:injective-divisible}.  For every nonzero \(a\in A\),
there is a homomorphism
\[
\chi_a:A\longrightarrow\mathbb Q/\mathbb Z
\]
with \(\chi_a(a)\ne0\).  Indeed, first define such a map on the cyclic
subgroup generated by \(a\): send a generator of a finite cyclic group of
order \(n\) to \(1/n+\mathbb Z\), and send an infinite cyclic generator to
\(1/2+\mathbb Z\).  Injectivity of \(\mathbb Q/\mathbb Z\) extends this map
to \(A\).

Now form the evaluation map
\[
A\longrightarrow
\prod_{\chi\in\operatorname{Hom}_{\mathbb Z}(A,\mathbb Q/\mathbb Z)}
\mathbb Q/\mathbb Z,
\qquad
a\longmapsto\bigl(\chi(a)\bigr)_\chi.
\]
It is injective because the coordinate \(\chi_a\) detects every nonzero
\(a\).  The product is divisible coordinatewise: an \(n\)-th root is chosen
in each coordinate.  Proposition~\ref{prop:injective-divisible} then makes
it injective as a \(\mathbb Z\)-module.
\end{proof}

\begin{remark}
The implication injective \(\Rightarrow\) divisible is elementary.  The
converse, and hence the equivalence used above, is established here through
Baer's criterion and its Zorn-lemma argument.  Thus the presentation relies
on the usual choice principles.
\end{remark}

\section{Vector Spaces, Span, and Finite Bases}
We now specialize to the case in which the scalar ring is a field \(F\).
An \(F\)-vector space is simply an \(F\)-module, so the notions of span,
linear independence, and basis from Sections~6.1 and~6.4 apply without
change. The additional strength of vector-space theory comes from the fact
that every nonzero scalar is invertible.

\begin{definition}
An \emph{\(F\)-vector space} is an \(F\)-module, and a \emph{subspace} is an
\(F\)-submodule.
\end{definition}

\begin{theorem}[Basis existence]
Every vector space has a basis.
\end{theorem}
\begin{proof}
Let \(\mathcal P\) be the set of linearly independent subsets of \(V\),
ordered by inclusion. It is nonempty because it contains the empty set. If
\(\mathcal C\) is a chain in \(\mathcal P\), then
\[
U=\bigcup_{S\in\mathcal C}S
\]
is linearly independent: a finite relation among distinct members of \(U\)
involves only finitely many members of the chain, hence all of them belong to
one member of \(\mathcal C\), where their coefficients vanish. Thus every
chain has an upper bound in \(\mathcal P\). Zorn's lemma gives a maximal
linearly independent set \(B\).

If \(v\notin\operatorname{span}(B)\), then \(B\cup\{v\}\) is linearly
independent: in a relation involving \(v\), a nonzero coefficient of \(v\)
would solve for \(v\) as a linear combination of members of \(B\), while a
zero coefficient leaves a relation in \(B\). This contradicts maximality.
Hence \(B\) spans \(V\), and therefore is a basis.
\end{proof}

\begin{remark}
It applies also to infinite-dimensional spaces. Its proof uses Zorn's lemma
and hence the usual Choice principles. The finite
replacement argument below is elementary and uses no such principle.
\end{remark}

\begin{theorem}[Replacement theorem]
Let
\[
G=\{w_1,\ldots,w_n\}
\]
span \(V\), and let
\[
L=\{v_1,\ldots,v_m\}
\]
be linearly independent. Then \(m\leq n\), and exactly \(n-m\) vectors from
\(G\) can be chosen so that, together with \(v_1,\ldots,v_m\), they span
\(V\). Equivalently, one may replace \(m\) vectors of the original spanning
list by \(v_1,\ldots,v_m\) and retain a spanning list of size \(n\).
\end{theorem}
\begin{proof}
We argue by induction on \(m\). For \(m=0\), the original list spans. Assume
the result for \(m-1\). The induction hypothesis gives a spanning list
\[
\{v_1,\ldots,v_{m-1},w_{i_1},\ldots,w_{i_{n-m+1}}\}.
\]
Write
\[
v_m=a_1v_1+\cdots+a_{m-1}v_{m-1}
+b_1w_{i_1}+\cdots+b_{n-m+1}w_{i_{n-m+1}}.
\]
Not every \(b_j\) is zero, for otherwise \(v_m\) belongs to the span of
\(v_1,\ldots,v_{m-1}\). Choose \(b_k\ne0\). Since \(F\) is a field,
\[
w_{i_k}
=b_k^{-1}\left(
v_m-\sum_{\ell=1}^{m-1}a_\ell v_\ell
-\sum_{j\ne k}b_jw_{i_j}
\right).
\]
Replacing \(w_{i_k}\) by \(v_m\) therefore preserves spanning. There is an
unreplaced \(w_{i_k}\) at every stage, so necessarily \(m\leq n\).
\end{proof}

\begin{corollary}[Extension and extraction]
If \(V\) has a finite spanning set, every finite linearly independent subset
extends to a basis, and every finite spanning set contains a basis.
\end{corollary}
\begin{proof}
Apply the replacement theorem to an independent list and a fixed finite
spanning list. Its retained members, together with the independent list, form
a finite spanning list; discard a retained member whenever it lies in the
span of the others. Conversely, applying the same deletion process to the
original spanning list leaves an independent spanning list.
\end{proof}

\begin{corollary}[Uniqueness of finite basis size]
Any two finite bases of a vector space have the same cardinality.
\end{corollary}
\begin{proof}
Apply the replacement theorem to either basis as an independent set and the
other as a spanning set. The two resulting inequalities are opposite.
\end{proof}

\begin{definition}
A vector space is \emph{finite-dimensional} if it has a finite basis. Its
\emph{dimension}, denoted by \(\dim_FV\), is the common cardinality of its
finite bases.
\end{definition}

Thus a basis expansion exists because a basis spans and is unique because a
basis is independent. An \emph{ordered basis} is a basis displayed as an
ordered finite list; this extra order becomes essential when coordinates are
introduced in Section~6.11.

\begin{proposition}[Subspaces of finite-dimensional spaces]
If \(V\) is finite-dimensional and \(W\leq V\), then \(W\) has a finite
basis and \(\dim_F W\leq\dim_FV\).
\end{proposition}
\begin{proof}
Begin with the empty independent list in \(W\), and adjoin an element of
\(W\) not in its span whenever one exists. The replacement theorem, applied
to the resulting independent list and a fixed basis of \(V\), shows that at
most \(\dim_FV\) choices can be made. When the process stops its list spans
\(W\), so it is a basis. The same theorem gives the inequality.
\end{proof}

\begin{corollary}[IBN for fields]
If
\[
F^m\cong F^n
\]
as \(F\)-vector spaces, then \(m=n\).
\end{corollary}
\begin{proof}
An isomorphism carries the standard basis of \(F^m\) to a basis of \(F^n\).
Uniqueness of finite basis size gives \(m=n\).
\end{proof}

\section{Invariant Basis Number and Finite Free Rank}
The preceding corollary proves IBN for fields. A finite free module over a
general ring needs a separate argument. This distinction is essential in
Chapter~7, where the free part of a module over a PID must have a uniquely
determined rank.

\begin{definition}
A ring \(R\) has \emph{invariant basis number} (IBN) if
\(R^m\cong R^n\) for finite \(m,n\) implies \(m=n\). When \(R\) has IBN,
the \emph{rank} of a finite free \(R\)-module is the integer \(n\) for which
it is isomorphic to \(R^n\).
\end{definition}

\begin{theorem}[Commutative rings have IBN]
Every nonzero commutative ring with identity has invariant basis number.
\end{theorem}
\begin{proof}
Suppose \(R^m\cong R^n\). Since \(R\) is nonzero, \((0)\) is proper; by the
\hyperref[thm:maximal-ideal-existence]{maximal-ideal theorem}, choose a
maximal ideal \(\mathfrak m\) of \(R\). An isomorphism \(f:R^m\to R^n\) carries
\(\mathfrak mR^m\) onto \(\mathfrak mR^n\), since
\(f(rx)=rf(x)\), and the same assertion for \(f^{-1}\) proves equality.
It consequently induces an isomorphism
\[
R^m/\mathfrak mR^m\ \cong\ R^n/\mathfrak mR^n.
\]
Coordinatewise reduction identifies these quotient modules with
\((R/\mathfrak m)^m\) and \((R/\mathfrak m)^n\). The quotient is a field
because \(\mathfrak m\) is maximal, so IBN for fields gives \(m=n\).
\end{proof}

The nonzero hypothesis is essential: for the zero ring,
\[
R^m=R^n=0
\]
for all finite \(m,n\), so invariant basis number fails.  Arbitrary
noncommutative rings need not have IBN, so no unqualified rank is attached to
their finite free modules.

\section{Linear Transformations and Dimension}
The module homomorphisms of vector spaces are the familiar linear
transformations. Finite dimension turns the abstract kernel--image picture
into a numerical conservation law.

\begin{definition}
A \emph{linear transformation} \(T:V\to W\) is an \(F\)-module
homomorphism: \(T(av+bw)=aT(v)+bT(w)\). Its kernel and image are the
submodules already defined in \S6.1, now called respectively its
\emph{kernel} and \emph{range}. It is injective, surjective, bijective, or
an isomorphism in the same sense as a module map. An endomorphism is a map
\(V\to V\), and an automorphism is an invertible endomorphism. The
endomorphisms form the ring \(\mathcal{L}_F(V)\); when the field is clear,
we write \(\mathcal{L}(V)\). Once
\(V\) is finite-dimensional, its \emph{nullity} and \emph{rank} are
\(\dim\ker T\) and \(\dim\operatorname{im}T\).
\end{definition}

The kernel and image are subspaces by Proposition~6.1, and the preceding
subspace theorem supplies their finite bases. Thus the numerical words in the
definition are legitimate before rank--nullity is invoked.

\begin{theorem}[Rank--nullity]
If \(T:V\to W\) is linear and \(V\) is finite-dimensional, then
\[\dim V=\dim\ker T+\dim\operatorname{im}T.\]
\end{theorem}
\begin{proof}
Choose a basis \(k_1,\ldots,k_a\) of \(\ker T\). Extend it to the basis
\[
k_1,\ldots,k_a,v_1,\ldots,v_b
\]
of \(V\). The vectors \(Tv_1,\ldots,Tv_b\)
span the image: applying \(T\) removes the \(k_i\)-part of any coordinate
expression. If
\[
\sum_i c_iTv_i=0,
\]
then
\[
\sum_i c_iv_i\in\ker T,
\]
and
uniqueness of the displayed basis expansion gives every \(c_i=0\). Hence
they form a basis of the image, so \(\dim V=a+b\).
\end{proof}

Thus a linear map between equal finite dimensions is injective exactly when
it is surjective, and hence exactly when it is an isomorphism: rank--nullity
shows that one of kernel and image has the necessary dimension precisely when
the other does.

\section{Matrices and Change of Basis}\label{sec:matrices-change-basis}
Once ordered bases are chosen, a linear transformation becomes a table of
coordinates. The matrix is a useful representative, but it is not the
transformation itself: changing coordinates changes its matrix.

Let \(\beta=(v_1,\ldots,v_n)\) be an ordered basis. Every \(v\in V\) has
one and only one expansion \(v=x_1v_1+\cdots+x_nv_n\); define its
\emph{coordinate column} by
\[
[v]_\beta=\begin{pmatrix}x_1\\x_2\\\vdots\\x_n\end{pmatrix}.
\]
The order matters: exchanging two basis vectors generally exchanges two
coordinates.

\begin{definition}
Let \(R\) be a ring. An \emph{\(m\times n\) matrix over \(R\)} is a
rectangular array
\[
A=(a_{ij})=
\begin{pmatrix}
a_{11}&a_{12}&\cdots&a_{1n}\\
a_{21}&a_{22}&\cdots&a_{2n}\\
\vdots&\vdots&\ddots&\vdots\\
a_{m1}&a_{m2}&\cdots&a_{mn}
\end{pmatrix},
\qquad
a_{ij}\in R.
\]
The set of all such matrices is \(M_{m\times n}(R)\), and
\[
M_n(R)=M_{n\times n}(R).
\]
Matrices are equal precisely when their corresponding entries are equal. For
matrices of the same size, define
\[
A+B=(a_{ij}+b_{ij}).
\]
If \(R\) is commutative, scalar multiplication is defined entrywise by
\[
rA=(ra_{ij}),
\]
making \(M_{m\times n}(R)\) an \(R\)-module.
\end{definition}

Matrices are arrays of scalars; they encode linear maps only after ordered
bases have been chosen. If
\[
A=(a_{ij})\in M_{m\times n}(R),
\]
its \emph{transpose} is the matrix
\[
A^t=
\begin{pmatrix}
a_{11}&a_{21}&\cdots&a_{m1}\\
a_{12}&a_{22}&\cdots&a_{m2}\\
\vdots&\vdots&\ddots&\vdots\\
a_{1n}&a_{2n}&\cdots&a_{mn}
\end{pmatrix}
\in M_{n\times m}(R),
\qquad
(A^t)_{ij}=a_{ji}.
\]
Thus transpose visibly interchanges rows and columns.

For \(A=(a_{ij})\in M_{m\times n}(F)\) and a column
\[
x=\begin{pmatrix}x_1\\\vdots\\x_n\end{pmatrix},
\]
define matrix--vector multiplication by
\[
(Ax)_i=\sum_{j=1}^n a_{ij}x_j.
\]
For instance,
\[
\begin{pmatrix}2&-1\\3&4\end{pmatrix}
\begin{pmatrix}5\\2\end{pmatrix}
=\begin{pmatrix}8\\23\end{pmatrix}.
\]

For
\[
A=(a_{ij})\in M_{m\times n}(R),
\qquad
B=(b_{jk})\in M_{n\times r}(R),
\]
define
\[
AB=(c_{ik})\in M_{m\times r}(R),
\qquad
c_{ik}=\sum_{j=1}^n a_{ij}b_{jk}.
\]
This is row-by-column multiplication, and its dimension rule is
\[
(m\times n)(n\times r)\longrightarrow(m\times r).
\]
Associativity follows by distributing the finite sums; no commutativity of
\(R\) is required. The \emph{identity matrix}
\[
I_n=(\delta_{ij})\in M_n(R)
\]
has entries \(1\) on the diagonal and \(0\) elsewhere. A square matrix \(A\)
is \emph{invertible} if some square matrix \(B\) satisfies \(AB=BA=I_n\);
then \(B\) is unique and is denoted by \(A^{-1}\).

\begin{proposition}[Transpose identities]
For compatible matrices over a commutative ring,
\[
(A^t)^t=A,
\qquad
(A+B)^t=A^t+B^t,
\qquad
(rA)^t=rA^t,
\qquad
(AB)^t=B^tA^t.
\]
\end{proposition}
\begin{proof}
The first three identities follow by comparing entries. For the product,
\[
\bigl((AB)^t\bigr)_{ij}
=(AB)_{ji}
=\sum_k a_{jk}b_{ki}
=\sum_k b_{ki}a_{jk}
=(B^tA^t)_{ij}.
\]
Commutativity is used in the third equality.
\end{proof}

Now let \(T:V\to W\), let
\[
\beta=(v_1,\ldots,v_n)
\]
be an ordered basis of \(V\), and let
\[
\gamma=(w_1,\ldots,w_m)
\]
be an ordered basis of \(W\). Write
\[
\begin{aligned}
T(v_1)&=a_{11}w_1+a_{21}w_2+\cdots+a_{m1}w_m,\\
T(v_2)&=a_{12}w_1+a_{22}w_2+\cdots+a_{m2}w_m,\\
&\ \vdots\\
T(v_n)&=a_{1n}w_1+a_{2n}w_2+\cdots+a_{mn}w_m.
\end{aligned}
\]
The \emph{matrix of \(T\) from \(\beta\) to \(\gamma\)} is
\[
[T]^\gamma_\beta=
\begin{pmatrix}
a_{11}&a_{12}&\cdots&a_{1n}\\
a_{21}&a_{22}&\cdots&a_{2n}\\
\vdots&\vdots&\ddots&\vdots\\
a_{m1}&a_{m2}&\cdots&a_{mn}
\end{pmatrix}.
\]
Equivalently,
\[
\boxed{\text{the \(j\)-th column of \([T]^\gamma_\beta\) is
\([T(v_j)]_\gamma\).}}
\]
Thus the matrix depends on the chosen ordered bases, though the linear map
does not. If
\[
v=x_1v_1+\cdots+x_nv_n,
\]
then linearity and the displayed expansions give
\[
T(v)
=\sum_{i=1}^m\left(\sum_{j=1}^n a_{ij}x_j\right)w_i.
\]
The inner sum is the \(i\)-th entry of
\([T]^\gamma_\beta[v]_\beta\), and hence
\[
[T(v)]_\gamma=[T]^\gamma_\beta[v]_\beta. \tag{6.1}
\]
For \(S:W\to X\), with an ordered basis \(\delta\) of \(X\), this immediately
explains the order of matrix multiplication:
\[
[S\circ T]^\delta_\beta=[S]^\delta_\gamma[T]^\gamma_\beta.
\]
Formula (6.1) also shows that a square matrix representing an endomorphism is
invertible exactly when that endomorphism is an automorphism.

\begin{example}[Constructing a representing matrix]
Let
\[
T:\mathbb R^2\longrightarrow\mathbb R^2,
\qquad
T(x,y)=(x+y,x-y),
\]
and choose the ordered bases
\[
\beta=((1,1),(1,0)),
\qquad
\gamma=((1,0),(1,1)).
\]
Then
\[
T(1,1)=(2,0)=2(1,0)+0(1,1),
\qquad
T(1,0)=(1,1)=0(1,0)+1(1,1).
\]
The coordinate columns therefore give
\[
[T]^\gamma_\beta=
\begin{pmatrix}2&0\\0&1\end{pmatrix}.
\]
For
\[
v=2(1,1)-(1,0)=(1,2),
\]
we have
\[
[v]_\beta=
\begin{pmatrix}2\\-1\end{pmatrix},
\qquad
[T(v)]_\gamma=
\begin{pmatrix}4\\-1\end{pmatrix}
=
\begin{pmatrix}2&0\\0&1\end{pmatrix}
\begin{pmatrix}2\\-1\end{pmatrix}.
\]
\end{example}

If \(\beta'\) is another ordered basis of \(V\), define the
change-of-coordinate matrix, with its direction displayed, by
\(P=[I_V]^\beta_{\beta'}\). Formula (6.1) for the identity says
\[
[v]_\beta=P[v]_{\beta'};
\]
thus \(P\) converts \(\beta'\)-coordinates into \(\beta\)-coordinates. The
reverse change-of-coordinate matrix
\(Q=[I_V]^{\beta'}_\beta\) satisfies \(PQ=QP=I\), so \(P^{-1}=Q\).
For an endomorphism \(T\), applying the composition formula to the two
identity maps gives
\[
[T]^{\beta'}_{\beta'}=P^{-1}[T]^\beta_\beta P.
\]
Matrices related in this way are \emph{similar}; they represent one operator
in different ordered bases.

Finally, the matrices \(E_{ij}\), having a single \(1\) in position
\((i,j)\), form a basis of \(M_{m\times n}(F)\). Hence this matrix space has
dimension \(mn\). For fixed ordered bases, \(T\mapsto[T]^\gamma_\beta\) is
a linear bijection
\[
\operatorname{Hom}_F(V,W)\longrightarrow M_{m\times n}(F)
\]
by the column construction above; this identification depends on the chosen
bases. In particular,
\(\dim_F\mathcal{L}_F(V)=(\dim_FV)^2\).

\begin{example}[A change of basis]
Let \(T:\mathbb R^2\to\mathbb R^2\) be given by
\(T(x,y)=(3x+y,y)\). In the standard basis its matrix is
\(A=\begin{pmatrix}3&1\\0&1\end{pmatrix}\). Put
\(v_1=(1,0)\), \(v_2=(1,-2)\), and let \(\beta=(v_1,v_2)\). The
change-of-coordinate matrix has these vectors as columns,
\[P=\begin{pmatrix}1&1\\0&-2\end{pmatrix},\qquad
P^{-1}=\begin{pmatrix}1&\tfrac12\\0&-\tfrac12\end{pmatrix}.
\]
Therefore
\[ [T]_\beta=P^{-1}AP=\begin{pmatrix}3&0\\0&1\end{pmatrix}. \]
The diagonal entries were not produced by a new operator: the chosen basis
consists of eigenvectors. This computation is the concrete meaning of
similarity and anticipates diagonalization.
\end{example}

\section{Determinants}\label{sec:determinants}
An endomorphism of an \(n\)-dimensional space either preserves oriented
\(n\)-dimensional volume up to a scalar or collapses some direction. The
determinant is that scalar. Its most useful definition is therefore
structural, not computational. Permutations and their signs, used below, were
proved in \hyperref[chap:groups-foundations]{Chapter~2}.

\begin{definition}
Let \(V\) be an \(n\)-dimensional vector space over \(F\), with ordered basis
\(\beta=(v_1,\ldots,v_n)\). A map \(\omega:V^n\to F\) is
\emph{\(n\)-multilinear} if it is linear in each argument separately, and
\emph{alternating} if it vanishes whenever two arguments agree. The
\emph{determinant form associated with \(\beta\)} is an alternating
\(n\)-multilinear map \(\omega_\beta\) satisfying
\[
\omega_\beta(v_1,\ldots,v_n)=1.
\]
\end{definition}

Multilinearity models how volume changes when one edge is changed; alternation
models the collapse caused by two equal directions; normalization fixes the
unit oriented parallelepiped determined by the chosen basis.

\begin{lemma}[Alternating sign-change lemma]\label{lem:alternating-sign-change}
Let \(\omega:V^n\to F\) be alternating and multilinear. Interchanging any
two arguments changes its value by a sign. More precisely, if every argument
other than those in two fixed positions is held fixed, then
\[
\omega(\ldots,x,\ldots,y,\ldots)
=-\omega(\ldots,y,\ldots,x,\ldots).
\]
\end{lemma}
\begin{proof}
Alternation gives
\[
0=\omega(\ldots,x+y,\ldots,x+y,\ldots).
\]
Expanding the two displayed positions by multilinearity gives four terms.
The terms with two copies of \(x\) or two copies of \(y\) vanish by
alternation. Thus
\[
0=\omega(\ldots,x,\ldots,y,\ldots)
 +\omega(\ldots,y,\ldots,x,\ldots),
\]
which is the asserted identity.
\end{proof}

Consequently, each transposition of arguments contributes a factor of
\(-1\). Chapter~2 proves that the parity of a permutation is well-defined,
so a permutation \(\sigma\in S_n\) changes the value of \(\omega\) by the
factor \(\operatorname{sgn}(\sigma)\). The lemma supplies the effect on an
alternating multilinear map; the earlier permutation theory supplies the
rule by which the signs of successive transpositions combine.

\begin{theorem}[Existence and uniqueness of the determinant form]\label{thm:det-form}
There is exactly one determinant form \(\omega_\beta\).
\end{theorem}
\begin{proof}
We first prove uniqueness. Write
\[
x_j=\sum_{i=1}^n a_{ij}v_i\qquad(1\leq j\leq n).
\]
Repeated multilinearity forces
\[
\omega(x_1,\ldots,x_n)
=\sum_{i_1,\ldots,i_n}
a_{i_1,1}\cdots a_{i_n,n}\,
\omega(v_{i_1},\ldots,v_{i_n}).
\]
Alternation kills every term with repeated indices. The surviving tuples are
exactly
\[
(i_1,\ldots,i_n)=(\sigma(1),\ldots,\sigma(n)),
\qquad \sigma\in S_n.
\]
By Lemma~\ref{lem:alternating-sign-change}, each adjacent transposition of
arguments contributes a factor of \(-1\). The Chapter~2 parity theory then
shows that the term indexed by \(\sigma\) has the factor
\(\operatorname{sgn}(\sigma)\), and normalization gives
\[
\omega_\beta(x_1,\ldots,x_n)
=\sum_{\sigma\in S_n}\operatorname{sgn}(\sigma)
a_{\sigma(1),1}\cdots a_{\sigma(n),n}. \tag{6.2}\label{eq:leibniz-form}
\]
Thus there can be at most one such map.

Conversely, define the right side of \eqref{eq:leibniz-form} to be
\(\omega_\beta(x_1,\ldots,x_n)\). Each coefficient in one column occurs
linearly, so the expression is multilinear in the corresponding argument.
If \(x_r=x_s\), pair a summand indexed by \(\sigma\) with the summand obtained
by exchanging \(\sigma(r)\) and \(\sigma(s)\). The products agree because
columns \(r\) and \(s\) agree, while the signs are opposite; every summand
therefore cancels. Finally, for \((x_1,\ldots,x_n)=(v_1,\ldots,v_n)\), only
the identity permutation contributes, and the value is \(1\). Hence the
candidate is a determinant form.
\end{proof}

Formula \eqref{eq:leibniz-form}, the \emph{Leibniz formula}, is consequently
not a separate definition: it is the forced coordinate expression for the
unique normalized alternating multilinear form.

\begin{definition}
For \(A\in M_n(F)\), with columns \(A_1,\ldots,A_n\), define
\[
\det(A)=\omega_e(A_1,\ldots,A_n),
\]
where \(e\) is the standard ordered basis of \(F^n\). Thus
\[
\det(A)=
\sum_{\sigma\in S_n}\operatorname{sgn}(\sigma)
a_{\sigma(1),1}\cdots a_{\sigma(n),n}.
\]
This is a column convention. The row convention follows below from transpose
invariance.
\end{definition}

\begin{definition}
For \(T\in\mathcal{L}_F(V)\), define \(\det(T)\) by
\[
\omega_\beta(Tx_1,\ldots,Tx_n)
=\det(T)\,\omega_\beta(x_1,\ldots,x_n)
\qquad(x_1,\ldots,x_n\in V).
\]
\end{definition}

The left side is alternating and multilinear in the \(x_i\). The expansion
used in the proof of Theorem~\ref{thm:det-form} shows that every such form is
its value on \(\beta\) times \(\omega_\beta\), so the displayed scalar exists
and is unique; evaluating at \(\beta\) shows that it equals
\(\omega_\beta(Tv_1,\ldots,Tv_n)\). The identity itself holds for every
tuple, so evaluating in another ordered basis proves that \(\det(T)\) is
independent of \(\beta\). Expanding the columns \(Tv_j\) in \(\beta\) gives
\[
\det(T)=\det([T]^\beta_\beta). \tag{6.3}\label{eq:operator-matrix-det}
\]
This separates the intrinsic operator from its coordinate matrix.

\begin{theorem}[Basic determinant identities]\label{thm:det-identities}
For endomorphisms \(S,T\) of a finite-dimensional vector space over a field,
\[
\det(I)=1,\qquad
\det(S\circ T)=\det(S)\det(T).
\]
For square matrices over \(F\),
\[
\det(AB)=\det(A)\det(B),
\qquad
\det(A^t)=\det(A).
\]
\end{theorem}
\begin{proof}
The defining identity for determinants gives
\[
\omega_\beta(STx_1,\ldots,STx_n)
=\det(S)\det(T)\omega_\beta(x_1,\ldots,x_n),
\]
which proves multiplicativity and \(\det(I)=1\). Matrix multiplicativity
follows from \eqref{eq:operator-matrix-det} and the composition formula
proved in \S6.11. Replacing \(\sigma\) by \(\sigma^{-1}\) in
\eqref{eq:leibniz-form} gives \(\det(A^t)=\det(A)\).
\end{proof}

Thus adding a multiple of one column to another leaves the determinant
unchanged, scaling a column scales the determinant by the same scalar, and
interchanging two columns changes its sign; these are immediate consequences
of multilinearity and alternation. Transpose invariance supplies the
corresponding row operations.

\begin{definition}
For \(A=(a_{ij})\in M_n(F)\), let \(A_{ij}\) be the matrix obtained by
deleting row \(i\) and column \(j\). Its \emph{cofactor} is
\[
C_{ij}=(-1)^{i+j}\det(A_{ij}).
\]
\end{definition}

\begin{theorem}[Laplace expansion and adjugate]\label{thm:adjugate}
For every row \(i\) and column \(j\),
\[
\det A=\sum_{k=1}^n a_{ik}C_{ik},
\qquad
\det A=\sum_{k=1}^n a_{kj}C_{kj}.
\]
If \(\operatorname{adj}(A)=(C_{ji})\), then
\[
A\operatorname{adj}(A)=\operatorname{adj}(A)A=(\det A)I_n. \tag{6.4}\label{eq:adjugate}
\]
\end{theorem}
\begin{proof}
Group the terms of the Leibniz formula according to the column selected in
row \(i\). Moving that selected entry to the first position requires
\(i+j\) parity changes, giving the first expansion; transpose invariance
gives the column version. The \((i,j)\)-entry of
\(A\operatorname{adj}(A)\) is the row-\(i\) cofactor expansion if \(i=j\).
If \(i\ne j\), it is the same expansion after replacing row \(j\) by row
\(i\), which is zero by alternation. The column calculation proves the
identity on the other side.
\end{proof}

\begin{corollary}[Invertibility criterion over a field]\label{cor:field-det-invertibility}
For \(A\in M_n(F)\), equivalently for an endomorphism of \(V\),
\[
A\text{ is invertible}\quad\Longleftrightarrow\quad\det(A)\ne0.
\]
\end{corollary}
\begin{proof}
If \(A\) is invertible, determinant multiplicativity gives
\[
1=\det(A)\det(A^{-1}).
\]
Conversely, if \(\det(A)\ne0\), then
\(\det(A)^{-1}\operatorname{adj}(A)\) is an inverse by
\eqref{eq:adjugate}.
\end{proof}

\begin{example}
For a \(2\times2\) matrix,
\[
\det\begin{pmatrix}a&b\\c&d\end{pmatrix}=ad-bc.
\]
Thus \(\det\begin{pmatrix}2&1\\4&2\end{pmatrix}=0\), while
\[
\det\begin{pmatrix}2&1\\1&1\end{pmatrix}=1,
\qquad
\begin{pmatrix}2&1\\1&1\end{pmatrix}^{-1}
=\begin{pmatrix}1&-1\\-1&2\end{pmatrix}.
\]
Laplace expansion and the adjugate are computational consequences of the
same structural determinant form, not competing definitions.
\end{example}

\subsection{Determinants over commutative rings}
The Leibniz expression uses only addition, multiplication, commutativity,
and the elements \(1\) and \(-1\). It therefore has a useful scope beyond
fields.

\begin{theorem}[Determinants over commutative rings]\label{thm:ring-determinants}
Let \(R\) be a commutative ring with identity. For \(A\in M_n(R)\), define
\[
\det(A)=
\sum_{\sigma\in S_n}\operatorname{sgn}(\sigma)
a_{\sigma(1),1}\cdots a_{\sigma(n),n}.
\]
On \(R^n\), this is the unique normalized alternating \(n\)-multilinear
form with respect to the standard ordered basis. The identities
\[
\det(I)=1,\qquad
\det(AB)=\det(A)\det(B),\qquad
\det(A^t)=\det(A)
\]
and the Laplace and adjugate identities \eqref{eq:adjugate} remain valid.
\end{theorem}
\begin{proof}
The uniqueness, cancellation, and multilinearity arguments in
Theorem~\ref{thm:det-form} use no division, only commutativity and the sign
of a permutation. The structural proof of multiplicativity and the cofactor
proof use the same operations, and hence carry over word for word.
\end{proof}

If \(P\) is an invertible change-of-basis matrix of a finite free
\(R\)-module \(F\), then
\[
\det(P)\det(P^{-1})=1.
\]
Consequently, for an endomorphism \(T:F\to F\), the value
\[
\det(T)=\det([T]^\beta_\beta)
\]
is independent of the ordered basis \(\beta\): two matrices are related by
\(P^{-1}[T]^\beta_\beta P\), and determinant multiplicativity cancels the
two unit determinants. Thus the preceding field construction extends exactly
to finite free modules, not to arbitrary finitely generated modules.

\begin{corollary}[Invertibility over a commutative ring]
For \(A\in M_n(R)\),
\[
A\text{ is invertible}\quad\Longleftrightarrow\quad\det(A)\in R^\times.
\]
\end{corollary}
\begin{proof}
An inverse of \(A\) makes \(\det(A)\det(A^{-1})=1\). Conversely, if
\(\det(A)\) is a unit, \((\det A)^{-1}\operatorname{adj}(A)\) is an inverse
by \eqref{eq:adjugate}.
\end{proof}

Over a field, nonzero and unit are equivalent. Over \(\mathbb Z\), they are
not: the \(1\times1\) matrix \((2)\) has nonzero determinant but is not
invertible. In particular, the field criterion involving linear dependence
must not be transferred unqualified to general commutative rings. For
\(A\in M_n(R)\), the polynomial
\(\chi_A(t)=\det(tI-A)\in R[t]\) is nevertheless defined, and the adjugate
identity over \(R[t]\) gives \(\chi_A(A)=0\). Chapter~7 uses this argument in
its operator-theoretic field setting.

\section{Eigenvalues, Eigenvectors, and Invariant Subspaces}
An eigenvector is a direction whose direction is not changed by an operator:
the operator merely scales it. Such directions are the first indication that
a complicated transformation might have a simple coordinate system.

\begin{definition}
For \(T:V\to V\), a nonzero \(v\) is an \emph{eigenvector} with
\emph{eigenvalue} \(\lambda\in F\) if \(Tv=\lambda v\). The
\emph{eigenspace} is \(E_\lambda=\ker(T-\lambda I)\). A subspace
\(W\leq V\) is \emph{\(T\)-invariant} if \(T(W)\subseteq W\).
\end{definition}

Because eigenspaces are kernels, they are subspaces before any dimension is
attached to them. Now let \(\beta\) be an ordered basis and
\(A=[T]^\beta_\beta\). Define the \emph{characteristic polynomial} by
\[\chi_T(x)=\det(xI-A)\in F[x].\]
This definition is independent of \(\beta\). Indeed, if
\(A'=P^{-1}AP\), then, in the polynomial ring \(F[x]\),
\[
xI-A'=P^{-1}(xI-A)P.
\]
Determinant multiplicativity, already proved above, gives
\(\det(xI-A')=\det(P^{-1})\det(xI-A)\det(P)=\det(xI-A)\).

\begin{proposition}
For \(\lambda\in F\), \(\lambda\) is an eigenvalue of \(T\) if and only if
\(\chi_T(\lambda)=0\).
\end{proposition}
\begin{proof}
The equation \(Tv=\lambda v\) for a nonzero \(v\) says that
\(T-\lambda I\) is not injective. In finite dimension it is therefore not
invertible. Its matrix is \(A-\lambda I\), so the determinant criterion
above says precisely that \(\det(A-\lambda I)=0\), equivalently
\(\det(\lambda I-A)=0\). This is \(\chi_T(\lambda)=0\). The reverse chain
of equivalences proves the converse.
\end{proof}

The eigenspace is a subspace because it is a kernel. It can have dimension
larger than one: for the identity operator, the unique eigenspace is all of
\(V\). By contrast, a rotation of the real plane through a nonzero angle less
than \(\pi\) has no real eigenvector at all. The field therefore matters:
the same rotation has complex eigenvalues, a point returned to in Chapter~7.

\section{Diagonalization}
The best possible matrix of an operator is diagonal: each coordinate evolves
independently. The question is therefore not whether a matrix happens to
look diagonal, but whether the underlying operator has a basis made of
eigenvectors.

\begin{theorem}[Diagonalization criterion]
For an endomorphism \(T:V\to V\), the following are equivalent:
\begin{enumerate}
\item \(T\) is represented by a diagonal matrix in some basis;
\item \(V\) has a basis of eigenvectors of \(T\);
\item
\[
V=\bigoplus_\lambda E_\lambda,
\]
where the sum ranges over the
eigenvalues that occur.
\end{enumerate}
\end{theorem}
\begin{proof}
The columns of a diagonal matrix say exactly that its basis vectors are
eigenvectors, proving (1) iff (2). Eigenvectors belonging to distinct
eigenvalues are independent. For if a shortest nontrivial relation
\[
\sum_{i=1}^rc_iv_i=0
\]
involved distinct eigenvalues \(\lambda_i\), then
applying \(T-\lambda_rI\) gives a shorter relation
\[
\sum_{i<r}c_i(\lambda_i-\lambda_r)v_i=0;
\]
its coefficients are nonzero
whenever the original ones are, a contradiction. Thus the sum of the
eigenspaces is direct. A basis of eigenvectors spans precisely this direct
sum, proving (2) iff (3).
\end{proof}

In particular, if \(\dim V=n\) and \(T\) has \(n\) distinct eigenvalues,
then it is diagonalizable. Diagonalization is powerful but not universal. The
next chapter explains the additional structure that remains when no eigenbasis
exists.

\begin{example}[Three first tests for diagonalization]
On \(\mathbb R^2\), the operator with matrix
\(\begin{pmatrix}3&1\\0&1\end{pmatrix}\) has distinct eigenvalues \(3\) and
\(1\), so the preceding worked calculation diagonalizes it. The identity
operator has the repeated eigenvalue \(1\), yet is diagonalizable because
every nonzero vector is an eigenvector. In contrast,
\[J=\begin{pmatrix}1&1\\0&1\end{pmatrix}\]
has characteristic polynomial \((x-1)^2\), but
\(\ker(J-I)\) is one-dimensional. It has only one independent eigenvector
and therefore is not diagonalizable. Thus a split characteristic polynomial
does not by itself supply an eigenbasis; the missing information is precisely
what the module structure in Chapter~7 records.
\end{example}

\section{Dual Spaces and Dual Maps}\label{sec:dual-spaces}
Coordinates record a vector after a basis has been chosen.  There is another
useful way to probe a vector: apply every scalar-valued linear measurement to
it.  This viewpoint is intrinsic, and it is the starting point for
bilinear forms and alternating forms.

\begin{definition}
A \emph{linear functional} on an \(F\)-vector space \(V\) is a linear map
\(\varphi:V\to F\).  The \emph{algebraic dual space} of \(V\) is
\[
V^\vee=\operatorname{Hom}_F(V,F).
\]
\end{definition}

The notation \(V^\vee\) is deliberately different from the notation
\(T^*\) used for the inner-product adjoint in Chapter~7. The construction
below is algebraic and requires no inner product.

\begin{proposition}[Dual basis]\label{prop:dual-basis}
Let \(\beta=(v_1,\ldots,v_n)\) be an ordered basis of a finite-dimensional
vector space \(V\).  There are unique linear functionals
\[
v_i^\vee:V\longrightarrow F
\]
such that
\[
v_i^\vee(v_j)=\delta_{ij}.
\]
Moreover,
\[
\beta^\vee=(v_1^\vee,\ldots,v_n^\vee)
\]
is a basis of \(V^\vee\).  In particular,
\[
\dim V^\vee=\dim V.
\]
\end{proposition}
\begin{proof}
For
\[
v=\sum_{j=1}^na_jv_j,
\]
define \(v_i^\vee(v)=a_i\).  The uniqueness of basis coordinates proves
that this is well-defined and linear, and gives the displayed values.
If
\[
\sum_{i=1}^nc_iv_i^\vee=0,
\]
evaluation at \(v_j\) gives \(c_j=0\), so the functionals are independent.
For \(\varphi\in V^\vee\), linearity gives
\[
\varphi=\sum_{i=1}^n\varphi(v_i)v_i^\vee.
\]
Thus they span and form a basis.
\end{proof}

Consequently \(V\) and \(V^\vee\) are isomorphic when \(V\) is
finite-dimensional, but this proof used the chosen basis.  In general there
is no preferred isomorphism
\[
V\cong V^\vee.
\]
An inner product can supply one, but that is additional structure.

\begin{definition}
For a linear map
\[
T:V\longrightarrow W,
\]
its \emph{dual map}, also called its algebraic transpose, is
\[
T^\vee:W^\vee\longrightarrow V^\vee,
\qquad
T^\vee(\varphi)=\varphi\circ T.
\]
\end{definition}

The direction reverses because a measurement on \(W\) can be pulled back
along \(T\) to a measurement on \(V\).

\begin{proposition}[Functorial identities and transpose matrix]
For linear maps \(T:V\to W\) and \(S:W\to U\),
\[
(\operatorname{id}_V)^\vee=\operatorname{id}_{V^\vee},
\qquad
(S\circ T)^\vee=T^\vee\circ S^\vee.
\]
If \(\beta\) and \(\gamma\) are ordered bases of \(V\) and \(W\), then
\[
[T^\vee]^{\beta^\vee}_{\gamma^\vee}
=\left([T]^\gamma_\beta\right)^t.
\]
\end{proposition}
\begin{proof}
For \(\varphi\in V^\vee\),
\[
(\operatorname{id}_V)^\vee(\varphi)
=\varphi\circ\operatorname{id}_V=\varphi.
\]
Likewise, for \(\psi\in U^\vee\),
\[
(S\circ T)^\vee(\psi)
=\psi\circ S\circ T
=T^\vee(S^\vee(\psi)).
\]
Write
\[
T(v_j)=\sum_{i=1}^ma_{ij}w_i,
\qquad
[T]^\gamma_\beta=(a_{ij}).
\]
For the \(i\)-th dual basis vector of \(W^\vee\),
\[
T^\vee(w_i^\vee)(v_j)=w_i^\vee(T(v_j))=a_{ij}.
\]
Hence
\[
T^\vee(w_i^\vee)=\sum_{j=1}^na_{ij}v_j^\vee.
\]
Its matrix therefore has \((j,i)\)-entry \(a_{ij}\), which is the transpose
of the matrix of \(T\).
\end{proof}

\begin{definition}
The \emph{evaluation map} is
\[
\iota_V:V\longrightarrow V^{\vee\vee},
\qquad
\iota_V(v)(\varphi)=\varphi(v).
\]
\end{definition}

\begin{proposition}[Double dual]\label{prop:double-dual}
The evaluation map \(\iota_V\) is linear and injective.  If \(V\) is
finite-dimensional, it is an isomorphism.
\end{proposition}
\begin{proof}
For \(a,b\in F\), \(v,w\in V\), and \(\varphi\in V^\vee\),
\[
\iota_V(av+bw)(\varphi)
=\varphi(av+bw)
=a\iota_V(v)(\varphi)+b\iota_V(w)(\varphi),
\]
so \(\iota_V\) is linear.  If \(v\ne0\), extend \(v\) to a basis and take
the corresponding first coordinate functional.  It does not vanish on
\(v\), so \(\iota_V(v)\ne0\).  Thus \(\iota_V\) is injective.  In finite
dimension, Proposition~\ref{prop:dual-basis} gives
\[
\dim V^{\vee\vee}=\dim V,
\]
so an injective map between these finite-dimensional spaces is surjective.
\end{proof}

Thus finite-dimensional vector spaces have a canonical isomorphism
\[
V\cong V^{\vee\vee},
\]
even though an isomorphism with the first dual is usually noncanonical.
This canonical character is expressed by the naturality identity
\begin{equation}\label{eq:double-dual-naturality}
\iota_W\circ T=T^{\vee\vee}\circ\iota_V
\qquad(T:V\to W).
\end{equation}
Indeed, for \(v\in V\) and \(\varphi\in W^\vee\),
\[
\begin{aligned}
\bigl(T^{\vee\vee}\iota_V(v)\bigr)(\varphi)
&=\iota_V(v)(\varphi\circ T)\\
&=(\varphi\circ T)(v)\\
&=\varphi(Tv)\\
&=\iota_W(Tv)(\varphi).
\end{aligned}
\]
Thus evaluation commutes with every linear map and requires no basis choice.

\begin{definition}
For a subspace \(W\leq V\), its \emph{annihilator} is
\[
W^\circ=\{\varphi\in V^\vee:\varphi(w)=0\text{ for every }w\in W\}.
\]
\end{definition}

\begin{proposition}
If \(V\) is finite-dimensional, then
\[
\dim W^\circ=\dim V-\dim W.
\]
\end{proposition}
\begin{proof}
Choose a basis \(w_1,\ldots,w_r\) of \(W\) and extend it to a basis
\[
(w_1,\ldots,w_r,u_1,\ldots,u_s)
\]
of \(V\).  Its dual basis shows that \(W^\circ\) has basis
\[
(u_1^\vee,\ldots,u_s^\vee).
\]
Since \(r+s=\dim V\), the result follows.
\end{proof}

\begin{theorem}[Dual maps, kernels, and images]\label{thm:dual-map-kernel-image}
Let \(T:V\to W\) be a linear map between finite-dimensional vector spaces.
Then
\[
\ker(T^\vee)=(\operatorname{im}T)^\circ,
\qquad
\operatorname{im}(T^\vee)=(\ker T)^\circ.
\]
Consequently,
\[
\operatorname{rank}(T^\vee)=\operatorname{rank}(T).
\]
\end{theorem}
\begin{proof}
For \(\varphi\in W^\vee\), every equivalence is direct:
\[
\begin{aligned}
\varphi\in\ker(T^\vee)
&\Longleftrightarrow T^\vee\varphi=0\\
&\Longleftrightarrow \varphi\circ T=0\\
&\Longleftrightarrow \varphi\text{ vanishes on }\operatorname{im}T\\
&\Longleftrightarrow \varphi\in(\operatorname{im}T)^\circ.
\end{aligned}
\]
Hence, using the annihilator dimension formula and rank--nullity,
\[
\begin{aligned}
\operatorname{rank}(T^\vee)
&=\dim W^\vee-\dim\ker(T^\vee)\\
&=\dim W-\bigl(\dim W-\dim\operatorname{im}T\bigr)\\
&=\operatorname{rank}T.
\end{aligned}
\]
If \(\varphi=T^\vee\psi=\psi\circ T\) and \(v\in\ker T\), then
\(\varphi(v)=\psi(Tv)=0\).  Thus
\(\operatorname{im}(T^\vee)\subseteq(\ker T)^\circ\).  The two spaces have
the same dimension, because
\[
\dim\operatorname{im}(T^\vee)=\operatorname{rank}T
=\dim V-\dim\ker T=\dim(\ker T)^\circ.
\]
The containment is therefore equality.
\end{proof}

In dual bases,
\([T^\vee]=[T]^t\).  The theorem therefore yields
\[
\operatorname{rank}(A^t)=\operatorname{rank}(A),
\]
giving the conceptual dual-space explanation of the row-rank equals
column-rank theorem proved earlier by elementary matrix methods.

\begin{proposition}[Quotient dual]\label{prop:quotient-dual}
If \(W\leq V\) and \(\pi:V\to V/W\) is the quotient map, then
\[
\pi^\vee:(V/W)^\vee\longrightarrow V^\vee
\]
is injective with image \(W^\circ\).  Hence
\[
(V/W)^\vee\cong W^\circ.
\]
\end{proposition}
\begin{proof}
Because \(\pi\) is surjective, \(\psi\circ\pi=0\) implies \(\psi=0\), so
\(\pi^\vee\) is injective.  Theorem~\ref{thm:dual-map-kernel-image} gives
\[
\operatorname{im}(\pi^\vee)=(\ker\pi)^\circ=W^\circ.
\]
\end{proof}

\begin{proposition}[Annihilator identities]\label{prop:annihilator-identities}
For subspaces \(U,W\leq V\), with \(V\) finite-dimensional,
\[
U\subseteq W\Longrightarrow W^\circ\subseteq U^\circ,
\qquad
(U+W)^\circ=U^\circ\cap W^\circ,
\]
\[
(U\cap W)^\circ=U^\circ+W^\circ,
\qquad
W^{\circ\circ}=W
\]
under the evaluation identification \(V\cong V^{\vee\vee}\).
\end{proposition}
\begin{proof}
A functional vanishing on \(W\) certainly vanishes on \(U\subseteq W\),
which proves reversal of inclusions.  A functional vanishes on \(U+W\)
exactly when it vanishes on both summands, proving the second identity.

Every element of \(U^\circ+W^\circ\) vanishes on \(U\cap W\), so
\[
U^\circ+W^\circ\subseteq(U\cap W)^\circ.
\]
The dimension formula for sums and intersections gives
\[
\begin{aligned}
\dim(U^\circ+W^\circ)
&=\dim U^\circ+\dim W^\circ-\dim(U^\circ\cap W^\circ)\\
&=(\dim V-\dim U)+(\dim V-\dim W)-\dim(U+W)^\circ\\
&=\dim V-\dim(U\cap W),
\end{aligned}
\]
which is the dimension of \((U\cap W)^\circ\); hence equality holds.
Finally, evaluation sends every \(w\in W\) to a functional on \(V^\vee\)
that vanishes on \(W^\circ\), so \(W\subseteq W^{\circ\circ}\).  Both spaces
have dimension \(\dim W\), and therefore they are equal.
\end{proof}

\section{Bilinear and Quadratic Forms}\label{sec:bilinear-quadratic-forms}
Duality turns a vector into a linear measurement.  A bilinear form instead
assigns a scalar to an ordered pair of vectors, linearly in each argument.
This is the algebraic framework behind dot products, but it does not include
positivity or conjugation.

\begin{definition}
A \emph{bilinear form} on an \(F\)-vector space \(V\) is a map
\[
B:V\times V\longrightarrow F
\]
that is linear in each variable.  It is \emph{symmetric} if
\[
B(x,y)=B(y,x),
\]
\emph{skew-symmetric} if
\[
B(x,y)=-B(y,x),
\]
and \emph{alternating} if
\[
B(v,v)=0
\qquad(v\in V).
\]
\end{definition}

Every alternating bilinear form is skew-symmetric, by expanding
\[
0=B(x+y,x+y).
\]
Conversely, a skew-symmetric form is alternating only when the
characteristic is not \(2\), in the sense of
Definition~\ref{def:ring-characteristic}. This distinction is essential in
characteristic \(2\).
In every ordered basis \(\beta\),
\[
B\text{ symmetric}\Longleftrightarrow[B]_\beta^t=[B]_\beta,
\qquad
B\text{ skew-symmetric}\Longleftrightarrow[B]_\beta^t=-[B]_\beta.
\]
Alternating always implies the second condition.  If
\(\operatorname{char}F\ne2\), the second condition forces
\(2B(v,v)=0\) and hence \(B(v,v)=0\).  In characteristic \(2\), however,
\(-A=A\), so the transpose equation merely says that \(A\) is symmetric; an
alternating form must additionally have zero diagonal.

\begin{proposition}[Matrix and congruence]\label{prop:bilinear-matrix}
For an ordered basis \(\beta=(v_1,\ldots,v_n)\), put
\[
[B]_\beta=(B(v_i,v_j)).
\]
Then
\[
B(x,y)=[x]_\beta^t[B]_\beta[y]_\beta.
\]
If
\[
P=[I_V]^\beta_{\beta'},
\]
so that
\[
[x]_\beta=P[x]_{\beta'},
\]
then
\[
[B]_{\beta'}=P^t[B]_\beta P.
\]
\end{proposition}
\begin{proof}
Write
\[
x=\sum_i a_iv_i,
\qquad
y=\sum_j b_jv_j.
\]
By bilinearity,
\[
B(x,y)=\sum_{i,j}a_iB(v_i,v_j)b_j,
\]
which is the stated matrix product.  The columns of \(P\) are the
\(\beta\)-coordinates of the vectors of \(\beta'\).  Substitution of
\[
[x]_\beta=P[x]_{\beta'},
\qquad
[y]_\beta=P[y]_{\beta'}
\]
into the coordinate formula gives
\[
B(x,y)
=[x]_{\beta'}^tP^t[B]_\beta P[y]_{\beta'},
\]
and uniqueness of a matrix representing the form in \(\beta'\) proves the
claim.
\end{proof}

This is \emph{congruence}, not similarity.  For an operator, one input is
converted into output coordinates and change of basis gives
\[
[T]_{\beta'}=P^{-1}[T]_\beta P.
\]
A bilinear form has two vector inputs, so both coordinates are substituted
and the formula has \(P^t\) on the left instead.

\begin{definition}
The \emph{left radical} and \emph{right radical} of \(B\) are, respectively,
\[
\operatorname{rad}_L(B)
=\{v\in V:B(v,w)=0\text{ for every }w\in V\},
\]
\[
\operatorname{rad}_R(B)
=\{w\in V:B(v,w)=0\text{ for every }v\in V\}.
\]
For a general bilinear form these are conceptually distinct.  The associated
linear maps are
\[
\Phi_B:V\longrightarrow V^\vee,
\qquad
\Phi_B(v)=B(v,-).
\]
and
\[
\Psi_B:V\longrightarrow V^\vee,
\qquad
\Psi_B(w)=B(-,w).
\]
\end{definition}

\begin{definition}
The \emph{rank of a bilinear form} is
\[
\operatorname{rank}B:=\operatorname{rank}\Phi_B.
\]
\end{definition}

\begin{proposition}[Rank and radicals of a bilinear form]
For every ordered basis \(\beta\),
\[
\operatorname{rank}B=\operatorname{rank}[B]_\beta.
\]
This number is independent of the basis.  Moreover,
\[
\dim\operatorname{rad}_L(B)=\dim V-\operatorname{rank}B,
\qquad
\dim\operatorname{rad}_R(B)=\dim V-\operatorname{rank}B.
\]
\end{proposition}
\begin{proof}
In \(\beta\) and its dual basis, the matrix of \(\Phi_B\) is
\([B]_\beta^t\).  Equality of row and column rank therefore gives
\[
\operatorname{rank}\Phi_B
=\operatorname{rank}[B]_\beta^t
=\operatorname{rank}[B]_\beta.
\]
Alternatively, basis independence follows directly from
\([B]_{\beta'}=P^t[B]_\beta P\) and invariance of rank under multiplication
by invertible matrices.  Rank--nullity applied to \(\Phi_B\) proves the
left-radical formula.  The right-radical formula follows in the same way from
\(\Psi_B\), whose matrix is \([B]_\beta\).
\end{proof}

\begin{proposition}[Nondegeneracy criterion]
For a finite-dimensional vector space, the following are equivalent:
\begin{enumerate}
\item \(\operatorname{rad}_L(B)=\{0\}\);
\item \(\operatorname{rad}_R(B)=\{0\}\);
\item \([B]_\beta\) is invertible for one, equivalently every, basis
\(\beta\).
\end{enumerate}
When these conditions hold, \(B\) is called \emph{nondegenerate}.
\end{proposition}
\begin{proof}
By their definitions,
\[
\ker\Phi_B=\operatorname{rad}_L(B),
\qquad
\ker\Psi_B=\operatorname{rad}_R(B).
\]
In the basis \(\beta\) and its dual, the matrices of these maps are
\[
[\Phi_B]_{\beta}^{\beta^\vee}=[B]_\beta^t,
\qquad
[\Psi_B]_{\beta}^{\beta^\vee}=[B]_\beta.
\]
Indeed, the \(j\)-th coordinate of \(\Phi_B(v_i)\) is \(B(v_i,v_j)\),
whereas that of \(\Psi_B(v_i)\) is \(B(v_j,v_i)\).  Thus either map is
injective exactly when \([B]_\beta\) is invertible.  Since
\[
\dim V=\dim V^\vee,
\]
injectivity is equivalent to bijectivity for both maps.  Finally, the
congruence formula in Proposition~\ref{prop:bilinear-matrix} shows that
matrix invertibility is independent of the chosen basis.
\end{proof}

Equivalently, in finite dimension,
\[
\begin{aligned}
B\text{ nondegenerate}
&\Longleftrightarrow \operatorname{rank}B=\dim V\\
&\Longleftrightarrow [B]_\beta\text{ is invertible}
\Longleftrightarrow \det[B]_\beta\ne0.
\end{aligned}
\]
A chosen nondegenerate bilinear form therefore supplies the isomorphism
\[
\Phi_B:V\xrightarrow{\ \sim\ }V^\vee,\qquad v\longmapsto B(v,-).
\]
This is precisely the extra structure missing from the bare vector space:
there is no preferred isomorphism \(V\cong V^\vee\), but a nondegenerate form
provides one.

For a symmetric bilinear form, \(B(v,w)=B(w,v)\), so the two radicals
coincide.  In that case we write their common value as \(\operatorname{rad}(B)\).

\begin{definition}
For a symmetric bilinear form \(B\) and a subspace \(W\leq V\), define
\[
W^{\perp_B}=\{v\in V:B(v,w)=0\text{ for every }w\in W\}.
\]
\end{definition}

\begin{proposition}[Orthogonal complements for a bilinear form]
If \(B\) is nondegenerate on finite-dimensional \(V\), then
\[
\dim W^{\perp_B}=\dim V-\dim W.
\]
Moreover,
\[
V=W\oplus W^{\perp_B}
\quad\Longleftrightarrow\quad
B|_W\text{ is nondegenerate}.
\]
\end{proposition}
\begin{proof}
The isomorphism \(\Phi_B:V\to V^\vee\) carries \(W^{\perp_B}\) onto the
annihilator \(W^\circ\).  The annihilator dimension formula gives the first
claim.  Also,
\[
W\cap W^{\perp_B}=\operatorname{rad}(B|_W).
\]
Thus the sum is direct exactly when \(B|_W\) is nondegenerate; in that case
the dimension formula shows that the direct sum has dimension \(\dim V\) and
therefore equals \(V\).  The converse is immediate from the displayed
intersection.
\end{proof}

Unlike the positive-definite case, nondegeneracy of \(B\) on \(V\) alone does
not guarantee this direct sum.  For example, on \(F^2\) with
\[
[B]=\begin{pmatrix}0&1\\1&0\end{pmatrix},
\]
the form is nondegenerate, but \(e_1\) is isotropic and
\(Fe_1=(Fe_1)^{\perp_B}\).  Thus isotropic vectors can make
\(W\cap W^{\perp_B}\ne0\).

\begin{definition}
Assume \(\operatorname{char}F\ne2\), in the sense of
Definition~\ref{def:ring-characteristic}. A quadratic form associated with a
symmetric bilinear form \(B\) is
\[
q(v)=B(v,v).
\]
\end{definition}

The form can be recovered from its quadratic function by polarization:
\[
B(x,y)
=\frac12\bigl(q(x+y)-q(x)-q(y)\bigr).
\]
Characteristic \(2\) requires a separate theory: a quadratic form must not
there be identified blindly with a symmetric bilinear form.

In coordinates, if \(A=[B]_\beta=A^t\) and
\(x=(x_1,\ldots,x_n)^t\), then
\[
q(x)=x^tAx
=\sum_i a_{ii}x_i^2+2\sum_{i<j}a_{ij}x_ix_j.
\]
Conversely, the polarization formula recovers the unique symmetric bilinear
form from \(q\).  If \(x=Py\) and \(P^tAP=D\), then
\[
q(Py)=(Py)^tA(Py)=y^tP^tAPy=y^tDy
=d_1y_1^2+\cdots+d_ny_n^2.
\]
Thus congruence is exactly the coordinate change that diagonalizes a
quadratic form.

\begin{example}[Completing the square]
Over a field of characteristic not \(2\), consider
\[
q(x,y)=x^2+4xy+3y^2.
\]
The cross term disappears visibly:
\[
q(x,y)=(x+2y)^2-y^2.
\]
Set \(u=x+2y\) and \(v=y\), equivalently
\[
\begin{pmatrix}x\\y\end{pmatrix}
=\begin{pmatrix}1&-2\\0&1\end{pmatrix}
\begin{pmatrix}u\\v\end{pmatrix}.
\]
Then \(q=u^2-v^2\).  This is the coordinate mechanism behind diagonalization,
not a classification of all quadratic forms.
\end{example}

\begin{theorem}[Diagonalization of symmetric bilinear forms]
Suppose that \(\operatorname{char}F\ne2\).  Every symmetric bilinear form on
a finite-dimensional \(F\)-vector space has a basis in which its matrix is
diagonal.  Equivalently, its quadratic form has coordinates
\[
q(x_1,\ldots,x_n)=a_1x_1^2+\cdots+a_nx_n^2.
\]
\end{theorem}
\begin{proof}
Induct on \(\dim V\).  If \(B=0\), there is nothing to prove.  Otherwise
some \(v\) satisfies \(B(v,v)\ne0\): if every such value were zero, then
polarization would give \(B=0\).  Let
\[
W=\{w\in V:B(v,w)=0\}.
\]
For \(x\in V\),
\[
x=\frac{B(x,v)}{B(v,v)}v+
\left(x-\frac{B(x,v)}{B(v,v)}v\right),
\]
and the parenthesized vector lies in \(W\).  Thus
\[
V=Fv\oplus W.
\]
The restriction of \(B\) to \(W\) is symmetric, so induction supplies a
diagonalizing basis of \(W\); adjoining \(v\) gives the desired basis.
\end{proof}

Over \(\mathbb R\), inner products are precisely positive-definite symmetric
bilinear forms. Over \(\mathbb C\), the inner products of Chapter~7 are
Hermitian: they are linear in the first variable and conjugate-linear in the
second. Thus Chapter~7 is not merely the complex special case of this
bilinear discussion.


\section*{Exercises}
\begin{exercise}[Basic] Verify directly that abelian groups are precisely
\(\mathbb Z\)-modules, and determine the submodules of
\(\mathbb Z/12\mathbb Z\).\end{exercise}
\begin{exercise}[Core] Prove the universal property of the direct sum for an
arbitrary family of modules.\end{exercise}
\begin{exercise}[Core] Let \(T:F^4\to F^3\) have rank \(2\). Find its
nullity and decide whether it can be injective or surjective.\end{exercise}
\begin{exercise}[Core] Compute the matrix of differentiation on polynomials
of degree at most \(3\), using the basis \(1,x,x^2,x^3\).\end{exercise}
\begin{exercise}[Core] Show that similar matrices have the same characteristic
polynomial and determinant.\end{exercise}
\begin{exercise}[Further] Decide which of the matrices
\(\begin{pmatrix}1&1\\0&1\end{pmatrix}\) and
\(\begin{pmatrix}1&0\\0&1\end{pmatrix}\) are diagonalizable.\end{exercise}
\begin{exercise}[Further, Challenging] Prove that an operator is
diagonalizable exactly when the sum of the dimensions of its eigenspaces is
\(\dim V\).\end{exercise}
\begin{exercise}[Core] For the standard basis of \(F^3\), compute its dual
basis and the matrix of the dual of a specified \(3\times3\) matrix.  Then
verify the congruence formula for the bilinear form whose matrix is
\[
\begin{pmatrix}2&1\\1&0\end{pmatrix}
\]
after the basis change \((e_1,e_2)\mapsto(e_1+e_2,e_2)\).
\end{exercise}

\chapter[PID Modules and Canonical Forms]{Modules over PIDs and Canonical Forms}\label{chap:pid-modules-canonical-forms}

Chapter~5 showed that diagonalization is possible exactly when an operator has
enough eigenvectors. This chapter asks what remains when it does not. The
answer comes from viewing a linear operator as multiplication by \(x\), and
from the remarkable fact that finitely generated modules over a principal
ideal domain have a complete decomposition theory. The same theorem therefore
classifies finitely generated abelian groups and produces rational and Jordan
canonical forms.

The first half of the chapter deliberately uses three different viewpoints.
Compatible ambient coordinates construct the decomposition, primary layers
recognize its pieces intrinsically, and Smith reduction expresses the result in computational
matrix language.  Keeping those tasks separate is the key to the proof
architecture.

Through the Smith normal form section, \(R\) is a principal ideal domain (PID).
A PID need not be a Euclidean domain. We use only principal ideals, gcds and
B\'ezout identities; no division-with-remainder argument is used over an
arbitrary PID. The polynomial ring \(F[x]\), used later, is Euclidean, so its
ordinary polynomial division algorithm is available there.

\section{Submodules of Free Modules over a PID}\label{sec:pid-submodules-free}
Over a field, every subspace has a basis. The analogous statement fails over
general rings, but a PID retains just enough divisibility control for finite
free modules to behave similarly.

\begin{theorem}[Submodules of finite free modules over a PID]
\label{thm:pid-submodules-compatible-bases}
Let \(R\) be a PID, let \(F\) be a free \(R\)-module of finite rank \(n\),
and let \(N\subseteq F\) be a submodule.
\begin{enumerate}
\item The module \(N\) is free of rank at most \(n\).
\item There exist a basis \(y_1,\ldots,y_n\) of \(F\), an integer
      \(m\leq n\), and nonzero elements \(a_1,\ldots,a_m\in R\) such that
\[
a_1\mid a_2\mid\cdots\mid a_m
\qquad\text{and}\qquad
N=Ra_1y_1\oplus\cdots\oplus Ra_my_m.
\]
\end{enumerate}
Part~(1) is the freeness conclusion. Part~(2) strengthens it by choosing
compatible coordinates in the ambient module; it makes no claim about an
arbitrarily prescribed basis of \(F\).
\end{theorem}
\begin{proof}
\emph{Proof of (1).}
Choose a basis of \(F\) and identify it with \(R^n\). We argue by induction
on \(n\). The assertion is clear for \(n=0\). Let
\[
\pi:R^n\longrightarrow R
\]
be projection onto the first coordinate, and set
\[
K=\ker(\pi|_N)=N\cap(0\oplus R^{n-1}).
\]
The image \(\pi(N)\) is an ideal, say \((d)\).

If \(\pi(N)=0\), then \(N\subseteq0\oplus R^{n-1}\). Induction makes \(N\)
free of rank at most \(n-1\), and hence at most \(n\).

Suppose instead that \(\pi(N)=(d)\ne0\). Choose \(x\in N\) with
\(\pi(x)=d\). For every \(y\in N\), there is \(r\in R\) such that
\(\pi(y)=rd\). Therefore
\[
\pi(y-rx)=rd-rd=0,
\]
so \(y-rx\in K\), and
\[
y=rx+(y-rx)\in Rx+K.
\]
This proves \(N=Rx+K\). The sum is direct: if \(ax+k=0\) with \(k\in K\),
then applying \(\pi\) gives
\[
ad=\pi(ax+k)=0.
\]
Because \(R\) is a domain and \(d\ne0\), we have \(a=0\), and then \(k=0\).

By induction, choose a basis \(k_1,\ldots,k_s\) of \(K\), with
\(s\leq n-1\). The decomposition above shows that
\(x,k_1,\ldots,k_s\) spans \(N\), and the directness calculation together
with independence of the \(k_i\) shows that this list is independent. Hence
it is a basis of \(N\), and
\[
\operatorname{rank}N=1+s\leq n.
\]

\medskip
\noindent\emph{Proof idea for (2).}
We find one invariant-factor direction at a time.  Among all linear
functionals \(\phi:F\to R\), choose one whose image \(\phi(N)\) is maximal
as an ideal.  Its generator becomes the first invariant factor.  We show that
a corresponding element of \(N\) is a scalar multiple of a primitive basis
vector, split this direction from both \(F\) and \(N\), and repeat inside a
free module of rank one smaller:
\[
\begin{gathered}
\text{choose the first coefficient}\longrightarrow
\text{find a primitive direction},\\
\text{split it off}\longrightarrow\text{induct}.
\end{gathered}
\]

\medskip
\noindent\emph{Proof of (2).}
We induct on \(n\).  The result is clear if \(n=0\), and if \(N=0\) we may
take \(m=0\) and any basis of \(F\).  Suppose therefore that \(n>0\) and
\(N\ne0\).

\emph{Stage 1: choose a maximal image ideal.}
For every \(\phi\in\operatorname{Hom}_R(F,R)\), the set \(\phi(N)\) is an
ideal.  A PID is Noetherian because all its ideals are principal, so the
maximal-member characterization of Noetherian rings proved in Chapter~4
gives a maximal member of the partially ordered set
\[
\Sigma=\{\phi(N):\phi\in\operatorname{Hom}_R(F,R)\},
\qquad\text{ordered by set inclusion}.
\]
Choose \(\psi:F\to R\) with \(\psi(N)=(a_1)\) maximal.  This ideal is
nonzero. Indeed, choose \(0\ne z\in N\) and a basis
\(x_1,\ldots,x_n\) of \(F\). Write
\[
z=r_1x_1+\cdots+r_nx_n.
\]
Some \(r_i\ne0\). Define the projection onto the \(i\)-th coordinate
relative to this basis by
\[
\pi_i:F\longrightarrow R,
\qquad
\pi_i\left(\sum_{j=1}^n r_jx_j\right)=r_i.
\]
Then \(\pi_i(z)=r_i\ne0\), so \(\pi_i(N)\ne0\) is a nonzero member of
\(\Sigma\). Recall the divisibility reversal
\[
(a)\subseteq(b)\quad\Longleftrightarrow\quad b\mid a.
\]
Thus enlarging an image ideal means finding a more fundamental divisor.

\emph{Stage 2: force divisibility of every scalar measurement.}
Choose \(z_1\in N\) with \(\psi(z_1)=a_1\).  We claim that
\[
a_1\mid\phi(z_1)
\qquad\text{for every }\phi:F\to R.
\]
If one scalar measurement of \(z_1\) were not divisible by \(a_1\), the
B\'ezout calculation below would produce a strictly larger image ideal.

Fix an arbitrary \(\phi:F\to R\). Let \((a_1,\phi(z_1))=(g)\), and choose
\(r,s\in R\) such that
\(g=ra_1+s\phi(z_1)\). Define \(\nu=r\psi+s\phi\). Then
\[
\nu(z_1)=r\psi(z_1)+s\phi(z_1)
=ra_1+s\phi(z_1)=g.
\]
Since \(z_1\in N\), this gives \(g\in\nu(N)\), and hence
\((g)\subseteq\nu(N)\). Also \(g\mid a_1\), so principal-ideal inclusion
reverses divisibility and gives \((a_1)\subseteq(g)\). Therefore
\[
(a_1)\subseteq(g)\subseteq\nu(N).
\]
The map \(\nu\) belongs to \(\operatorname{Hom}_R(F,R)\), so \(\nu(N)\)
is a member of \(\Sigma\). Since \((a_1)=\psi(N)\) was chosen maximal in
\(\Sigma\), the inclusion \((a_1)\subseteq\nu(N)\) forces
\(\nu(N)=(a_1)\). Consequently \((a_1)=(g)\), so \(g\) is associate to
\(a_1\). Since \(g\mid\phi(z_1)\), it follows that
\(a_1\mid\phi(z_1)\).

\emph{Stage 3: identify a primitive direction.}
Choose any basis \(x_1,\ldots,x_n\) of \(F\), and define the projection onto
the \(i\)-th coordinate relative to this basis by
\[
\pi_i:F\longrightarrow R,
\qquad
\pi_i\left(\sum_{j=1}^n r_jx_j\right)=r_i.
\]
Since the preceding step applies to every \(\pi_i\), write
\(\pi_i(z_1)=a_1b_i\) and set
\[
y_1=\sum_{i=1}^n b_ix_i.
\]
The coordinate expansion of \(z_1\) now gives
\[
z_1=\sum_i\pi_i(z_1)x_i
=\sum_i a_1b_ix_i=a_1y_1.
\]
Moreover,
\[
a_1=\psi(z_1)=\psi(a_1y_1)=a_1\psi(y_1).
\]
Because \(R\) is a domain and \(a_1\ne0\), we have \(\psi(y_1)=1\).
Thus \(y_1\) has an \(R\)-linear functional taking it to \(1\). This is the
key primitive property: it makes \(Ry_1\) split off from \(F\).

\emph{Stage 4: split \(F\) and \(N\) simultaneously.}
For every \(x\in F\),
\[
x=\psi(x)y_1+\bigl(x-\psi(x)y_1\bigr).
\]
Indeed,
\[
\psi\bigl(x-\psi(x)y_1\bigr)
=\psi(x)-\psi(x)\psi(y_1)=0,
\]
so the second term lies in \(\ker\psi\), and hence
\(F=Ry_1+\ker\psi\). If \(ry_1\in\ker\psi\), then
\[
0=\psi(ry_1)=r\psi(y_1)=r.
\]
Thus \(Ry_1\cap\ker\psi=0\), and therefore
\[
F=Ry_1\oplus\ker\psi.
\]
If \(z\in N\), then \(\psi(z)=ra_1\) for some \(r\in R\), so
\(z-ra_1y_1\in N\), and
\[
\psi(z-ra_1y_1)=ra_1-ra_1\psi(y_1)=0.
\]
Therefore
\[
z=ra_1y_1+(z-ra_1y_1),
\]
with the second summand in \(N\cap\ker\psi\). This proves
\(N=Ra_1y_1+(N\cap\ker\psi)\). If
\(ra_1y_1\in N\cap\ker\psi\), then
\[
0=\psi(ra_1y_1)=ra_1.
\]
Since \(R\) is a domain and \(a_1\ne0\), we have \(r=0\). Hence the sum is
direct:
\[
N=Ra_1y_1\oplus(N\cap\ker\psi).
\]
One relation direction has now been completely separated from all the others.

\emph{Stage 5: induct and recover the divisibility chain.}
The module \(\ker\psi\) is a submodule of \(F\), hence free by part~(1).
Since \(\operatorname{rank}F=n\), \(\operatorname{rank}(Ry_1)=1\), and
\(F=Ry_1\oplus\ker\psi\),
Proposition~\ref{prop:finite-direct-sum-rank} gives
\[
n=1+\operatorname{rank}(\ker\psi).
\]
Thus \(\operatorname{rank}(\ker\psi)=n-1\), which justifies the induction
parameter. Apply induction to
\(N\cap\ker\psi\subseteq\ker\psi\).  We obtain basis vectors
\(y_2,\ldots,y_n\) and nonzero elements \(a_2,\ldots,a_m\) with
\[
a_2\mid\cdots\mid a_m,
\qquad
N\cap\ker\psi=Ra_2y_2\oplus\cdots\oplus Ra_my_m.
\]
If \(m\geq2\), define \(\phi:F\to R\) on the new basis by
\[
\phi(y_1)=1,\qquad \phi(y_2)=1,\qquad
\phi(y_i)=0\quad(i\geq3).
\]
Then
\[
\phi(a_1y_1)=a_1,
\qquad
\phi(a_2y_2)=a_2.
\]
Hence \(a_1,a_2\in\phi(N)\), and in particular
\((a_1)\subseteq\phi(N)\). Because
\(\phi\in\operatorname{Hom}_R(F,R)\), the ideal \(\phi(N)\) is a member of
\(\Sigma\). Maximality of \((a_1)\) therefore gives
\(\phi(N)=(a_1)\).
Therefore \(a_2\in(a_1)\), or \(a_1\mid a_2\), as required.
\end{proof}

\begin{corollary}[Rank--nullity for finite free modules over a PID]
\label{cor:pid-rank-nullity}
Let \(R\) be a PID and let \(f:F\to G\) be an \(R\)-linear map between
finite free \(R\)-modules. Then \(\ker f\) and \(\operatorname{im}f\) are
finite free, the short exact sequence
\[
0\longrightarrow\ker f\longrightarrow F
\xrightarrow{f}\operatorname{im}f\longrightarrow0
\]
splits, and
\[
F\cong\ker f\oplus\operatorname{im}f,
\qquad
\operatorname{rank}F
=\operatorname{rank}(\ker f)+\operatorname{rank}(\operatorname{im}f).
\]
More precisely, after choosing a section
\(s:\operatorname{im}f\to F\), the literal internal decomposition is
\(F=\ker f\oplus s(\operatorname{im}f)\), with
\(s(\operatorname{im}f)\cong\operatorname{im}f\).
\end{corollary}
\begin{proof}
The inclusions \(\ker f\leq F\) and \(\operatorname{im}f\leq G\), together
with Theorem~\ref{thm:pid-submodules-compatible-bases}(1), show that both
modules are finite free. In particular \(\operatorname{im}f\) is projective,
so Theorem~\ref{thm:projective-characterizations} splits the displayed short
exact sequence. The Splitting Lemma gives the stated internal decomposition,
and Proposition~\ref{prop:finite-direct-sum-rank} gives the rank formula.
\end{proof}

One should not conclude similarly that
\(G\cong\operatorname{im}f\oplus\operatorname{coker}f\): although
\(2\mathbb Z\leq\mathbb Z\) is free, it is not a direct summand.

\section{Structure Theorem: Existence in Invariant-Factor Form}
\label{sec:pid-structure-existence}
A finitely generated module is built from generators subject to relations.
The compatible-basis conclusion of Theorem~\ref{thm:pid-submodules-compatible-bases}
separates those relations inside a finite free module; taking the quotient
then reads off the answer coordinate by coordinate.

\begin{theorem}[Structure theorem: existence]\label{thm:pid-structure-existence}
Let \(M\) be a finitely generated module over a PID \(R\).
\begin{enumerate}
\item There exist an integer \(r\geq0\) and nonzero nonunits
\(a_1,\ldots,a_m\) such that
\[
a_1\mid a_2\mid\cdots\mid a_m
\quad\text{and}\quad
M\cong R/(a_1)\oplus\cdots\oplus R/(a_m)\oplus R^r.
\]
\item The module \(M\) is free if and only if it is torsion-free.
\item In the decomposition above,
\[
T(M)\cong R/(a_1)\oplus\cdots\oplus R/(a_m),
\qquad M/T(M)\cong R^r.
\]
\end{enumerate}
\end{theorem}
\begin{proofidea}
Choose generators of \(M\) and realize it as a quotient \(F/N\) of a finite
free module. Apply Theorem~\ref{thm:pid-submodules-compatible-bases}(2) to
choose compatible coordinates for \(N\subseteq F\). The quotient becomes
coordinatewise, producing the invariant factors. Once the free and torsion
summands are visible, parts (2) and (3) follow immediately.
\end{proofidea}
\begin{proof}
Choose a surjection \(F=R^n\twoheadrightarrow M\) and let \(N\) be its
kernel. By Theorem~\ref{thm:pid-submodules-compatible-bases}(2), there are a basis
\(y_1,\ldots,y_n\) of \(F\) and nonzero \(a_1\mid\cdots\mid a_s\) such that
\[
N=Ra_1y_1\oplus\cdots\oplus Ra_sy_s.
\]
Define
\[
\Phi:F\longrightarrow
R/(a_1)\oplus\cdots\oplus R/(a_s)\oplus R^{n-s}
\]
by
\[
\Phi\left(\sum_{i=1}^n r_iy_i\right)
=\bigl(r_1+(a_1),\ldots,r_s+(a_s),r_{s+1},\ldots,r_n\bigr).
\]
Coordinatewise addition and scalar multiplication show directly that
\(\Phi\) is \(R\)-linear. It is surjective: choose representatives in \(R\)
for the first \(s\) coordinates of any target tuple and use its remaining
coordinates directly. Moreover,
\[
\Phi\left(\sum_i r_iy_i\right)=0
\]
if and only if \(r_i\in(a_i)\) for \(1\leq i\leq s\) and \(r_i=0\) for
\(i>s\). Thus \(r_i=a_ic_i\) for \(i\leq s\), and
\[
\ker\Phi=Ra_1y_1\oplus\cdots\oplus Ra_sy_s=N.
\]
The First Isomorphism Theorem therefore gives
\[
M\cong F/N
\cong R/(a_1)\oplus\cdots\oplus R/(a_s)\oplus R^{n-s}.
\]
If some \(a_i\) is a unit, then \(R/(a_i)=0\); omit every such zero summand.
Relabel the remaining divisibility chain as \(a_1,\ldots,a_m\).  Its cyclic
factors are nonzero and the free rank is \(r=n-s\), which proves part (1).

If \(M\) is free, it is torsion-free because \(R\) is a domain.  Conversely,
each nonzero summand \(R/(a_i)\) in part (1) is torsion, since
\[
a_i\bigl(1+(a_i)\bigr)=0.
\]
Thus a torsion-free \(M\) has no cyclic torsion summands and is free, proving
part (2).

Finally, set \(C=R/(a_1)\oplus\cdots\oplus R/(a_m)\). If \(m>0\), the
divisibility \(a_i\mid a_m\) makes \(a_m\) annihilate every cyclic summand,
hence every element of \(C\); if \(m=0\), then \(C=0\). Thus
\(C\subseteq T(M)\). Conversely, write an element of
\(M\cong C\oplus R^r\) as \(x=c+u\). If \(x\) is torsion, some
\(0\ne d\in R\) satisfies \(d(c+u)=0\). Its free coordinate gives
\(du=0\). Since \(R^r\) is torsion-free, \(u=0\), so every torsion element
lies in \(C\). Therefore \(T(M)=C\), and quotienting gives
\(M/T(M)\cong R^r\). This proves part (3).
\end{proof}

\section{Elementary Divisors and Primary Decomposition}
Invariant factors group the torsion into a divisibility chain.  Elementary
divisors reorganize the same module in the opposite direction: they separate
the independent prime-power phenomena.  This is still an existence argument;
intrinsic uniqueness comes afterward.

\begin{proposition}[Elementary-divisor existence]
The torsion part of a finitely generated PID module is a direct sum of modules
\(R/(p^e)\), where \(p\) is a prime element and \(e\geq1\).
\end{proposition}
\begin{proof}
Every PID is a UFD by Theorem~\ref{thm:pid-ufd}. Factor an invariant factor as
\[
d=u\prod_i p_i^{e_i},
\]
with distinct \(p_i\). The ideals
\((p_i^{e_i})\) are pairwise comaximal: a common divisor of two distinct
prime powers is a unit, so a B\'ezout identity gives their sum as \(R\).
The Chinese Remainder Theorem therefore gives
\[R/(d)\cong\bigoplus_iR/(p_i^{e_i}).\]
Apply this to every invariant factor.
\end{proof}

Fix one representative from each associate class of prime elements of \(R\).
When \(R=F[x]\), we take the monic irreducible polynomials as representatives.
For such a representative \(p\), define the \(p\)-primary component of \(M\) by
\[
M_p:=M[p^\infty]
=\{m\in M:p^em=0\text{ for some }e\geq1\}.
\]
Grouping the elementary-divisor summands by these representatives gives
submodules \(E_p\) and a decomposition
\[
T(M)=\bigoplus_p E_p.
\]
Each \(E_p\) is killed by a power of \(p\), so
\(E_p\subseteq M_p\). Conversely, suppose
\(m\in M_p\) and write \(m=\sum_qm_q\), with \(m_q\in E_q\).
For \(q\ne p\), some powers \(p^e\) and \(q^f\) kill \(m_q\). Since
distinct prime representatives are relatively prime, B\'ezout gives
\(up^e+vq^f=1\), and therefore
\(m_q=up^em_q+vq^fm_q=0\). Thus \(m=m_p\in E_p\), proving
\(E_p=M_p\). Hence the intrinsic primary decomposition is
\[
T(M)=\bigoplus_p M_p,
\]
where \(p\) runs over the fixed representatives and only finitely many
summands are nonzero. For example, the CRT
decomposition \(\mathbb Z/12\cong\mathbb Z/4\oplus\mathbb Z/3\) separates
the \(2\)-primary and \(3\)-primary behavior.

\section{Why Is the Decomposition Unique?}
Existence asked a constructive question: how can we choose coordinates that
split a module into standard pieces?  Uniqueness asks a different question:
how can those pieces be recognized from \(M\) without referring to chosen
coordinates?

The free part is immediate from Theorem~\ref{thm:pid-structure-existence}(3).
The torsion submodule \(T(M)\) is intrinsic, so
\[
M/T(M)\cong R^r.
\]
Invariant basis number from Chapter~5 therefore makes \(r\) intrinsic.  The
real problem is to recover the torsion blocks.

There is a useful first clue. Suppose a nonzero torsion module \(M\) has an
invariant-factor decomposition
\[
M\cong R/(a_1)\oplus\cdots\oplus R/(a_k),
\qquad a_1\mid\cdots\mid a_k.
\]
The annihilator of a direct sum is the intersection of the annihilators of
its summands: a scalar annihilates the direct sum exactly when it annihilates
every coordinate. The divisibility chain therefore gives
\[
\operatorname{Ann}(M)=(a_1)\cap\cdots\cap(a_k)=(a_k).
\]
Thus the largest invariant factor is visible intrinsically.  It is tempting
to remove a displayed copy of \(R/(a_k)\) and repeat.  But an abstract
isomorphism of direct sums need not identify the last displayed summand on
one side with the last displayed summand on the other.  Cancellation cannot
be assumed silently.  We instead seek stronger intrinsic invariants that
recover every block at once.

Different primes describe independent torsion phenomena, so the intrinsic
submodules \(M_p\) tell us to fix \(p\) and study \(M_p\), solving uniqueness
one prime at a time.

\section{Primary Layers and Uniqueness}\label{sec:primary-layers}
Let \(P\) be a finitely generated \(p\)-primary module.  Elementary-divisor
existence shows that one common power of \(p\) annihilates \(P\), so its powers
of \(p\) form a finite filtration
\[
P\supseteq pP\supseteq p^2P\supseteq\cdots\supseteq0.
\]
What disappears between \(p^{j-1}P\) and \(p^jP\)?  The quotient
\(p^{j-1}P/p^jP\) measures exactly that loss.  Multiplication by \(p\) kills
the quotient, so it is naturally a vector space over \(R/(p)\).  This
motivates the intrinsic layer dimensions
\[
a_j(P)=\dim_{R/(p)}\bigl(p^{j-1}P/p^jP\bigr)
\qquad(j\geq1).
\]
At the first layer, every cyclic \(p\)-power block is still present, so
\(a_1(P)\) counts all blocks. At the second layer only blocks of exponent at
least two survive, so \(a_2(P)\) counts those blocks; at the third layer only
blocks of exponent at least three survive, and so on. The successive layers
therefore record the block heights without choosing a decomposition.

First consider one block \(C=R/(p^4)\).  Its filtration is
\[
C\supseteq pC\supseteq p^2C\supseteq p^3C\supseteq p^4C=0,
\]
and every nonzero successive quotient is one-dimensional over \(R/(p)\).
Its layer sequence is therefore \(1,1,1,1,0,\ldots\).  Similarly,
\(R/(p^2)\) contributes \(1,1,0,\ldots\), while \(R/(p)\) contributes
\(1,0,\ldots\).

For the multi-block example
\[
P=R/(p^4)\oplus R/(p^2)\oplus R/(p^2),
\]
the Young-diagram bookkeeping appears numerically as
\[
\begin{array}{c|ccccc}
 &a_1&a_2&a_3&a_4&a_5\\ \hline
R/(p^4)&1&1&1&1&0\\
R/(p^2)&1&1&0&0&0\\
R/(p^2)&1&1&0&0&0\\ \hline
P&3&3&1&1&0.
\end{array}
\]
Representing each cyclic summand by a column whose height is its exponent,
\(a_j\) counts columns reaching height \(j\), while \(a_j-a_{j+1}\) counts
columns ending there.  The picture reveals the recovery mechanism before the
formal proof.

\begin{lemma}[Primary layers]\label{lem:primary-layers}
If
\[
P\cong\bigoplus_iR/(p^{e_i}),
\]
then
\[
a_j(P)=\#\{i:e_i\geq j\}.
\]
\end{lemma}
\begin{proof}
We separate the four mechanisms in the formula.

\emph{Vector-space structure.}
Because \(p\) is a nonzero prime element of the PID \(R\),
Proposition~\ref{prop:pid-nonzero-prime-maximal} shows that \(R/(p)\) is a
field.
Define \((r+(p))(z+p^jP)=rz+p^jP\) for \(z\in p^{j-1}P\).  This is
well-defined: changing \(z\) changes \(rz\) by an element of \(p^jP\), and
if \(r-r'=ps\), then \((r-r')z\in p^jP\).  Thus the quotient is an
\(R/(p)\)-module, hence a vector space.

\emph{Intrinsic character.}
If \(f:P\to Q\) is an isomorphism, then
\(f(p^kP)=p^kf(P)=p^kQ\) for every \(k\). Hence
\[
\overline f_j:p^{j-1}P/p^jP\longrightarrow p^{j-1}Q/p^jQ,
\qquad z+p^jP\longmapsto f(z)+p^jQ
\]
is well-defined and \(R/(p)\)-linear. The map induced in the same way by
\(f^{-1}\) is its inverse. Thus corresponding layers are isomorphic, and
\(a_j(P)\) depends only on \(P\).

\emph{Additivity.}
For modules \(P_1,P_2\), powers of \(p\) commute with their direct sum. There
is therefore a natural \(R/(p)\)-linear isomorphism
\[
\frac{p^{j-1}(P_1\oplus P_2)}{p^j(P_1\oplus P_2)}
\cong
\frac{p^{j-1}P_1}{p^jP_1}\oplus
\frac{p^{j-1}P_2}{p^jP_2}.
\]
Explicitly, it is given by
\[
(z_1,z_2)+p^j(P_1\oplus P_2)
\longmapsto (z_1+p^jP_1,z_2+p^jP_2)
\]
and its inverse is obtained by pairing representatives. Iterating this binary
isomorphism gives the corresponding statement for every finite direct sum;
Corollary~\ref{cor:finite-direct-sum-dimension} then gives
\[
a_j(P_1\oplus\cdots\oplus P_n)=\sum_{i=1}^n a_j(P_i).
\]

\emph{Contribution of one block.}
For \(R/(p^e)\), the \(j\)-th layer is zero when \(e<j\).  When \(e\geq j\),
it is generated by the class of \(p^{j-1}\), whose annihilator in that layer
is \((p)\); it is one copy of \(R/(p)\).  Additivity proves the formula.
\end{proof}

The decisive recovery formula is
\begin{equation}\label{eq:primary-layer-recovery}
m_j:=a_j(P)-a_{j+1}(P)
=\#\{\text{summands }R/(p^j)\text{ in }P\}.
\end{equation}
It converts intrinsic row lengths of the filtration into column heights,
hence into elementary divisors.

\begin{theorem}[Uniqueness: elementary-divisor form]
\label{thm:uniqueness-elementary-divisor}
For each prime representative \(p\) and each \(j\geq1\), the multiplicity of
\(R/(p^j)\) in a finitely generated PID module is uniquely determined by the
module. Hence the elementary-divisor decomposition is unique up to associates
and order; the free rank is unique as well.
\end{theorem}
\begin{proof}
The intrinsic module \(T(M)\) and
Theorem~\ref{thm:pid-structure-existence}(3) give
\(M/T(M)\cong R^r\), so IBN determines \(r\). Each primary component
\(M_p\) is intrinsic. On it, the intrinsic sequence \(a_j\)
determines through \eqref{eq:primary-layer-recovery} the multiplicity of every
block \(R/(p^j)\). Thus all elementary divisors are unique, prime by prime.
\end{proof}

Elementary divisors determine invariant factors by a canonical regrouping
procedure. Here is that regrouping explicitly. Suppose the elementary divisors are
\[
p,\quad p^3,\quad q^2,\quad q^4,\quad q^5,
\]
for nonassociate primes \(p,q\).  Arrange each prime's powers by increasing
exponent and pad the shorter list on the left by units:
\[
\begin{array}{ccc}
1&p&p^3\\
q^2&q^4&q^5
\end{array}.
\]
Multiplying down the columns gives
\[
q^2,\qquad pq^4,\qquad p^3q^5,
\]
and visibly \(q^2\mid pq^4\mid p^3q^5\).  Factoring an invariant-factor
chain reverses this construction. Thus elementary-divisor uniqueness already
provides one route to invariant-factor uniqueness; after establishing
cancellation, we will also give a direct proof in invariant-factor language.

\section{Cancellation and Consequences}
Cancellation is now a consequence of the intrinsic classification, rather
than an unproved step inside its proof.

\begin{corollary}[Cyclic prime-power cancellation]
Let \(M,N\) be finitely generated \(p\)-primary modules and let \(e\geq1\).
If
\[
M\oplus R/(p^e)\cong N\oplus R/(p^e),
\]
then \(M\cong N\).
\end{corollary}
\begin{proof}
Theorem~\ref{thm:uniqueness-elementary-divisor} says that the two sides have
the same multiplicity of every elementary divisor. Removing the one common
occurrence of \(p^e\) leaves identical elementary-divisor multisets for \(M\)
and \(N\).
\end{proof}

\begin{corollary}[Cyclic cancellation]\label{cor:cyclic-cancellation}
Let \(M,N\) be finitely generated torsion \(R\)-modules and let \(0\ne d\in R\).
If
\[
M\oplus R/(d)\cong N\oplus R/(d),
\]
then \(M\cong N\).
\end{corollary}
\begin{proof}
If \(d\) is a unit, the added cyclic module is zero and there is nothing to
prove. Otherwise factor \(d\), up to a unit, into powers of distinct fixed
prime representatives. The Chinese Remainder Theorem expresses \(R/(d)\) as
the direct sum of the corresponding prime-power cyclic modules.
Theorem~\ref{thm:uniqueness-elementary-divisor} says that the two sides of the
given isomorphism have the same multiset of elementary divisors. Remove the
common elementary divisors contributed by \(R/(d)\). The modules \(M\) and
\(N\) then have identical elementary-divisor decompositions, so \(M\cong N\).
\end{proof}

\begin{theorem}[Uniqueness: invariant-factor form]
\label{thm:uniqueness-invariant-factor}
Suppose a finitely generated torsion \(R\)-module has two decompositions
\[
M\cong R/(a_1)\oplus\cdots\oplus R/(a_r)
\cong R/(b_1)\oplus\cdots\oplus R/(b_s),
\]
where
\[
a_1\mid\cdots\mid a_r,
\qquad b_1\mid\cdots\mid b_s,
\]
and every displayed cyclic factor is nonzero. Then \(r=s\) and, after
matching the two divisibility chains, \((a_i)=(b_i)\) for every \(i\).
\end{theorem}
\begin{proofidea}
The annihilator identifies the largest invariant factor. Cyclic cancellation
removes that factor, and induction repeats the argument on the remaining
chains.
\end{proofidea}
\begin{proof}
We induct on \(r+s\). If \(r=0\), then \(M=0\); because every displayed
factor in the second decomposition is nonzero, necessarily \(s=0\). The same
argument applies with the two decompositions reversed, establishing the base
case.

Now suppose \(r,s>0\). The annihilator of a finite direct sum is the
intersection of the annihilators of its summands, because a scalar
annihilates the direct sum exactly when it annihilates every coordinate. The
divisibility chains give
\[
\operatorname{Ann}(M)
=(a_1)\cap\cdots\cap(a_r)=(a_r),
\qquad
\operatorname{Ann}(M)
=(b_1)\cap\cdots\cap(b_s)=(b_s).
\]
Thus \((a_r)=(b_s)\), so \(R/(a_r)\cong R/(b_s)\). Put
\[
A=\bigoplus_{i<r}R/(a_i),
\qquad
B=\bigoplus_{j<s}R/(b_j).
\]
After identifying the last two factors with the common cyclic type
\(R/(a_r)\), the original isomorphisms give
\[
A\oplus R/(a_r)\cong B\oplus R/(a_r).
\]
Corollary~\ref{cor:cyclic-cancellation} yields \(A\cong B\). The total number
of displayed factors has decreased, so induction gives \(r-1=s-1\) and
\((a_i)=(b_i)\) for \(i<r\). Together with \((a_r)=(b_s)\), this proves
\(r=s\) and equality of all matched ideals.
\end{proof}

The two uniqueness theorems reflect the two normal forms. Primary layers
recover all prime-power block multiplicities simultaneously and prove
Theorem~\ref{thm:uniqueness-elementary-divisor}. In contrast, the direct
invariant-factor argument follows the chain
\[
\text{annihilator}\longrightarrow\text{largest invariant factor}
\longrightarrow\text{cancellation}\longrightarrow\text{induction},
\]
detecting and removing one factor at a time. Consequently, invariant factors
are unique up to associates, while elementary divisors are unique up to
associates and order.

\begin{remark}
The annihilator alone does not single out a preferred embedded cyclic summand
inside an abstract direct sum, so cancellation cannot be assumed. The
primary-layer argument is the natural intrinsic proof of uniqueness in
elementary-divisor form. Once cyclic cancellation has been proved from that
classification, the annihilator argument gives a particularly short inductive
proof of uniqueness in invariant-factor form.
\end{remark}

\section{Finitely Generated Abelian Groups}
The familiar classification of finitely generated abelian groups is now a
corollary, rather than a separate theorem: abelian groups are precisely
\(\mathbb Z\)-modules.

\begin{corollary}[Finitely generated abelian groups]
Every finitely generated abelian group is uniquely isomorphic to
\[
\mathbb Z^r\oplus\mathbb Z/(d_1)\oplus\cdots\oplus\mathbb Z/(d_k),
\qquad 1<d_1\mid\cdots\mid d_k,
\]
or, equivalently, to a direct sum of a free abelian group and cyclic
prime-power groups.
\end{corollary}

Indeed, abelian groups are precisely \(\mathbb Z\)-modules and
\(\mathbb Z\) is a PID, so this is the PID module structure theorem applied
to \(R=\mathbb Z\).  Factoring the invariant factors into prime powers gives
the elementary-divisor form.

\begin{corollary}[Finiteness and order]
A finitely generated abelian group is finite if and only if its free rank is
zero.  In that case,
\[
\left|\mathbb Z/(d_1)\oplus\cdots\oplus\mathbb Z/(d_k)\right|
=d_1\cdots d_k.
\]
\end{corollary}
\begin{proof}
A nonzero free summand \(\mathbb Z\) makes the group infinite.  With no free
summand, the group is a finite Cartesian product of finite cyclic groups, so
its order is the product of their orders.
\end{proof}

\begin{example}[Free and torsion parts]
The group
\[
G=\mathbb Z^2\oplus\mathbb Z/6\mathbb Z\oplus\mathbb Z/30\mathbb Z
\]
has free rank \(2\), torsion subgroup of order \(180\), and is infinite.
Its elementary divisors are
\[
2,\ 3,\ 2,\ 3,\ 5
\]
with prime powers grouped by their two invariant-factor columns; more
transparently,
\[
G_{\mathrm{tors}}\cong
\mathbb Z/2\oplus\mathbb Z/3\oplus
\mathbb Z/2\oplus\mathbb Z/3\oplus\mathbb Z/5.
\]
\end{example}

\begin{example}[Converting the two forms]
\(\mathbb Z/12\mathbb Z\oplus\mathbb Z/18\mathbb Z\) has
elementary divisors \(\mathbb Z/4\mathbb Z\), \(\mathbb Z/3\mathbb Z\),
\(\mathbb Z/2\mathbb Z\), and \(\mathbb Z/9\mathbb Z\); regrouping gives
the invariant factors
\[
\mathbb Z/6\mathbb Z\oplus\mathbb Z/36\mathbb Z.
\]
Its order is \(6\cdot36=216\), the same as \(12\cdot18\).
\end{example}

\begin{example}[Classification at a fixed order]
Every abelian group of order \(72=2^3 3^2\) is obtained by independently
choosing a partition of \(3\) for the \(2\)-primary part and a partition of
\(2\) for the \(3\)-primary part.  The three possibilities
\[
\mathbb Z/8,\qquad \mathbb Z/4\oplus\mathbb Z/2,\qquad
(\mathbb Z/2)^3
\]
for the \(2\)-part combine with the two possibilities
\[
\mathbb Z/9,\qquad(\mathbb Z/3)^2
\]
for the \(3\)-part.  Thus there are exactly six isomorphism types, one for
each pairing of the displayed choices.  For example, pairing the middle and
second choices gives
\[
\mathbb Z/4\oplus\mathbb Z/2\oplus
\mathbb Z/3\oplus\mathbb Z/3.
\]
No search through group operations is needed.
\end{example}

\section{Smith Normal Form}
Theorem~\ref{thm:pid-submodules-compatible-bases}(2) gives the conceptual
existence argument, while the preceding uniqueness theorems identify the
elementary-divisor and invariant-factor data intrinsically. Smith normal form
is the matrix manifestation of the same structure. We first explain why it exists, then
show how B\'ezout row and column operations compute it.

For matrices \(A,B\in M_{n\times m}(R)\), say that \(A\) and \(B\) are
\emph{equivalent} if \(B=UAV\) for some
\(U\in\operatorname{GL}_n(R)\) and \(V\in\operatorname{GL}_m(R)\).
Here \(U\) changes the basis of the target \(R^n\), while \(V\) changes the
basis of the source \(R^m\). This differs fundamentally from similarity,
\(B=P^{-1}AP\), which concerns square matrices representing one endomorphism
after one change of basis. Thus matrix equivalence is not operator similarity.
This distinction matters when Smith reduction is applied to \(xI-A\), whereas
rational canonical form classifies \(A\) up to similarity.

\begin{definition}
A matrix over a PID is in \emph{Smith normal form} if it is
\[\operatorname{diag}(d_1,\ldots,d_t,0,\ldots,0),\qquad
d_1\mid\cdots\mid d_t,\quad d_i\ne0.\]
\end{definition}

The elementary invertible row and column operations are: interchange two
rows or columns; multiply one row or column by a unit; and add a scalar
multiple of one row or column to another. Each is multiplication on the
appropriate side by an elementary matrix. Its inverse is respectively the
same interchange, multiplication by the inverse unit, or addition of the
negative scalar multiple, so each operation preserves matrix equivalence.

\begin{theorem}[Smith normal form]
For every matrix \(A\) over a PID, there are invertible matrices \(U,V\)
such that \(UAV\) is in Smith normal form.
\end{theorem}
\begin{proofidea}
Regard the matrix as a homomorphism \(A:R^m\to R^n\). Choose compatible
target coordinates for \(\operatorname{im}(A)\), then split the surjection from the source
onto that free image.  Compatible source and target bases make the matrix
diagonal.
\end{proofidea}
\begin{proof}
Apply Theorem~\ref{thm:pid-submodules-compatible-bases}(2) to
\(\operatorname{im}(A)\subseteq R^n\).  There is a target basis
\(y_1,\ldots,y_n\) such that
\[
\operatorname{im}(A)=Rd_1y_1\oplus\cdots\oplus Rd_ty_t,
\qquad d_1\mid\cdots\mid d_t.
\]
The displayed generators form a basis of \(\operatorname{im}(A)\).  Because
this module is free, it is projective by Chapter~5, so the surjection
\[
R^m\twoheadrightarrow\operatorname{im}(A)
\]
splits.  Choose a section \(s\), put \(x_i=s(d_iy_i)\), and choose a basis
of \(\ker A\), which is free by Section~\ref{sec:pid-submodules-free}.  The
splitting gives
\[
R^m=s(\operatorname{im}(A))\oplus\ker A,
\]
so \(x_1,\ldots,x_t\), followed by the kernel basis, is a basis of the
source.  In these bases,
\[
A(x_i)=d_iy_i\quad(1\leq i\leq t),
\qquad A(\ker A)=0,
\]
and the matrix is \(\operatorname{diag}(d_1,\ldots,d_t,0,\ldots,0)\).
The two changes of basis give the required invertible \(U,V\).
\end{proof}

\subsection*{Constructive viewpoint: B\'ezout row and column reduction}
The structural proof explains why Smith form exists.  The next argument
explains how to reach it through explicit invertible operations.

\begin{proposition}[Constructive PID diagonal reduction]
Every finite matrix over a PID can be transformed to Smith normal form by
invertible row and column operations.
\end{proposition}
\begin{proofidea}
Improve the upper-left pivot until it divides every entry; prove termination
by ACC on principal ideals; clear its row and column; then induct on the
remaining submatrix.
\end{proofidea}
\begin{proof}
If the matrix is zero, there is nothing to prove.  Otherwise move a nonzero
entry \(d\) to the upper-left corner.  Suppose first that an entry \(b\) in
the first column is not divisible by \(d\).  Let \(g\) generate \((d,b)\),
write \(d=gd'\), \(b=gb'\), and choose \(u,v\in R\) with \(ud+vb=g\).  Then
\[
\begin{pmatrix}u&v\\-b'&d'\end{pmatrix}
\begin{pmatrix}d\\b\end{pmatrix}
=\begin{pmatrix}g\\0\end{pmatrix},
\qquad
\det\begin{pmatrix}u&v\\-b'&d'\end{pmatrix}=ud'+vb'=1.
\]
This invertible row operation replaces the pivot by \(g\).  Analogous right
multiplication on two columns handles a nondivisible first-row entry.

Repeat these B\'ezout improvements until the pivot divides every entry in its
row and column, then clear the rest of that row and column.  The matrix has
the form
\[
\begin{pmatrix}d&0\\0&A_1\end{pmatrix}.
\]
If an interior entry \(b\) of \(A_1\) is not divisible by \(d\), add its row
to the first row.  The first column remains unchanged, while \(b\) enters the
first row.  Apply to the pivot column and the column containing \(b\) the
two-column operation
\[
\begin{pmatrix}u&-b'\\v&d'\end{pmatrix},
\qquad ud+vb=g,\quad d=gd',\quad b=gb'.
\]
Indeed,
\[
(d,b)\begin{pmatrix}u&-b'\\v&d'\end{pmatrix}=(g,0),
\]
and its determinant is \(1\).  The new pivot generates \((d,b)\), and
\(d\nmid b\) implies
\[
(d)\subsetneq(d,b).
\]
Every such improvement strictly enlarges the principal pivot ideal.
Chapter~4 proves that a PID is Noetherian and that ACC terminates ascending
ideal chains, so this process stops.

At termination, clear the first row and column once more.  No interior entry
can fail to be divisible by the pivot, or the preceding step would improve it
again.  Thus the first diagonal entry \(d_1\) divides every entry of the
remaining matrix.  Apply induction to that submatrix.  Its operations form
only linear combinations of its entries, so divisibility by \(d_1\) survives,
and the resulting diagonal satisfies
\[
d_1\mid d_2\mid\cdots\mid d_t.
\]
Every operation used is invertible; no Euclidean division is required.
\end{proof}

Over an arbitrary PID this is an explicit terminating reduction procedure,
with termination supplied by ACC. Effectiveness additionally requires the
ability to compute gcds and B\'ezout coefficients. Over Euclidean domains such
as \(\mathbb Z\) and \(F[x]\), the Euclidean algorithm supplies these, so the
procedure is a computational algorithm in the usual sense.

Presentations explain what this computation does to a module.  Chosen
generators and relations give an exact sequence
\[
\begin{tikzcd}[column sep=large]
R^m\arrow[r,"A"]&R^n\arrow[r,two heads]&M\arrow[r]&0,
\end{tikzcd}
\]
where the columns of \(A\) record the relations and
\(M\cong\operatorname{coker}(A)\).  Row operations change the target basis,
column operations change the source basis, and neither changes the cokernel.
Smith reduction therefore separates the defining relations coordinatewise.

\begin{corollary}[Kernel and cokernel from Smith form]
\label{cor:smith-kernel-cokernel}
Suppose \(A:R^m\to R^n\) has Smith form
\(\operatorname{diag}(d_1,\ldots,d_t,0,\ldots,0)\). Then
\[
\operatorname{rank}(\operatorname{im}A)=t,
\qquad \ker A\cong R^{m-t},
\]
and
\[
\operatorname{coker}A\cong
R/(d_1)\oplus\cdots\oplus R/(d_t)\oplus R^{n-t}.
\]
Unit \(d_i\) contribute zero cyclic summands and may be omitted. Thus Smith
reduction diagonalizes the relations in a finite presentation.
\end{corollary}
\begin{proof}
Equivalent matrices represent the same map after source and target
isomorphisms. For the diagonal map the assertions follow coordinatewise;
multiplication by each nonzero \(d_i\) is injective because a PID is a domain.
\end{proof}

Over a field every nonzero Smith entry is a unit and can be normalized to
\(1\). Thus Smith form becomes
\(\operatorname{diag}(1,\ldots,1,0,\ldots,0)\): ordinary rank normal form is
the field case of the divisibility data carried by Smith form over a PID.

If \(A\) is \(n\times n\) and has full rank, determinant multiplicativity
gives \(\det A\sim d_1\cdots d_n\). For an integral matrix with
\(\det A\ne0\), the cokernel formula therefore gives
\[
\lvert\mathbb Z^n/A\mathbb Z^n\rvert=\lvert\det A\rvert.
\]

\noindent\textbf{Worked example.} Over \(\mathbb Z\), elementary operations
give
\[
\begin{pmatrix}2&4\\6&8\end{pmatrix}
\xrightarrow{R_2\mapsto R_2-3R_1}
\begin{pmatrix}2&4\\0&-4\end{pmatrix}
\xrightarrow{C_2\mapsto C_2-2C_1}
\begin{pmatrix}2&0\\0&-4\end{pmatrix}.
\]
Multiplying the second row by \(-1\) yields
\[
\operatorname{diag}(2,4).
\]
Consequently its cokernel is
\[
\mathbb Z/2\mathbb Z\oplus\mathbb Z/4\mathbb Z.
\]
Equivalently, the abelian group with two generators and relation columns
\((2,6)^t\) and \((4,8)^t\) has precisely these invariant factors.  This is
the promised passage from a presentation matrix to the classification of its
cokernel.

\noindent\textbf{A genuine B\'ezout pivot improvement.}
For the column matrix \(A=(6,15)^t\), the pivot \(6\) does not divide \(15\).
Since \((-2)6+15=3\),
\[
\begin{pmatrix}-2&1\\-5&2\end{pmatrix}
\begin{pmatrix}6\\15\end{pmatrix}
=\begin{pmatrix}3\\0\end{pmatrix},
\qquad
\det\begin{pmatrix}-2&1\\-5&2\end{pmatrix}=1.
\]
This invertible row operation replaces the pivot by
\(\gcd(6,15)=3\), displaying the B\'ezout improvement that repeated
subtraction can conceal.

For a matrix, let \(D_j(A)\) be the ideal generated by its \(j\times j\)
minors, with \(D_j(A)=R\) for \(j\leq0\) and \(D_j(A)=0\) when \(j\)
exceeds the possible minor size. Elementary row and column operations preserve
each \(D_j(A)\), because each new minor is an \(R\)-linear combination of old
minors, and the inverse operation gives the reverse inclusion.  In Smith form
with \(t\) nonzero entries,
\[
D_j(A)=(d_1\cdots d_j)\quad(1\leq j\leq t),
\qquad D_j(A)=0\quad(j>t).
\]
If \(D_j(A)=(\Delta_j)\), then, up to associates,
\(d_1\sim\Delta_1\) and
\(\Delta_j\sim\Delta_{j-1}d_j\) successively. Concretely, \(d_1\) is
associate to the gcd of all entries, \(d_1d_2\) to the gcd of all
\(2\times2\) minors, and so forth.

\begin{corollary}[Uniqueness of Smith normal form]
\label{cor:smith-uniqueness}
Suppose
\[
\operatorname{diag}(d_1,\ldots,d_t,0,\ldots,0),\qquad
\operatorname{diag}(e_1,\ldots,e_s,0,\ldots,0)
\]
are Smith forms equivalent to the same matrix, with both nonzero lists ordered
by divisibility. Then \(t=s\) and \(d_i\sim e_i\) for every \(i\).
With fixed associate representatives, such as positive integers in
\(\mathbb Z\) or monic polynomials in \(F[x]\), the Smith form is literally
unique.
\end{corollary}
\begin{proof}
Equivalent matrices have the same determinantal ideals. The integer \(t\) is
the largest \(j\geq0\) for which \(D_j(A)\ne0\), with \(D_0(A)=R\) by
convention. Hence equivalent matrices have the same \(t\). Then
\((d_1\cdots d_j)=(e_1\cdots e_j)\) for \(1\leq j\leq t\); cancellation in the
domain, applied successively, gives \((d_j)=(e_j)\).
\end{proof}

\subsection*{Alternative uniqueness via Smith normal form and Fitting ideals}
Smith uniqueness compares equivalent matrices. An abstract module, however,
can have many finite presentations, so another step is required before minors
become module invariants.

Let
\[
R^m\xrightarrow{A}R^n\longrightarrow M\longrightarrow0
\]
be a finite presentation. For every integer \(k\), define
\[
\operatorname{Fitt}_k(M)=D_{n-k}(A),
\]
using \(D_j(A)=R\) for \(j\leq0\) and \(D_j(A)=0\) beyond the possible
minor size.

\begin{theorem}[Presentation independence of Fitting ideals]
\label{thm:fitting-presentation-independent}
The ideal \(\operatorname{Fitt}_k(M)\) is independent of the chosen finite
presentation. Hence every Fitting ideal is intrinsic to \(M\).
\end{theorem}
\begin{proof}
We compare separately the relations and the module generators.

First fix a surjection \(R^n\twoheadrightarrow M\). If the columns of two
matrices \(A\) and \(A'\) generate its kernel, each column of either matrix is
a linear combination of columns of the other. Thus \(A'=AC\) and \(A=A'C'\)
for suitable matrices. Fix a \(j\times j\) minor of \(AC\). After selecting
its \(j\) rows, each of its columns is an \(R\)-linear combination of the
columns of \(A\) restricted to those rows. Multilinearity expands its
determinant as an \(R\)-linear combination of determinants formed from \(j\)
restricted columns of \(A\). Terms with a repeated column vanish; every
remaining determinant is, up to sign, a \(j\times j\) minor of \(A\).
Therefore \(D_j(AC)\subseteq D_j(A)\), so
\(D_j(A')\subseteq D_j(A)\). The reverse factorization gives equality, and
the ideals do not depend on the chosen finite generating list of relations.

Now adjoin to module generators \(y_1,\ldots,y_n\) one redundant generator
\(y=\sum c_i y_i\). If \(c=(c_1,\ldots,c_n)^t\), a relation matrix for the
enlarged list is
\[
B=\begin{pmatrix}A&-c\\0&1\end{pmatrix}.
\]
Indeed, \((u,s)\) is a relation on the enlarged list exactly when
\(u+sc\in\operatorname{im}A\), which is exactly when
\((u,s)=B(z,s)\) for some \(z\).
Adding \(c_i\) times the last row to row \(i\) makes \(B\) equivalent to
\(A\oplus(1)\). A Laplace expansion of any \(j\times j\) minor of \(A\)
along one row expresses it as an \(R\)-linear combination of
\((j-1)\times(j-1)\) minors of \(A\). Hence
\(D_j(A)\subseteq D_{j-1}(A)\). A \(j\)-minor of the block matrix either avoids the unit
block or uses it, so
\[
D_j(A\oplus(1))=D_j(A)+D_{j-1}(A)=D_{j-1}(A),
\]
because \(D_j(A)\subseteq D_{j-1}(A)\). Therefore
\[
D_{(n+1)-k}(B)=D_{n-k}(A),
\]
so adjoining a redundant module generator does not change
\(\operatorname{Fitt}_k\).

Finally, take two finite generating lists for \(M\) and adjoin every member
of each list to the other. Both become the same combined list, one redundant
generator at a time. For that common surjection, the first paragraph compares
any two finite relation lists. Such lists exist because its kernel is a
submodule of a finite free module over the PID \(R\), hence is finitely
generated by Theorem~\ref{thm:pid-submodules-compatible-bases}. Thus the two
original presentations give the same ideals.
\end{proof}

\begin{theorem}[Invariant factors from Fitting ideals]
\label{thm:fitting-invariant-factors}
Suppose existence gives
\[
M\cong R^r\oplus R/(d_1)\oplus\cdots\oplus R/(d_t),
\qquad d_1\mid\cdots\mid d_t,
\]
where the \(d_i\) are nonzero nonunits. Then the free rank and all associate
classes \((d_i)\) are determined by \(M\).
\end{theorem}
\begin{proof}
The free rank \(r\) is intrinsic because
\(M/T(M)\cong R^r\) and IBN determines \(r\). A diagonal presentation with
\(r+t\) generators has relation matrix consisting of a zero \(r\)-row block
above \(\operatorname{diag}(d_1,\ldots,d_t)\). Consequently
\[
\begin{aligned}
\operatorname{Fitt}_k(M)&=0 &&(k<r),\\
\operatorname{Fitt}_{r+t-j}(M)&=(d_1\cdots d_j)
&&(1\leq j\leq t),\\
\operatorname{Fitt}_{r+t}(M)&=R.
\end{aligned}
\]
Because every \(d_i\) is a nonunit, \(r+t\) is the least index at which the
Fitting ideal becomes \(R\); hence \(t\) is intrinsic. Reading the middle
line from \(j=1\) upward gives
\[
(d_1),\quad(d_1d_2),\quad\ldots,\quad(d_1\cdots d_t).
\]
Successive cancellation recovers each \((d_j)\). This proves invariant-factor
uniqueness without choosing or identifying a summand of \(M\). Factoring the
recovered invariant factors into prime powers then recovers the elementary
divisors as well.
\end{proof}

The three uniqueness arguments illuminate different structure. Primary
layers use the intrinsic \(p\)-primary filtration and recover all prime-power
block multiplicities at once; this remains the conceptual main proof.
Annihilator plus cancellation detects the largest invariant factor, removes
it, and inducts. Smith form and Fitting ideals pass from a finite presentation
to minors, prove those ideals presentation-independent, and recover invariant
factors. These are complementary rather than circular: the Fitting argument
uses only existence of a decomposition, not its uniqueness, and the Smith
uniqueness corollary rests directly on determinantal ideals. Fitting ideals
also make sense beyond PIDs, where they help measure how many generators a
module needs locally.

\section{Linear Operators as \texorpdfstring{\(F[x]\)}{F[x]}-Modules}
\label{sec:linear-operators-fx-modules}
Chapter~5 showed what generalized eigenvectors and Jordan chains look like.
It did not yet explain why enough compatible chains exist or why their block
sizes are intrinsic. Why do the canonical blocks exist, and why are their
sizes unique? The structural answer is to let \(x\) act as \(T\). Then \(V\)
becomes a finitely generated torsion \(F[x]\)-module, and the PID structure
theorem answers both questions: the blocks are its cyclic summands, whose
types are intrinsic.

Let \(T:V\to V\) be linear over a field \(F\). Define
\[p(x)\cdot v=p(T)v,\qquad
p(T)=a_0I+a_1T+\cdots+a_nT^n.\]
Because \((pq)(T)=p(T)q(T)\), the module axioms hold.
Write \(V_T\) when the operator defining this \(F[x]\)-module structure
needs to be displayed.

Recall from Section~\ref{sec:matrices-change-basis} the direct sums of
operators and their block-diagonal matrices. Direct calculation on pairs
gives, for every \(p\in F[x]\),
\[
p(T_1\oplus T_2)=p(T_1)\oplus p(T_2),
\]
and the identity map on the underlying direct sum is an \(F[x]\)-module
isomorphism
\[
(V_1\oplus V_2)_{T_1\oplus T_2}
\cong (V_1)_{T_1}\oplus(V_2)_{T_2}.
\]

\begin{proposition}
A subspace \(W\subseteq V\) is \(T\)-invariant if and only if it is an
\(F[x]\)-submodule for this action.
\end{proposition}
\begin{proof}
If \(W\) is a submodule, then \(Tw=x\cdot w\in W\). Conversely, if
\(T(W)\subseteq W\), every polynomial in \(T\) preserves \(W\), so \(W\)
is an \(F[x]\)-submodule.
\end{proof}

\begin{proposition}[Intertwiners are module maps]\label{prop:intertwiner-module-map}
Let \(T:V\to V\) and \(S:W\to W\).  An \(F\)-linear map
\(\phi:V\to W\) is \(F[x]\)-linear from \(V_T\) to \(W_S\) if and only if
\[
\phi T=S\phi.
\]
Consequently \(V_T\cong W_S\) as \(F[x]\)-modules if and only if there is
an invertible intertwining map between \(T\) and \(S\).  In coordinates this
is precisely similarity of their matrices.
\end{proposition}
\begin{proof}
If \(\phi\) is \(F[x]\)-linear, then
\[
\phi(Tv)=\phi(x\cdot v)=x\cdot\phi(v)=S\phi(v).
\]
Conversely, \(\phi T=S\phi\) implies \(\phi T^j=S^j\phi\) by induction,
and linear combinations give
\(\phi p(T)=p(S)\phi\) for every polynomial \(p\).  Thus \(\phi\) is
\(F[x]\)-linear.  If \(\phi\) is invertible and its matrix is \(P\), the
intertwining equation becomes \(PA=BP\), or \(B=PAP^{-1}\), exactly the
change-of-basis relation.
\end{proof}

\begin{proposition}[Finite-dimensional operator modules are torsion]
\label{prop:operator-module-torsion}
If \(V\) is finite-dimensional, then \(V_T\) is a finitely generated torsion
\(F[x]\)-module.
\end{proposition}
\begin{proof}
An \(F\)-basis \(v_1,\ldots,v_n\) generates \(V_T\), because constants in
\(F[x]\) already produce every \(F\)-linear combination.  For each \(v_i\),
the \(n+1\) vectors
\[
v_i,Tv_i,\ldots,T^nv_i
\]
are linearly dependent, so a nonzero polynomial \(f_i\) satisfies
\(f_i(T)v_i=0\).  The nonzero product \(f_1\cdots f_n\) annihilates every
basis vector and hence all of \(V\).  In particular every vector is torsion.
\end{proof}

For a monic polynomial
\[
d(x)=x^m+a_{m-1}x^{m-1}+\cdots+a_0
\]
define its \emph{companion matrix} by
\[
C_d=
\begin{pmatrix}
0&0&\cdots&0&-a_0\\
1&0&\cdots&0&-a_1\\
0&1&\cdots&0&-a_2\\
\vdots&&\ddots&\vdots&\vdots\\
0&0&\cdots&1&-a_{m-1}
\end{pmatrix}.
\]
The next proposition derives this matrix directly from the first polynomial
relation in an operator orbit.

\begin{proposition}[Cyclic subspaces, companion matrices, and annihilators]
\label{prop:cyclic-subspace-annihilator}
Let \(0\ne v\in V\), and let
\[
W=F[x]v=\operatorname{span}_F\{v,Tv,T^2v,\ldots\}.
\]
If \(\dim W=m\), then
\[
\beta=(v,Tv,\ldots,T^{m-1}v)
\]
is a basis of \(W\). There is a unique monic polynomial
\[
d(x)=x^m+a_{m-1}x^{m-1}+\cdots+a_1x+a_0
\]
generating \(\operatorname{Ann}_{F[x]}(v)\), and
\[
F[x]v\cong F[x]/(d),
\qquad
[T|_W]_\beta=C_d.
\]
\end{proposition}
\begin{proof}
Let \(k\) be the least index for which
\(v,Tv,\ldots,T^kv\) are linearly dependent. Then the first \(k\) vectors
are independent, and the first dependence expresses \(T^kv\) in their span.
Applying \(T\) repeatedly shows that every later orbit vector remains in that
span. Thus those \(k\) vectors form a basis of \(W\), so \(k=\dim W=m\).
Equivalently, \(v,Tv,\ldots,T^{m-1}v\) are independent, and adjoining
\(T^mv\) produces the first dependence.

Consequently there are scalars \(c_0,\ldots,c_m\in F\) such that
\[
c_mT^mv+c_{m-1}T^{m-1}v+\cdots+c_1Tv+c_0v=0.
\]
Here \(c_m\ne0\): otherwise this would already be a dependence among the
first \(m\) orbit vectors. Because \(F\) is a field, \(c_m\) is invertible.
We intentionally divide the entire relation by \(c_m\), normalizing its
leading coefficient to \(1\). With \(a_i=c_i/c_m\), the relation becomes
\[
T^mv+a_{m-1}T^{m-1}v+\cdots+a_1Tv+a_0v=0,
\]
or equivalently
\[
T^mv=-a_0v-a_1Tv-\cdots-a_{m-1}T^{m-1}v.
\]
This deliberate normalization makes
\(d(x)=x^m+a_{m-1}x^{m-1}+\cdots+a_1x+a_0\) monic, matching the convention
for generators of ideals in \(F[x]\).

On the orbit basis, the action is
\[
T(v)=Tv,\quad T(Tv)=T^2v,\quad\ldots,\quad
T(T^{m-2}v)=T^{m-1}v,
\]
while the last basis vector satisfies
\[
T(T^{m-1}v)=-a_0v-a_1Tv-\cdots-a_{m-1}T^{m-1}v.
\]
Reading these coordinate columns gives
\[
[T|_W]_\beta=C_d=
\begin{pmatrix}
0&0&\cdots&0&-a_0\\
1&0&\cdots&0&-a_1\\
0&1&\cdots&0&-a_2\\
\vdots&&\ddots&\vdots&\vdots\\
0&0&\cdots&1&-a_{m-1}
\end{pmatrix}.
\]
Thus a companion matrix is precisely the matrix of \(T\) on a cyclic
\(T\)-invariant subspace in its natural orbit basis.

The normalized relation says \(d(T)v=0\), so
\(d\in\operatorname{Ann}_{F[x]}(v)\). No nonzero polynomial of degree less
than \(m\) can annihilate \(v\), since it would give a dependence among
\(v,Tv,\ldots,T^{m-1}v\). If \(d_v\) is the monic generator of
\(\operatorname{Ann}(v)\), then \(d_v\mid d\), while the preceding
independence gives \(\deg d_v\geq m\). Hence \(d_v=d\) and
\(\operatorname{Ann}(v)=(d)\).

Finally, the surjective module map
\[
F[x]\longrightarrow F[x]v,\qquad p\longmapsto p(T)v,
\]
has kernel \((d)\). The First Isomorphism Theorem gives
\(F[x]v\cong F[x]/(d)\), making the concrete orbit calculation and the
abstract cyclic-module description identical.
\end{proof}

\begin{theorem}[The \(xI-A\) presentation]
\label{thm:operator-cokernel-presentation}
Let \(V\) have ordered basis \(\beta=(v_1,\ldots,v_n)\), let
\(T\in\mathcal L(V)\), and put \(A=[T]_\beta\). Then, as \(F[x]\)-modules,
\[
V_T\cong\operatorname{coker}\bigl(xI_n-A:F[x]^n\to F[x]^n\bigr).
\]
If the Smith form over \(F[x]\) is
\[
\operatorname{diag}(1,\ldots,1,d_1,\ldots,d_r),
\qquad d_1\mid\cdots\mid d_r,
\]
with the \(d_i\) monic nonunits, then they are precisely the invariant factors
of \(V_T\).
\end{theorem}
\begin{proof}
Send the standard basis vector \(e_i\) of \(F[x]^n\) to \(v_i\). This gives
a surjective \(F[x]\)-linear map to \(V_T\). If \(A=(a_{ij})\), column \(j\)
of \(xI_n-A\) maps to
\[
xv_j-\sum_i a_{ij}v_i=Tv_j-Tv_j=0,
\]
so the map factors through the cokernel. Conversely, send
\(\sum_i c_iv_i\) to the class of the constant vector \((c_i)\). In the
cokernel, the displayed column relations say
\(x[e_j]=\sum_i a_{ij}[e_i]\), so this inverse candidate respects the action
of \(x\), hence every polynomial. The two induced maps are inverse on the
displayed generators.

Smith equivalence preserves the cokernel. Corollary~\ref{cor:smith-kernel-cokernel}
therefore identifies it with the direct sum of the \(F[x]/(d_i)\); unit
entries contribute zero modules. Also
\(\det(xI_n-A)=\chi_T(x)\) is a nonzero monic polynomial, so \(xI_n-A\)
has full rank and its Smith form has no zero diagonal entries. The remaining
nonunit entries are exactly the invariant factors.
\end{proof}

Thus
\[
\operatorname{SNF}_{F[x]}(xI-A)
\longrightarrow\text{invariant factors of }V_T
\longrightarrow\text{rational canonical form}.
\]
The matrix \(xI-A\) is a matrix over \(F[x]\), and its Smith equivalence is
the computational bridge from presentations to similarity classification.

The cyclic-subspace proposition is the first half of the canonical-form dictionary:
a cyclic module \(F[x]/(d)\) is a companion block.  The special case
\(d=(x-\lambda)^k\) is the Jordan-chain block described below.

\[
\begin{array}{c|c}
\text{Linear algebra}&F[x]\text{-module theory}\\ \hline
T\text{-invariant subspace}&F[x]\text{-submodule}\\
\text{cyclic \(T\)-subspace generated by }v&F[x]v\\
p(T)v=0&p(x)\text{ annihilates }v\\
\text{similarity of operators}&F[x]\text{-module isomorphism}\\
\text{generalized eigenspace for }\lambda&(x-\lambda)\text{-primary component}\\
\text{companion block }C_d&F[x]/(d)\\
\text{Jordan block }J_k(\lambda)&F[x]/((x-\lambda)^k)
\end{array}
\]
This dictionary synthesizes the correspondences proved in this section and
turns the linear-algebra questions into the PID classification problem.

\begin{definition}
For an irreducible polynomial \(p\in F[x]\), the \emph{\(p\)-primary
component} of \(V_T\) is
\[
V_p=\{v\in V:p(T)^Nv=0\text{ for some }N\geq1\}.
\]
We also write this component as \((V_T)_p\) when the ambient operator module
needs to be emphasized.
\end{definition}

\begin{proposition}[Primary components and generalized eigenspaces]
\label{prop:primary-generalized-eigenspace}
For \(p(x)=x-\lambda\),
\[
V_p=G_\lambda(T).
\]
Thus the concrete generalized eigenspace of Chapter~5 is exactly the
\((x-\lambda)\)-primary component of the operator module.
\end{proposition}
\begin{proof}
The defining condition \(p(T)^Nv=0\) is precisely
\((T-\lambda I)^Nv=0\).  Taking all such vectors gives the same set in the
two definitions.
\end{proof}

If \(p_1,\ldots,p_s\) are the distinct monic irreducibles occurring in the
elementary divisors of \(V_T\), the PID primary decomposition gives
\[
V_T=(V_T)_{p_1}\oplus\cdots\oplus(V_T)_{p_s}.
\]

\begin{proposition}[Jordan-chain dictionary]
\label{prop:jordan-chain-cyclic-module}
A Jordan chain \(v_1,\ldots,v_k\) for \(\lambda\) spans a cyclic
\(F[x]\)-module isomorphic to
\[
F[x]/((x-\lambda)^k).
\]
Conversely, every such cyclic summand supplies a Jordan chain of length
\(k\).
\end{proposition}
\begin{proof}
The last vector \(v_k\) generates the chain because
\[
(x-\lambda)^{k-j}\cdot v_k=v_j.
\]
It is killed by \((x-\lambda)^k\), but not by a smaller positive power,
so if \(\operatorname{Ann}(v_k)=(d)\) with \(d\) monic, then
\(d\mid(x-\lambda)^k\). Hence \(d=(x-\lambda)^j\) for some \(j\leq k\).
No smaller positive power kills \(v_k\), so \(j=k\) and
\(\operatorname{Ann}(v_k)=((x-\lambda)^k)\). Conversely, in
\(F[x]/((x-\lambda)^k)\), use the ordered basis
\[
(x-\lambda)^{k-1},(x-\lambda)^{k-2},\ldots,x-\lambda,1.
\]
Multiplication by \(x-\lambda\) sends each basis vector to the preceding
one and kills the first.  Multiplication by \(x\) therefore has matrix
\(J_k(\lambda)\), and the basis is a Jordan chain in the convention of
Chapter~5.
\end{proof}

\section{Minimal Polynomial and Cayley--Hamilton}
The characteristic polynomial introduced in Chapter~5 detects eigenvalues.
The minimal polynomial records the smallest polynomial relation satisfied by
the whole operator and will control the sizes of canonical-form blocks.  In
module language it is the global relation shared by every vector, whereas
\(\operatorname{Ann}(v)\) records only the relation on one cyclic subspace.

\begin{lemma}
If \(V\) is finite-dimensional, then
\(\operatorname{Ann}_{F[x]}(V)=\{p\in F[x]:p(T)=0\}\) is nonzero.
\end{lemma}
\begin{proof}
The proof of Proposition~\ref{prop:operator-module-torsion} constructs, from
an \(F\)-basis \(v_1,\ldots,v_n\), nonzero polynomials \(f_i\) that
annihilate the individual basis vectors.  Their product annihilates every
basis vector and hence all of \(V\).
\end{proof}

Since \(F[x]\) is a PID, this annihilator has a unique monic generator,
called the \emph{minimal polynomial} \(m_T\):
\[
\operatorname{Ann}_{F[x]}(V_T)=(m_T).
\]
It is the monic polynomial of least degree satisfying \(m_T(T)=0\), and
every polynomial annihilating \(T\) is a multiple of it.

For \(v\in V\), write \(\operatorname{Ann}_{F[x]}(v)=(d_v)\) with \(d_v\)
monic. Since every global annihilator kills \(v\),
\((m_T)\subseteq(d_v)\), equivalently \(d_v\mid m_T\). If
\(v_1,\ldots,v_n\) is an \(F\)-basis and
\(\operatorname{Ann}(v_i)=(d_i)\), then a polynomial kills \(V\) exactly
when it kills every basis vector. Therefore
\[
\operatorname{Ann}(V_T)=\bigcap_{i=1}^n\operatorname{Ann}(v_i),
\qquad
m_T=\operatorname{lcm}(d_1,\ldots,d_n),
\]
with monic representatives. The minimal polynomial is thus the least common
polynomial relation satisfied simultaneously in every direction of \(V\).

\begin{proposition}[Eigenvalues and the minimal polynomial]
\label{prop:eigenvalue-minimal-polynomial}
For \(\lambda\in F\),
\[
\lambda\text{ is an eigenvalue of }T
\quad\Longleftrightarrow\quad m_T(\lambda)=0
\quad\Longleftrightarrow\quad (x-\lambda)\mid m_T.
\]
\end{proposition}
\begin{proof}
If \(Tv=\lambda v\) for some \(v\ne0\), then
\(0=m_T(T)v=m_T(\lambda)v\), so \(m_T(\lambda)=0\). The factor theorem gives
the second equivalence in one direction. Conversely, write
\(m_T=(x-\lambda)q\). If \(T-\lambda I\) were invertible, then
\[
0=m_T(T)=(T-\lambda I)q(T)
\]
would imply \(q(T)=0\), contradicting the minimal degree of \(m_T\). Thus
\(T-\lambda I\) is singular and \(\lambda\) is an eigenvalue.
\end{proof}

\begin{proposition}[Polynomial invariants of a direct sum]
\label{prop:direct-sum-polynomials}
Let \(T_i\in\mathcal L(V_i)\), \(i=1,2\), where the \(V_i\) are
finite-dimensional. Then
\[
m_{T_1\oplus T_2}
=\operatorname{lcm}(m_{T_1},m_{T_2}),
\qquad
\chi_{T_1\oplus T_2}=\chi_{T_1}\chi_{T_2},
\]
where the least common multiple is chosen monic.
\end{proposition}
\begin{proof}
A polynomial annihilates \(T_1\oplus T_2\) exactly when it annihilates both
summands.  Hence
\[
(m_{T_1\oplus T_2})
=(m_{T_1})\cap(m_{T_2}),
\]
whose monic generator in the PID \(F[x]\) is the least common multiple.
This also follows directly from
\(p(T_1\oplus T_2)=p(T_1)\oplus p(T_2)\). Choose bases and write
\(A_i=[T_i]\). In the concatenated basis,
\([T_1\oplus T_2]=[T_1]\oplus[T_2]\). Therefore
Chapter~5's block-determinant formula gives
\[
\det(xI-(A_1\oplus A_2))
=\det(xI-A_1)\det(xI-A_2).
\]
\end{proof}

\begin{theorem}[Cayley--Hamilton]\label{thm:cayley-hamilton}
Every linear operator satisfies its characteristic polynomial:
\[\chi_T(T)=0.\]
\end{theorem}
\begin{proofidea}
The matrix \(TI-A\) records the relations describing \(T\) on a basis.
Multiplying this relation matrix by its adjugate collapses it to its
determinant, forcing \(\chi_T(T)\) to kill every basis vector.
\end{proofidea}
\begin{proof}
Choose a basis \(\beta=(v_1,\ldots,v_n)\), write
\(A=(a_{ij})=[T]_\beta\), and regard the basis as the row vector
\(\mathbf v=(v_1,\ldots,v_n)\). For a matrix of operators \(Q=(q_{ij})\),
the \(j\)-th entry of \(\mathbf vQ\) means
\(\sum_i q_{ij}(v_i)\). Since column \(j\) of \(A\) gives the coordinates of
\(Tv_j\), the coordinate relations are exactly
\begin{equation}\label{eq:cayley-hamilton-relation-matrix}
\mathbf v(TI_n-A)=0.
\end{equation}

Every entry of \(TI_n-A\) lies in the subring
\[
F[T]=\{p(T):p\in F[x]\}\subseteq\operatorname{End}_F(V).
\]
Here each scalar entry \(a_{ij}\in F\) is identified with the scalar operator
\(a_{ij}I_V\in F[T]\), so every entry of \(TI_n-A\) genuinely belongs to
this one coefficient ring.
This ring is commutative because all its elements are polynomials in the same
operator. Determinants and adjugates therefore obey their usual identities
over \(F[T]\):
\[
(TI_n-A)\operatorname{adj}(TI_n-A)
=\det(TI_n-A)I_n.
\]
Multiplying \eqref{eq:cayley-hamilton-relation-matrix} on the right gives
\[
0=\mathbf v(TI_n-A)\operatorname{adj}(TI_n-A)
=\mathbf v\det(TI_n-A).
\]
Thus \(\det(TI_n-A)\) kills every basis vector.

It remains to identify this determinant. By definition,
\(\chi_T(x)=\det(xI_n-A)\). The evaluation homomorphism
\[
F[x]\longrightarrow F[T],\qquad p(x)\longmapsto p(T),
\]
takes \(\det(xI_n-A)\) to \(\det(TI_n-A)\), because a determinant is a
finite sum of products of entries. Hence
\[
\det(TI_n-A)=\chi_T(T).
\]
This operator kills every basis vector and therefore all of \(V\), proving
\(\chi_T(T)=0\).
\end{proof}

\begin{corollary}[Minimal polynomial divides characteristic polynomial]
\label{cor:minimal-divides-characteristic}
For every endomorphism of a finite-dimensional vector space,
\[
m_T\mid\chi_T,
\qquad \deg m_T\leq\dim V.
\]
\end{corollary}
\begin{proof}
Cayley--Hamilton says \(\chi_T\in\operatorname{Ann}(V_T)=(m_T)\), which is
exactly the divisibility assertion. Since \(\deg\chi_T=\dim V\), the degree
bound follows. Equivalently, polynomial division
\(\chi_T=qm_T+r\), \(\deg r<\deg m_T\), and evaluation at \(T\) give
\(r(T)=0\); minimality forces \(r=0\).
\end{proof}

If \(\chi_T\) splits over \(F\), then
Corollary~\ref{cor:minimal-divides-characteristic} shows that \(m_T\), as a
divisor of a split polynomial, also splits into linear factors. Every
irreducible \(p_i\) occurring in the operator module divides its annihilator
generator \(m_T\), so \(p_i=x-\lambda_i\) for some \(\lambda_i\in F\).
Proposition~\ref{prop:eigenvalue-minimal-polynomial} identifies the distinct
\(\lambda_i\) as the eigenvalues of \(T\), and
Proposition~\ref{prop:primary-generalized-eigenspace} turns the primary
decomposition from Section~\ref{sec:linear-operators-fx-modules} into
\[
V=G_{\lambda_1}(T)\oplus\cdots\oplus G_{\lambda_s}(T).
\]
This proves the generalized-eigenspace decomposition promised in Chapter~5,
using only results already established.

\section{Rational Canonical Form}
Rational canonical form exists over every field; it is the field-independent
classification of linear operators. Jordan form is available only when the
relevant irreducible polynomials split into linear factors over the base
field. Thus Jordan form is a split-polynomial refinement of rational form.

Recall from Section~\ref{sec:linear-operators-fx-modules} that the companion
matrix \(C_d\) is multiplication by \(x\) on the cyclic module
\(F[x]/(d)\) in the basis \(1,x,\ldots,x^{n-1}\). It is not an artificial
matrix recipe: it is the coordinate form of one generator subject to the
single polynomial relation \(d(x)=0\).

\begin{theorem}[Characteristic polynomial of a companion matrix]
\label{thm:companion-characteristic}
For a monic polynomial
\[
d(x)=x^n+a_{n-1}x^{n-1}+\cdots+a_0,
\]
the companion matrix \(C_d\) has characteristic polynomial
\[
\chi_{C_d}(x)=d(x).
\]
\end{theorem}
\begin{proof}
Write \(D_n(x)=\det(xI-C_d)\). Explicitly,
\[
\begin{pmatrix}
x&0&\cdots&0&a_0\\
-1&x&\cdots&0&a_1\\
0&-1&\ddots&0&a_2\\
\vdots&\ddots&\ddots&x&\vdots\\
0&\cdots&0&-1&x+a_{n-1}
\end{pmatrix}.
\]
Expand along the first row. Only the entries \(x\) and \(a_0\) contribute.
Deleting the first row and column leaves the companion-type determinant for
\[
x^{n-1}+a_{n-1}x^{n-2}+\cdots+a_2x+a_1;
\]
call it \(D_{n-1}(x)\). Hence the first contribution is
\(xD_{n-1}(x)\). The minor of \(a_0\) is triangular with diagonal entries
\(-1\), so its determinant is \((-1)^{n-1}\). Its cofactor sign is
\((-1)^{1+n}\), and
\[
(-1)^{1+n}(-1)^{n-1}=1.
\]
Thus
\[
D_n(x)=xD_{n-1}(x)+a_0.
\]
The case \(n=1\) is \(D_1(x)=x+a_0\). Induction now gives
\[
D_n(x)=x^n+a_{n-1}x^{n-1}+\cdots+a_1x+a_0=d(x).
\]
\end{proof}

The minimal polynomial of \(C_d\) is also \(d\), independently of the
determinant calculation. The matrix acts as multiplication by \(x\) on
\(F[x]/(d)\), whose class of \(1\) is cyclic. We have \(d(C_d)(1)=0\), while
no nonzero polynomial of degree less than \(n\) annihilates \(1\), because
such a polynomial cannot belong to the ideal \((d)\). Because the class of
\(1\) generates the entire \(F[x]\)-module, a polynomial annihilates the
whole module if and only if it annihilates that class. Therefore
\(m_{C_d}=d\).

\begin{lemma}[Characteristic polynomials and invariant subspaces]
\label{lem:invariant-subspace-characteristic-divisibility}
If \(W\leq V\) is \(T\)-invariant, then
\[
\chi_{T|_W}(x)\mid\chi_T(x).
\]
\end{lemma}
\begin{proof}
Choose a basis of \(W\) and extend it to a basis of \(V\). In this adapted
basis the matrix of \(T\) has block form
\[
\begin{pmatrix}[T|_W]&*\\0&C\end{pmatrix}.
\]
The block-triangular determinant formula,
Theorem~\ref{thm:block-triangular-determinant}, gives
\[
\chi_T(x)=\chi_{T|_W}(x)\det(xI-C).
\]
\end{proof}

\begin{corollary}[Polynomial invariants of a cyclic subspace]
\label{cor:cyclic-subspace-polynomials}
Use the notation of Proposition~\ref{prop:cyclic-subspace-annihilator}. Then
\[
[T|_W]_\beta=C_d,
\qquad
m_{T|_W}=\chi_{T|_W}=d,
\qquad
d\mid\chi_T.
\]
If \(v\) is cyclic for all of \(V\), so \(W=V\) and \(n=\dim V\), then
\[
\beta=(v,Tv,\ldots,T^{n-1}v),\qquad
[T]_\beta=C_d,\qquad
m_T=\chi_T=d.
\]
\end{corollary}
\begin{proof}
The proposition gives the natural orbit basis and
\([T|_W]_\beta=C_d\). The direct companion determinant computation gives
\(\chi_{T|_W}=d\). Since \(v\) generates \(W\), a polynomial annihilates
\(T|_W\) on all of \(W\) exactly when it annihilates \(v\); hence
\(\operatorname{Ann}_{F[x]}(W)=\operatorname{Ann}_{F[x]}(v)=(d)\) and
\(m_{T|_W}=d\). Finally, \(W\) is \(T\)-invariant, so
Lemma~\ref{lem:invariant-subspace-characteristic-divisibility} gives
\(d=\chi_{T|_W}\mid\chi_T\). Taking \(W=V\) yields the last display.
\end{proof}

\Needspace{10\baselineskip}
\begin{remark}[Alternative proof: cyclic invariant subspaces]
Here is a classical linear-algebra proof of Cayley--Hamilton. Fix
\(0\ne v\in V\) and put
\[
W=F[x]v.
\]
Corollary~\ref{cor:cyclic-subspace-polynomials} gives both
\[
\chi_{T|_W}(T)v=0
\qquad\text{and}\qquad
\chi_{T|_W}\mid\chi_T.
\]
Write
\(\chi_T=q\chi_{T|_W}\). Hence
\[
\chi_T(T)v=q(T)\chi_{T|_W}(T)v=0.
\]
The zero vector is automatic, and \(v\) was otherwise arbitrary, so
\(\chi_T(T)=0\).
\end{remark}

\begin{theorem}[Rational canonical form]\label{thm:rational-canonical-form}
Let \(T\) be an endomorphism of a finite-dimensional vector space over an
arbitrary field \(F\).  There are unique monic nonconstant polynomials
\[
d_1\mid d_2\mid\cdots\mid d_r
\]
such that, in a suitable basis,
\[
[T]=C_{d_1}\oplus\cdots\oplus C_{d_r}.
\]
The invariant factors \(d_i\) classify \(T\) up to similarity.  Moreover,
\[
m_T=d_r,
\qquad
\chi_T=d_1\cdots d_r.
\]
\end{theorem}
\begin{proof}
Proposition~\ref{prop:operator-module-torsion} makes \(V_T\) a finitely
generated torsion module over the PID \(F[x]\).  The structure theorem gives
its unique invariant factors and a decomposition
\[
V_T\cong\bigoplus_{i=1}^rF[x]/(d_i).
\]
The cyclic-subspace proposition turns multiplication by \(x\) on each
summand into \(C_{d_i}\).  The intertwiner proposition says that isomorphic
operator modules are exactly similar operators, proving both existence and
classification.

The annihilator of the direct sum is generated by the least common multiple
of the \(d_i\), which is \(d_r\) because they form a divisibility chain.
Thus \(m_T=d_r\).  The direct-sum characteristic-polynomial formula and
Theorem~\ref{thm:companion-characteristic} give
\(\chi_T=d_1\cdots d_r\).
\end{proof}

\begin{corollary}[Irreducible factors of the polynomial invariants]
\label{cor:minimal-characteristic-factors}
The polynomials \(m_T\) and \(\chi_T\) have exactly the same distinct monic
irreducible factors. Consequently, \(m_T\) splits over \(F\) if and only if
\(\chi_T\) splits over \(F\).
\end{corollary}
\begin{proof}
In rational canonical form, \(m_T=d_r\) and
\(\chi_T=d_1\cdots d_r\). Every monic irreducible factor of \(\chi_T\)
occurs in some \(d_i\); since \(d_i\mid d_r=m_T\), that irreducible factor
divides \(m_T\). Conversely, every monic irreducible factor of
\(m_T=d_r\) occurs in the product \(\chi_T=d_1\cdots d_r\).
\end{proof}

\begin{definition}
A vector \(v\in V\) is a \emph{cyclic vector} for \(T\) if
\(V=F[x]v\). The operator \(T\) is \emph{cyclic} if it has a cyclic vector.
\end{definition}

\begin{corollary}[Cyclic-operator criterion]
\label{cor:cyclic-operator-criterion}
For an endomorphism of a nonzero finite-dimensional vector space, the
following are equivalent:
\begin{enumerate}
\item \(T\) is cyclic;
\item its rational canonical form consists of one companion block;
\item \(m_T=\chi_T\).
\end{enumerate}
\end{corollary}
\begin{proof}
If \(V=F[x]v\), then
\(V_T\cong F[x]/\operatorname{Ann}(v)\), so its invariant-factor
decomposition has one summand and its rational form has one companion block.
Conversely, the class of \(1\) generates the module belonging to one
companion block, so such an operator is cyclic. A single block has both
minimal and characteristic polynomial equal to its defining polynomial.

Finally, if the rational form has invariant factors
\(d_1\mid\cdots\mid d_r\), then
\(m_T=d_r\) and \(\chi_T=d_1\cdots d_r\). All \(d_i\) are nonconstant, so
equality of these monic polynomials forces
\(\sum_i\deg d_i=\deg d_r\), hence \(r=1\).
\end{proof}

\begin{remark}[Cayley--Hamilton from the PID structure theorem]
The rational canonical form gives a second structural derivation. If
\(d_1\mid\cdots\mid d_r\) are the invariant factors, then
\[
V_T\cong\bigoplus_iF[x]/(d_i),\qquad
m_T=d_r,\qquad
\chi_T=d_1\cdots d_r.
\]
The last formula uses the direct companion determinant computation and the
block-diagonal determinant formula. Since \(d_r\mid d_1\cdots d_r\), we have
\(m_T\mid\chi_T\). Therefore \(m_T(T)=0\) implies
\(\chi_T(T)=0\). This argument is noncircular: the companion characteristic
polynomial was proved directly by cofactor expansion, without invoking
Cayley--Hamilton.
\end{remark}

\begin{remark}[Three views of Cayley--Hamilton]
The adjugate proof is a matrix identity from determinant algebra. The cyclic
invariant-subspace proof uses the characteristic polynomial of each local
cyclic restriction and its divisibility into \(\chi_T\). The PID proof says structurally that
\(\chi_T\) is a multiple of the annihilator generator \(m_T\):
\[
\begin{array}{c}
\text{adjugate identity}\longrightarrow\text{matrix proof},\\
\text{cyclic invariant subspaces}\longrightarrow\text{local restrictions and divisibility},\\
\text{PID structure theorem}\longrightarrow\text{module-theoretic proof}.
\end{array}
\]
Each proof exposes a different reason for the same theorem.
\end{remark}

\begin{example}[Rational form without real Jordan form]
Over \(\mathbb R\), let
\[
A=\begin{pmatrix}0&-1\\1&0\end{pmatrix}.
\]
Then \(\chi_A(x)=x^2+1\), and \(A^2=-I\) while no linear polynomial
annihilates \(A\), so \(m_A(x)=x^2+1\). Moreover,
\[
e_1=\begin{pmatrix}1\\0\end{pmatrix},
\qquad Ae_1=\begin{pmatrix}0\\1\end{pmatrix}
\]
form a basis, so \(e_1\) is cyclic and the operator module is explicitly
\(\mathbb R[x]/(x^2+1)\). Its rational canonical form is the companion
matrix of \(x^2+1\), which is exactly \(A\) in the convention above. There
is no real Jordan form because \(A\) has no real eigenvalues. Over
\(\mathbb C\), the factorization \(x^2+1=(x-i)(x+i)\) makes the operator
diagonalizable with eigenvalues \(i\) and \(-i\).
\end{example}

\begin{example}[From invariant factors to rational canonical form]
Suppose an \(F[x]\)-module has invariant factors
\[d_1=x-1,\qquad d_2=(x-1)(x^2+1).\]
Then it is the direct sum \(F[x]/(d_1)\oplus F[x]/(d_2)\). The first
summand contributes the \(1\times1\) companion matrix \((1)\), and the
second contributes the companion matrix of
\(x^3-x^2+x-1\). Their block diagonal sum is the rational canonical form.
At once we read
\[m_T=d_2=(x-1)(x^2+1),\qquad
\chi_T=d_1d_2=(x-1)^2(x^2+1).\]
The example illustrates the division of labor: invariant factors describe
the cyclic summands, while matrices make that description visible in a chosen
basis.
\end{example}

\begin{example}[One matrix through both canonical forms]
Over \(\mathbb C\), begin with the actual matrix
\[
A=\begin{pmatrix}1&1&0\\0&1&0\\0&0&-1\end{pmatrix}.
\]
Let \(T\) be its operator on \(\mathbb C^3\), and give the space its
\(\mathbb C[x]\)-module structure from
Section~\ref{sec:linear-operators-fx-modules}. On the first two standard
basis vectors, \((T-I)e_2=e_1\) and \((T-I)e_1=0\), so this plane is the
cyclic module \(\mathbb C[x]/((x-1)^2)\). The last basis vector gives
\(\mathbb C[x]/(x+1)\). Thus the elementary divisors are
\((x-1)^2\) and \(x+1\). They are coprime, so CRT combines them into the
single invariant factor
\[
d=(x-1)^2(x+1)=x^3-x^2-x+1.
\]
Consequently \(\chi_T=d\) and \(m_T=d\), the rational canonical form is
the companion matrix of \(d\),
\[
\begin{pmatrix}0&0&-1\\1&0&1\\0&1&1\end{pmatrix},
\]
and the Jordan form is \(J_2(1)\oplus(-1)\). The original matrix happens
already to be in Jordan form, but the calculation shows how its module
invariants produce both canonical descriptions rather than merely recognize
one by inspection.
\end{example}

\begin{example}[A noncanonical matrix through the classification]
Over \(\mathbb R\), consider
\[
A=\begin{pmatrix}2&0&1\\0&3&-1\\0&0&2\end{pmatrix}.
\]
Triangularity gives \(\chi_A(x)=(x-2)^2(x-3)\).  For \(N=A-2I\),
\[
\ker N=\operatorname{span}(e_1),
\qquad
\ker N^2=\operatorname{span}(e_1,e_2+e_3).
\]
Thus the generalized \(2\)-eigenspace has dimension \(2\), and
\[
N(e_2+e_3)=e_1,\qquad Ne_1=0
\]
is a Jordan chain of length \(2\).  The ordinary \(3\)-eigenspace is
\(\operatorname{span}(e_2)\).  Hence the elementary divisors are
\((x-2)^2\) and \(x-3\), so
\[
m_A=(x-2)^2(x-3)=\chi_A.
\]
Because the elementary divisors are coprime, they combine into the single
invariant factor
\[
d=(x-2)^2(x-3)=x^3-7x^2+16x-12.
\]
The rational canonical form is
\[
C_d=\begin{pmatrix}0&0&12\\1&0&-16\\0&1&7\end{pmatrix},
\]
while the Jordan canonical form is \(J_2(2)\oplus(3)\).

The Jordan basis \(\beta=(e_1,e_2+e_3,e_2)\) gives an explicit change of
basis.  With
\[
P=\begin{pmatrix}1&0&0\\0&1&1\\0&1&0\end{pmatrix}
\]
whose columns are the vectors of \(\beta\), one has
\[
P^{-1}AP=J_2(2)\oplus(3).
\]
This example starts from a matrix in neither canonical form and uses
generalized eigenspaces, kernel growth, a Jordan chain, elementary divisors,
and invariant factors to reach both forms.
\end{example}

\section{Jordan Canonical Form}
When the relevant polynomials split, elementary divisors refine the rational
blocks into the familiar Jordan blocks. This is the structural completion of
the diagonalization discussion in Chapter~5. The two canonical forms mirror
the two forms of the PID structure theorem:
\[
\begin{array}{ccl}
\text{invariant factors}&\longleftrightarrow&\text{rational companion blocks},\\
\text{elementary divisors}&\longleftrightarrow&\text{Jordan blocks}.
\end{array}
\]
Accordingly, Theorem~\ref{thm:uniqueness-invariant-factor} gives uniqueness
of rational canonical form, while
Theorem~\ref{thm:uniqueness-elementary-divisor} gives uniqueness of the
Jordan blocks.

\begin{theorem}[Jordan canonical form]\label{thm:jordan-canonical-form}
Let \(T\) be an endomorphism of a finite-dimensional \(F\)-vector space and
suppose that \(\chi_T\) splits over \(F\).  Then \(V\) has a basis in which
\[
[T]=J_{k_1}(\lambda_1)\oplus\cdots\oplus J_{k_s}(\lambda_s).
\]
The blocks correspond to the elementary divisors
\[
(x-\lambda_1)^{k_1},\ldots,(x-\lambda_s)^{k_s},
\]
and their multiset is unique up to ordering.  Equivalently, \(V\) has a
basis obtained by concatenating Jordan chains.
\end{theorem}
\begin{proof}
When \(\chi_T\) splits, every irreducible factor in the elementary-divisor
decomposition of \(V_T\) is \(x-\lambda\).  Thus the PID structure theorem
gives a direct sum of cyclic modules
\[
F[x]/((x-\lambda)^k).
\]
Proposition~\ref{prop:jordan-chain-cyclic-module} supplies a Jordan-chain
basis and the block \(J_k(\lambda)\) on each summand.  Uniqueness of the
elementary divisors makes the multiset of blocks unique immediately; no
separate linear-algebraic uniqueness proof is needed.
\end{proof}

\begin{corollary}[Diagonalizability criterion]
\label{cor:diagonalizable-minimal-polynomial}
An endomorphism \(T\) is diagonalizable over \(F\) if and only if \(m_T\) is
a product of distinct linear factors over \(F\).
\end{corollary}
\begin{proof}
In Jordan form, diagonalizability says exactly that every block has size
\(1\). The elementary divisor belonging to a block \(J_k(\lambda)\) is
\((x-\lambda)^k\), and \(m_T\) is the least common multiple of all elementary
divisors. Thus all block sizes are \(1\) exactly when \(m_T\) splits into
distinct linear factors; for the reverse implication,
Corollary~\ref{cor:minimal-characteristic-factors} first ensures that
\(\chi_T\) splits, so Jordan form exists. The splitting condition is
essential: real rotation
of the plane has no real eigenbasis.
\end{proof}

\begin{corollary}[Largest blocks and kernel growth]
\label{cor:jordan-kernel-growth}
Assume that \(\chi_T\) splits over \(F\). For each eigenvalue \(\lambda\),
the exponent of \(x-\lambda\) in \(m_T\)
is the size of the largest Jordan block for \(\lambda\). Moreover, put
\[
K_j=\ker(T-\lambda I)^j,
\qquad K_0=0,
\qquad c_j=\dim K_j-\dim K_{j-1}.
\]
Then
\[
c_j=\#\{\text{Jordan blocks for \(\lambda\) of size at least \(j\)}\},
\]
and therefore
\[
c_j-c_{j+1}
=\#\{\text{Jordan blocks for \(\lambda\) of size exactly \(j\)}\}.
\]
Finally,
\[
\dim G_\lambda(T)=
\text{the algebraic multiplicity of \(\lambda\)}.
\]
\end{corollary}
\begin{proof}
The least common multiple of the elementary divisors contains
\((x-\lambda)^k\) to the largest exponent occurring among the
\(\lambda\)-blocks, proving the first assertion. On a block of size \(k\)
for \(\lambda\), the kernel of the \(j\)-th power has dimension
\(\min(j,k)\), so its contribution to \(c_j\) is \(1\) exactly when
\(k\geq j\). On a block for a different eigenvalue, \(T-\lambda I\) is
invertible and the contribution is zero. Summing over blocks proves both
counting formulas.

The generalized eigenspace is the direct sum of all \(\lambda\)-blocks, so
its dimension is their total size. The same total is the exponent of
\(x-\lambda\) in \(\chi_T\), hence the algebraic multiplicity.
\end{proof}

Chapter~5 established these kernel-growth formulas conditionally from a
displayed Jordan decomposition. The Jordan theorem supplies the required
decomposition for every split characteristic polynomial, so the bookkeeping
is now unconditional under precisely that hypothesis.

\begin{example}[From elementary divisors to Jordan form]
Over \(\mathbb C\), suppose the elementary divisors are
\[(x-2)^3,\qquad (x-2),\qquad (x+1)^2.\]
There are therefore two Jordan blocks for eigenvalue \(2\), of sizes \(3\)
and \(1\), and one block for eigenvalue \(-1\), of size \(2\). The Jordan
form is
\[J_3(2)\oplus J_1(2)\oplus J_2(-1).\]
Its characteristic polynomial is \((x-2)^4(x+1)^2\), whereas its minimal
polynomial is \((x-2)^3(x+1)^2\): the largest block for each eigenvalue
determines the corresponding exponent in the minimal polynomial.
\end{example}

\begin{remark}[How the polynomial invariants fit together]
If \(d_1\mid\cdots\mid d_r\) are invariant factors, then
\(\chi_T=d_1\cdots d_r\) and \(m_T=d_r\). Factoring the \(d_i\)'s and
using CRT yields elementary divisors. When these factors are powers of linear
polynomials, their exponents are exactly the Jordan-block sizes. Thus the
characteristic polynomial records total algebraic multiplicities, the minimal
polynomial records the largest required exponents, invariant factors produce
rational canonical form over every field, and elementary divisors produce
Jordan form when splitting is available.
\[
\begin{gathered}
\text{PID structure theorem}\\
\big\downarrow\\
\text{invariant factors}\\
\big\downarrow\\
\text{rational canonical form over every field},\\[1.5ex]
\text{factor into prime powers}\\
\big\downarrow\\
\text{elementary divisors}\\
\big\downarrow\\
\text{Jordan blocks when }p=x-\lambda.
\end{gathered}
\]
Rational canonical form is the field-independent classification.  Jordan
form is its split-polynomial refinement.
\end{remark}

\section*{Exercises}
\begin{exercise}[Basic] Classify
\[
\mathbb Z/8\mathbb Z\oplus\mathbb Z/12\mathbb Z
\]
in invariant-factor and elementary-divisor form.
\end{exercise}
\begin{exercise}[Core] Convert
\[
\mathbb Z/4\oplus\mathbb Z/8\oplus\mathbb Z/9\oplus\mathbb Z/3
\]
from elementary-divisor form to invariant-factor form, and convert
\(\mathbb Z/12\oplus\mathbb Z/60\) in the opposite direction.
\end{exercise}
\begin{exercise}[Core] List the abelian groups of order \(360=2^3 3^2 5\)
up to isomorphism by choosing partitions prime by prime.  Check the order of
each group from its invariant factors.\end{exercise}
\begin{exercise}[Core] Find the Smith normal form of
\(\begin{pmatrix}2&4&6\\4&10&14\end{pmatrix}\) over \(\mathbb Z\), regarded
as a homomorphism \(\mathbb Z^3\to\mathbb Z^2\), and determine its
cokernel.\end{exercise}
\begin{exercise}[Core] For \(T(x,y)=(-y,x)\) on \(\mathbb R^2\), find the
minimal polynomial and explain why no real Jordan form exists.\end{exercise}
\begin{exercise}[Core] Let \(T:V\to V\), \(S:W\to W\), and let
\(\phi:V\to W\) be \(F\)-linear.  Verify directly for
\(p(x)=2x^3-x+1\) that \(\phi T=S\phi\) implies
\(\phi p(T)=p(S)\phi\), and then prove the assertion for every
\(p\in F[x]\).\end{exercise}
\begin{exercise}[Core] For a vector \(v\), suppose
\(v,Tv,T^2v,T^3v\) are independent and
\[
T^4v=2v-Tv+3T^3v.
\]
Find \(\operatorname{Ann}(v)\) and the companion matrix of \(T\) on
\(F[x]v\).\end{exercise}
\begin{exercise}[Core] Prove the minimal-polynomial formula for a direct sum
of three operators by applying Proposition~\ref{prop:direct-sum-polynomials}
twice.  Give an example showing why the answer is an lcm rather than a
product.\end{exercise}
\begin{exercise}[Core] For
\[
A=\begin{pmatrix}1&1\\0&1\end{pmatrix}\in M_2(F),
\]
compute the Smith normal form of \(xI-A\) over \(F[x]\), recover the invariant
factors of \(V_A\), and determine the rational canonical form.\end{exercise}
\begin{exercise}[Further] Show that the invariant factors of a diagonalizable
operator are square-free.\end{exercise}
\begin{exercise}[Further] First suppose an operator has
\[
\dim\ker(T-\lambda I)^j=2,4,5,6,6,\ldots,
\]
and recover the sizes of all its Jordan blocks for \(\lambda\).  In a
separate example with different block data, construct explicit chains for
the matrix
\[
\begin{pmatrix}
\lambda&1&0&0\\0&\lambda&0&0\\
0&0&\lambda&1\\0&0&0&\lambda
\end{pmatrix}.
\]\end{exercise}
\begin{exercise}[Further] Decide whether the matrices
\[
\begin{pmatrix}1&1&0\\0&1&0\\0&0&2\end{pmatrix},
\qquad
\begin{pmatrix}1&0&0\\0&1&1\\0&0&2\end{pmatrix}
\]
are similar.  Justify the answer using elementary divisors rather than an
attempt to solve directly for a change-of-basis matrix.\end{exercise}
\begin{exercise}[Further, Challenging] Determine the rational and Jordan
canonical forms of an operator whose elementary divisors are
\((x-1)^3,(x-1),(x+2)^2\).\end{exercise}
\begin{exercise}[Looking Ahead] For a Jordan matrix, read off
\(\chi_T\), \(m_T\), and the dimensions of all generalized eigenspaces.
For a rational canonical matrix, read off \(\chi_T\) and \(m_T\) from its
companion blocks and explain how factorization over a splitting field predicts
the Jordan blocks.\end{exercise}

\chapter[Inner Product Spaces]{Inner Product Spaces and Spectral Theory}\label{chap:inner-products-spectral-theory}

Chapters~6 and~7 develop linear maps algebraically: kernels, images,
eigenvalues, and canonical forms do not require a notion of length. An inner
product adds geometry. It lets us ask stronger questions: when is a direct
sum orthogonal, when is a matrix a genuine change of orthonormal coordinates,
and when can an eigenbasis be chosen orthonormally? The answers lead to the
finite-dimensional spectral theorem and singular value decomposition.

Throughout this chapter \(V\) and \(W\) are finite-dimensional vector spaces
over \(F=\mathbb R\) or \(F=\mathbb C\). In the complex case our convention is
that the inner product is linear in the first variable and conjugate-linear in
the second. This convention agrees with the one used for adjoints below.

\section{Inner Products and Norms}

Length and angle in \(\mathbb R^n\) are encoded by the dot product. The
following axioms retain its useful properties while incorporating the complex
conjugation that makes a complex squared length real and nonnegative.

\begin{definition}
An \emph{inner product} on \(V\) is a map
\[
\langle\ ,\ \rangle:V\times V\longrightarrow F
\]
such that, for all \(u,v,w\in V\) and \(a,b\in F\),
\[
\langle au+bv,w\rangle
=a\langle u,w\rangle+b\langle v,w\rangle,
\]
\[
\langle v,u\rangle=\overline{\langle u,v\rangle},
\]
and
\[
\langle v,v\rangle\geq0,
\qquad
\langle v,v\rangle=0\Longleftrightarrow v=0.
\]
The associated \emph{norm} is
\[
\lVert v\rVert=\sqrt{\langle v,v\rangle}.
\]
\end{definition}

Conjugate symmetry turns first-variable linearity into
\[
\langle u,av+bw\rangle
=\overline a\langle u,v\rangle+\overline b\langle u,w\rangle.
\]
In particular,
\[
\lVert av\rVert=|a|\lVert v\rVert.
\]
Thus the complex convention is not cosmetic: it is what makes a scalar act
on length by its ordinary absolute value.

\begin{example}[Standard inner products]
On \(F^n\), the standard inner product is
\[
\langle u,v\rangle=\sum_{i=1}^n u_i\overline{v_i}.
\]
On \(M_{m,n}(F)\), the formula
\[
\langle A,B\rangle=\operatorname{tr}(AB^*)
=\sum_{i=1}^m\sum_{j=1}^n a_{ij}\overline{b_{ij}},
\qquad B^*=\overline B^{\,t},
\]
is an inner product. On the vector space \(P_n(F)\) of polynomials of degree
at most \(n\),
\[
\langle p,q\rangle=\int_0^1p(t)\overline{q(t)}\,dt
\]
is also an inner product. Positivity in the last example follows because a
nonzero polynomial is nonzero on some interval, so its squared absolute value
has strictly positive integral.
\end{example}

\begin{example}[Why conjugation is necessary]
The formula
\[
\sum_{i=1}^n u_iv_i
\]
is not an inner product on \(\mathbb C^n\): for \(u=(1,i)\), it gives
\[
1+i^2=0
\]
although \(u\ne0\). The standard formula instead gives
\[
\langle u,u\rangle=|1|^2+|i|^2=2.
\]
Conjugation is exactly what turns the complex dot product into a genuine
squared length.
\end{example}

\begin{theorem}[Cauchy--Schwarz inequality]
For all \(u,v\in V\),
\[
|\langle u,v\rangle|\leq\lVert u\rVert\lVert v\rVert.
\]
Equality holds, for nonzero \(u\) and \(v\), if and only if they are linearly
dependent.
\end{theorem}
\begin{proof}
If \(v=0\), the inequality is immediate. Otherwise put
\[
c=\frac{\langle u,v\rangle}{\lVert v\rVert^2}.
\]
Positivity gives
\[
0\leq\lVert u-cv\rVert^2.
\]
Using the linearity convention in the two variables, the right-hand side
expands to
\[
\lVert u\rVert^2
-\frac{|\langle u,v\rangle|^2}{\lVert v\rVert^2}.
\]
Multiplication by \(\lVert v\rVert^2\) proves the inequality. Equality
holds exactly when \(u-cv=0\), which is exactly linear dependence when both
vectors are nonzero.
\end{proof}

\begin{proposition}[Elementary norm identities]
For \(u,v\in V\),
\[
\lVert u+v\rVert\leq\lVert u\rVert+\lVert v\rVert,
\]
\[
u\perp v
\quad\Longrightarrow\quad
\lVert u+v\rVert^2=\lVert u\rVert^2+\lVert v\rVert^2,
\]
and
\[
\lVert u+v\rVert^2+\lVert u-v\rVert^2
=2\lVert u\rVert^2+2\lVert v\rVert^2.
\]
The last identity is the \emph{parallelogram identity}.
\end{proposition}
\begin{proof}
Expanding gives
\[
\lVert u+v\rVert^2
=\lVert u\rVert^2+2\operatorname{Re}\langle u,v\rangle+\lVert v\rVert^2.
\]
Cauchy--Schwarz bounds this by
\[
(\lVert u\rVert+\lVert v\rVert)^2,
\]
which proves the triangle inequality after taking nonnegative square roots.
When \(u\perp v\), the middle term vanishes, giving Pythagoras. Adding the
corresponding expansions for \(u+v\) and \(u-v\) cancels the two cross terms
and proves the parallelogram identity.
\end{proof}

\begin{proposition}[Polarization]
The norm induced by an inner product determines the inner product. More
precisely, over \(\mathbb R\),
\[
\langle u,v\rangle
=\frac14\bigl(\lVert u+v\rVert^2-\lVert u-v\rVert^2\bigr),
\]
whereas over \(\mathbb C\), with our first-variable-linear convention,
\[
\langle u,v\rangle
=\frac14\bigl(\lVert u+v\rVert^2-\lVert u-v\rVert^2\bigr)
+\frac{i}{4}\bigl(\lVert u+iv\rVert^2-\lVert u-iv\rVert^2\bigr).
\]
\end{proposition}
\begin{proof}
The first difference is
\[
\lVert u+v\rVert^2-\lVert u-v\rVert^2
=4\operatorname{Re}\langle u,v\rangle.
\]
Over \(\mathbb C\), direct expansion of the second difference gives
\[
\lVert u+iv\rVert^2-\lVert u-iv\rVert^2
=4\operatorname{Im}\langle u,v\rangle.
\]
The displayed formulas follow by combining real and imaginary parts. The
sign of the second term reflects the convention that the first, rather than
the second, variable is linear.
\end{proof}

\section{Orthonormal Sets and Gram--Schmidt}

Orthogonality makes coordinates behave like independent directions: all cross
terms disappear. Gram--Schmidt shows that this convenience is available in
every finite-dimensional inner-product space.

\begin{definition}
Vectors \(u,v\in V\) are \emph{orthogonal}, written \(u\perp v\), if
\[
\langle u,v\rangle=0.
\]
A set is \emph{orthogonal} if distinct members are orthogonal and
\emph{orthonormal} if, in addition, every member has norm one. For a subset
\(S\subseteq V\), its \emph{orthogonal complement} is
\[
S^\perp=\{v\in V:\langle v,s\rangle=0\text{ for every }s\in S\}.
\]
\end{definition}

\begin{proposition}[Orthogonal coordinates]\leavevmode\\
Every orthogonal set of nonzero vectors is independent. If
\[
(e_1,\ldots,e_n)
\]
is an orthonormal basis of \(V\), then every \(v\in V\) has the expansion
\[
v=\sum_{j=1}^n\langle v,e_j\rangle e_j,
\]
and
\[
\lVert v\rVert^2=\sum_{j=1}^n|\langle v,e_j\rangle|^2.
\]
\end{proposition}
\begin{proof}
If
\[
\sum_{j=1}^na_je_j=0
\]
for an orthogonal set of nonzero vectors, take the inner product with
\(e_k\). This gives
\[
a_k\lVert e_k\rVert^2=0,
\]
so every \(a_k=0\). For an orthonormal basis, write
\[
v=\sum_{j=1}^na_je_j.
\]
Taking the inner product with \(e_k\) gives \(a_k=\langle v,e_k\rangle\).
Pythagoras applied to this orthogonal sum gives the norm formula.
\end{proof}

\begin{corollary}[Bessel and Parseval]
If \(e_1,\ldots,e_r\) is an orthonormal set in \(V\), then
\[
\sum_{j=1}^r|\langle v,e_j\rangle|^2\leq\lVert v\rVert^2
\qquad(v\in V).
\]
If the set is an orthonormal basis, equality holds; this equality is
\emph{Parseval's identity} in the finite-dimensional setting.
\end{corollary}
\begin{proof}
Put
\[
w=\sum_{j=1}^r\langle v,e_j\rangle e_j.
\]
Then \(v-w\) is orthogonal to every \(e_j\), hence to \(w\). Pythagoras
gives
\[
\lVert v\rVert^2=\lVert w\rVert^2+\lVert v-w\rVert^2
=\sum_{j=1}^r|\langle v,e_j\rangle|^2+\lVert v-w\rVert^2.
\]
This proves Bessel's inequality. If the \(e_j\)'s form a basis, then
\(w=v\), so equality holds.
\end{proof}

\begin{theorem}[Gram--Schmidt orthogonalization]
Let
\[
(v_1,\ldots,v_n)
\]
be a linearly independent list. Define first
\[
u_1=v_1,
\qquad
e_1=\frac{u_1}{\lVert u_1\rVert},
\]
then
\[
u_2=v_2-\langle v_2,e_1\rangle e_1,
\qquad
e_2=\frac{u_2}{\lVert u_2\rVert},
\]
then
\[
u_3=v_3-\langle v_3,e_1\rangle e_1
-\langle v_3,e_2\rangle e_2,
\qquad
e_3=\frac{u_3}{\lVert u_3\rVert}.
\]
In general,
\[
u_k=v_k-\sum_{j=1}^{k-1}\langle v_k,e_j\rangle e_j,
\qquad
e_k=\frac{u_k}{\lVert u_k\rVert}.
\]
Every \(u_k\) is nonzero, the resulting list is orthonormal, and
\[
\operatorname{span}(e_1,\ldots,e_k)
=\operatorname{span}(v_1,\ldots,v_k)
\qquad(1\leq k\leq n).
\]
\end{theorem}
\begin{proof}
Assume that \(e_1,\ldots,e_{k-1}\) have been constructed. For \(i<k\),
\[
\langle u_k,e_i\rangle
=\langle v_k,e_i\rangle
-\sum_{j=1}^{k-1}\langle v_k,e_j\rangle\langle e_j,e_i\rangle
=0.
\]
Thus \(u_k\) is orthogonal to the preceding \(e_i\)'s. If \(u_k=0\), the
definition would put \(v_k\) in the span of \(e_1,\ldots,e_{k-1}\), which,
by the induction hypothesis, is the span of \(v_1,\ldots,v_{k-1}\). This
contradicts independence. Hence \(e_k\) is defined and has norm one.

The formula defining \(u_k\) shows that \(u_k\) lies in the span of the first
\(k\) vectors \(v_i\). Conversely, it can be rearranged to express \(v_k\)
in the span of \(u_k,e_1,\ldots,e_{k-1}\). The two successive spans are
therefore equal. Induction proves every assertion.
\end{proof}

\begin{corollary}[Gram--Schmidt fixes orthogonal lists]
If \(v_1,\ldots,v_k\) is an orthogonal list of nonzero vectors, then
Gram--Schmidt applied to this ordered list gives
\[
e_i=\frac{v_i}{\lVert v_i\rVert}
\qquad(1\leq i\leq k).
\]
In particular, Gram--Schmidt fixes every ordered orthonormal set.
\end{corollary}
\begin{proof}
At the \(i\)-th step, the preceding vectors \(e_j\) are scalar multiples of
\(v_j\).  Orthogonality gives
\[
\langle v_i,e_j\rangle=0
\qquad(j<i),
\]
so every projection term in the Gram--Schmidt formula vanishes.  Thus
\(u_i=v_i\), and normalizing gives the displayed formula.  If the list is
already orthonormal, then \(\lVert v_i\rVert=1\), so \(e_i=v_i\).
\end{proof}

\begin{corollary}[Orthonormal bases]\leavevmode\\
Each finite-dimensional inner-product space admits an orthonormal basis.
Every orthonormal set can be enlarged to one.
\end{corollary}
\begin{proof}
Apply Gram--Schmidt to any basis. For an orthonormal set, first extend it to
an ordinary basis using the finite-dimensional basis theorem of Chapter~5,
then apply Gram--Schmidt to the additional vectors. The first members are
unchanged because they are already orthonormal.
\end{proof}

\begin{example}[Gram--Schmidt in \(\mathbb R^2\)]
Starting with
\[
v_1=(1,1),
\qquad
v_2=(1,0),
\]
we obtain
\[
e_1=\frac1{\sqrt2}(1,1),
\]
and
\[
u_2=(1,0)-\left\langle(1,0),e_1\right\rangle e_1
=\left(\frac12,-\frac12\right).
\]
Thus
\[
e_2=\frac1{\sqrt2}(1,-1).
\]
The original and orthonormal lists span the same plane, but the latter makes
lengths, coordinates, and projections immediate.
\end{example}

\section{Orthogonal Complements and Projections}

An arbitrary direct-sum complement of a subspace is seldom canonical. The
orthogonal complement is canonical because it is determined by the inner
product alone; it separates a vector into a component in the subspace and a
perpendicular error.

\begin{proposition}[Orthogonal complements]
If \(W\leq V\), then \(W^\perp\) is a subspace. In finite dimensions,
\[
V=W\oplus W^\perp,
\qquad
\dim W+\dim W^\perp=\dim V,
\qquad
(W^\perp)^\perp=W.
\]
\end{proposition}
\begin{proof}
Closure of \(W^\perp\) under linear combinations follows from linearity in
the first variable. Choose an orthonormal basis
\[
(e_1,\ldots,e_r)
\]
of \(W\). For \(v\in V\), define
\[
P_Wv=\sum_{i=1}^r\langle v,e_i\rangle e_i.
\]
Then \(P_Wv\in W\), while, for each \(i\),
\[
\langle v-P_Wv,e_i\rangle=0.
\]
Thus \(v-P_Wv\in W^\perp\), proving \(V=W+W^\perp\). A vector in the
intersection is orthogonal to itself and is therefore zero, so the sum is
direct. Taking dimensions gives the formula. Since \(W\subseteq
(W^\perp)^\perp\) and the two spaces have the same dimension, they are
equal.
\end{proof}

\begin{definition}
The linear operator
\[
P_W:V\longrightarrow V,
\qquad
P_Wv=\sum_{i=1}^r\langle v,e_i\rangle e_i,
\]
where \(e_1,\ldots,e_r\) is any orthonormal basis of \(W\), is the
\emph{orthogonal projection onto \(W\)}; its image is \(W\).
\end{definition}

The preceding proof shows that the formula depends only on the unique
decomposition
\[
v=P_Wv+(v-P_Wv),
\qquad
P_Wv\in W,
\quad v-P_Wv\in W^\perp,
\]
and therefore not on the chosen orthonormal basis.

\begin{theorem}[Best approximation]
For \(v\in V\) and \(w\in W\),
\[
\lVert v-w\rVert^2
=\lVert v-P_Wv\rVert^2+\lVert P_Wv-w\rVert^2.
\]
Consequently \(P_Wv\) is the unique point of \(W\) nearest to \(v\).
\end{theorem}
\begin{proof}
The vector \(v-P_Wv\) lies in \(W^\perp\), whereas \(P_Wv-w\) lies in
\(W\). They are orthogonal and their sum is \(v-w\), so Pythagoras gives the
formula. The second summand is nonnegative and vanishes exactly when
\(w=P_Wv\), proving uniqueness.
\end{proof}

\begin{example}[Projection onto a line]
Let
\[
W=\operatorname{span}(1,1,1)\subseteq\mathbb R^3.
\]
The unit vector \(e=(1,1,1)/\sqrt3\) gives
\[
P_W(1,2,3)=\langle(1,2,3),e\rangle e=2(1,1,1).
\]
The error
\[
(-1,0,1)
\]
has dot product zero with \((1,1,1)\), visibly confirming the orthogonal
decomposition.
\end{example}

\section{Linear Functionals and Adjoints}

The algebraic dual \(V^\vee\), developed in Section~6.15, consists of all
linear measurements on \(V\). In finite-dimensional inner-product spaces,
the Riesz theorem says that every such measurement is obtained by taking an
inner product with one uniquely determined vector. It is this theorem that
makes the adjoint an intrinsic operator rather than merely a matrix operation.

\begin{theorem}[Finite-dimensional Riesz representation]
For every linear functional
\[
\varphi:V\longrightarrow F,
\]
there is a unique \(u\in V\) such that
\[
\varphi(v)=\langle v,u\rangle
\qquad(v\in V).
\]
\end{theorem}
\begin{proof}
Choose an orthonormal basis \(e_1,\ldots,e_n\). Put
\[
u=\sum_{j=1}^n\overline{\varphi(e_j)}e_j.
\]
If
\[
v=\sum_{j=1}^nx_je_j,
\]
then
\[
\langle v,u\rangle
=\sum_{j=1}^nx_j\varphi(e_j)
=\varphi(v).
\]
If both \(u\) and \(u'\) represent \(\varphi\), then
\[
\langle v,u-u'\rangle=0
\]
for every \(v\). Taking \(v=u-u'\) proves \(u=u'\).
\end{proof}

\begin{theorem}[Adjoint]
For every linear map
\[
T:V\longrightarrow W,
\]
there is a unique linear map
\[
T^*:W\longrightarrow V
\]
satisfying
\[
\langle Tv,w\rangle=\langle v,T^*w\rangle
\qquad(v\in V,\ w\in W).
\]
\end{theorem}
\begin{proof}
For fixed \(w\in W\), the function
\[
v\longmapsto\langle Tv,w\rangle
\]
is a linear functional on \(V\). Riesz representation supplies a unique
vector, denoted \(T^*w\), satisfying the displayed equality. For
\(a,b\in F\) and \(w_1,w_2\in W\), conjugate-linearity in the second
variable gives
\[
\langle Tv,aw_1+bw_2\rangle
=\overline a\langle Tv,w_1\rangle
+\overline b\langle Tv,w_2\rangle
=\langle v,aT^*w_1+bT^*w_2\rangle.
\]
Uniqueness in Riesz representation proves the linearity of \(T^*\). The
same uniqueness proves uniqueness of the adjoint.
\end{proof}

\noindent When \(V=W\), this is precisely the adjoint of the linear operator
\(T:V\to V\) considered below.  The point of allowing two spaces here is
that the adjoint reverses the direction of a map:
\[
T:V\longrightarrow W,
\qquad
T^*:W\longrightarrow V.
\]

\begin{proposition}[Adjoint identities and matrices]
The familiar adjoint identities remain valid for maps between possibly
different inner-product spaces, provided all sums and compositions are
defined.  Thus, if \(S,T:V\to W\) and \(a\in F\), then
\[
(S+T)^*=S^*+T^*,
\qquad
(aT)^*=\overline a\,T^*,
\qquad
(T^*)^*=T.
\]
If \(T:V\to W\) and \(S:W\to U\), then \(ST:V\to U\) and
\[
(ST)^*=T^*S^*:U\to V.
\]
If \(T:V\to W\) is an isomorphism, then
\[
(T^{-1})^*=(T^*)^{-1}.
\]
Here \(T^{-1}:W\to V\), so \((T^{-1})^*:V\to W\).  Likewise,
\(T^*:W\to V\), so \((T^*)^{-1}:V\to W\).
Relative to orthonormal bases, the matrix of \(T^*\) is
\[
[T^*]=[T]^*=\overline{[T]}^{\,t}.
\]
\end{proposition}
\begin{proof}
For the sum identity, the defining equation and linearity in the first
variable give
\[
\langle(S+T)v,w\rangle
=\langle v,(S^*+T^*)w\rangle.
\]
For the scalar identity, our convention gives
\[
\langle aTv,w\rangle
=a\langle Tv,w\rangle
=a\langle v,T^*w\rangle
=\langle v,\overline a\,T^*w\rangle.
\]
Over \(\mathbb R\), \(\overline a=a\).  For composition, if
\(T:V\to W\) and \(S:W\to U\), then
\[
\langle STv,w\rangle
=\langle Tv,S^*w\rangle
=\langle v,T^*S^*w\rangle.
\]
Finally,
\[
\langle T^*w,v\rangle_V
=\overline{\langle v,T^*w\rangle_V}
=\overline{\langle Tv,w\rangle_W}
=\langle w,Tv\rangle_W,
\]
so uniqueness shows that \((T^*)^*=T\).  For an isomorphism \(T:V\to W\),
\[
(T^{-1}T)^*=T^*(T^{-1})^*=I_V
\]
and
\[
(TT^{-1})^*=(T^{-1})^*T^*=I_W.
\]
Thus \((T^{-1})^*\) is the inverse of \(T^*\).
For orthonormal coordinate columns \(x,y\) and a matrix \(A=[T]\),
\[
\langle Ax,y\rangle
=x^tA^t\overline y
=\langle x,\overline A^{\,t}y\rangle.
\]
Thus the matrix of the adjoint is the conjugate transpose. Over
\(\mathbb R\), this is the ordinary transpose.
\end{proof}

\begin{definition}[Hermitian and symmetric matrices]
For \(A\in M_n(\mathbb C)\), its \emph{conjugate transpose} is
\[
A^*=\overline A^{\,\mathsf T}.
\]
The matrix \(A\) is \emph{Hermitian} if
\[
A^*=A.
\]
Equivalently, its entries satisfy
\[
a_{ij}=\overline{a_{ji}}
\qquad(1\leq i,j\leq n).
\]
In particular, every diagonal entry of a Hermitian matrix is real.  For
example,
\[
\begin{pmatrix}2&i\\-i&3\end{pmatrix}
\]
is Hermitian.  Over \(\mathbb R\), conjugation has no effect, and the
corresponding condition \(A^{\mathsf T}=A\) defines a \emph{symmetric}
matrix.
\end{definition}

\begin{remark}[Self-adjoint operators and matrices]
If \(\beta\) is an orthonormal basis of a finite-dimensional inner-product
space, then
\[
[T^*]_\beta=[T]_\beta^*.
\]
Consequently, over \(\mathbb C\),
\[
T\text{ is self-adjoint}
\quad\Longleftrightarrow\quad
[T]_\beta\text{ is Hermitian};
\]
over \(\mathbb R\), the corresponding statement uses symmetric matrices.
The orthonormality of \(\beta\) matters: in an arbitrary basis, the matrix
of a self-adjoint operator need not literally be Hermitian or symmetric.
\end{remark}

\begin{theorem}[Kernels, images, and adjoints]
For \(T:V\to W\),
\[
\ker T^*=(\operatorname{im}T)^\perp,
\qquad
\operatorname{im}T^*=(\ker T)^\perp.
\]
Consequently,
\[
\operatorname{rank}T=\operatorname{rank}T^*,
\]
and the equation \(Tv=w\) is solvable if and only if
\[
w\perp\ker T^*.
\]
\end{theorem}
\begin{proof}
For \(w\in W\),
\[
\begin{aligned}
w\in\ker T^*
&\Longleftrightarrow
\langle v,T^*w\rangle=0
&&\text{for every }v\in V \\
&\Longleftrightarrow
\langle Tv,w\rangle=0
&&\text{for every }v\in V.
\end{aligned}
\]
This is exactly \(w\in(\operatorname{im}T)^\perp\), proving the first
identity. Applying it to \(T^*\), using \((T^*)^*=T\) and taking a second
orthogonal complement in finite dimensions gives
\[
\operatorname{im}T^*=(\ker T)^\perp.
\]
The rank equality follows either by dimensions of these complements or by
rank--nullity. Finally,
\[
w\in\operatorname{im}T
\Longleftrightarrow
w\in(\ker T^*)^\perp,
\]
which is the asserted solvability criterion.
\end{proof}

\begin{proposition}[Characterization of orthogonal projections]
For a linear operator \(P:V\to V\), the following are equivalent:
\begin{enumerate}
\item \(P=P_W\) for some subspace \(W\);
\item \(P^2=P\) and \(P^*=P\).
\end{enumerate}
In this situation,
\[
W=\operatorname{im}P,
\qquad
\ker P=W^\perp.
\]
\end{proposition}
\begin{proof}
The formula for \(P_W\) shows directly that \(P_W^2=P_W\). If
\(v,w\in V\), expansion in an orthonormal basis of \(W\) gives
\[
\langle P_Wv,w\rangle=\langle v,P_Ww\rangle,
\]
so \(P_W^*=P_W\).

Conversely, suppose \(P^2=P=P^*\). The idempotence gives
\[
V=\operatorname{im}P\oplus\ker P.
\]
If \(x=Pu\in\operatorname{im}P\) and \(y\in\ker P\), then
\[
\langle x,y\rangle=\langle Pu,y\rangle
=\langle u,Py\rangle=0.
\]
Thus \(\ker P\subseteq(\operatorname{im}P)^\perp\). The direct-sum
dimension formula makes the two subspaces equal, so \(P\) is the orthogonal
projection onto its image.
\end{proof}

An ordinary algebraic projection only requires \(P^2=P\). It need not be
self-adjoint, and its kernel need not be perpendicular to its image. The
self-adjoint condition is exactly what distinguishes orthogonal projection.

\section{Self-Adjoint, Normal, Unitary, and Orthogonal Operators}

Adjoints identify the maps compatible with inner-product geometry. These
operator classes are important not merely for their names: normality is the
precise complex hypothesis for orthonormal diagonalization.

\begin{definition}
For an endomorphism \(T:V\to V\), we say that \(T\) is \emph{self-adjoint}
if
\[
T=T^*,
\]
\emph{skew-adjoint} if
\[
T^*=-T,
\]
and \emph{normal} if
\[
TT^*=T^*T.
\]
It is \emph{positive} if it is self-adjoint and
\[
\langle Tv,v\rangle\geq0
\qquad(v\in V),
\]
and \emph{positive-definite} if equality occurs only for \(v=0\). A complex
operator is \emph{unitary}, and a real operator is \emph{orthogonal}, if
\[
T^*T=I.
\]
\end{definition}

Self-adjoint matrices are real symmetric or complex Hermitian. The terms
unitary and orthogonal distinguish the base field because the latter is the
real specialization of the former.

\begin{proposition}[Isometries and unitary operators]
For an endomorphism \(T\) of a finite-dimensional inner-product space, the
following are equivalent:
\begin{enumerate}
\item \(T\) is unitary or orthogonal, as appropriate to the field;
\item \(T^*T=I\);
\item
\[
\langle Tu,Tv\rangle=\langle u,v\rangle
\qquad(u,v\in V);
\]
\item
\[
\lVert Tv\rVert=\lVert v\rVert
\qquad(v\in V).
\]
\end{enumerate}
In this case \(T\) is invertible and
\[
T^{-1}=T^*.
\]
Relative to an orthonormal basis, its matrix \(A\) satisfies
\[
A^*A=AA^*=I.
\]
Thus the columns and rows of \(A\) are orthonormal bases.
\end{proposition}
\begin{proof}
The equivalence of \(1\) and \(2\) is the definition. Moreover,
\[
\langle Tu,Tv\rangle=\langle u,T^*Tv\rangle,
\]
so \(2\) implies \(3\), and \(3\) implies \(4\). Conversely, the polarization
identity shows that \(4\) determines the same inner products, so \(4\) implies
\(3\). If \(3\) holds, then
\[
\langle u,(T^*T-I)v\rangle=0
\]
for every \(u,v\); taking \(u=(T^*T-I)v\) proves \(2\).

Condition \(4\) makes \(T\) injective. Finite dimensionality then makes it
surjective, so it is invertible; from \(T^*T=I\) we obtain \(T^*=T^{-1}\),
and hence \(TT^*=I\). The matrix identities follow in orthonormal bases.
Their entries state exactly that both the columns and rows are orthonormal.
\end{proof}

\begin{proposition}[Basic consequences of the adjoint]
Every operator \(T^*T\) is self-adjoint and positive. If \(T\) is normal,
then
\[
\lVert Tv\rVert=\lVert T^*v\rVert
\qquad(v\in V),
\]
and
\[
\ker T=\ker T^*.
\]
Self-adjoint and unitary or orthogonal operators are normal.
\end{proposition}
\begin{proof}
The adjoint identities give
\[
(T^*T)^*=T^*T,
\]
and
\[
\langle T^*Tv,v\rangle=\langle Tv,Tv\rangle\geq0.
\]
If \(T\) is normal, then
\[
\lVert Tv\rVert^2
=\langle T^*Tv,v\rangle
=\langle TT^*v,v\rangle
=\lVert T^*v\rVert^2.
\]
This equality proves the kernel identity. Self-adjointness makes the two
products identical, while for a unitary or orthogonal operator they both
equal \(I\), proving the final assertion.
\end{proof}

\begin{proposition}[Eigenvalue geometry]
Let \(T\) be an endomorphism of a finite-dimensional inner-product space.
\begin{enumerate}
\item If \(T\) is self-adjoint, every eigenvalue of \(T\) is real.
\item If \(T\) is unitary or orthogonal, every eigenvalue \(\lambda\) has
\[
|\lambda|=1.
\]
\item If \(T\) is normal on a complex inner-product space and
\[
Tv=\lambda v,
\]
then
\[
T^*v=\overline\lambda v.
\]
Consequently, eigenvectors of a normal operator belonging to distinct
eigenvalues are orthogonal.
\end{enumerate}
\end{proposition}
\begin{proof}
For a self-adjoint \(T\) and a nonzero eigenvector \(v\),
\[
\lambda\langle v,v\rangle
=\langle Tv,v\rangle
=\langle v,Tv\rangle
=\overline\lambda\langle v,v\rangle,
\]
so \(\lambda=\overline\lambda\). If \(T\) is unitary or orthogonal,
\[
\lVert v\rVert=\lVert Tv\rVert=|\lambda|\lVert v\rVert,
\]
proving the second claim. For normal \(T\), the preceding proposition gives
\[
\lVert(T-\lambda I)v\rVert
=\lVert(T^*-\overline\lambda I)v\rVert.
\]
The left side is zero, hence the third assertion follows. If also
\[
Tw=\mu w,
\]
then
\[
\lambda\langle v,w\rangle
=\langle Tv,w\rangle
=\langle v,T^*w\rangle
=\langle v,\overline\mu w\rangle
=\mu\langle v,w\rangle.
\]
Thus \(\langle v,w\rangle=0\) when \(\lambda\ne\mu\).
\end{proof}

For a normal operator on a complex inner-product space, the last assertion
may be recorded as
\[
\lambda\ne\mu
\quad\Longrightarrow\quad
E_\lambda\perp E_\mu.
\]
Here the proof above applies to arbitrary vectors in the two eigenspaces,
not only to a particular choice of eigenvectors.

\begin{example}[Why real eigenvalues are not enough]
On \(\mathbb C^2\) with its standard inner product, consider
\[
A=\begin{pmatrix}1&1\\0&2\end{pmatrix}.
\]
Its eigenvalues are \(1\) and \(2\), with eigenvectors
\[
v_1=(1,0),
\qquad
v_2=(1,1),
\]
respectively.  Thus \(A\) is diagonalizable, and all of its eigenvalues are
real.  Nevertheless,
\[
A^*=\begin{pmatrix}1&0\\1&2\end{pmatrix}\ne A,
\qquad
A^*A=\begin{pmatrix}1&1\\1&5\end{pmatrix}
\ne
\begin{pmatrix}2&2\\2&4\end{pmatrix}=AA^*,
\]
so \(A\) is neither self-adjoint nor normal.  Also,
\(\langle v_1,v_2\rangle=1\), showing that eigenvectors for distinct
eigenvalues of an arbitrary diagonalizable operator need not be orthogonal.

This example rules out both false implications
\[
\begin{gathered}
\text{real eigenvalues}\ \not\Longrightarrow\ \text{self-adjoint},\\
\text{diagonalizable with real eigenvalues}\ \not\Longrightarrow\ \text{self-adjoint}.
\end{gathered}
\]
It also illustrates the danger of applying Gram--Schmidt globally to an
eigenbasis.  Starting with \((v_1,v_2)\), the process gives
\[
e_1=(1,0),
\qquad
e_2=v_2-\langle v_2,e_1\rangle e_1=(0,1).
\]
But
\[
Ae_2=(1,2)
\]
is not a scalar multiple of \(e_2\), so the resulting orthonormal basis is
not an eigenbasis.
\end{example}

\begin{remark}[Local, not global, Gram--Schmidt]
If
\[
V=\bigoplus_\lambda E_\lambda
\]
is an eigenspace decomposition, Gram--Schmidt can safely be applied inside
each individual eigenspace: linear combinations of vectors in
\(E_\lambda\) remain eigenvectors with eigenvalue \(\lambda\).  Its union
is an orthonormal basis of \(V\) precisely when the distinct eigenspaces are
mutually orthogonal.  The example shows why applying the process across
nonorthogonal eigenspaces can fail.  For a normal complex operator, the
orthogonality just proved makes the local procedure work globally; if an
eigenbasis is already orthonormal, the preceding corollary says there is
nothing for Gram--Schmidt to change.
\end{remark}

\begin{remark}[The real and complex cases]
A self-adjoint operator on a finite-dimensional real inner-product space has
real eigenvalues and an orthonormal eigenbasis, as the real spectral theorem
below states.  A real normal operator need not be diagonalizable over
\(\mathbb R\).  Indeed,
\[
R_\theta=
\begin{pmatrix}
\cos\theta&-\sin\theta\\
\sin\theta&\cos\theta
\end{pmatrix}
\]
is orthogonal and hence normal, but for
\(\theta\notin\{0,\pi\}\) it has no real eigenvalues.  The real analogue
of unitary diagonalization for general normal operators is orthogonal block
diagonalization, with real \(1\times1\) blocks and \(2\times2\)
rotation-type blocks.  Thus the full normal spectral theorem is naturally a
complex statement.
\end{remark}

\section{Spectral Theorems}

The spectral theorem identifies exactly when ordinary diagonalization can be
made geometric. Its finite-dimensional proof requires no Hilbert-space
theory, but the complex result uses the Fundamental Theorem of Algebra to
produce an eigenvalue.

\begin{remark}[The Fundamental Theorem of Algebra]
The Fundamental Theorem of Algebra says that every nonconstant polynomial in
\(\mathbb C[x]\) has a root in \(\mathbb C\). We use it only to ensure that
the characteristic polynomial of a complex endomorphism has a root. A
Galois-theoretic proof appears later in Chapter~11; no complex analysis is
needed elsewhere in this chapter.
\end{remark}

\begin{theorem}[Complex spectral theorem]
For an endomorphism \(T\) of a finite-dimensional complex inner-product
space, the following are equivalent:
\begin{enumerate}
\item \(T\) is normal;
\item \(V\) has an orthonormal basis of eigenvectors of \(T\);
\item the matrix of \(T\) in some orthonormal basis is diagonal.
\end{enumerate}
Equivalently, \(T\) is unitarily diagonalizable.
\end{theorem}
\begin{proof}
The equivalence of \(2\) and \(3\) is the definition of a matrix in a chosen
basis. A diagonal matrix commutes with its conjugate transpose, so \(2\)
implies \(1\). Conversely, suppose \(T\) is normal and argue by induction on
\(\dim V\). The case \(\dim V=0\) is clear. The Fundamental Theorem of
Algebra supplies an eigenvalue \(\lambda\), and let
\[
E=\ker(T-\lambda I).
\]
If \(z\in E^\perp\) and \(e\in E\), then
\[
\langle Tz,e\rangle
=\langle z,T^*e\rangle
=\langle z,\overline\lambda e\rangle=0.
\]
Thus \(E^\perp\) is \(T\)-invariant. The same calculation with \(T^*\)
shows that it is \(T^*\)-invariant, so the restriction of \(T\) to
\(E^\perp\) remains normal. It has smaller dimension unless \(E=V\), and
the induction hypothesis supplies an orthonormal eigenbasis of \(E^\perp\).
An orthonormal basis of \(E\), obtained by Gram--Schmidt, together with that
basis gives an orthonormal eigenbasis of \(V\).
\end{proof}

\begin{corollary}[Real spectral theorem]
For an endomorphism \(T\) of a finite-dimensional real inner-product space,
the following are equivalent:
\begin{enumerate}
\item \(T\) is self-adjoint;
\item \(V\) has an orthonormal basis of real eigenvectors of \(T\);
\item \(T\) is orthogonally diagonalizable.
\end{enumerate}
\end{corollary}
\begin{proof}
The last two conditions are equivalent by coordinates, and either one plainly
implies self-adjointness. Conversely, complexify \(V\) and extend the real
inner product to the usual Hermitian inner product on \(V_\mathbb C\). The
complexified operator is self-adjoint, hence normal; the complex spectral
theorem applies. Its eigenvalues are real by the preceding eigenvalue
proposition. If
\[
T_\mathbb C(x+iy)=\lambda(x+iy),
\]
then
\[
Tx=\lambda x,
\qquad
Ty=\lambda y.
\]
Thus the complex eigenspaces are complexifications of real eigenspaces.
They span \(V_\mathbb C\), so the real eigenspaces span \(V\). Distinct
eigenspaces are orthogonal, and Gram--Schmidt within each one produces a
real orthonormal eigenbasis.
\end{proof}

\begin{remark}[Spectral-theorem implications]
For a finite-dimensional complex inner-product space, the main implications
are
\[
\text{self-adjoint}
\ \Longrightarrow\
\text{normal}
\ \Longrightarrow\
\text{unitarily diagonalizable}
\ \Longrightarrow\
\text{diagonalizable}.
\]
The reverse implications fail without additional hypotheses.  More exactly,
\[
T\text{ is normal}
\quad\Longleftrightarrow\quad
T\text{ has an orthonormal eigenbasis},
\]
and
\[
\begin{gathered}
T\text{ is self-adjoint}\\
\Longleftrightarrow\\
T\text{ has an orthonormal eigenbasis with real eigenvalues}.
\end{gathered}
\]
The latter equivalence explains the geometric content that is missing from
ordinary diagonalizability.  It also gives the useful converse
\[
T\text{ normal with real spectrum}
\quad\Longrightarrow\quad
T\text{ self-adjoint}:
\]
in an orthonormal eigenbasis, its diagonal matrix has real entries and hence
equals its conjugate transpose.  At the matrix level, this says
\[
\begin{gathered}
A\text{ is Hermitian}\\
\Longleftrightarrow\\
A\text{ is unitarily diagonalizable with real eigenvalues}.
\end{gathered}
\]
\end{remark}

\begin{example}[Normal diagonalization]
The matrix
\[
\begin{pmatrix}2&0\\0&i\end{pmatrix}
\]
is normal and already unitarily diagonal. By contrast,
\[
\begin{pmatrix}1&a\\0&1\end{pmatrix}
\]
is diagonalizable only when \(a=0\). For \(a\ne0\), it does not commute
with its conjugate transpose and is not normal. The spectral theorem
therefore excludes precisely the Jordan-block obstruction described in
Chapter~6.
\end{example}

\section{Singular Value Decomposition}

Eigenvalue diagonalization applies only to endomorphisms.  Singular value
decomposition applies to a map between two possibly different spaces.  Its
conceptual path is
\[
\begin{gathered}
\text{adjoint}\longrightarrow T^*T
\longrightarrow\text{positive self-adjoint operator}\\
\longrightarrow\text{spectral theorem}\longrightarrow\text{SVD}.
\end{gathered}
\]
Here the adjoint constructed in the preceding section has the defining
property
\[
\langle Tv,w\rangle_W=\langle v,T^*w\rangle_V
\qquad(v\in V,\ w\in W).
\]
Thus, although \(T\) need not be an endomorphism, the two composites are
endomorphisms of different spaces:
\[
T^*T:V\longrightarrow V,
\qquad
TT^*:W\longrightarrow W.
\]

\begin{proposition}[The positive operators associated to a map]
For a linear map \(T:V\to W\), both \(T^*T\) and \(TT^*\) are self-adjoint
and positive.  In particular, all their eigenvalues are real and
nonnegative.
\end{proposition}
\begin{proof}
The generalized adjoint identities just proved give
\[
(T^*T)^*=T^*(T^*)^*=T^*T,
\qquad
(TT^*)^*=(T^*)^*T^*=TT^*.
\]
Moreover, for \(v\in V\) and \(w\in W\),
\[
\langle T^*Tv,v\rangle_V=\langle Tv,Tv\rangle_W
=\lVert Tv\rVert^2\geq0,
\]
and
\[
\langle TT^*w,w\rangle_W=\langle T^*w,T^*w\rangle_V
=\lVert T^*w\rVert^2\geq0.
\]
If, for example, \(T^*Tv=\lambda v\) with \(v\ne0\), then
\[
\lambda\lVert v\rVert^2=\langle T^*Tv,v\rangle_V
=\lVert Tv\rVert^2.
\]
The left side is real because \(T^*T\) is self-adjoint (or by the
eigenvalue geometry proposition), and the right side is nonnegative.
Hence \(\lambda\) is real and nonnegative.  The same argument applies to
\(TT^*\).
\end{proof}

\begin{proposition}[Nonzero spectra]
The operators \(T^*T\) and \(TT^*\) have the same nonzero eigenvalues, with
the same multiplicities.
\end{proposition}
\begin{proof}
Let \(\lambda>0\).  The map \(T\) sends the \(\lambda\)-eigenspace of
\(T^*T\) into the \(\lambda\)-eigenspace of \(TT^*\), since
\[
TT^*(Tv)=T(T^*Tv)=\lambda Tv.
\]
It is injective on this eigenspace: if \(Tv=0\), then
\(\lambda v=T^*Tv=0\).  Conversely, if \(TT^*w=\lambda w\), then
\[
T\left(\frac{T^*w}{\lambda}\right)=w.
\]
Moreover,
\[
T^*T\left(\frac{T^*w}{\lambda}\right)
=\frac{T^*(TT^*w)}{\lambda}
=\lambda\left(\frac{T^*w}{\lambda}\right),
\]
so this preimage belongs to the \(\lambda\)-eigenspace of \(T^*T\).
Thus \(T\) is an isomorphism between the two \(\lambda\)-eigenspaces.
Since both operators are self-adjoint and hence diagonalizable, the
dimensions of these eigenspaces are their multiplicities.
\end{proof}

\begin{definition}
List the eigenvalues of \(T^*T\), with multiplicity, as
\[
\lambda_1,\ldots,\lambda_n,
\qquad n=\dim V.
\]
The numbers
\[
\sigma_i=\sqrt{\lambda_i}
\]
are the \emph{singular values} of \(T\).  They are well-defined nonnegative
real numbers by the preceding proposition.  The nonzero singular values are
called the \emph{positive singular values}.
\end{definition}

The norm identity also records exactly where the zero singular values occur:
\[
\ker(T^*T)=\ker T.
\]
Indeed, \(Tv=0\) plainly implies \(T^*Tv=0\), while \(T^*Tv=0\) implies
\(\lVert Tv\rVert^2=0\), hence \(Tv=0\).  Consequently, the number of
positive singular values, counted with multiplicity, is
\[
\dim V-\dim\ker T=\operatorname{rank}T.
\]

\begin{theorem}[Singular value decomposition]
Let \(T:V\to W\) be a linear map of finite-dimensional real or complex
inner-product spaces.  Write
\[
n=\dim V,\qquad m=\dim W,\qquad r=\operatorname{rank}T.
\]
There are orthonormal bases
\[
(v_1,\ldots,v_n)\quad\text{of }V,
\qquad
(u_1,\ldots,u_m)\quad\text{of }W,
\]
and positive singular values \(\sigma_1,\ldots,\sigma_r\) such that
\[
Tv_i=\sigma_i u_i\quad(1\leq i\leq r),
\qquad
Tv_i=0\quad(r<i\leq n).
\]
Equivalently, relative to these bases the matrix of \(T\) is the rectangular
diagonal matrix \(\Sigma=(\Sigma_{ij})\) of size \(m\times n\), where
\[
\Sigma_{ij}=
\begin{cases}
\sigma_i,& i=j\leq r,\\
0,&\text{otherwise}.
\end{cases}
\]
\end{theorem}
\begin{proof}
By the real or complex spectral theorem, the positive self-adjoint operator
\(T^*T:V\to V\) has an orthonormal eigenbasis.  Arrange it so that
\[
T^*Tv_i=\sigma_i^2v_i\quad(1\leq i\leq r),
\]
where \(\sigma_i>0\), and so that \(v_{r+1},\ldots,v_n\) span the
zero-eigenspace.  The kernel identity above says that this zero-eigenspace
is \(\ker T\), so \(Tv_i=0\) for \(i>r\).

For \(1\leq i\leq r\), define
\[
u_i=\frac{Tv_i}{\sigma_i}\in W.
\]
These vectors are orthonormal.  Indeed, since the singular values are real,
\[
\begin{aligned}
\langle u_i,u_j\rangle_W
&=\frac{1}{\sigma_i\sigma_j}\langle Tv_i,Tv_j\rangle_W\\
&=\frac{1}{\sigma_i\sigma_j}\langle T^*Tv_i,v_j\rangle_V\\
&=\frac{\sigma_i}{\sigma_j}\langle v_i,v_j\rangle_V
=\delta_{ij}.
\end{aligned}
\]
Extend \(u_1,\ldots,u_r\) to an orthonormal basis
\(u_1,\ldots,u_m\) of \(W\).  The defining equation for the \(u_i\)'s and
the kernel calculation give the asserted formulas for \(Tv_i\), and these
formulas say exactly that the matrix is \(\Sigma\).
\end{proof}

For matrices, work in fixed orthonormal coordinate bases and write \(A=[T]\).
Let \(V\) and \(U\) be the matrices whose
columns are respectively the coordinate columns of \(v_1,\ldots,v_n\) and
\(u_1,\ldots,u_m\).  These matrices are unitary over \(\mathbb C\) and
orthogonal over \(\mathbb R\).  Since the matrix in the singular-vector
bases is \(\Sigma\),
\[
\Sigma=U^*AV,
\qquad\text{and hence}\qquad
A=U\Sigma V^*.
\]
Over \(\mathbb R\), this reads
\[
A=U\Sigma V^{\mathsf T}.
\]

\subsection{Spectral decomposition and SVD}

The case \(T:V\to V\) makes especially transparent both what SVD inherits
from the spectral theorem and what information it discards.

\begin{proposition}[Spectral decomposition]
Let \(T:V\to V\) be self-adjoint, let
\(\lambda_1,\ldots,\lambda_s\) be its distinct eigenvalues, and let
\(P_{\lambda_j}\) be the orthogonal projection onto
\(E_{\lambda_j}\).  Then
\[
T=\sum_{j=1}^s\lambda_jP_{\lambda_j},
\qquad
I=\sum_{j=1}^sP_{\lambda_j},
\]
and
\[
P_{\lambda_i}P_{\lambda_j}=0\quad(i\ne j),
\qquad
P_{\lambda_j}^*=P_{\lambda_j}.
\]
\end{proposition}
\begin{proof}
The real or complex spectral theorem gives the orthogonal direct sum
\[
V=\bigoplus_{j=1}^sE_{\lambda_j}.
\]
For \(v=\sum_jv_j\), with \(v_j\in E_{\lambda_j}\), the projections satisfy
\(P_{\lambda_j}v=v_j\).  Hence
\[
Tv=\sum_j\lambda_jv_j
=\sum_j\lambda_jP_{\lambda_j}v,
\]
and the remaining identities follow from the orthogonal direct sum and the
fact that orthogonal projections are self-adjoint.
\end{proof}

This formula is the coordinate-free form of unitary or orthogonal
diagonalization.  It also describes the SVD of a self-adjoint operator.  If
\(v_1,\ldots,v_n\) is an orthonormal eigenbasis with
\[
Tv_i=\lambda_iv_i,
\]
then \(T^*T=T^2\), so
\[
T^*Tv_i=\lambda_i^2v_i.
\]
Thus the corresponding singular values are
\[
\sigma_i=|\lambda_i|.
\]
For \(\lambda_i\ne0\), set
\[
u_i=\operatorname{sgn}(\lambda_i)v_i.
\]
Then \(u_i\) is a unit vector and
\[
Tv_i=\sigma_i u_i.
\]
When \(\lambda_i=0\), take \(u_i=v_i\); then again
\(Tv_i=0=\sigma_i u_i\).  The \(u_i\)'s form an orthonormal basis, so this
is an SVD.  The ordinary spectral decomposition retains the sign of every
real eigenvalue, whereas SVD replaces it by its absolute value.

If \(T\) is positive semidefinite as well as self-adjoint, then
\(\lambda_i\geq0\), and hence \(\sigma_i=\lambda_i\).  In this case one
may take \(u_i=v_i\): the SVD and spectral decomposition essentially
coincide.  At the matrix level, a Hermitian positive semidefinite matrix has
\[
A=U\Lambda U^*,
\qquad
\Lambda=\operatorname{diag}(\lambda_1,\ldots,\lambda_n),
\qquad
\lambda_i\geq0,
\]
and this is simultaneously a spectral decomposition and an SVD.

\begin{remark}[A brief polar-decomposition viewpoint]
For an endomorphism \(T\), the positive operator
\[
|T|=(T^*T)^{1/2}
\]
is defined by applying the spectral theorem to \(T^*T\).  Its eigenvalues
are the singular values of \(T\).  When \(T\) is self-adjoint, \(|T|\)
replaces each eigenvalue \(\lambda\) by \(|\lambda|\).  This is the
starting point for polar decomposition and another direct bridge between the
spectral theorem and SVD.
\end{remark}

\begin{example}[An SVD calculation]
Consider the real \(3\times2\) matrix
\[
A=\begin{pmatrix}1&1\\1&1\\0&0\end{pmatrix}.
\]
We first compute the operator on the domain:
\[
A^{\mathsf T}A
=\begin{pmatrix}1&1&0\\1&1&0\end{pmatrix}
\begin{pmatrix}1&1\\1&1\\0&0\end{pmatrix}
=\begin{pmatrix}2&2\\2&2\end{pmatrix}.
\]
Its characteristic polynomial is
\[
\det\begin{pmatrix}2-\lambda&2\\2&2-\lambda\end{pmatrix}
=\lambda(\lambda-4),
\]
so its eigenvalues are \(4\) and \(0\).  Solving the two eigenspace
equations gives
\[
(A^{\mathsf T}A-4I)\begin{pmatrix}x\\y\end{pmatrix}=0
\quad\Longleftrightarrow\quad x=y,
\qquad
A^{\mathsf T}A\begin{pmatrix}x\\y\end{pmatrix}=0
\quad\Longleftrightarrow\quad x=-y.
\]
Normalizing vectors in these two directions gives the orthonormal eigenbasis
\[
v_1=\frac1{\sqrt2}\begin{pmatrix}1\\1\end{pmatrix},
\qquad
v_2=\frac1{\sqrt2}\begin{pmatrix}1\\-1\end{pmatrix}.
\]
Thus the singular values are \(\sigma_1=\sqrt4=2\) and
\(\sigma_2=\sqrt0=0\).  The vectors \(v_1,v_2\) are the right singular
vectors; in particular, \(v_2\) spans \(\ker A\).

For the positive singular value, the corresponding left singular vector is
\[
u_1=\frac{Av_1}{2}
=\frac12\begin{pmatrix}\sqrt2\\\sqrt2\\0\end{pmatrix}
=\begin{pmatrix}1/\sqrt2\\1/\sqrt2\\0\end{pmatrix}.
\]
Complete it to an orthonormal basis of \(\mathbb R^3\), for example by
choosing
\[
u_2=\begin{pmatrix}1/\sqrt2\\-1/\sqrt2\\0\end{pmatrix},
\qquad
u_3=\begin{pmatrix}0\\0\\1\end{pmatrix}.
\]
Therefore
\[
U=\begin{pmatrix}
1/\sqrt2&1/\sqrt2&0\\
1/\sqrt2&-1/\sqrt2&0\\
0&0&1
\end{pmatrix},
\qquad
\Sigma=\begin{pmatrix}2&0\\0&0\\0&0\end{pmatrix},
\qquad
V=\begin{pmatrix}
1/\sqrt2&1/\sqrt2\\
1/\sqrt2&-1/\sqrt2
\end{pmatrix}.
\]
Both \(U\) and \(V\) are orthogonal.  Finally, direct multiplication gives
\[
U\Sigma V^{\mathsf T}
=2u_1v_1^{\mathsf T}
=2\begin{pmatrix}1/\sqrt2\\1/\sqrt2\\0\end{pmatrix}
\begin{pmatrix}1/\sqrt2&1/\sqrt2\end{pmatrix}
=\begin{pmatrix}1&1\\1&1\\0&0\end{pmatrix}=A.
\]
\end{example}

\begin{remark}[Why SVD matters]
The SVD is also a practical tool.  It underlies least-squares problems and
the Moore--Penrose pseudoinverse, best low-rank approximation, and principal
component analysis for dimensionality reduction.  Keeping only the largest
singular values can compress images and other data, while discarding small
ones is useful in denoising and signal processing.  Singular values also
measure numerical conditioning, so SVD is central in numerical linear
algebra and in many applications across computer science and engineering.
\end{remark}

\section*{Exercises}
\begin{exercise}[Basic] Verify that the standard formula on \(\mathbb C^n\)
is an inner product under the convention of this chapter.\end{exercise}
\begin{exercise}[Core] Use the complex polarization identity to recover the
standard inner product of two specified vectors in \(\mathbb C^2\) from their
norms.\end{exercise}
\begin{exercise}[Core] Apply Gram--Schmidt to
\[
(1,1,0),\ (1,0,1),\ (0,1,1)
\]
in \(\mathbb R^3\).\end{exercise}
\begin{exercise}[Core] Find the orthogonal projection of \((1,2,3)\) onto
the span of \((1,1,1)\), and verify its best-approximation property for an
arbitrary point on that line.\end{exercise}
\begin{exercise}[Core] Compute the adjoint of
\[
\begin{pmatrix}1&i\\0&2\end{pmatrix}
\]
on \(\mathbb C^2\), and determine its kernel.\end{exercise}
\begin{exercise}[Further] Prove that a positive-definite operator is
invertible.\end{exercise}
\begin{exercise}[Further, Challenging] Find an SVD of
\[
\begin{pmatrix}1&1\\0&0\end{pmatrix}
\]
and identify its nonzero singular subspaces.\end{exercise}

\chapter{Tensor Products of Modules}\label{chap:tensor-products}

Direct sums place modules side by side. Tensor products solve a different
problem: they turn expressions that are additive in two variables into
ordinary homomorphisms from one abelian group. The price is that sidedness now
matters. We begin over an arbitrary ring and only later specialize to the
more familiar commutative case.

\section{Balanced Maps and Binary Tensor Products}\label{sec:binary-tensors}
\begin{definition}
Let \(M_R\) be a right \(R\)-module, \({}_RN\) a left \(R\)-module, and
\(A\) an abelian group. A map \(b:M\times N\to A\) is \emph{\(R\)-balanced}
if
\[
\begin{aligned}
b(m+m',n)&=b(m,n)+b(m',n),\\
b(m,n+n')&=b(m,n)+b(m,n'),\\
b(mr,n)&=b(m,rn).
\end{aligned}
\]
\end{definition}

The last equation is the balancing condition. It says that the scalar
inserted between the two variables may be moved across their interface. This
is precisely what makes a product such as \(m\otimes n\) meaningful when the
two module actions lie on opposite sides.

\begin{theorem}[Tensor-product universal property]\label{thm:binary-tensor}
There is an abelian group \(M\otimes_RN\), with an \(R\)-balanced map
\[
\tau:M\times N\longrightarrow M\otimes_RN,
\qquad (m,n)\longmapsto m\otimes n,
\]
such that every \(R\)-balanced map \(b:M\times N\to A\) has a unique
homomorphism \(\bar b:M\otimes_RN\to A\) satisfying
\(\bar b(m\otimes n)=b(m,n)\).
\end{theorem}
\[
\begin{tikzcd}[column sep=large, row sep=large]
M\times N \arrow[r, "\tau"] \arrow[dr, "b"'] &
M\otimes_R N \arrow[d, "\bar b"] \\
& A
\end{tikzcd}
\]
\begin{proof}
Let \(G\) be the free abelian group on symbols \([m,n]\), and quotient by
the subgroup generated by
\[
\begin{aligned}
[m+m',n]-[m,n]-[m',n],\qquad
[m,n+n']-[m,n]-[m,n'],\\
[mr,n]-[m,rn].
\end{aligned}
\]
Call the quotient \(M\otimes_RN\), and write \(m\otimes n\) for the class of
\([m,n]\). The displayed relations make \(\tau\) balanced. Given balanced
\(b\), the map from \(G\) sending \([m,n]\) to \(b(m,n)\) kills every
listed relation, so it descends to \(\bar b\). Since pure tensors generate
the quotient, this homomorphism is unique.
\end{proof}

If two groups satisfy this universal property, applying each universal
property to the other universal balanced map gives mutually inverse
homomorphisms. Thus \(M\otimes_RN\) is unique up to a \emph{canonical
isomorphism}. It is not a Cartesian product, and an element need not be a
single pure tensor; it is a finite sum of pure tensors.

\section{Bimodules and Tensor Actions}\label{sec:bimodules-tensors}
\begin{definition}
For rings \(R,S\), an \emph{\((R,S)\)-bimodule} \({}_RM_S\) is an abelian
group carrying a left \(R\)-action and a right \(S\)-action such that
\[
(rm)s=r(ms)
\qquad(r\in R,\ m\in M,\ s\in S).
\]
\end{definition}

Every ring is an \((R,R)\)-bimodule over itself. A left \(R\)-module is an
\((R,\mathbb Z)\)-bimodule and a right \(R\)-module is an
\((\mathbb Z,R)\)-bimodule. Bimodules retain an action on each outside of a
tensor product.

\begin{proposition}[Induced outside actions]\label{prop:tensor-actions}
Let \({}_SM_R\), \({}_RN_T\) be bimodules. Then \(M\otimes_RN\) is an
\((S,T)\)-bimodule under
\[
s(m\otimes n)=(sm)\otimes n,
\qquad
(m\otimes n)t=m\otimes(nt).
\]
\end{proposition}
\begin{proof}
For fixed \(s\), the map \((m,n)\mapsto(sm)\otimes n\) is balanced; indeed,
\[
s(mr)\otimes n=((sm)r)\otimes n=(sm)\otimes(rn).
\]
The universal property therefore gives a well-defined endomorphism. The
argument for \(t\) is analogous. The module axioms follow on pure tensors,
and the actions commute because
\[
(s(m\otimes n))t=(sm)\otimes(nt)=s((m\otimes n)t).
\]
\end{proof}

\section{Finite Associativity}\label{sec:tensor-associativity}
Associativity of tensor products is a canonical isomorphism, not literal
equality. Only finitely many tensor factors are considered.

\begin{theorem}[Associativity]\label{thm:tensor-associativity}
For bimodules \({}_SM_R\), \({}_RN_T\), and \({}_TP\), there is a canonical
isomorphism of left \(S\)-modules
\[
(M\otimes_RN)\otimes_TP
\ \cong\
M\otimes_R(N\otimes_TP),
\]
given by
\[
(m\otimes n)\otimes p\longmapsto m\otimes(n\otimes p).
\]
\end{theorem}
\begin{proof}
The proposed formula respects the \(R\)-balancing relation because
\[
(mr\otimes n)\otimes p
\longmapsto mr\otimes(n\otimes p)
=m\otimes(rn\otimes p).
\]
It respects the outer \(T\)-balancing relation because
\[
(m\otimes nt)\otimes p
\longmapsto m\otimes(nt\otimes p)
=m\otimes(n\otimes tp).
\]
Additivity is immediate, so the two universal properties produce a
well-defined homomorphism. The reversed formula
\[
m\otimes(n\otimes p)\longmapsto(m\otimes n)\otimes p
\]
satisfies the same checks. Both composites fix pure generators, hence are
the identity.
\end{proof}

\section{Tensor Products and Arbitrary Direct Sums}\label{sec:tensor-direct-sums}
This is the intended infinite-index construction: tensor product commutes
with arbitrary \emph{direct sums}, not arbitrary direct products. The
finite-support condition makes the following formula meaningful.

\begin{theorem}[Arbitrary direct sums]\label{thm:tensor-direct-sums}
For right \(R\)-modules \((M_i)_{i\in I}\) and a left \(R\)-module \(N\),
\[
\left(\bigoplus_{i\in I}M_i\right)\otimes_RN
\ \cong\
\bigoplus_{i\in I}(M_i\otimes_RN).
\]
For a right \(R\)-module \(M\) and left \(R\)-modules \((N_i)_{i\in I}\),
\[
M\otimes_R\left(\bigoplus_{i\in I}N_i\right)
\ \cong\
\bigoplus_{i\in I}(M\otimes_RN_i).
\]
\end{theorem}
\begin{proof}
For the first map, send \((m_i)\otimes n\) to the finite sum
\[
\sum_{i\in I}m_i\otimes n.
\]
It is balanced, so the tensor universal property makes it a homomorphism.
For each canonical injection
\[
\iota_i:M_i\longrightarrow\bigoplus_{i\in I}M_i,
\]
the balanced map \((m,n)\mapsto\iota_i(m)\otimes n\) induces
\[
M_i\otimes_RN\longrightarrow
\left(\bigoplus_{i\in I}M_i\right)\otimes_RN.
\]
The direct-sum universal property combines these maps into the reverse
homomorphism. The two maps are inverse on pure tensors and on the injected
summands, hence everywhere. The second isomorphism follows by the identical
argument with the variables reversed.
\end{proof}

\section{The Commutative-Ring Specialization}\label{sec:commutative-tensors}
Suppose now that \(R\) is commutative. Every left \(R\)-module becomes
canonically an \((R,R)\)-bimodule by defining \(mr=rm\). Thus
\(M\otimes_RN\) is again an \(R\)-module, and
\[
r(m\otimes n)=(rm)\otimes n=m\otimes(rn).
\]

\begin{proposition}[Symmetry]\label{prop:tensor-symmetry}
For \(R\)-modules \(M,N\), there is a canonical isomorphism
\[
M\otimes_RN\cong N\otimes_RM,
\qquad
m\otimes n\longmapsto n\otimes m.
\]
\end{proposition}
\begin{proof}
The map \((m,n)\mapsto n\otimes m\) is balanced because
\[
n\otimes(mr)=(rn)\otimes m.
\]
The universal property gives the displayed map; swapping the factors again
gives its inverse. Tensor products are symmetric only up to canonical
isomorphism, not by literal equality.
\end{proof}

\section{Finite Multilinear Tensor Products}\label{sec:multilinear-tensors}
For a finite list \(M_1,\ldots,M_n\) of \(R\)-modules and an \(R\)-module
\(A\), a map
\[
b:M_1\times\cdots\times M_n\longrightarrow A
\]
is \emph{\(R\)-multilinear} if it is additive and respects scalar
multiplication in each variable. Define
\[
M_1\otimes_R\cdots\otimes_RM_n
\]
to be the quotient of the free \(R\)-module on tuples
\((m_1,\ldots,m_n)\) by the additivity and scalar-linearity relations in each
slot. Write the class as \(m_1\otimes\cdots\otimes m_n\).

\begin{theorem}[Finite multilinear universal property]\label{thm:finite-tensors}
For every finite \(n\), \(R\)-module homomorphisms
\[
M_1\otimes_R\cdots\otimes_RM_n\longrightarrow A
\]
are in bijection with \(R\)-multilinear maps
\[
M_1\times\cdots\times M_n\longrightarrow A.
\]
For \(n=3\), the direct tensor product is canonically isomorphic to both
\[
(M_1\otimes_RM_2)\otimes_RM_3
\quad\text{and}\quad
M_1\otimes_R(M_2\otimes_RM_3).
\]
\end{theorem}
\begin{proof}
The quotient construction makes the universal tuple map multilinear. A
multilinear map from the product kills its defining relations and therefore
factors uniquely through the quotient, proving the bijection. For \(n=3\),
the map
\[
m_1\otimes m_2\otimes m_3
\longmapsto
(m_1\otimes m_2)\otimes m_3
\]
is well-defined by multilinearity and has the evident inverse on pure
tensors. Theorem~\ref{thm:tensor-associativity} gives the other
parenthesization. Induction on \(n\) proves the finite case.
\end{proof}

This justifies suppressing parentheses in finite tensor products. Repeated
associativity and symmetry also let every permutation of finitely many factors
induce a canonical isomorphism. No countable or arbitrary infinite tensor
product is being defined here.

\section{Tensoring Homomorphisms and Concrete Calculations}\label{sec:tensor-maps-calculations}
If \(f:M_R\to M'_R\) and \(g:{}_RN\to{}_RN'\) are module homomorphisms, then
\[
(f\otimes g)(m\otimes n)=f(m)\otimes g(n)
\]
defines a homomorphism \(f\otimes g:M\otimes_RN\to M'\otimes_RN'\).
Indeed, the displayed expression is balanced, so the universal property
proves well-definedness. On pure tensors,
\[
(f'\otimes g')\circ(f\otimes g)=(f'f)\otimes(g'g),
\qquad
1_M\otimes1_N=1_{M\otimes_RN}.
\]

Return to a commutative ring \(R\). The following calculations anchor the
abstract construction.

\begin{proposition}\label{prop:tensor-calculations}
There are canonical isomorphisms
\[
R\otimes_RM\cong M,
\qquad
R^m\otimes_RR^n\cong R^{mn}.
\]
If \(I\) is an ideal, then
\[
(R/I)\otimes_RM\cong M/IM.
\]
\end{proposition}
\begin{proof}
The balanced map \((r,m)\mapsto rm\) induces \(R\otimes_RM\to M\), with
inverse \(m\mapsto1\otimes m\). For finite free modules, \(e_i\otimes f_j\)
map to the \((i,j)\)-th standard vector; expanding both factors proves that
this map and its evident inverse are mutually inverse. Finally,
\((r+I,m)\mapsto rm+IM\) is balanced and well-defined because changing \(r\)
by an element of \(I\) changes the value by an element of \(IM\). Its inverse
is \(m+IM\mapsto(1+I)\otimes m\); balancing kills every generator \(im\).
\end{proof}

\begin{example}
Taking \(R=\mathbb Z\) and \(I=(m)\) gives
\[
\mathbb Z/m\mathbb Z\otimes_{\mathbb Z}\mathbb Z/n\mathbb Z
\cong
\mathbb Z/\gcd(m,n)\mathbb Z.
\]
Indeed, the preceding proposition identifies this with the quotient of
\(\mathbb Z/n\mathbb Z\) by the subgroup generated by \(m\).
\end{example}

\section{The Tensor--Hom Correspondence}
The universal property of a tensor product can be tested against arbitrary
abelian groups.  This produces the basic tensor--Hom correspondence without
assuming commutativity of the ring.

\begin{theorem}[Tensor--Hom]\label{thm:tensor-hom}
Let \(M_R\) be a right \(R\)-module and let \({}_RN\) be a left \(R\)-module.
Let \(A\) be an abelian group.  Give
\[
\operatorname{Hom}_{\mathbb Z}(M,A)
\]
the left \(R\)-module structure
\[
(r\cdot f)(m)=f(mr).
\]
Then there is an isomorphism of abelian groups
\[
\operatorname{Hom}_{\mathbb Z}(M\otimes_RN,A)
\cong
\operatorname{Hom}_R\!\left(N,\operatorname{Hom}_{\mathbb Z}(M,A)\right).
\]
\end{theorem}
\begin{proof}
Given \(u:M\otimes_RN\to A\), set
\[
\widehat u(n)(m)=u(m\otimes n).
\]
This is additive in \(n\).  Balancing gives
\[
\widehat u(rn)(m)=u(m\otimes rn)=u(mr\otimes n)
=(r\cdot\widehat u(n))(m),
\]
so \(\widehat u\) is \(R\)-linear.  Conversely, for an \(R\)-linear map
\(v:N\to\operatorname{Hom}_{\mathbb Z}(M,A)\), the rule
\[
(m,n)\longmapsto v(n)(m)
\]
is additive in each variable and balanced, since
\[
v(n)(mr)=(r\cdot v(n))(m)=v(rn)(m).
\]
It therefore induces \(M\otimes_RN\to A\).  Evaluating each composite on a
pure tensor shows that the two constructions are inverse.
\end{proof}

\begin{proposition}[Compatibility of Tensor--Hom]\label{prop:tensor-hom-compat}
The correspondence in Theorem~\ref{thm:tensor-hom} has the following
explicit compatibilities.  If \(f:M'\to M\) is a map of right modules and
\(u:M\otimes_RN\to A\), then \(u\circ(f\otimes1_N)\) corresponds to
\[
n\longmapsto\widehat u(n)\circ f.
\]
If \(h:A\to A'\) is a homomorphism of abelian groups, then \(h\circ u\)
corresponds to
\[
n\longmapsto h\circ\widehat u(n).
\]
Finally, if \(\ell:N'\to N\) is a map of left modules, then
\(u\circ(1_M\otimes\ell)\) corresponds to \(\widehat u\circ\ell\).
\end{proposition}
\begin{proof}
For example, the first map sends \(n\) and \(m'\) to
\[
u(f(m')\otimes n)=\widehat u(n)(f(m')).
\]
The second and third identities follow in exactly the same way by evaluating
on \(m\otimes n\) and \(m\otimes n'\), respectively.  The displayed
formulas also show that the resulting maps have the required module
linearity.
\end{proof}

\begin{proposition}[Commutative-ring Tensor--Hom]\label{prop:comm-tensor-hom}
If \(R\) is commutative and \(M,N,P\) are \(R\)-modules, then
\[
\operatorname{Hom}_R(M\otimes_RN,P)
\cong
\operatorname{Hom}_R\!\left(M,\operatorname{Hom}_R(N,P)\right)
\cong
\operatorname{Hom}_R\!\left(N,\operatorname{Hom}_R(M,P)\right).
\]
\end{proposition}
\begin{proof}
An \(R\)-linear map \(M\otimes_RN\to P\) is the same, by the tensor-product
universal property, as an \(R\)-balanced bilinear map \(M\times N\to P\).
Sending such a map \(b\) to
\[
m\longmapsto\bigl(n\longmapsto b(m,n)\bigr)
\]
gives the first isomorphism, with inverse obtained by evaluation.  Switching
the two variables gives the second.  Commutativity is what identifies the
two sidedness conventions and makes each Hom group an \(R\)-module.
\end{proof}

\begin{remark}
The preceding compatibility identities are often summarized by saying that
the Tensor--Hom correspondence is natural in its variables.  The explicit
sidedness in Theorem~\ref{thm:tensor-hom} remains essential when \(R\) is
not commutative.
\end{remark}

\section{Right Exactness}\label{sec:tensor-right-exactness}
Tensor product preserves cokernels, which is stronger than merely preserving
surjections.  Indeed, an exact sequence
\[
A\longrightarrow B\longrightarrow C\longrightarrow0
\]
says that \(C\cong B/\operatorname{im}(A)\), so tensoring preserves this
quotient/cokernel even though it may destroy injectivity of \(A\to B\).

\begin{theorem}[Right exactness]\label{thm:tensor-right-exact}
Let
\[
A_R\xrightarrow{u}B_R\xrightarrow{v}C_R\longrightarrow0
\]
be exact, and let \({}_RN\) be a left \(R\)-module. Then
\[
A\otimes_RN\xrightarrow{u\otimes1}
B\otimes_RN\xrightarrow{v\otimes1}
C\otimes_RN\longrightarrow0
\]
is exact.
\end{theorem}
\begin{proof}
We prove that \(v\otimes1\) is a cokernel of \(u\otimes1\).  Let \(X\) be an
abelian group and let
\[
F:B\otimes_RN\longrightarrow X
\]
satisfy \(F(u\otimes1)=0\).  Theorem~\ref{thm:tensor-hom} identifies \(F\)
with
\[
\widehat F:N\longrightarrow\operatorname{Hom}_{\mathbb Z}(B,X).
\]
By Proposition~\ref{prop:tensor-hom-compat}, the vanishing condition says
\[
\widehat F(n)\circ u=0
\]
for every \(n\in N\).  Since \(C\) is the cokernel of \(u\), each
\(\widehat F(n)\) factors uniquely as
\[
\widehat F(n)=G_n\circ v
\]
for a homomorphism \(G_n:C\to X\).  The assignment is additive because
\[
\bigl(G_{n+n'}-G_n-G_{n'}\bigr)\circ v=0,
\]
so uniqueness of factorization through the cokernel gives
\[
G_{n+n'}=G_n+G_{n'}.
\]
For the left \(R\)-action on \(\operatorname{Hom}_{\mathbb Z}(C,X)\),
\[
(r\cdot G_n)\circ v
=r\cdot(G_n\circ v)
=r\cdot\widehat F(n)
=\widehat F(rn).
\]
Uniqueness similarly gives \(G_{rn}=r\cdot G_n\).  Therefore
\[
n\longmapsto G_n
\]
is an \(R\)-linear map \(N\to\operatorname{Hom}_{\mathbb Z}(C,X)\).
Applying Theorem~\ref{thm:tensor-hom} backward gives a unique
\[
G:C\otimes_RN\longrightarrow X
\]
with \(G(v(b)\otimes n)=F(b\otimes n)\).  Hence
\[
F=G\circ(v\otimes1),
\]
and uniqueness follows because pure tensors generate.  Thus \(v\otimes1\)
is the cokernel of \(u\otimes1\), which is exactly the asserted exactness.
\end{proof}

\begin{corollary}[Tensoring a quotient]\label{cor:tensor-quotient}
For a map \(u:A_R\to B_R\) and a left \(R\)-module \({}_RN\), there is a
canonical isomorphism
\[
\bigl(B/\operatorname{im}u\bigr)\otimes_RN
\cong
\bigl(B\otimes_RN\bigr)/\operatorname{im}(u\otimes1).
\]
\end{corollary}
\begin{proof}
Apply Theorem~\ref{thm:tensor-right-exact} to the quotient map
\[
B\longrightarrow B/\operatorname{im}u.
\]
The two modules displayed are cokernels of the same map \(u\otimes1\), and
the unique map induced by the quotient is the required canonical
isomorphism.
\end{proof}

No injectivity statement is implied. Tensoring
\[
2\mathbb Z\hookrightarrow\mathbb Z
\]
with \(\mathbb Z/2\mathbb Z\) gives a zero map from a nonzero module. Later
homological algebra measures this failure; the next definition isolates when
the failure does not occur.

\section{Flat Modules}
\begin{definition}
A left \(R\)-module \(N\) is \emph{flat} if tensoring with \(N\) takes every
short exact sequence of right \(R\)-modules to a short exact sequence.
\end{definition}
Equivalently, tensoring with \(N\) preserves injections, since
Theorem~\ref{thm:tensor-right-exact} already gives right exactness.

\begin{proposition}\label{prop:projective-flat}
Every free module is flat, and every projective module is flat.  Thus
\[
\text{free}\Longrightarrow\text{projective}\Longrightarrow\text{flat}.
\]
\end{proposition}
\begin{proof}
If \(F\) is free on \(X\), then, for every right \(R\)-module \(A\),
\[
A\otimes_R F\cong\bigoplus_{x\in X}A.
\]
Consequently tensoring a short exact sequence of right modules with \(F\) is a direct sum of
short exact sequences and remains exact.  Thus free modules are flat.

By Theorem~\ref{thm:projective-characterizations}, a projective module \(P\)
is a direct summand of a free module \(F=P\oplus P'\).  Tensoring distributes
over this finite direct sum, so a tensor map involving \(F\) is the direct
sum of the corresponding tensor maps involving \(P\) and \(P'\).  If the
first is injective, each summand map is injective; if the middle sequence is
exact, exactness holds componentwise.  Since \(F\) is flat, \(P\) is flat.
\end{proof}
The converses need not hold in general.

\section{Exterior Powers and Exterior Algebra}\label{sec:exterior-algebra}
Finite tensor products represent multilinear maps.  Exterior powers impose
one further relation: a multilinear expression should vanish whenever two
inputs agree.  This packages alternating forms, oriented volume, and the
determinant into one construction.

Throughout this section \(R\) is a commutative ring and \(M\) is an
\(R\)-module.  A multilinear map
\[
f:M^k\longrightarrow N
\]
is \emph{alternating} if
\[
f(m_1,\ldots,m_k)=0
\]
whenever two arguments coincide.  This is the same alternating condition
used for determinants in \hyperref[sec:determinants]{Section~6.12}; it is
not to be confused with skew-symmetry in characteristic \(2\).

\begin{theorem}[Exterior-power universal property]\label{thm:exterior-power}
For every integer \(k\geq1\), there is an \(R\)-module
\[
\bigwedge_R^kM
\]
and an alternating multilinear map
\[
\omega_k:M^k\longrightarrow\bigwedge_R^kM,
\qquad
(m_1,\ldots,m_k)\longmapsto m_1\wedge\cdots\wedge m_k,
\]
such that every alternating multilinear map
\[
f:M^k\longrightarrow N
\]
has a unique linear factorization
\[
\widetilde f:\bigwedge_R^kM\longrightarrow N.
\]
\end{theorem}
\[
\begin{tikzcd}[column sep=large, row sep=large]
M^k \arrow[r, "\omega_k"] \arrow[dr, "f"'] &
\bigwedge_R^k M \arrow[d, "\widetilde f"] \\
& N
\end{tikzcd}
\]
\begin{proof}
By Theorem~\ref{thm:finite-tensors}, form the finite tensor power
\[
M^{\otimes k}
=M\otimes_R\cdots\otimes_RM.
\]
Let \(J\) be the submodule generated by all pure tensors having two equal
entries, and define
\[
\bigwedge_R^kM=M^{\otimes k}/J.
\]
The quotient of the universal multilinear tensor map is \(\omega_k\), and
it is alternating by construction.  If \(f\) is alternating, its associated
linear map
\[
M^{\otimes k}\longrightarrow N
\]
vanishes on every generator of \(J\), so it descends uniquely through the
quotient.  This gives the required factorization.  Any two modules with this
property are canonically isomorphic by applying their universal properties
to one another's universal alternating maps.
\end{proof}

\begin{proposition}[Alternating wedge identities]\label{prop:wedge-signs}
If two entries agree, then
\[
m_1\wedge\cdots\wedge m_k=0.
\]
Interchanging two entries changes sign; consequently
\[
m_{\sigma(1)}\wedge\cdots\wedge m_{\sigma(k)}
=\operatorname{sgn}(\sigma)\,
m_1\wedge\cdots\wedge m_k
\]
for every \(\sigma\in S_k\).
\end{proposition}
\begin{proof}
The first assertion is part of the construction.  Holding all other entries
fixed, alternation and multilinearity give
\[
0=\omega_k(\ldots,x+y,\ldots,x+y,\ldots).
\]
The repeated-\(x\) and repeated-\(y\) terms vanish, leaving the sum of the
two wedges with \(x\) and \(y\) interchanged.  Thus a transposition
contributes \(-1\).  The sign theory of Chapter~2 then gives the formula for
an arbitrary permutation.
\end{proof}

\noindent\textbf{Example.} If \(M=R^3\) has basis \(e_1,e_2,e_3\), then
\[
\bigwedge_R^2M
\]
has basis
\[
e_1\wedge e_2,\qquad e_1\wedge e_3,\qquad e_2\wedge e_3.
\]
Terms \(e_i\wedge e_i\) vanish, and reversing two basis vectors changes the
sign, so these are precisely the increasing-index wedges.

\begin{theorem}[Basis of an exterior power]\label{thm:exterior-basis}
If \(M\) is finite free with ordered basis
\[
(e_1,\ldots,e_n),
\]
then the wedges
\[
e_{i_1}\wedge\cdots\wedge e_{i_k},
\qquad
1\leq i_1<\cdots<i_k\leq n,
\]
form a basis of
\[
\bigwedge_R^kM.
\]
Hence
\[
\operatorname{rank}_R\bigl(\bigwedge_R^kM\bigr)=\binom nk.
\]
\end{theorem}
\begin{proof}
Multilinearity and Proposition~\ref{prop:wedge-signs} show that these wedges
span: expand every argument in the basis, discard repeated-index terms, and
reorder the rest.  To prove independence, fix an increasing \(k\)-tuple
\[
I=(i_1,\ldots,i_k).
\]
For \(m_j=\sum_ra_{rj}e_r\), define
\[
f_I(m_1,\ldots,m_k)=
\det(a_{i_rj})_{1\leq r,j\leq k}.
\]
The determinant theory of
\hyperref[sec:determinants]{Section~6.12} shows that \(f_I\) is alternating
and multilinear, so it induces a linear map from \(\bigwedge_R^kM\) to
\(R\).  On the displayed wedges it has value \(1\) for the tuple \(I\) and
\(0\) for every other increasing tuple.  Applying these maps to a linear
relation proves that every coefficient is zero.
\end{proof}

We set
\[
\bigwedge_R^0M=R,
\qquad
\bigwedge_R^1M=M.
\]
The basis theorem gives
\[
\bigwedge_R^kM=0
\qquad(k>n)
\]
when \(M\) is finite free of rank \(n\), and it shows that
\[
\bigwedge_R^nM
\]
is free of rank one.

\begin{proposition}[Induced exterior maps]\label{prop:exterior-maps}
An \(R\)-linear map
\[
T:M\longrightarrow N
\]
induces a unique linear map
\[
\bigwedge^kT:\bigwedge_R^kM\longrightarrow\bigwedge_R^kN
\]
satisfying
\[
\bigwedge^kT(m_1\wedge\cdots\wedge m_k)
=Tm_1\wedge\cdots\wedge Tm_k.
\]
Moreover,
\[
\bigwedge^k(\operatorname{id}_M)
=\operatorname{id}_{\bigwedge_R^kM},
\qquad
\bigwedge^k(S\circ T)
=\bigl(\bigwedge^kS\bigr)\circ\bigl(\bigwedge^kT\bigr).
\]
\end{proposition}
\begin{proof}
The map
\[
(m_1,\ldots,m_k)\longmapsto Tm_1\wedge\cdots\wedge Tm_k
\]
is alternating and multilinear.  By Theorem~\ref{thm:exterior-power}, it
therefore induces a unique map.  The two displayed identities agree
on every pure wedge, and pure wedges generate the relevant exterior powers.
\end{proof}

\begin{theorem}[Top exterior action]
Let \(V\) be \(n\)-dimensional over \(F\), and let \(T:V\to V\).  Then
\[
\bigwedge^nT
\]
acts on the one-dimensional space
\[
\bigwedge_F^nV
\]
by multiplication by \(\det(T)\).  Equivalently,
\[
Tv_1\wedge\cdots\wedge Tv_n
=\det(T)\,
\bigl(v_1\wedge\cdots\wedge v_n\bigr).
\]
\end{theorem}
\begin{proof}
Fix an ordered basis \((v_1,\ldots,v_n)\).  The linear functional on
\(\bigwedge_F^nV\) that takes \(v_1\wedge\cdots\wedge v_n\) to \(1\)
corresponds, by the exterior universal property, to a normalized alternating
\(n\)-linear form on \(V\).  By the uniqueness theorem for determinants in
\hyperref[sec:determinants]{Section~6.12}, its value on
\[
(Tv_1,\ldots,Tv_n)
\]
is \(\det(T)\).  Since the top exterior power is one-dimensional, this is
exactly the scalar by which \(\bigwedge^nT\) acts.
\end{proof}

This gives a later, equivalent characterization of the determinant already
developed from normalized alternating multilinear forms in
\hyperref[sec:determinants]{Section~6.12}.  Since
\[
\dim_F\bigwedge_F^nV=1,
\]
the determinant of \(T\) is the unique scalar \(\lambda\in F\) such that
\[
\bigwedge^nT
=\lambda\operatorname{id}_{\bigwedge_F^nV}.
\]
Equivalently, for every ordered basis \((v_1,\ldots,v_n)\) of \(V\),
\[
Tv_1\wedge\cdots\wedge Tv_n
=\det(T)\bigl(v_1\wedge\cdots\wedge v_n\bigr).
\]
Thus this top-exterior-power viewpoint describes the same invariant as the
Leibniz and alternating-multilinear descriptions, rather than defining a
new determinant.

For \(A\in M_n(F)\), regarded as an endomorphism of \(F^n\), the matrix of
\[
\bigwedge^nA
\]
with respect to the one-element basis
\[
e_1\wedge\cdots\wedge e_n
\]
of \(\bigwedge_F^nF^n\) is the \(1\times1\) matrix
\[
(\det A).
\]
Moreover, the composition law proved in Proposition~\ref{prop:exterior-maps}
makes determinant multiplicativity conceptually immediate:
\[
\bigwedge^n(S\circ T)
=\bigl(\bigwedge^nS\bigr)\circ\bigl(\bigwedge^nT\bigr)
\]
implies
\[
\det(S\circ T)=\det(S)\det(T).
\]
This is a retrospective explanation of the multiplicativity already proved
in Chapter~6, not a replacement for that earlier proof.

\section{The Exterior Algebra and Alternating Forms}\label{sec:exterior-forms}
The individual exterior powers fit together into a graded algebra.  This
uses an arbitrary direct sum of already constructed modules; it does
\emph{not} introduce an infinite tensor product.

\begin{definition}
The \emph{exterior algebra} of \(M\) is
\[
\bigwedge M
=\bigoplus_{k\geq0}\bigwedge_R^kM.
\]
For pure wedges, define the product by concatenation:
\[
\begin{aligned}
&(m_1\wedge\cdots\wedge m_p)
\wedge
(n_1\wedge\cdots\wedge n_q)\\
&\qquad =
m_1\wedge\cdots\wedge m_p
\wedge n_1\wedge\cdots\wedge n_q.
\end{aligned}
\]
\end{definition}

The right-hand side is alternating separately in the \(m\)'s and in the
\(n\)'s, so the two exterior universal properties make this a well-defined
bilinear map
\[
\bigwedge_R^pM\times\bigwedge_R^qM
\longrightarrow
\bigwedge_R^{p+q}M.
\]
Pure wedges generate each factor, so associativity follows from associativity
of concatenation on pure wedges.  The degree-\(k\) summand is the \(k\)-th
graded piece.

\begin{proposition}[Graded commutativity]
For homogeneous elements
\[
\alpha\in\bigwedge_R^pM,
\qquad
\beta\in\bigwedge_R^qM,
\]
one has
\[
\alpha\wedge\beta=(-1)^{pq}\beta\wedge\alpha.
\]
\end{proposition}
\begin{proof}
For pure wedges, moving the block of \(p\) entries past the block of \(q\)
entries requires \(pq\) transpositions.  Proposition~\ref{prop:wedge-signs}
gives the sign.  Bilinearity extends the identity to arbitrary homogeneous
elements.  The argument uses alternation itself, so it remains valid in
characteristic \(2\), where the displayed sign may equal \(1\).
\end{proof}

\begin{theorem}[Alternating forms as exterior powers of the dual]
Let \(V\) be a finite-dimensional vector space over \(F\).  There is a
natural isomorphism from
\[
\bigwedge_F^kV^\vee
\]
onto the space of alternating \(k\)-linear maps
\[
V^k\longrightarrow F.
\]
On a decomposable element it is given by
\[
\varphi_1\wedge\cdots\wedge\varphi_k
\longmapsto
\bigl((v_1,\ldots,v_k)\mapsto
\det(\varphi_i(v_j))_{1\leq i,j\leq k}\bigr).
\]
\end{theorem}
\begin{proof}
The displayed determinant is alternating and multilinear in the vectors, so
the exterior universal property makes the assignment well-defined and
linear.  If \((v_1,\ldots,v_n)\) is a basis of \(V\), then the wedges of its
dual basis indexed by increasing \(k\)-tuples map to the corresponding
coordinate alternating forms.  These forms are linearly independent and
span by evaluating an alternating form on basis tuples.  Thus the map is an
isomorphism.
\end{proof}

This identifies the algebra developed here with the pointwise algebra behind
differential forms.  On a smooth manifold, a differential \(k\)-form assigns
to each point \(p\) an alternating \(k\)-linear form on the tangent space
\[
T_pM,
\]
that is, an element of
\[
\bigwedge^k(T_pM)^\vee.
\]
The wedge product sends a \(p\)-form and a \(q\)-form at \(p\) to an element
of
\[
\bigwedge^{p+q}(T_pM)^\vee.
\]
No manifold theory is needed here; the point is that exterior algebra is the
linear-algebraic model that makes those geometric objects possible.


\section*{Exercises}
\begin{exercise}[Basic] Verify that multiplication \(R\times R\to R\) is
balanced when the first copy is the regular right module and the second the
regular left module.\end{exercise}
\begin{exercise}[Core] Give an \((R,R)\)-bimodule structure on \(R^n\), and
verify the induced actions of Proposition~\ref{prop:tensor-actions}.
\end{exercise}
\begin{exercise}[Core] Construct explicitly the associativity isomorphism
for \(\mathbb Z/2\mathbb Z\), \(\mathbb Z/3\mathbb Z\), and
\(\mathbb Z/4\mathbb Z\).\end{exercise}
\begin{exercise}[Core] Prove the second arbitrary-direct-sum isomorphism in
Theorem~\ref{thm:tensor-direct-sums} directly.\end{exercise}
\begin{exercise}[Core] Compute \((R/(a))\otimes_R(R/(b))\) for a commutative
ring \(R\).\end{exercise}
\begin{exercise}[Further] Verify the universal property of the three-factor
tensor product directly from the quotient construction.\end{exercise}
\begin{exercise}[Further, Challenging] Show that tensoring the inclusion
\(2\mathbb Z\hookrightarrow\mathbb Z\) with \(\mathbb Z/2\mathbb Z\) is not
injective.\end{exercise}
\begin{exercise}[Core] If \(M\) is free with basis \(e_1,e_2,e_3\), give a
basis for
\[
\bigwedge_R^2M.
\]
For a diagonal endomorphism of \(F^3\), compute its induced map on this
exterior power and recover its determinant from the top exterior power.
\end{exercise}

\part{Fields and Galois Theory}
\chapter{Field Extensions}

Many equations require a larger field before they can be solved:
\(\mathbb Q(\sqrt2)\) contains a root of \(x^2-2\), and \(\mathbb C\)
contains a root of \(x^2+1\) over \(\mathbb R\).  Extensions also let us
place several related quantities in a single field and, later, enlarge the
scalars on which a vector space is considered.  This chapter makes adjoining
elements precise and establishes the results on which separability and Galois
theory later depend.

\section{Extensions and Degree}\label{sec:extensions-degree}
The vector-space language developed in Chapter~5 is therefore immediately
available: once \(F\subseteq E\), scalar multiplication by elements of \(F\)
makes \(E\) into an \(F\)-vector space, and degree is simply its dimension.
For instance, the common field \(\mathbb Q(\sqrt2,\sqrt3)\) is where the two
square roots can be added and multiplied while still retaining
\(\mathbb Q\) as the base field.

\begin{definition}
An \emph{extension field} \(E/F\), written \(F\subseteq E\), is a field
\(E\) containing \(F\) as a subfield.  Then \(E\) is an \(F\)-vector space.
Its \emph{degree} is
\[
[E:F]=\dim_F E.
\]
The extension is \emph{finite} when this dimension is finite.
\; If
\[
F\subseteq K\subseteq E,
\]
then \(K\) is an \emph{intermediate field} of \(E/F\).
\end{definition}

This definition deliberately reuses familiar linear algebra.  The field
operations in \(E\) are richer than vector-space operations, but the smaller
field \(F\) may still be used as a field of scalars.  Thus the degree measures
how many independent \(F\)-directions must be added in passing from \(F\) to
\(E\).  For example,
\[
[\mathbb Q(\sqrt2):\mathbb Q]=2,
\qquad
[\mathbb C:\mathbb R]=2,
\]
with bases \(1,\sqrt2\) and \(1,i\), respectively.  For elements in a common
ambient field, \(F(S)\) denotes the smallest subfield containing \(F\cup S\).
The inclusion \(\mathbb Q\subseteq\mathbb R\), by contrast, has infinite
degree: no finite list of real numbers spans all of \(\mathbb R\) over
\(\mathbb Q\).  The degree is therefore a genuine measure of the amount of
new scalar information in the extension, not merely a count of symbols.
The infinite extension \(F(t)/F\), formed by adjoining an indeterminate, will
appear in Section~\ref{sec:algebraic-transcendental}; its powers
\(1,t,t^2,\ldots\) cannot fit into any finite
basis.

The Tower Law is the numerical bookkeeping principle for successive
adjunctions.  It says that a basis for the top field is obtained by first
choosing coordinates over the middle field and then expanding each of those
coordinates over the base field.

\begin{theorem}[Tower Law]\label{thm:tower-law}
If
\[
F\subseteq E\subseteq L
\]
and \([E:F]\), \([L:E]\) are finite, then
\[
[L:F]=[L:E][E:F].
\]
\end{theorem}
\begin{proof}
Let \(e_1,\ldots,e_m\) be an \(F\)-basis of \(E\) and
\(u_1,\ldots,u_n\) an \(E\)-basis of \(L\).  The products \(e_i u_j\) span:
expand first in the \(u_j\), then expand each coefficient in the \(e_i\).
If
\[
\sum_{i,j}c_{ij}e_i u_j=0,
\]
then grouping by \(u_j\) and using their independence gives
\[
\sum_i c_{ij}e_i=0
\]
for each \(j\).  Independence of the \(e_i\) makes all \(c_{ij}\) zero.
Thus the products are a basis of size \(mn\).
\end{proof}

\begin{corollary}
If \(F\subseteq E\subseteq L\) and \([L:F]\) is finite, then
\([E:F]\) divides \([L:F]\).  In particular, a prime-degree extension has
no proper intermediate field.
\end{corollary}
\begin{proof}
Apply the Tower Law and factor the prime integer \([L:F]\).
\end{proof}

\section{Algebraic and Transcendental Elements}
\label{sec:algebraic-transcendental}
\begin{definition}
For an extension \(E/F\) and \(\alpha\in E\), the element \(\alpha\) is
\emph{algebraic over \(F\)} if \(f(\alpha)=0\) for some nonzero
\(f\in F[x]\); otherwise it is \emph{transcendental}.  The extension is
\emph{algebraic} when each of its elements is algebraic.
\end{definition}

The evaluation homomorphism
\[
\operatorname{ev}_\alpha:F[x]\longrightarrow E,
\qquad
f\longmapsto f(\alpha)
\]
collects all relations on \(\alpha\).  If \(\alpha\) is algebraic, its
nonzero kernel has a unique monic generator, the \emph{minimal polynomial}
\(m_{\alpha,F}\).

Thus the minimal polynomial is not an arbitrary equation satisfied by
\(\alpha\): it is the irreducible relation of least degree, and it records
exactly the algebraic information lost when the indeterminate \(x\) is
replaced by \(\alpha\).  The next proposition explains why every other
polynomial relation is a consequence of this one.

\begin{proposition}[Minimal-polynomial properties]\label{thm:minimal-polynomial}
The minimal polynomial \(m_{\alpha,F}\) is irreducible and divides every
polynomial in \(F[x]\) which vanishes at \(\alpha\).
\end{proposition}
\begin{proof}
If \(m=gh\), then \(g(\alpha)h(\alpha)=0\), so one factor belongs to the
evaluation kernel.  Minimal positive degree forces the other factor to be a
unit.  For a polynomial \(f\) vanishing at \(\alpha\), divide by \(m\):
\[
f=qm+r,
\qquad
\deg r<\deg m.
\]
Then \(r(\alpha)=0\), so \(r=0\) by minimality.
\end{proof}

\begin{definition}
\(F[\alpha]\) is the subring consisting of polynomial expressions in
\(\alpha\), while \(F(\alpha)\) is the smallest subfield containing \(F\)
and \(\alpha\).
\end{definition}

\begin{theorem}[Simple algebraic extensions]\label{thm:simple-algebraic}
If \(\alpha\) is algebraic and \(d=\deg m_{\alpha,F}\), then
\[
F[\alpha]\cong F[x]/(m_{\alpha,F})=F(\alpha).
\]
The elements
\[
1,\alpha,\ldots,\alpha^{d-1}
\]
are an \(F\)-basis of \(F(\alpha)\), so
\[
[F(\alpha):F]=d.
\]
\end{theorem}
\begin{proof}
Evaluation and the First Isomorphism Theorem give the quotient.  Since the
minimal polynomial is irreducible, that quotient is a field, and therefore
\(F[\alpha]=F(\alpha)\).  Division by the minimal polynomial gives spanning
by the displayed powers.  A relation among them is a smaller nonzero
polynomial in the evaluation kernel, which is impossible.
\end{proof}

\begin{example}[Two square roots]\label{ex:two-square-roots}
The simple-extension theorem now lets us complete the calculation suggested
in Section~\ref{sec:extensions-degree}.  The polynomial \(x^2-3\) has no root in
\(\mathbb Q(\sqrt2)\).  Indeed, every element of this field has the form
\(a+b\sqrt2\) with \(a,b\in\mathbb Q\).  If
\[
(a+b\sqrt2)^2=3,
\]
comparison of the rational and \(\sqrt2\) parts gives
\[
a^2+2b^2=3,
\qquad
2ab=0.
\]
If \(a=0\), then \(b^2=3/2\), impossible in \(\mathbb Q\); if \(b=0\), then
\(a^2=3\), also impossible in \(\mathbb Q\).  Thus \(x^2-3\) is
irreducible over \(\mathbb Q(\sqrt2)\), so it is the minimal polynomial of
\(\sqrt3\) over that field.  The theorem gives
\[
[\mathbb Q(\sqrt2,\sqrt3):\mathbb Q(\sqrt2)]=2.
\]
The Tower Law therefore yields
\[
[\mathbb Q(\sqrt2,\sqrt3):\mathbb Q]
=2\cdot2=4.
\]
This is the first synthesis of irreducibility, minimal polynomials, simple
extension degree, and the Tower Law.  A second square root increases the
degree precisely when it was absent from the first extension.
\end{example}

\begin{proposition}[The transcendental case]
If \(t\) is transcendental over \(F\), then
\[
F[t]\cong F[x],
\qquad
F(t)\cong F(x).
\]
In particular \(F(t)/F\) is infinite.
\end{proposition}
\begin{proof}
The evaluation kernel is zero.  Passing from the polynomial ring to its
field of fractions gives the second isomorphism, and the powers of \(t\) are
linearly independent over \(F\).
\end{proof}

The algebraic and transcendental cases have opposite shapes.  An algebraic
element has a power that can be reduced to a combination of finitely many
earlier powers, whereas the powers of a transcendental element never satisfy
such a relation.  Passing from \(F[t]\) to \(F(t)\) adds fractions, just as
passing from \(\mathbb Z\) to \(\mathbb Q\) adds quotients; it does not create
a polynomial relation for \(t\).
\[
\begin{array}{c|c|c}
& \text{algebraic }\alpha & \text{transcendental }t\\ \hline
\text{polynomial relations} & \text{generated by }m_{\alpha,F} & \text{none}\\
\text{associated field} & F(\alpha)\cong F[x]/(m_{\alpha,F}) & F(t)\cong F(x)\\
\text{degree over }F & \text{finite} & \text{infinite}
\end{array}
\]

\begin{example}
Eisenstein's criterion makes \(x^3-2\) the minimal polynomial of
\(\sqrt[3]2\) over \(\mathbb Q\).  Thus
\[
\mathbb Q(\sqrt[3]2)
=\mathbb Q\oplus\mathbb Q\sqrt[3]2\oplus\mathbb Q\sqrt[3]4.
\]
In contrast, \(F(x)\) contains rational functions not in \(F[x]\).
\end{example}

\begin{definition}
An extension is \emph{finitely generated} if it has the form
\[
F(\alpha_1,\ldots,\alpha_r).
\]
\end{definition}

\begin{proposition}\label{prop:algebraic-towers}
Finite extensions are algebraic, finitely generated algebraic extensions are
finite, and algebraicity is transitive through towers.
\end{proposition}
\begin{proof}
First suppose \([E:F]=n\), and let \(a\in E\).  The \(n+1\) elements
\[
1,a,\ldots,a^n
\]
are linearly dependent over \(F\).  Their dependence relation is a nonzero
polynomial over \(F\) which kills \(a\), proving that \(E/F\) is algebraic.

Next let
\[
E=F(\alpha_1,\ldots,\alpha_r)
\]
with every \(\alpha_i\) algebraic over \(F\).  Each \(\alpha_i\) remains
algebraic over the intermediate field already constructed, because its
polynomial over \(F\) is also a polynomial over that field.  Thus every
successive extension is finite by Theorem~\ref{thm:simple-algebraic}; repeated
use of the Tower Law makes \(E/F\) finite.

Finally suppose \(F\subseteq E\subseteq L\), \(E/F\) is algebraic, and
\(\alpha\in L\) is algebraic over \(E\).  Write a monic relation
\[
\alpha^n+a_{n-1}\alpha^{n-1}+\cdots+a_0=0,
\qquad
a_i\in E,
\]
and put
\[
K=F(a_0,\ldots,a_{n-1}).
\]
The elements \(a_i\) are algebraic over \(F\), so \(K/F\) is finite by the
second part.  The displayed relation shows that \(\alpha\) is algebraic over
\(K\); hence \(K(\alpha)/K\) is finite.  The Tower Law gives
\[
[K(\alpha):F]<\infty.
\]
Its first part now says that \(\alpha\) is algebraic over \(F\).  This proves
transitivity.
\end{proof}

\section{Root Adjunction and Splitting Fields}
When an irreducible polynomial has no root in the current field, this is not
an obstruction to doing algebra with a root.  The quotient construction below
creates the smallest possible field in which a formal symbol is required to
satisfy that polynomial.

\begin{theorem}[Root adjunction]\label{thm:root-adjunction}
If \(f\in F[x]\) is irreducible, then
\[
F[x]/(f)
\]
is a field containing \(F\), and \(x+(f)\) is a root of \(f\).  Its minimal
polynomial is \(f\), and its degree over \(F\) is \(\deg f\).
\end{theorem}
\begin{proof}
The quotient is the formal way to impose the single relation \(f(x)=0\).
Because \(f\) is irreducible, the
\hyperref[thm:irreducible-maximal-field]{irreducible-quotient theorem of
Chapter~4} makes the principal ideal \((f)\) maximal; hence the quotient is a
field rather than merely a ring.
Constants embed as their residue classes, and the class
\[
\alpha=x+(f)
\]
is the formal adjoined root since \(f(\alpha)=f(x)+(f)=0\).  The evaluation
kernel is \((f)\); the remaining claims follow from
Theorem~\ref{thm:simple-algebraic}.
\end{proof}

\begin{example}[A root made concrete]
Taking \(f=x^2-2\in\mathbb Q[x]\) gives
\[
\mathbb Q[x]/(x^2-2).
\]
If \(\alpha=x+(x^2-2)\), then \(\alpha^2=2\), and every residue class has
a unique form
\[
a+b\alpha,
\qquad a,b\in\mathbb Q.
\]
For instance,
\[
(a+b\alpha)(c+d\alpha)=(ac+2bd)+(ad+bc)\alpha.
\]
Sending \(\alpha\) to the usual real number \(\sqrt2\) identifies this
quotient with \(\mathbb Q(\sqrt2)\).  The quotient construction is useful
because it works before a convenient ambient field, such as \(\mathbb R\),
has been chosen.
\end{example}

\begin{lemma}[Extension of embeddings]\label{lem:root-extension}
Let \(\sigma:K\to\Omega\) be an \(F\)-embedding.  Let \(\alpha\) be
algebraic over \(K\), with minimal polynomial \(m\).  Extensions
\[
\widetilde\sigma:K(\alpha)\longrightarrow\Omega
\]
of \(\sigma\) correspond bijectively to roots in \(\Omega\) of the
coefficientwise image \(\sigma(m)\).
\end{lemma}
\begin{proof}
Write
\[
m(x)=a_0+a_1x+\cdots+a_nx^n,
\qquad
\sigma(m)(x)=\sigma(a_0)+\sigma(a_1)x+\cdots+\sigma(a_n)x^n.
\]
For a root \(\beta\) of this transported polynomial, the map
\[
\sum a_ix^i\longmapsto\sum\sigma(a_i)\beta^i
\]
has \(m\) in its kernel and descends to a homomorphism from
\[
K[x]/(m)\cong K(\alpha)
\]
to \(\Omega\).  This is well-defined precisely because the defining relation
\(m(\alpha)=0\) is respected by the chosen root \(\beta\).  The source is a
field, so it is injective.  It is uniquely determined by its values on \(K\)
and \(\alpha\).  Conversely an extension
sends \(m(\alpha)=0\) to
\[
\sigma(m)(\widetilde\sigma(\alpha))=0.
\]
\end{proof}

\begin{definition}
A \emph{splitting field} of a nonconstant polynomial \(f\in F[x]\) is an
extension generated over \(F\) by the roots of \(f\), in which \(f\) splits
into linear factors.
\end{definition}

\begin{theorem}[Splitting fields]\label{thm:splitting-existence}
Every nonconstant polynomial in \(F[x]\) has a splitting field, unique up to
an \(F\)-isomorphism.
\end{theorem}
\begin{proof}
For existence, induct on \(n=\deg f\).  There is nothing to prove for
\(n=1\).  Otherwise choose an irreducible factor \(p\) of \(f\), and use
Theorem~\ref{thm:root-adjunction} to obtain a field \(K/F\) containing a root
\(\alpha\) of \(p\).  In \(K[x]\), division gives
\[
f(x)=(x-\alpha)g(x),
\qquad
\deg g=n-1.
\]
By induction, \(g\) has a splitting field \(L/K\).  Then \(f\) splits in
\(L\), and \(L\) is generated over \(F\) by \(\alpha\) together with the
roots of \(g\), which are precisely all roots of \(f\).

For uniqueness, let \(L=F(\alpha_1,\ldots,\alpha_r)\) and \(L'\) be
splitting fields of \(f\).  Start with \(1_F:F\to L'\).  Suppose that an
embedding has been defined on \(F(\alpha_1,\ldots,\alpha_{j-1})\).  The
minimal polynomial of \(\alpha_j\) over this field divides \(f\), after the
coefficients have been transported.  Since \(f\) splits in \(L'\), that
transported minimal polynomial has a root in \(L'\).  Lemma
\ref{lem:root-extension} extends the embedding over \(\alpha_j\).

After the last step we have an \(F\)-embedding \(\sigma:L\to L'\).  If
\[
f(x)=c\prod_i(x-\alpha_i)
\]
in \(L[x]\), then applying \(\sigma\), which fixes the coefficients of
\(f\), gives
\[
f(x)=c\prod_i(x-\sigma(\alpha_i))
\]
in \(L'[x]\).  Thus the images are the complete roots of \(f\), with
multiplicities.  Since \(L'\) is generated by those roots, it lies in
\(\sigma(L)\).  Therefore \(\sigma\) is onto and is an isomorphism.
\end{proof}

\begin{example}
The splitting field of \(x^2-2\) is \(\mathbb Q(\sqrt2)\).  The three roots
of \(x^3-2\) are
\[
\sqrt[3]2,\qquad\zeta_3\sqrt[3]2,\qquad\zeta_3^2\sqrt[3]2,
\]
so its splitting field is
\[
\mathbb Q(\sqrt[3]2,\zeta_3).
\]
The first adjunction supplies one real root, but it does not supply the two
conjugate nonreal roots.  The primitive cube root \(\zeta_3\) is precisely
what rotates \(\sqrt[3]2\) into those roots, so adjoining it produces the
whole splitting field.  This is the first indication that the symmetries of
all roots together, rather than a single chosen root, are what Galois theory
will study.  The polynomial \((x-1)(x+2)\) already splits over
\(\mathbb Q\), so its splitting field is \(\mathbb Q\).  In characteristic
\(p\), the polynomial \(x^p-x\) splits in \(\mathbb F_p\), because each
\(a\in\mathbb F_p\) satisfies \(a^p=a\), and it has degree \(p\).
\end{example}

\section{Algebraically Closed Fields and Embeddings}
\begin{definition}
A field \(\Omega\) is \emph{algebraically closed} if every nonconstant
polynomial in \(\Omega[x]\) has a root in \(\Omega\).
\end{definition}
\begin{proposition}
A field is algebraically closed if and only if every nonconstant polynomial
splits into linear factors.
\end{proposition}
\begin{proof}
Repeatedly factor out a root and induct on degree.  The converse is immediate.
\end{proof}

\begin{definition}
An \emph{\(F\)-embedding} \(\sigma:E\to\Omega\) is a field homomorphism
between extensions of \(F\) that fixes \(F\) pointwise.  Its images of an
algebraic element are called its \emph{conjugates}.
\end{definition}

\begin{proposition}[Embeddings and conjugates]\label{prop:embeddings-roots}
If \(\Omega\) is algebraically closed and \(\alpha\) is algebraic over \(F\),
then
\[
\sigma\longmapsto\sigma(\alpha)
\]
is a bijection from \(F\)-embeddings
\[
F(\alpha)\longrightarrow\Omega
\]
to roots of \(m_{\alpha,F}\) in \(\Omega\).
\end{proposition}
\begin{proof}
An embedding sends the equation \(m_{\alpha,F}(\alpha)=0\) to the required
root relation.  Thus every image is a root of the minimal polynomial; these
are exactly the possible algebraic realizations of \(\alpha\) compatible with
the base field.  Conversely, a root gives the unique embedding by
Lemma~\ref{lem:root-extension}.  The two constructions are inverse because
an embedding of \(F(\alpha)\) is determined by its value on \(\alpha\).
\end{proof}

\begin{example}
The conjugates of \(\sqrt2\) over \(\mathbb Q\) are
\[
\sqrt2,\qquad-\sqrt2.
\]
There are correspondingly two \(\mathbb Q\)-embeddings of
\(\mathbb Q(\sqrt2)\) into \(\mathbb C\).  The conjugates of
\(\sqrt[3]2\) are its three complex cube roots; Chapter~10 will explain why
the number of distinct conjugates can be smaller than the degree in positive
characteristic.
\end{example}

More explicitly, the three \(\mathbb Q\)-embeddings of
\(\mathbb Q(\sqrt[3]2)\) into \(\mathbb C\) are determined by the three
possible images
\[
\sqrt[3]2,
\qquad
\zeta_3\sqrt[3]2,
\qquad
\zeta_3^2\sqrt[3]2.
\]
Only the first has image equal to the original real cubic field.  The other
two land in different conjugate subfields of its splitting field.  This is
why embeddings are the right preliminary notion: automorphisms are the
special embeddings whose image remains in the same field.

The next extension theorem is the standard place where Zorn's Lemma enters
field theory.  Start with every partial extension of the prescribed embedding,
order these choices by further extension, take the union along a chain, and
observe that a maximal partial embedding can be enlarged unless its domain is
already all of \(E\).  The proof records the details because this pattern will
return when algebraic closures and normal extensions are compared.

\begin{theorem}[Extension of embeddings]\label{thm:embedding-extension}
If \(E/K\) is algebraic, \(\Omega\) is algebraically closed, and
\(\sigma:K\to\Omega\) is an \(F\)-embedding, then \(\sigma\) extends to an
\(F\)-embedding \(E\to\Omega\).
\end{theorem}
\begin{proof}
Let \(\mathcal P\) consist of pairs \((M,\tau)\) with
\[
K\subseteq M\subseteq E
\]
and \(\tau:M\to\Omega\) an \(F\)-embedding extending \(\sigma\).  Define
\[
(M,\tau)\leq(N,\rho)
\]
when \(M\subseteq N\) and \(\rho|_M=\tau\).  This is a set: the possible
domains are subsets of the fixed set \(E\), and possible maps are subsets of
\(E\times\Omega\).

Let \(\mathcal C\) be a chain.  Its domains have union
\[
M_{\mathcal C}=\bigcup_{(M,\tau)\in\mathcal C}M.
\]
For two elements of this union, one chain member contains both, so it
contains their sum, product, and inverses.  Hence the union is a field.  The
maps agree on overlaps because the pairs are comparable; their union is
therefore a well-defined \(F\)-embedding
\(\tau_{\mathcal C}:M_{\mathcal C}\to\Omega\).  It is an upper bound.
Zorn's Lemma supplies a maximal pair \((M,\tau)\).

If \(M\ne E\), choose \(\alpha\in E\setminus M\).  Since \(E/M\) is
algebraic, \(\alpha\) has a minimal polynomial over \(M\).  The coefficientwise
image of that polynomial has a root in the algebraically closed field
\(\Omega\), and Lemma~\ref{lem:root-extension} extends \(\tau\) to
\(M(\alpha)\), contradicting maximality.  Thus \(M=E\).
\end{proof}

\section{Algebraic Closures}
\begin{definition}
An \emph{algebraic closure} \(\overline F/F\) is an algebraic extension
\(\overline F\) that is algebraically closed.
\end{definition}

An algebraic closure is a sufficiently large, but still algebraic, universe
in which every polynomial over the starting field can be completely split.
The construction below first enlarges a field so that every current polynomial
has a root, then repeats, and finally retains the elements algebraic over the
original field.  That final restriction is what keeps the closure algebraic
over \(F\).

\begin{theorem}[Existence of algebraic closures]\label{thm:algebraic-closure-existence}
Every field has an algebraic closure.
\end{theorem}
\begin{proof}
We use a construction that works within ordinary sets throughout.  For any
field \(K\), let \(S_K\) be the set of nonconstant polynomials in \(K[x]\).
Introduce one indeterminate \(X_f\) for each \(f\in S_K\), and form
\[
R_K=K[X_f:f\in S_K].
\]
Let \(I_K\) be the ideal generated by the elements
\[
f(X_f),
\qquad
f\in S_K.
\]
This ideal is proper.  Indeed, any alleged relation \(1\in I_K\) involves
only finitely many variables \(X_{f_1},\ldots,X_{f_r}\).  Adjoin roots one at
a time by Theorem~\ref{thm:root-adjunction}; in the resulting field every
\(f_i(X_{f_i})\) vanishes, while \(1\) does not.  Thus no such relation can
exist.

By the \hyperref[thm:maximal-ideal-existence]{maximal-ideal theorem}, choose
a maximal ideal \(\mathfrak m_K\supseteq I_K\), and set
\[
K^+=R_K/\mathfrak m_K.
\]
The quotient is a field containing a copy of \(K\): the ideal
\(\mathfrak m_K\cap K\) is an ideal of the field \(K\), and it cannot equal
\(K\), since then \(1\in\mathfrak m_K\).  Hence
\[
\mathfrak m_K\cap K=(0),
\]
so the composite \(K\to R_K\to K^+\) is injective.  Every nonconstant
polynomial \(f\in K[x]\) has the root \(X_f+\mathfrak m_K\) in \(K^+\).

Starting with \(F_0=F\), recursively construct
\[
F_0\subseteq F_1\subseteq F_2\subseteq\cdots,
\qquad
F_{n+1}=F_n^+,
\]
and set
\[
\Omega=\bigcup_{n\geq0}F_n.
\]
This union is a field.  If \(g\in\Omega[x]\) is nonconstant, its finitely
many coefficients lie in some \(F_n\), and then \(g\) has a root in
\(F_{n+1}\subseteq\Omega\).  Hence \(\Omega\) is algebraically closed.

Finally let \(\overline F\) be the subfield of \(\Omega\) consisting of
elements algebraic over \(F\).  This set is a field: if \(\alpha,\beta\) are
algebraic over \(F\), then \(F(\alpha,\beta)/F\) is finite.  Hence
\[
\alpha+\beta,\qquad \alpha-\beta,\qquad \alpha\beta,
\qquad \alpha^{-1}\quad(\alpha\ne0)
\]
belong to a finite extension of \(F\), and are therefore algebraic over \(F\).
If \(h\in\overline F[x]\) is nonconstant, its coefficients lie in a finite
algebraic extension \(K/F\) inside \(\overline F\).  A root \(b\in\Omega\)
of \(h\) is algebraic over \(K\), hence algebraic over \(F\); so
\(b\in\overline F\).  Therefore \(\overline F\) is algebraically closed and
is an algebraic closure of \(F\).
\end{proof}

\begin{theorem}[Uniqueness of algebraic closures]
Any two algebraic closures of \(F\) are \(F\)-isomorphic.
\end{theorem}
\begin{proof}
For closures \(\overline F,\overline F'\), Theorem
\ref{thm:embedding-extension} extends \(1_F\) to an embedding
\[
\sigma:\overline F\longrightarrow\overline F'.
\]
The image is algebraically closed: if a polynomial has coefficients in
\(\sigma(\overline F)\), transport its coefficients back through \(\sigma\),
take a root in \(\overline F\), and transport that root forward.  Now take
\(\beta\in\overline F'\).  It is algebraic over \(\sigma(\overline F)\).
Its irreducible minimal polynomial has a root in the algebraically closed
field \(\sigma(\overline F)\), so it is linear.  Hence
\(\beta\in\sigma(\overline F)\).  Therefore the image is all of
\(\overline F'\).
\end{proof}

\begin{remark}
An algebraic closure is thus unique only up to an isomorphism fixing \(F\).
There is usually no distinguished isomorphism between two constructions, so one fixes
a chosen algebraic closure when it is useful to speak of all conjugates in a
single ambient field.
\end{remark}

\section*{Exercises}
\begin{exercise}[Basic] Prove directly that
\[
[\mathbb Q(\sqrt2,\sqrt3):\mathbb Q]=4.
\]
\end{exercise}
\begin{exercise}[Basic] Find the minimal polynomial of \(1+\sqrt2\) over
\(\mathbb Q\).
\end{exercise}
\begin{exercise}[Basic] Show that \(\mathbb Q(\sqrt2)=\mathbb Q(1+\sqrt2)\).
\end{exercise}
\begin{exercise}[Basic] Prove that \(\mathbb Q(\sqrt2)\cap\mathbb Q(\sqrt3)
=\mathbb Q\).
\end{exercise}
\begin{exercise}[Basic] Show that \(F[t]\) is not a field when \(t\) is
transcendental over \(F\), whereas \(F(t)\) is a field.
\end{exercise}
\begin{exercise}[Core] Show that if \(\alpha\) has odd degree over \(F\),
then \(F(\alpha^2)=F(\alpha)\).
\end{exercise}
\begin{exercise}[Core] Let \(\alpha\) have minimal polynomial
\(x^4+2x^2+2\) over \(F\).  Write \(\alpha^7\) as an \(F\)-linear combination
of \(1,\alpha,\alpha^2,\alpha^3\).
\end{exercise}
\begin{exercise}[Core] Prove that every element of a finite extension
\(E/F\) has minimal polynomial of degree at most \([E:F]\).
\end{exercise}
\begin{exercise}[Core] If \(F\subseteq E\subseteq L\), with
\([L:F]=12\), list the possible degrees of \(E/F\).
\end{exercise}
\begin{exercise}[Core] Construct the splitting field of \(x^4-2\) over
\(\mathbb Q\) and determine its degree.
\end{exercise}
\begin{exercise}[Core] Prove that an intermediate field of a degree-six
extension cannot have degree four over the base.
\end{exercise}
\begin{exercise}[Further] Show that the number of \(F\)-embeddings of
\(F(\alpha)\) into an algebraic closure is at most \(\deg m_{\alpha,F}\).
\end{exercise}
\begin{exercise}[Further] Let \(f\in F[x]\) be irreducible and let
\(\alpha,\beta\) be roots of \(f\) in an algebraically closed field.  Use the
root-extension lemma to construct an \(F\)-isomorphism
\[
F(\alpha)\longrightarrow F(\beta).
\]
\end{exercise}
\begin{exercise}[Further] Prove that a field which is algebraic over an
algebraically closed field is that field itself.
\end{exercise}
\begin{exercise}[Further] Show that \(\mathbb C\) is an algebraic closure of
\(\mathbb R\), using the Fundamental Theorem of Algebra as an external input.
A Galois-theoretic proof of that theorem appears later in Chapter~11.
\end{exercise}
\begin{exercise}[Looking Ahead] Show that every finite extension of a finite
field is a splitting field of a polynomial over that finite field.
\end{exercise}
\begin{exercise}[Looking Ahead] Let \(f\in F[x]\) be irreducible.  Show that
every \(F\)-embedding of \(F(\alpha)\) into an algebraic closure sends
\(\alpha\) to a root of \(f\), and compare this with the automorphisms of a
splitting field of \(f\).
\end{exercise}

\chapter[Separability, Normality, and Finite Fields]{Separability, Normality, and Finite Fields}
\chaptermark{Separability and Finite Fields}

Embeddings count genuinely different conjugate roots.  Separability measures
that distinction; normality says that no conjugate root escapes the field.
Only after both ideas are clear will we define a Galois extension.

\section{Separable Polynomials and Extensions}
\begin{definition}
A nonconstant polynomial \(f\in F[x]\) is \emph{separable} if its roots in a
splitting field are distinct.  An algebraic element is \emph{separable} over
\(F\) when its minimal polynomial is separable; an algebraic extension is
separable when each of its elements is separable.
\end{definition}

The definition applies to arbitrary polynomials.  For an irreducible
polynomial it asks whether the possible conjugates of one root occur only
once.  This is the distinction needed for counting embeddings.

\begin{proposition}[Derivative criterion]\label{prop:separable-criterion}
For nonconstant \(f\in F[x]\),
\[
f\text{ has a repeated root in }\overline F
\quad\Longleftrightarrow\quad
\gcd(f,f')\ne1.
\]
For irreducible \(f\), this is equivalent to \(f'=0\).
\end{proposition}
\begin{proof}
Write \(f=(x-a)^rg\) in a splitting field, with \(g(a)\ne0\).  The product
rule gives
\[
f'=(x-a)^{r-1}\bigl(rg+(x-a)g'\bigr).
\]
Thus \(a\) is repeated exactly when \(r\ge2\), equivalently when
\(f(a)=f'(a)=0\).  A common divisor of \(f\) and \(f'\) is therefore
equivalent, after splitting, to a common linear factor, hence to a repeated
root.  For irreducible \(f\), the gcd is either \(1\) or an associate of
\(f\); the latter can occur only if \(f'=0\).
\end{proof}

\begin{remark}
In characteristic \(p>0\), \(f'=0\) exactly when every exponent appearing
in \(f\) is divisible by \(p\), so
\[
f(x)=g(x^p).
\]
This is the source of inseparability.  For example,
\[
x^p-t\in\mathbb F_p(t)[x]
\]
has derivative zero and is irreducible because a rational function whose
\(p\)-th power is \(t\) would have every zero and pole order divisible by
\(p\), whereas \(t\) has a simple zero.  In an algebraic closure it is
\[
x^p-t=(x-t^{1/p})^p;
\]
it has one distinct root, with multiplicity \(p\).
\end{remark}

\begin{example}[A separable polynomial with repeated-looking coefficients]\label{ex:sep-cubic}
Over \(\mathbb F_2\), the polynomial
\[
x^3+x+1
\]
has derivative \(x^2+1\).  It has no root in \(\mathbb F_2\), hence is
irreducible, and it cannot share a root with its derivative.  It is therefore
separable.  The contrast with \(x^2-t\) over \(\mathbb F_2(t)\), whose
derivative is zero, shows that characteristic alone is not the issue:
inseparability comes from the special \(p\)-th-power shape.
\end{example}

\section{Embeddings, Perfect Fields, and Primitive Elements}

The next three ideas answer complementary questions about a finite extension.
Separability asks whether the degree is fully visible through distinct
conjugates; normality, introduced later, asks whether those conjugates remain
inside the field; and the primitive element theorem says that a finite
separable extension can often be described by one well-chosen generator.
The comparison is useful to keep in view:
\[
\begin{array}{c|c}
\text{property} & \text{what it controls}\\ \hline
\text{separable} & \text{distinct embeddings and conjugate roots}\\
\text{normal} & \text{whether all conjugates stay in the extension}\\
\text{Galois} & \text{both conditions, hence internal symmetries}
\end{array}
\]

\begin{definition}
A field \(F\) is \emph{perfect} if every algebraic extension of \(F\) is
separable, equivalently if every irreducible polynomial in \(F[x]\) is
separable.
\end{definition}
\begin{proposition}
Every characteristic-zero field is perfect.  A field \(F\) of characteristic
\(p>0\) is perfect if and only if Frobenius
\[
F\longrightarrow F,\qquad a\longmapsto a^p
\]
is surjective.  In particular every finite field is perfect.
\end{proposition}
\begin{proof}
In characteristic zero a nonzero irreducible polynomial cannot have zero
derivative, so the criterion proves separability.  In characteristic \(p\),
an inseparable irreducible polynomial has derivative zero and is of the form
\(g(x^p)\).  Surjectivity of Frobenius makes each coefficient a \(p\)-th
power, so this is a \(p\)-th power polynomial, impossible for a nonconstant
irreducible polynomial.  Conversely, if \(a\) has no \(p\)-th root, then
\(x^p-a\) has derivative zero and an irreducible factor which is inseparable.
Frobenius is injective on every field and hence bijective on a finite field.
\end{proof}

Before turning this count into a definition, notice the strategy: build the
field one algebraic generator at a time.  At each step an embedding has only
as many possible extensions as the minimal polynomial has roots.  The Tower
Law then turns this local bound into a global one.

\begin{theorem}[Embedding-count inequality]\label{thm:embedding-count}
For finite \(E/F\),
\[
\#\operatorname{Hom}_F(E,\overline F)\le [E:F],
\]
\end{theorem}
\begin{proof}
Choose generators
\[
F=F_0\subseteq F_1\subseteq\cdots\subseteq F_r=E,
\qquad
F_i=F_{i-1}(\alpha_i).
\]
An embedding \(F_{i-1}\to\overline F\) extends across \(\alpha_i\) only by
sending it to a root of the transported minimal polynomial.  By the
root-extension lemma, there are at most
\([F_i:F_{i-1}]\) choices.  Successively extending embeddings therefore
gives
\[
\#\operatorname{Hom}_F(E,\overline F)
\leq\prod_i[F_i:F_{i-1}]=[E:F].
\]
Equality at every step holds exactly when each transported minimal polynomial
has its full number of distinct roots.  We analyze that equality after
introducing separable degree.
\end{proof}

\begin{definition}
For a finite extension, the number
\[
[E:F]_s=\#\operatorname{Hom}_F(E,\overline F)
\]
is its \emph{separable degree}.  Thus
\[
[E:F]_s\leq[E:F],
\]
by Theorem~\ref{thm:embedding-count}.
\end{definition}

The ordinary degree is multiplicative in a tower by the Tower Law.  The
following result says that the count of genuinely distinct embeddings has the
same multiplicative behavior.  Its proof is a refined version of the
preceding counting argument: first choose an embedding of the middle field,
then count its possible continuations.

\begin{proposition}[Multiplicativity of separable degree]\label{prop:separable-degree}
If
\[
F\subseteq K\subseteq E
\]
is a tower of finite extensions, then
\[
[E:F]_s=[E:K]_s[K:F]_s.
\]
\end{proposition}
\begin{proof}
Choose a tower of simple extensions
\[
K=K_0\subseteq K_1\subseteq\cdots\subseteq K_r=E,
\qquad K_i=K_{i-1}(\alpha_i),
\]
and let \(d_i\) be the number of distinct roots of the minimal polynomial
of \(\alpha_i\) over \(K_{i-1}\).  The root-extension lemma shows that an
embedding of \(K_{i-1}\) has exactly \(d_i\) extensions to \(K_i\).
Indeed, the transported minimal polynomial is obtained by applying the
embedding to its coefficients.  It has the same degree and its derivative
is the transported derivative, so its gcd with its derivative has the same
degree; hence it has exactly the same number \(d_i\) of distinct roots.

It follows that every \(F\)-embedding \(\sigma:K\to\overline F\) has
\[
\prod_{i=1}^r d_i=[E:K]_s
\]
extensions to \(E\).  The equality with \([E:K]_s\) is obtained by applying
the same step-by-step count beginning with the identity embedding of \(K\).
Counting first the possible restrictions to \(K\), then their extensions,
proves the formula.
\end{proof}

\begin{theorem}[Separable degree criterion]\label{thm:separable-degree}
For a finite extension \(E/F\),
\[
[E:F]_s=[E:F]
\]
if and only if \(E/F\) is separable.
\end{theorem}
\begin{proof}
Choose a simple generating tower as in the proof of
Proposition~\ref{prop:separable-degree}.  The embedding count there gives
\[
[E:F]_s=\prod_i d_i,
\qquad
[E:F]=\prod_i[K_i:K_{i-1}].
\]
If \(E/F\) is separable, the minimal polynomial of \(\alpha_i\) over
\(K_{i-1}\) divides its separable minimal polynomial over \(F\), so it is
separable.  Thus \(d_i=[K_i:K_{i-1}]\) for every \(i\), proving equality.

Conversely, equality of the two products forces equality in every factor,
so each \(\alpha_i\) is separable over \(K_{i-1}\).  To see that every
\(\gamma\in E\) is separable over \(F\), set \(L=F(\gamma)\).  By
Proposition~\ref{prop:separable-degree},
\[
[E:F]_s=[E:L]_s[L:F]_s.
\]
Both factors are bounded above by their ordinary degrees, while their
product is \([E:F]=[E:L][L:F]\).  Hence
\[
[L:F]_s=[L:F].
\]
For the simple extension \(L=F(\gamma)\), this says precisely that the
minimal polynomial of \(\gamma\) has its full number of distinct roots.
Thus \(\gamma\) is separable.
\end{proof}

\begin{proposition}[Separability in towers]\label{prop:separable-towers}
For a finite tower
\[
F\subseteq K\subseteq E,
\]
\(E/F\) is separable if and only if both \(E/K\) and \(K/F\) are separable.
The same statement holds for arbitrary algebraic towers.
\end{proposition}
\begin{proof}
If \(E/F\) is separable, elements of \(K\) are already separable over
\(F\), and a minimal polynomial over \(K\) divides the corresponding
separable minimal polynomial over \(F\); hence both subextensions are
separable.  Conversely, if both are separable, the multiplicativity formula
and Theorem~\ref{thm:separable-degree} give
\[
[E:F]_s=[E:K][K:F]=[E:F].
\]
Thus \(E/F\) is separable.  For arbitrary algebraic towers, apply the finite
statement as follows.  The forward direction is immediate for \(K/F\), since
\(K\subseteq E\); and, for \(\alpha\in E\), the minimal polynomial over
\(K\) divides the separable minimal polynomial over \(F\), proving that
\(E/K\) is separable.

Conversely, suppose \(K/F\) and \(E/K\) are separable, and choose
\(\alpha\in E\).  The minimal polynomial
\[
m_{\alpha,K}(x)
\]
has only finitely many coefficients.  Let \(K_0\) be the subfield of \(K\)
generated over \(F\) by those coefficients.  Then \(K_0/F\) is finite; it is
separable because every element of the subfield \(K_0\subseteq K\) is already
separable over \(F\).  The polynomial \(m_{\alpha,K}\) lies in \(K_0[x]\) and is
separable because \(E/K\) is separable.  Hence the minimal polynomial of
\(\alpha\) over \(K_0\), which divides \(m_{\alpha,K}\), is separable.
Thus \(K_0(\alpha)/K_0\) is finite separable.  The finite tower statement,
applied to
\[
F\subseteq K_0\subseteq K_0(\alpha),
\]
shows that \(\alpha\) is separable over \(F\).  Since \(\alpha\) was
arbitrary, \(E/F\) is separable.
\end{proof}

\begin{lemma}[Cyclicity lemma]\label{lem:finite-mult-cyclic}
Every finite subgroup of the multiplicative group of a field is cyclic.
\end{lemma}
\begin{proof}
Let \(G\) be such a subgroup and let \(m\) be its exponent.  Every element
of \(G\) is a root of \(x^m-1\), so \(|G|\leq m\).  Conversely, by the
structure theorem for finite abelian groups, \(G\) has an element of order
equal to its exponent \(m\), hence \(m\leq|G|\).  Therefore \(m=|G|\), and
an element of order \(m\) generates \(G\).
\end{proof}

The primitive element theorem is not merely a convenient notation theorem.
It reduces a finite separable extension to a single minimal polynomial, so
that its embeddings and conjugates can be studied through one element.  Over
an infinite field, the proof chooses a linear combination
\(\alpha+c\beta\) avoiding the finitely many values of \(c\) that would make
two embeddings coincide; over a finite field, cyclicity of the multiplicative
group supplies a generator instead.

\begin{theorem}[Primitive element theorem]
Every finite separable extension is simple.
\end{theorem}
\begin{proof}
First suppose \(F\) is infinite and \(E=F(\alpha,\beta)\).  List the
\([E:F]\) \(F\)-embeddings of \(E\) into an algebraic closure as
\(\sigma_i\), which exist by Theorem~\ref{thm:separable-degree}.  For
each pair of distinct images, the equality
\[
\sigma_i(\alpha)+c\sigma_i(\beta)
=\sigma_j(\alpha)+c\sigma_j(\beta)
\]
excludes at most one value of \(c\in F\).  Choose \(c\) outside this finite
set and put \(\gamma=\alpha+c\beta\).  The conjugates of \(\gamma\) are then
\([E:F]\) distinct values.  They are values of distinct restrictions of the
\(\sigma_i\) to \(F(\gamma)\): if two restrictions agreed, they would give
the same displayed value.  Hence
\[
\#\operatorname{Hom}_F(F(\gamma),\overline F)\ge [E:F].
\]
The embedding-count theorem gives the reverse bound by
\([F(\gamma):F]\), while the Tower Law gives
\([F(\gamma):F]\le [E:F]\).  Thus equality holds throughout.
Since \(F(\gamma)\subseteq E\), the Tower Law gives \(E=F(\gamma)\).
Induction reduces finitely many generators to this case.

If \(F\) is finite, then \(E^\times\) is a finite subgroup of a field and
is cyclic by Lemma~\ref{lem:finite-mult-cyclic}.  A generator \(\gamma\) of
\(E^\times\) has \(E=F(\gamma)\).
\end{proof}

\begin{example}[A primitive element]
The extension
\[
\mathbb Q(\sqrt2,\sqrt3)/\mathbb Q
\]
is finite and separable.  The proof above predicts a primitive element; here
one may take \(\gamma=\sqrt2+\sqrt3\).  Since
\[
\gamma^2=5+2\sqrt6,
\]
we recover \(\sqrt6\), and then
\[
\sqrt2=\frac{\gamma^2-1}{2\gamma},
\qquad
\sqrt3=\gamma-\sqrt2.
\]
Thus \(\mathbb Q(\sqrt2,\sqrt3)=\mathbb Q(\sqrt2+\sqrt3)\).
\end{example}

\section{Normal Extensions and Galois Extensions}
\begin{definition}
An algebraic extension \(E/F\) is \emph{normal} if every irreducible
polynomial in \(F[x]\) having one root in \(E\) splits over \(E\).  It is
\emph{Galois} if it is algebraic, normal, and separable.
\end{definition}

The guiding slogan is:
\[
\text{normal means that no conjugate root escapes.}
\]
Thus normality is a condition on the \emph{whole} field, whereas separability
is a condition on the distinctness of roots.  The normality criteria below
translate this root-theoretic statement into the embedding language developed
above.

\begin{theorem}[Normality criteria]\label{thm:normality-criteria}
For finite \(E/F\), the following are equivalent:
\begin{enumerate}
\item \(E/F\) is normal;
\item every \(F\)-embedding \(E\to\overline F\) has image \(E\);
\item \(E\) is a splitting field over \(F\) of a collection of polynomials.
\end{enumerate}
\end{theorem}
\begin{proof}
Assume first that \(E/F\) is normal.  If \(\sigma:E\to\overline F\) is an
\(F\)-embedding and \(\alpha\in E\), then \(\sigma(\alpha)\) is a conjugate
root of \(m_{\alpha,F}\).  Normality makes this root belong to \(E\); hence
\(\sigma(E)\subseteq E\).  Both fields have degree \([E:F]\) over \(F\), so
\(\sigma(E)=E\).

Conversely, suppose every embedding has image \(E\).  If an irreducible
\(f\in F[x]\) has a root \(\alpha\in E\), every root \(\beta\) of \(f\) gives
an embedding \(F(\alpha)\to\overline F\) by Chapter~10.  Extend it to \(E\)
by Theorem~\ref{thm:embedding-extension}.  Its image is \(E\), so
\(\beta\in E\).  Thus \(f\) splits in \(E\), and \(E/F\) is normal.

Finally, suppose \(E\) is a splitting field over \(F\).  Every
\(F\)-embedding of \(E\) fixes the defining coefficients and permutes the
corresponding roots.  Its image is consequently contained in \(E\); equality
follows from degrees, and the already proved implication shows that \(E/F\)
is normal.  If \(E/F\) is
finite normal, choose generators \(\alpha_1,\ldots,\alpha_r\).  The splitting
field in \(E\) of the finitely many minimal polynomials
\(m_{\alpha_i,F}\) contains every generator and is therefore \(E\).
\end{proof}

\begin{example}[Normality and separability are independent]
\(\mathbb Q(\sqrt[3]2)/\mathbb Q\) is separable because characteristic-zero
fields are perfect, but it is not normal: the two nonreal roots
\[
\omega\sqrt[3]2,\qquad\omega^2\sqrt[3]2
\]
of \(x^3-2\) are absent.  Its normal closure is obtained by adjoining a
primitive cube root of unity \(\omega\).  On the other hand, put
\[
E=\mathbb F_p(t^{1/p}),
\qquad
F=\mathbb F_p(t).
\]
For every \(\alpha\in E\), the Frobenius identity gives
\[
\alpha^p\in F.
\]
Thus the minimal polynomial of \(\alpha\) divides
\[
x^p-\alpha^p=(x-\alpha)^p
\]
in an algebraic closure and has only one distinct root.  Every irreducible
polynomial over \(F\) having a root in \(E\) consequently splits in \(E\) as
a power of a linear factor, so \(E/F\) is normal.  It is not separable,
because the generator \(t^{1/p}\) has derivative-zero minimal polynomial
\(x^p-t\).
\end{example}

The following three familiar extensions summarize the distinctions.  The
middle example has all of its distinct conjugates but not all of them inside
the displayed field; the last loses embeddings because its minimal polynomial
has a repeated root.
\[
\begin{array}{c|c|c|c|c|c}
\text{extension} & \text{degree} & \text{\(F\)-embeddings}
& \text{separable} & \text{normal} & \text{Galois}\\ \hline
\mathbb Q(\sqrt2)/\mathbb Q & 2 & 2 & \text{yes} & \text{yes} & \text{yes}\\
\mathbb Q(\sqrt[3]2)/\mathbb Q & 3 & 3 & \text{yes} & \text{no} & \text{no}\\
\mathbb F_p(t^{1/p})/\mathbb F_p(t) & p & 1 & \text{no} & \text{yes} & \text{no}
\end{array}
\]

\begin{proposition}[Finite Galois criteria]\label{prop:finite-galois-criteria}
A finite extension is Galois precisely when it is the splitting field of a
separable polynomial, equivalently when
\[
|\operatorname{Aut}_F(E)|=[E:F].
\]
\end{proposition}
\begin{proof}
A splitting field is normal by Theorem~\ref{thm:normality-criteria}.
When its defining polynomial is separable, the tower criterion gives
separability.
Conversely, if \(E/F\) is Galois, choose generators
\[
E=F(\alpha_1,\ldots,\alpha_r).
\]
Normality makes every \(m_{\alpha_i,F}\) split in \(E\), and separability
makes each of these polynomials separable.  The product of the distinct
polynomials among them is separable, and its splitting field contains every
generator and is contained in \(E\); hence it is \(E\).

If \(E/F\) is normal and separable, Theorem~\ref{thm:separable-degree} gives
\[
\#\operatorname{Hom}_F(E,\overline F)=[E:F].
\]
Theorem~\ref{thm:normality-criteria} says every such embedding has image
\(E\), so all of them are automorphisms and
\[
|\operatorname{Aut}_F(E)|=[E:F].
\]
Conversely, if this automorphism count holds, then
\[
[E:F]
=|\operatorname{Aut}_F(E)|
\leq\#\operatorname{Hom}_F(E,\overline F)
\leq[E:F].
\]
Equality holds throughout.  Theorem~\ref{thm:separable-degree} gives
separability, and every \(F\)-embedding is already one of the automorphisms.
Its image is therefore \(E\), so Theorem~\ref{thm:normality-criteria} gives
normality.  Thus \(E/F\) is Galois.
\end{proof}

\phantomsection\label{thm:finite-galois-closure}
\begin{applicationtheorem}[Finite Galois closures]
If \(E/F\) is finite separable, then there is a finite Galois extension
\(L/F\) containing \(E\).
\end{applicationtheorem}
\begin{proof}
Choose generators
\[
E=F(\alpha_1,\ldots,\alpha_r),
\]
and let \(f_i\) be the minimal polynomial of \(\alpha_i\) over \(F\).  Each
\(f_i\) is separable because \(E/F\) is separable.  Let \(L\) be the
splitting field over \(F\) of the product of the distinct \(f_i\)'s.  A
splitting field of finitely many polynomials is finite, and \(L/F\) is
normal.  The product is separable, so its splitting field is separable.
Thus \(L/F\) is finite Galois, and it contains every \(\alpha_i\), hence
contains \(E\).
\end{proof}

Such an \(L\) is called a finite Galois closure (or normal Galois closure)
of \(E/F\).  Uniqueness is not needed here.

\section{Finite Fields}
\begin{theorem}[Finite fields]\label{thm:finite-fields}
For each prime power \(q=p^r\), there is, up to isomorphism, one field
\(\mathbb F_q\).  It is the splitting field of
\[
x^q-x
\]
over \(\mathbb F_p\), and \(\mathbb F_q^\times\) is cyclic of order \(q-1\).
Its subfields are exactly \(\mathbb F_{p^d}\) for \(d\mid r\).
\end{theorem}
\begin{proof}
Let \(q=p^r\).  In an algebraic closure of \(\mathbb F_p\), the roots of
\[
x^q-x
\]
form a field: in characteristic \(p\),
\[
(a+b)^q=a^q+b^q,\qquad(ab)^q=a^qb^q,
\]
and nonzero roots remain roots on inversion.  The derivative is \(-1\), so
there are \(q\) distinct roots.  This gives a field with \(q\) elements.

Conversely, if \(K\) has \(q\) elements, then \(K^\times\) has order \(q-1\),
so \(a^{q-1}=1\) for every nonzero \(a\), and \(a^q=a\) for all \(a\in K\).
Thus \(K\) is exactly the root set of \(x^q-x\), proving uniqueness up to
isomorphism and showing that it is a splitting field.  Cyclicity of
\(\mathbb F_q^\times\) is Lemma~\ref{lem:finite-mult-cyclic}.

For subfields, start with an arbitrary field
\[
\mathbb F_p\subseteq K\subseteq\mathbb F_{p^r}
\]
and put \(d=[K:\mathbb F_p]\).  Then \(|K|=p^d\), so every element of \(K\)
satisfies \(x^{p^d}-x=0\).  Thus \(K\) is the unique field with \(p^d\)
elements in the chosen algebraic closure, namely \(\mathbb F_{p^d}\), and
the Tower Law gives \(d\mid r\).  Conversely, if \(d\mid r\), every root of
\[
x^{p^d}-x
\]
lies in \(\mathbb F_{p^r}\), because
\[
x^{p^d}-x\mid x^{p^r}-x
\]
when \(d\mid r\).  To justify the divisibility, write \(r=dk\).  If
\(a^{p^d}=a\), iterating this equality \(k\) times gives
\[
a^{p^r}=a.
\]
Thus every root of \(x^{p^d}-x\) is a root of \(x^{p^r}-x\).  The former
polynomial has derivative \(-1\), hence has no repeated roots, and this
root containment proves the divisibility.  Its \(p^d\) roots form the
required subfield.  This proves both the existence of each listed subfield
and the exhaustion of all subfields.
\end{proof}

\begin{example}[The field with four elements]
The polynomial
\[
x^2+x+1
\]
has no root in \(\mathbb F_2\), so
\[
\mathbb F_4\cong\mathbb F_2[\alpha]/(\alpha^2+\alpha+1).
\]
Its elements are
\[
0,\quad1,\quad\alpha,\quad\alpha+1,
\]
and the relation \(\alpha^2=\alpha+1\) makes multiplication concrete:
\[
(\alpha+1)^2=\alpha^2+1=\alpha.
\]
The Frobenius automorphism fixes \(0,1\) and exchanges the two nontrivial
elements:
\[
\alpha^2=\alpha+1,\qquad(\alpha+1)^2=\alpha.
\]
\end{example}

\begin{example}[The field with eight elements]
By Example~\ref{ex:sep-cubic}, the polynomial \(x^3+x+1\) has no root in
\(\mathbb F_2\) and is irreducible.  Thus
\[
\mathbb F_8\cong\mathbb F_2[\beta]/(\beta^3+\beta+1).
\]
Every element has a unique form
\[
a+b\beta+c\beta^2,
\qquad a,b,c\in\mathbb F_2,
\]
and multiplication is reduced using \(\beta^3=\beta+1\).  For instance,
\[
\beta(\beta^2+1)=\beta^3+\beta=1.
\]
Since \(\beta\ne0,1\) and \(\mathbb F_8^\times\) has prime order \(7\),
Lagrange's theorem implies that \(\beta\) has multiplicative order \(7\).
\end{example}

\begin{proposition}[Frobenius has exact order]\label{prop:frobenius-order}
For \(q=p^r\),
\[
\operatorname{Gal}(\mathbb F_{q^n}/\mathbb F_q)\cong C_n,
\]
generated by Frobenius \(x\mapsto x^q\), whose order is exactly \(n\).
\end{proposition}
\begin{proof}
Frobenius fixes \(\mathbb F_q\), and its \(n\)-th power is the identity on
\(\mathbb F_{q^n}\).  If its \(d\)-th power were the identity for
\(0<d<n\), every element would satisfy \(x^{q^d}=x\), placing the field
among the \(q^d\) roots of \(x^{q^d}-x\), an impossibility since
\(q^n>q^d\).  Thus its \(n\) powers are distinct.  They are all
\(\mathbb F_q\)-automorphisms, and the embedding count bounds the whole
automorphism group by \(n\).
\end{proof}

These cyclic finite Galois groups will become a running example in
Chapter~12.  There the compatible restrictions among the groups
\(\operatorname{Gal}(\mathbb F_{q^n}/\mathbb F_q)\) assemble into the
absolute Galois group of \(\mathbb F_q\).

\begin{example}[Subfields of \(\mathbb F_{p^{12}}\)]
The divisors of \(12\) give all subfields:
\[
\mathbb F_p,\quad
\mathbb F_{p^2},\quad
\mathbb F_{p^3},\quad
\mathbb F_{p^4},\quad
\mathbb F_{p^6},\quad
\mathbb F_{p^{12}}.
\]
For example,
\[
\mathbb F_{p^3}\cap\mathbb F_{p^4}=\mathbb F_p,
\]
since an element in the intersection lies in the unique subfield whose
degree divides both \(3\) and \(4\).
\end{example}

\section{Cyclotomic Extensions}

Cyclotomic fields are the natural meeting point of roots of unity, splitting
fields, and explicit Galois groups.  The multiplicative order of a root of
unity supplies its arithmetic label, while the way automorphisms change that
root will later be read directly from the unit group
\((\mathbb Z/n\mathbb Z)^\times\).

\begin{definition}
A \emph{primitive \(n\)-th root of unity} is an element of multiplicative
order \(n\).  The \(n\)-th cyclotomic polynomial \(\Phi_n(x)\) is the monic
polynomial whose roots are the primitive \(n\)-th roots of unity.
\end{definition}
The finite cyclicity lemma shows that the roots of \(x^n-1\), in
characteristic not dividing \(n\), form a cyclic group.  Partitioning them
by their order gives
\[
x^n-1=\prod_{d\mid n}\Phi_d(x).
\]
Induction in this formula, together with Gauss's Lemma, gives
\(\Phi_n\in\mathbb Z[x]\): after the factors for proper divisors have been
constructed as monic integral polynomials, their product divides the monic
polynomial \(x^n-1\) in \(\mathbb Q[x]\), and polynomial division followed by
Gauss's Lemma makes the monic quotient integral.  The number of primitive
roots gives
\[
\deg\Phi_n=\varphi(n).
\]
For example,
\[
\begin{aligned}
\Phi_1&=x-1, & \Phi_2&=x+1, & \Phi_3&=x^2+x+1,\\
\Phi_4&=x^2+1, & \Phi_5&=x^4+x^3+x^2+x+1,\\
\Phi_6&=x^2-x+1, & \Phi_8&=x^4+1.
\end{aligned}
\]

\begin{theorem}[Cyclotomic irreducibility]
\(\Phi_n(x)\) is irreducible over \(\mathbb Q\).
\end{theorem}
\begin{proof}
Let \(\zeta\) be primitive and let \(g\in\mathbb Z[x]\) be its monic
irreducible polynomial.  Write
\[
\Phi_n=gh
\]
in \(\mathbb Z[x]\).  Let \(p\nmid n\) be prime.  Suppose that \(\zeta^p\)
is not a root of \(g\).  It is a primitive root and hence a root of \(h\),
so \(h(\zeta^p)=0\).  Therefore \(\zeta\) is a root of \(h(x^p)\), and
minimality of \(g\) gives
\[
g(x)\mid h(x^p)
\]
in \(\mathbb Q[x]\), hence in \(\mathbb Z[x]\) by Gauss's Lemma.  Reducing
modulo \(p\), Frobenius gives
\[
\overline{h(x^p)}=\overline{h(x)}^{\,p}.
\]
Thus an irreducible factor of \(\bar g\) also divides \(\bar h\), so
\(\overline{\Phi_n}=\bar g\bar h\) has a repeated irreducible factor.
But \(\Phi_n\) divides \(x^n-1\), whereas
\[
(x^n-1)'=nx^{n-1}
\]
is relatively prime to \(x^n-1\) modulo \(p\).  This contradicts the
derivative criterion.  Hence \(\zeta^p\) is a root of \(g\) for every prime
\(p\nmid n\).  Let \(a>0\) with \((a,n)=1\), and factor
\[
a=p_1\cdots p_s.
\]
No \(p_i\) divides \(n\).  Repeatedly applying the preceding conclusion
shows that \(\zeta^a\) is a root of \(g\).  Thus every primitive root is a
root of \(g\), and consequently \(g=\Phi_n\).
\end{proof}

If \(\zeta_n\) is primitive, irreducibility gives
\[
[\mathbb Q(\zeta_n):\mathbb Q]=\varphi(n).
\]
It is the splitting field of \(x^n-1\), hence is finite Galois.  The map
\[
(\mathbb Z/n\mathbb Z)^\times\longrightarrow
\operatorname{Gal}(\mathbb Q(\zeta_n)/\mathbb Q),
\qquad a\longmapsto(\zeta_n\mapsto\zeta_n^a)
\]
is an isomorphism.  Indeed, \(\zeta_n^a\) is primitive exactly when
\((a,n)=1\), and then it is a root of the minimal polynomial \(\Phi_n\).
The root-extension lemma therefore gives a unique \(\mathbb Q\)-embedding
sending \(\zeta_n\) to \(\zeta_n^a\); finite dimensionality makes it an
automorphism.  Composition corresponds to multiplication of exponents, and
the map is injective because an automorphism is determined by \(\zeta_n\).
Conversely every automorphism sends \(\zeta_n\) to a primitive root, so it
has this form.  Roots of unity will reappear in the later discussion of
radicals.

\begin{example}[Small cyclotomic Galois groups]
For \(n=3\) and \(n=4\),
\[
(\mathbb Z/n\mathbb Z)^\times\cong C_2,
\]
so both \(\mathbb Q(\zeta_3)\) and
\(\mathbb Q(\zeta_4)=\mathbb Q(i)\) are quadratic Galois extensions.  For
\(n=5\),
\[
(\mathbb Z/5\mathbb Z)^\times\cong C_4,
\]
whereas
\[
(\mathbb Z/8\mathbb Z)^\times\cong C_2\oplus C_2.
\]
Thus the degree-four fields \(\mathbb Q(\zeta_5)\) and
\(\mathbb Q(\zeta_8)\) have nonisomorphic Galois groups.
\end{example}

\section*{Exercises}
\subsection*{Basic}
\begin{exercise}[Basic] Determine whether \(x^4-2x^2+1\) is separable over
\(\mathbb Q\), and explain the answer using the derivative criterion.
\end{exercise}
\begin{exercise}[Basic] Use the derivative criterion to decide whether
\[
x^p-x\in\mathbb F_p[x]
\]
is separable.
\end{exercise}
\begin{exercise}[Basic] Show that a finite field is perfect directly from
the fact that its Frobenius map is bijective.
\end{exercise}

\subsection*{Core}
\begin{exercise}[Core] Test normality and separability of
\(\mathbb Q(\sqrt[3]2)/\mathbb Q\).
\end{exercise}
\begin{exercise}[Core] Let \(E/F\) be finite.  Explain why
\[
[E:F]_s=[E:F]
\]
when \(E/F\) is generated by a tower of separable simple extensions.
\end{exercise}
\begin{exercise}[Core] Find a primitive element for
\[
\mathbb Q(\sqrt2,\sqrt5)/\mathbb Q.
\]
\end{exercise}
\begin{exercise}[Core] Verify that \(x^3+x+1\) is irreducible over
\(\mathbb F_2\), and use it to list the elements of \(\mathbb F_8\).
\end{exercise}
\begin{exercise}[Core] Determine the subfields of \(\mathbb F_{p^{12}}\).
\end{exercise}
\begin{exercise}[Core] Prove that every finite extension of a finite field
is Galois.
\end{exercise}
\begin{exercise}[Core] Compute the Frobenius orbit of a nonzero element of
\(\mathbb F_8\) not in \(\mathbb F_2\).
\end{exercise}

\subsection*{Further}
\begin{exercise}[Further] Show that
\[
\Phi_{10}(x)=x^4-x^3+x^2-x+1.
\]
\end{exercise}
\begin{exercise}[Further] Show that \(\mathbb Q(\zeta_5)\) has Galois group
isomorphic to \(C_4\).
\end{exercise}
\begin{exercise}[Further] Determine
\(\operatorname{Gal}(\mathbb Q(\zeta_8)/\mathbb Q)\).
\end{exercise}
\begin{exercise}[Further] Prove that the splitting field of
\[
x^4-2
\]
over \(\mathbb Q\) is normal and separable, and determine its degree.
\end{exercise}

\subsection*{Looking Ahead}
\begin{exercise}[Challenging] Let \(E/F\) be finite and separable.  Use the
embedding-count theorem to show that \(E/F\) is normal exactly when every
\(F\)-embedding \(E\to\overline F\) is an automorphism of \(E\).
\end{exercise}
\begin{exercise}[Looking Ahead] If \(E/F\) is finite Galois, show that
\[
F\subseteq E^{\operatorname{Aut}_F(E)}.
\]
The reverse inclusion is the starting point of the Galois correspondence.
\end{exercise}

\chapter[Galois Theory]{Galois Theory, Limits, and Profinite Groups}

Galois theory replaces the study of intermediate fields by the study of the
symmetries of the larger field.  In the finite case this is an exact
dictionary.  In the infinite case the same dictionary survives once the
symmetry group remembers agreement at every finite stage through a topology.

\section{Galois Groups and Fixed Fields}

\begin{definition}
For an extension \(E/F\), its \emph{Galois group} is
\[
\operatorname{Gal}(E/F)=\operatorname{Aut}_F(E).
\]
If \(H\leq\operatorname{Gal}(E/F)\), its \emph{fixed field} is
\[
E^H=\{x\in E:\sigma(x)=x\text{ for every }\sigma\in H\}.
\]
For an intermediate field \(F\subseteq K\subseteq E\), write
\[
\operatorname{Gal}(E/K)=\{\sigma\in\operatorname{Gal}(E/F):
\sigma|_K=1_K\}.
\]
\end{definition}

The fixed field records the elements on which the chosen symmetries are
invisible; it will be the inverse operation to taking an automorphism group.

\begin{lemma}\label{lem:fixed-field-intermediate}
The set \(E^H\) is an intermediate field.  Furthermore,
\[
K_1\subseteq K_2\Longrightarrow
\operatorname{Gal}(E/K_2)\leq\operatorname{Gal}(E/K_1),
\qquad
H_1\leq H_2\Longrightarrow E^{H_2}\subseteq E^{H_1}.
\]
\end{lemma}
\begin{proof}
Every element of \(H\) fixes \(0\) and \(1\).  If \(x,y\in E^H\), then each
\(\sigma\in H\) fixes \(x-y\), \(xy\), and \(x^{-1}\) when \(x\ne0\).
Thus \(E^H\) is a subfield containing \(F\).  The two inclusions say,
respectively, that fixing more elements imposes more conditions and that
asking more automorphisms to fix an element imposes more conditions.
\end{proof}

\begin{example}
Complex conjugation has fixed field \(\mathbb R\) in \(\mathbb C\).  If \(d\)
is squarefree, the nonidentity automorphism of \(\mathbb Q(\sqrt d)\) sends
\(\sqrt d\) to \(-\sqrt d\) and fixes precisely \(\mathbb Q\).  In
\(\mathbb Q(\zeta_5)\), the automorphisms are
\(\zeta_5\mapsto\zeta_5^a\), \(a\in(\mathbb Z/5\mathbb Z)^\times\); its
unique order-two subgroup fixes \(\mathbb Q(\sqrt5)\).
\end{example}

The following linear-algebra fact is the key input in Artin's theorem.
Distinct field homomorphisms cannot satisfy an accidental linear relation as
functions.

\begin{lemma}[Artin independence lemma]\label{lem:artin-independence}
Let \(\sigma_1,\ldots,\sigma_n:E\to L\) be distinct field homomorphisms.
Then they are linearly independent over \(L\) as functions \(E\to L\).
\end{lemma}
\begin{proof}
Suppose that a nontrivial relation
\[
a_1\sigma_1+\cdots+a_n\sigma_n=0
\]
exists, and choose one with the fewest nonzero coefficients.  Omit zero
terms, so every \(a_i\ne0\) and \(n\ge2\).  Choose \(x\in E\) with
\(\sigma_1(x)\ne\sigma_n(x)\).  Evaluating at \(xy\), then subtracting
\(\sigma_n(x)\) times the original relation evaluated at \(y\), yields
\[
\sum_{i=1}^{n-1}a_i\bigl(\sigma_i(x)-\sigma_n(x)\bigr)\sigma_i(y)=0
\quad\text{for all }y\in E.
\]
The coefficient of \(\sigma_n\) has vanished while that of \(\sigma_1\) is
nonzero.  This is a shorter nontrivial relation, a contradiction.
\end{proof}

The proof of Artin's theorem has a useful two-sided shape.  Independence of
the automorphisms first produces enough linearly independent elements to give
the lower bound \(|G|\leq[E:E^G]\).  Orbit polynomials then show that every
finitely generated subfield has degree at most \(|G|\), producing the reverse
bound.  The final embedding count identifies the automorphisms.  This is the
bridge from a concrete group action to the field--subgroup dictionary.

\begin{theorem}[Artin fixed-field theorem]\label{thm:artin-fixed-field}
If \(G\) is a finite group of automorphisms of a field \(E\), then
\[
[E:E^G]=|G|,
\qquad
\operatorname{Gal}(E/E^G)=G.
\]
\end{theorem}
\begin{proof}
Put \(F=E^G\) and \(n=|G|\).  Artin independence implies that there are
\(x_1,\ldots,x_n\in E\) for which the evaluation matrix
\[
\bigl(\sigma_i(x_j)\bigr)_{1\leq i,j\leq n}
\]
is invertible.  Indeed, consider the evaluation vectors
\[
v(x)=(\sigma_1(x),\ldots,\sigma_n(x))\in E^n.
\]
If these did not span \(E^n\), a nonzero linear functional
\((a_1,\ldots,a_n)\) would vanish on every \(v(x)\), giving
\[
a_1\sigma_1+\cdots+a_n\sigma_n=0,
\]
contrary to Artin independence.  Choose \(n\) evaluation vectors forming a
basis, and use them as the columns of the displayed matrix.  The \(x_j\) are
therefore \(F\)-linearly independent, so \(n\leq[E:F]\).

For the reverse inequality, let \(\alpha\in E\) and form
\[
P_\alpha(X)=\prod_{\sigma\in G}(X-\sigma(\alpha)).
\]
Every \(\tau\in G\) permutes these factors, so the coefficients belong to
\(F\).  The minimal polynomial of \(\alpha\) over \(F\) divides
\(P_\alpha\), and is separable of degree at most \(n\).  Given finitely many
elements of \(E\), the field they generate over \(F\) is finite separable;
the primitive element theorem makes it \(F(\alpha)\), hence of degree at most
\(n\).  Therefore \(E\) has no \(n+1\) linearly independent elements over
\(F\), and \([E:F]\leq n\).

Finally \(G\leq\operatorname{Gal}(E/F)\), while the embedding-count
inequality (Theorem~\ref{thm:embedding-count}) gives
\[
|\operatorname{Gal}(E/F)|\leq[E:F]=|G|.
\]
Therefore the inclusion is equality.
\end{proof}

\begin{corollary}\label{cor:finite-automorphism-consequences}
If \(E/F\) is finite Galois, then
\[
E^{\operatorname{Gal}(E/F)}=F.
\]
Conversely, if \(G\) is a finite subgroup of \(\operatorname{Aut}(E)\), then
\(E/E^G\) is finite Galois with Galois group \(G\).
\end{corollary}
\begin{proof}
The first assertion follows by combining Artin's theorem with the finite
Galois criterion (Proposition~\ref{prop:finite-galois-criteria}).  The second
follows from the same two results.
\end{proof}

\section{The Fundamental Theorem of Finite Galois Theory}

\begin{theorem}[Finite Galois correspondence]\label{thm:finite-galois}
Let \(E/F\) be finite Galois and put \(G=\operatorname{Gal}(E/F)\).  The maps
\[
K\longmapsto\operatorname{Gal}(E/K),
\qquad
H\longmapsto E^H
\]
are inverse inclusion-reversing bijections between intermediate fields
\(F\subseteq K\subseteq E\) and subgroups \(H\leq G\).  Moreover,
\[
[E:K]=|\operatorname{Gal}(E/K)|,
\qquad
[K:F]=[G:\operatorname{Gal}(E/K)].
\]
\end{theorem}
\begin{proof}
For an intermediate \(K\), separability and normality descend from \(E/F\)
to \(E/K\): a minimal polynomial over \(K\) divides the corresponding
polynomial over \(F\), so it is separable and splits in \(E\).  Thus \(E/K\)
is finite Galois, and Corollary~\ref{cor:finite-automorphism-consequences}
gives
\[
E^{\operatorname{Gal}(E/K)}=K.
\]
For \(H\leq G\), Artin's theorem applied to the finite group \(H\) gives
\[
\operatorname{Gal}(E/E^H)=H.
\]
These identities prove the bijection.  The first degree formula is Artin's
theorem applied to \(E/K\), and the Tower Law gives
\[
[K:F]=\frac{[E:F]}{[E:K]}
=\frac{|G|}{|\operatorname{Gal}(E/K)|}.
\]
\end{proof}

The reversal is the central visual idea: a larger field has fewer
automorphisms that fix it.
\[
\begin{tikzcd}[column sep=large,row sep=small]
E \arrow[d,phantom,"\longleftrightarrow" description] & \{1\} \\
K \arrow[u] \arrow[d,phantom,"\longleftrightarrow" description] & H \arrow[u] \\
F \arrow[u] & G \arrow[u]
\end{tikzcd}
\]

\begin{proposition}[Normality refinement]\label{prop:finite-normality-refinement}
For an intermediate field \(K\),
\[
K/F\text{ is Galois}
\quad\Longleftrightarrow\quad
\operatorname{Gal}(E/K)\trianglelefteq G.
\]
When these conditions hold, restriction induces
\[
G/\operatorname{Gal}(E/K)\cong\operatorname{Gal}(K/F).
\]
\end{proposition}
\begin{proof}
Write \(H=\operatorname{Gal}(E/K)\).  First,
\[
K/F\text{ is normal}\quad\Longleftrightarrow\quad
\sigma(K)=K\text{ for every }\sigma\in G. \tag{12.1}
\]
Normality makes every conjugate of an element of \(K\) remain in \(K\).
Conversely, extend any \(F\)-embedding of \(K\) to \(E\) by
Theorem~\ref{thm:embedding-extension}; normality of \(E/F\) makes the
extension an automorphism, and stability then proves normality of \(K/F\).

For \(\sigma\in G\),
\[
\sigma H\sigma^{-1}=\operatorname{Gal}(E/\sigma(K)).
\]
The finite correspondence and (12.1) show that \(H\) is normal exactly when
\(K/F\) is normal.  Every intermediate extension is separable, so this is
equivalent to Galoisness.

When \(K\) is stable, restriction
\[
\operatorname{res}:G\longrightarrow\operatorname{Gal}(K/F)
\]
is well-defined and has kernel \(H\).  An \(F\)-automorphism of \(K\) extends
to an \(F\)-embedding of \(E\), and normality of \(E/F\) makes the extension
an automorphism.  Thus restriction is onto, and the First Isomorphism Theorem
gives the quotient isomorphism.
\end{proof}

\begin{example}[The full lattice for a cubic splitting field]
Put
\[
r=\sqrt[3]2,
\qquad
E=\mathbb Q(r,\zeta_3).
\]
The roots of \(x^3-2\) are \(r,\zeta_3r,\zeta_3^2r\), and
\(\operatorname{Gal}(E/\mathbb Q)\cong S_3\) acts by permuting them.  Let
\[
K_1=\mathbb Q(r),\qquad
K_2=\mathbb Q(\zeta_3r),\qquad
K_3=\mathbb Q(\zeta_3^2r).
\]
The intermediate fields have the following lattice; each displayed
intermediate field is contained in \(E\), and the arrows are inclusions.
\[
\begin{tikzcd}[column sep=large,row sep=large]
&& E &&\\
K_1 \arrow[urr] & K_2 \arrow[ur] & K_3 \arrow[u] &
\mathbb Q(\sqrt{-3}) \arrow[ul] \\
&& \mathbb Q \arrow[ull] \arrow[ul] \arrow[u] \arrow[ur] &&
\end{tikzcd}
\]
The order-two subgroup stabilizing the indicated root fixes the corresponding
\(K_i\), so the three order-two subgroups correspond to the three cubic
subfields.  The alternating subgroup \(A_3\) fixes
\(\mathbb Q(\sqrt{-3})=\mathbb Q(\zeta_3)\), and the trivial and full
subgroups correspond to \(E\) and \(\mathbb Q\), respectively:
\[
\begin{array}{c|c|c}
\text{field (degree)} & \text{subgroup} & \text{order}\\ \hline
E\ (6) & \{1\} & 1\\
K_1,\ K_2,\ K_3\ (3) & \text{three root stabilizers} & 2\\
\mathbb Q(\sqrt{-3})\ (2) & A_3 & 3\\
\mathbb Q\ (1) & S_3 & 6
\end{array}
\]
The three order-two subgroups are not normal.  Accordingly the three cubic
fields are not Galois over \(\mathbb Q\), whereas
\(\mathbb Q(\sqrt{-3})/\mathbb Q\) is Galois.  This is a useful contrast
between the field lattice itself and the additional normality information
carried by the group lattice.
\end{example}

\begin{example}[A biquadratic contrast]
For
\[
E=\mathbb Q(\sqrt2,\sqrt3),
\]
every \(\mathbb Q\)-automorphism independently changes the signs of
\(\sqrt2\) and \(\sqrt3\).  Hence
\[
\operatorname{Gal}(E/\mathbb Q)\cong C_2\oplus C_2.
\]
Its three order-two subgroups fix, respectively,
\[
\mathbb Q(\sqrt2),\qquad
\mathbb Q(\sqrt3),\qquad
\mathbb Q(\sqrt6).
\]
Unlike the \(S_3\) example, every subgroup is normal because the group is
abelian.  Thus all three quadratic intermediate fields are Galois over
\(\mathbb Q\).  The two examples show that the same-looking degree pattern
can have very different normality behavior.
\end{example}

\begin{example}[A cyclic correspondence]
\(\operatorname{Gal}(\mathbb Q(\zeta_5)/\mathbb Q)\cong C_4\).  Its unique
subgroup of order two fixes \(\mathbb Q(\sqrt5)\).  All subgroups are normal,
so every intermediate field is Galois over \(\mathbb Q\).
\end{example}

\subsection*{Optional Application: A Galois-Theoretic Proof of the Fundamental Theorem of Algebra}

Chapter~8 used the Fundamental Theorem of Algebra to obtain an eigenvalue of a
complex operator.  We now discharge that dependence using the field and group
theory developed so far.  Apart from the standard elementary properties of
the real numbers, including square roots of nonnegative reals, the only
substantive real-analysis input is the Intermediate Value Theorem.  The
needed facts are developed in \emph{Lecture Notes on Mathematical Analysis}
\cite{MA}: polynomials are continuous (\S5.2), and the Intermediate Value
Theorem is Theorem~5.14 (\S5.4).  Hence every real polynomial of odd degree
has a real root.  Indeed, if
\[
p(x)=a_{2m+1}x^{2m+1}+\cdots,
\qquad
a_{2m+1}\ne0,
\]
then, for sufficiently large \(R>0\), leading-term behavior makes \(p(R)\)
and \(p(-R)\) have opposite signs.  Continuity and the Intermediate Value
Theorem therefore give a point \(c\in(-R,R)\) with \(p(c)=0\).

\begin{applicationtheorem}[Fundamental Theorem of Algebra]
Every nonconstant polynomial in \(\mathbb C[x]\) has a root in \(\mathbb C\).
Equivalently, \(\mathbb C\) is algebraically closed.
\end{applicationtheorem}
\begin{proof}
We first show that every finite extension of \(\mathbb R\) has degree a power
of \(2\).  Let \(M/\mathbb R\) be finite.  Characteristic zero makes
\(M/\mathbb R\) separable, so the
\hyperref[thm:finite-galois-closure]{finite Galois closure theorem} supplies
a finite Galois extension \(L/\mathbb R\) containing \(M\).  Put
\(G=\operatorname{Gal}(L/\mathbb R)\).
If an odd prime divided \(|G|\), let \(P\) be a Sylow \(2\)-subgroup.  Then
\[
[G:P]
\]
is odd and greater than one.  By the finite Galois correspondence,
\(L^P/\mathbb R\) has this odd degree.  The primitive element theorem writes
\(L^P=\mathbb R(\beta)\), so the minimal polynomial of \(\beta\) has odd
degree.  It has a root in \(\mathbb R\) by the stated external fact, which
contradicts irreducibility unless the degree is one.  Thus \(G\) is a
\(2\)-group.  Since
\[
[M:\mathbb R]\mid[L:\mathbb R]=|G|,
\]
the claimed power-of-two conclusion follows.

Fix an algebraic closure \(\overline{\mathbb R}\) and an element
\(i\in\overline{\mathbb R}\) with \(i^2=-1\), so that
\(\mathbb C=\mathbb R(i)\).  Every nontrivial quadratic extension of
\(\mathbb R\) is this field.  Indeed, if \(\mathbb R(\alpha)/\mathbb R\) is
quadratic, the minimal polynomial of \(\alpha\) is an irreducible quadratic
with negative discriminant.  Completing the square therefore produces an
element of the extension whose square is \(-1\): after division by the
leading coefficient, a polynomial \(x^2+bx+c\) gives
\[
\left(\alpha+\frac b2\right)^2=\frac{b^2-4c}{4}<0,
\]
and rescaling by the positive real square root of its negative gives a square
root of \(-1\).  It is \(i\) or \(-i\), so the extension is
\(\mathbb R(i)=\mathbb C\).

We next show directly that every element of \(\mathbb C\) has a square root.
For \(z=a+bi\), put
\[
r=\sqrt{a^2+b^2}.
\]
If \(r+a>0\), set
\[
x=\sqrt{\frac{r+a}{2}},
\qquad
y=\frac{b}{2x}.
\]
Then \(x^2-y^2=a\) and \(2xy=b\), so \((x+iy)^2=z\).  If \(r+a=0\), then
\(b=0\) and \(a\leq0\), and \((i\sqrt{-a})^2=z\).  Consequently the
quadratic formula shows that every quadratic polynomial over \(\mathbb C\)
has a root.  Hence \(\mathbb C\) has no nontrivial quadratic extension.

We claim that no finite Galois extension \(L/\mathbb R\) has degree greater
than \(2\).  Otherwise its group \(G\) is a \(2\)-group of order greater than
two.  Every nontrivial finite \(2\)-group has a normal subgroup of index two:
for order two this is clear, and in general the center theorem and Cauchy's
theorem supply a central subgroup of order two; induction in the quotient
then gives the asserted subgroup.  Choose a normal subgroup \(H\trianglelefteq
G\) of index two.  The finite correspondence makes
\[
L^H/\mathbb R
\]
a quadratic extension, hence \(L^H=\mathbb C\).  Now \(H=
\operatorname{Gal}(L/\mathbb C)\) is a nontrivial \(2\)-group, so it has a
subgroup \(J\) of index two.  The finite correspondence for \(L/\mathbb C\)
then makes
\[
L^J/\mathbb C
\]
a quadratic extension, a contradiction.  Therefore \([L:\mathbb R]\leq2\).

Finally, let \(\alpha\in\overline{\mathbb R}\).  The finite extension
\(\mathbb R(\alpha)/\mathbb R\) is separable, so the
\hyperref[thm:finite-galois-closure]{finite Galois closure theorem} gives a
finite Galois extension containing it, whose degree is at most two by the
preceding argument.  The Tower Law therefore
gives \([\mathbb R(\alpha):\mathbb R]\leq2\), so it is either \(\mathbb R\)
or, in the quadratic case, \(\mathbb C\).  Thus every \(\alpha\) belongs to \(\mathbb C\), and
\(\overline{\mathbb R}=\mathbb C\).  Hence \(\mathbb C\) is algebraically
closed.
\end{proof}

\begin{unremark}
Chapter~8 met the Fundamental Theorem of Algebra through complex linear
algebra, while its classical proof belongs to complex analysis.  The present
proof instead reveals field extensions, finite \(2\)-groups, and the Galois
correspondence behind the same theorem.
\end{unremark}

\section{Radicals and Classical Constructions}

\begin{definition}
A \emph{radical extension} is an extension contained in a tower
\[
F=K_0\subseteq K_1\subseteq\cdots\subseteq K_r,
\]
where \(K_i=K_{i-1}(\alpha_i)\) and
\(\alpha_i^{n_i}\in K_{i-1}\) for some \(n_i>0\).  A polynomial is
\emph{solvable by radicals} if its splitting field is contained in a radical
extension of its coefficient field.
\end{definition}

Roots of unity govern the conjugates of an \(n\)-th root, which explains the
following honest part of the radicals--groups connection.

\begin{proposition}[A root-of-unity radical direction]\label{prop:radical-solvable}
Let
\[
F=K_0\subseteq K_1\subseteq\cdots\subseteq K_r
\]
be a tower such that every \(K_i/F\) is finite Galois and
\[
K_i=K_{i-1}(\alpha_i),
\qquad
\alpha_i^{n_i}\in K_{i-1}.
\]
Assume that \(K_{i-1}\) contains all \(n_i\)-th roots of unity.  Then
\[
\operatorname{Gal}(K_r/F)
\]
is solvable.
\end{proposition}
\begin{proof}
If \(\alpha_i=0\), then \(K_i=K_{i-1}\), so the relevant Galois group is
trivial.  Otherwise define
\[
\chi_i:\operatorname{Gal}(K_i/K_{i-1})\longrightarrow\mu_{n_i},
\qquad
\chi_i(\sigma)=\frac{\sigma(\alpha_i)}{\alpha_i}.
\]
The equation \(\alpha_i^{n_i}\in K_{i-1}\) shows that \(\chi_i(\sigma)\)
is an \(n_i\)-th root of unity.  Since \(\mu_{n_i}\subseteq K_{i-1}\), every
\(\sigma\) fixes it pointwise; consequently
\[
\chi_i(\sigma\tau)
=\frac{\sigma\tau(\alpha_i)}{\alpha_i}
=\frac{\sigma\bigl(\chi_i(\tau)\alpha_i\bigr)}{\alpha_i}
=\chi_i(\tau)\chi_i(\sigma).
\]
The target is abelian, so this is the homomorphism law.  Moreover,
\(\chi_i(\sigma)=1\) means that \(\sigma\) fixes \(\alpha_i\); since
\(K_i=K_{i-1}(\alpha_i)\), it is then the identity.  Thus \(\chi_i\) is
injective.  Therefore \(\operatorname{Gal}(K_i/K_{i-1})\) is isomorphic to
a subgroup of the abelian group \(\mu_{n_i}\), and hence is abelian.

Put
\[
G_i=\operatorname{Gal}(K_r/K_i).
\]
The hypothesis that \(K_i/F\) is Galois makes \(G_i\) normal in
\(G_0=\operatorname{Gal}(K_r/F)\).  The finite normality refinement
(Proposition~\ref{prop:finite-normality-refinement}) gives
\[
G_{i-1}/G_i\cong\operatorname{Gal}(K_i/K_{i-1}).
\]
Consequently
\[
G_0\trianglerighteq G_1\trianglerighteq\cdots\trianglerighteq G_r=1
\]
is a normal series with abelian factors.  Thus \(G_0\) is solvable.
\end{proof}

\begin{remark}
The classical theorem in characteristic zero says that a polynomial is
solvable by radicals if and only if its Galois group is solvable.  The converse
requires Kummer theory, which is outside these notes.  The proposition above
captures the group-theoretic mechanism after a radical tower has been placed
inside a suitable Galois tower; it does not prove the converse.  Accordingly,
we do not claim a proof here of the general quintic theorem.
\end{remark}

\begin{definition}
A real number is \emph{constructible over \(\mathbb Q\)} if it belongs to a
tower of subfields of \(\mathbb R\)
\[
\mathbb Q=K_0\subseteq K_1\subseteq\cdots\subseteq K_r\subseteq\mathbb R
\]
whose successive degrees are at most two.  This is the field-theoretic form
of straightedge-and-compass
construction: line intersections use field operations and circle
intersections require a quadratic equation.
\end{definition}

The algebraic definition records a necessary consequence of the ordinary
geometric construction rules.  Indeed, the intersection of two lines has
coordinates obtained by field operations from the preceding coordinates.
After eliminating one variable, the intersection of a line and a circle, or
of two circles, is obtained by solving a quadratic equation.  Hence every
ordinary straightedge-and-compass construction produces coordinates in a
tower whose successive degrees are at most two.  This one direction is all
that is needed for the impossibility results below.

\begin{proposition}[Degree obstruction]\label{prop:construction-degree}
If \(\alpha\) is constructible over \(\mathbb Q\), then
\[
[\mathbb Q(\alpha):\mathbb Q]
\]
divides a power of \(2\).
\end{proposition}
\begin{proof}
The construction operations give a tower with successive degrees at most two.
The Tower Law makes the degree of its top field a power of two, and the degree
of the subfield \(\mathbb Q(\alpha)\) divides that degree.
\end{proof}

\begin{corollary}
Doubling the cube and trisecting a general angle are impossible with
straightedge and compass.
\end{corollary}
\begin{proof}
\(\sqrt[3]2\) has irreducible minimal polynomial \(X^3-2\).  For angle
trisection, a \(60^\circ\) angle is constructible.  If arbitrary angle
trisection were possible, then \(20^\circ\) would be constructible.  But
\(\cos20^\circ\) satisfies
\[
8X^3-6X-1=0,
\]
which is irreducible over \(\mathbb Q\) by the rational-root test.  Both
degrees are \(3\), so Proposition~\ref{prop:construction-degree} applies.
Squaring the circle uses the transcendence of \(\pi\), outside the present
scope.
\end{proof}

\section{Direct and Inverse Limits}

Infinite Galois theory is assembled from finite stages.  Direct limits join
objects that grow; inverse limits record compatible information at all finite
quotients.  The common indexing language is designed so that a finite number
of comparisons can always be made at one later stage.

\begin{definition}
A partially ordered set \(I\) is \emph{directed} if it is nonempty and every
pair \(i,j\) has an upper bound \(k\).  Examples are \(\mathbb N\), finite
subsets of a set ordered by inclusion, and finite intermediate extensions
ordered by inclusion.  Repeatedly applying the pair condition gives a common
upper bound for every finite collection.  This is exactly what lets a direct
limit turn a finite sum of representatives into one representative later on.
\end{definition}

\begin{definition}
A \emph{direct system} of \(R\)-modules over \(I\) is a family
\((M_i)_{i\in I}\) and maps
\[
\phi_i^j:M_i\longrightarrow M_j\qquad(i\leq j)
\]
with \(\phi_i^i=1_{M_i}\) and
\[
\phi_j^k\phi_i^j=\phi_i^k
\qquad(i\leq j\leq k).
\]
For \(i\leq j\leq k\), this says that the following triangle commutes:
\[
\begin{tikzcd}[column sep=large,row sep=large]
M_i \arrow[rr,"\phi_i^k"] \arrow[dr,"\phi_i^j"'] && M_k \\
& M_j \arrow[ur,"\phi_j^k"'] &
\end{tikzcd}
\]
Thus passing from \(i\) to \(k\) does not depend on whether one pauses at
\(j\).
\end{definition}

\Needspace{24\baselineskip}
\begin{definition}
A \emph{direct limit} of a direct system is an \(R\)-module
\(\varinjlim_iM_i\), together with maps
\[
\iota_i:M_i\longrightarrow\varinjlim_iM_i,
\]
such that \(\iota_j\phi_i^j=\iota_i\) for \(i\leq j\), and satisfying the
following universal property: whenever maps \(f_i:M_i\to N\) satisfy
\[
f_j\phi_i^j=f_i,
\]
there is a unique map
\[
u:\varinjlim_iM_i\longrightarrow N
\]
with \(u\iota_i=f_i\) for every \(i\).
\end{definition}

\[
\begin{tikzcd}[column sep=large,row sep=large]
M_i \arrow[rr,"\phi_i^j"] \arrow[dr,"\iota_i"'] \arrow[ddr,bend right=35,"f_i"']&&
M_j \arrow[dl,"\iota_j"] \arrow[ddl,bend left=35,"f_j"]\\
&\varinjlim_iM_i \arrow[d,dashed,"u"]\\
&N
\end{tikzcd}
\]

\begin{proposition}[Uniqueness of direct limits]
Two direct limits of the same direct system are uniquely isomorphic by an
isomorphism commuting with every canonical map.
\end{proposition}
\begin{proof}
Let \((M,\iota_i)\) and \((M',\iota_i')\) be direct limits.  The universal
property of \(M\), applied to the compatible family \((\iota_i')\), gives a
unique \(u:M\to M'\) with \(u\iota_i=\iota_i'\).  Similarly there is a unique
\(v:M'\to M\) with \(v\iota_i'=\iota_i\).  Both \(vu\) and \(1_M\) have the
same composites with every \(\iota_i\), so uniqueness gives \(vu=1_M\).
Likewise \(uv=1_{M'}\).
\end{proof}

\begin{proposition}[Construction of direct limits]
Every direct system of \(R\)-modules has a direct limit.  It is
\[
\varinjlim_{i\in I}M_i=
\left(\bigoplus_{i\in I}M_i\right)\bigg/D,
\]
where \(D\) is generated by
\[
e_i(m)-e_j(\phi_i^j(m))\qquad(i\leq j,\ m\in M_i).
\]
\end{proposition}
\begin{proof}
Let \(\iota_i(m)\) be the class of \(e_i(m)\).  The displayed relations give
\(\iota_j\phi_i^j=\iota_i\).  A compatible family \(f_i:M_i\to N\) defines
\[
\sum_i e_i(m_i)\longmapsto\sum_i f_i(m_i).
\]
The sums are finite because the source is a direct sum, and each defining
relation maps to
\[
f_i(m)-f_j(\phi_i^j(m))=0.
\]
Hence the map descends to a map \(u\) from the quotient.  It satisfies
\(u\iota_i=f_i\).  Since every quotient class is a finite sum of classes from
the summands, these identities also force uniqueness.
\end{proof}

\begin{proposition}[Maps induced on direct limits]\label{prop:direct-limit-maps}
Suppose \((M_i,\phi_i^j)\) and \((N_i,\psi_i^j)\) are direct systems and maps
\(\rho_i:M_i\to N_i\) satisfy
\[
\rho_j\phi_i^j=\psi_i^j\rho_i
\qquad(i\leq j).
\]
Then there is a unique map
\[
\rho:\varinjlim_iM_i\longrightarrow\varinjlim_iN_i
\]
with \(\rho\iota_i=\kappa_i\rho_i\), where \(\kappa_i\) are the canonical maps
for the second limit.
\end{proposition}
\begin{proof}
The maps \(\kappa_i\rho_i\) are compatible, since
\[
\kappa_j\rho_j\phi_i^j
=\kappa_j\psi_i^j\rho_i
=\kappa_i\rho_i.
\]
The universal property of \(\varinjlim_iM_i\) gives exactly the asserted map.
\end{proof}

\begin{proposition}[Element criteria]\label{prop:direct-limit-elements}
Every element of a direct limit has the form \(\iota_i(m)\) for some
\(i\in I\) and \(m\in M_i\).  In addition,
\[
\iota_i(m)=0
\quad\Longleftrightarrow\quad
\phi_i^j(m)=0\text{ for some }j\geq i,
\]
and
\[
\iota_i(m)=\iota_j(n)
\quad\Longleftrightarrow\quad
\phi_i^k(m)=\phi_j^k(n)\text{ for some }k\geq i,j.
\]
\end{proposition}
\begin{proof}
A quotient element is represented by
\[
\sum_{\ell=1}^re_{i_\ell}(m_\ell).
\]
Choose \(k\) above every \(i_\ell\).
Compatibility gives
\[
\sum_{\ell=1}^r\iota_{i_\ell}(m_\ell)
=\iota_k\left(\sum_{\ell=1}^r\phi_{i_\ell}^k(m_\ell)\right),
\]
which proves the first assertion.  The reverse implications in the displayed
criteria are compatibility.  For the forward implications, the asserted
zero or equality in the quotient is a finite sum of defining relations.
Choose an index above all of the finitely many stages occurring in this
expression.  Mapping every summand to that stage cancels the relations and
gives the desired zero or equality.
\end{proof}

\begin{example}[A directed union written as fractions]
Consider the direct system indexed by \(n\geq0\)
\[
\mathbb Z\xrightarrow{\times2}\mathbb Z\xrightarrow{\times2}\mathbb Z
\xrightarrow{\times2}\cdots.
\]
Its direct limit is naturally
\[
\mathbb Z[1/2]
=\left\{\frac{m}{2^n}:m\in\mathbb Z,\ n\geq0\right\}.
\]
Indeed, send the integer \(m\) at stage \(n\) to \(m/2^n\).  The transition
map sends \(m\) to \(2m\), and
\[
\frac{2m}{2^{n+1}}=\frac{m}{2^n},
\]
so these maps are compatible.  They induce a map from the direct limit to
\(\mathbb Z[1/2]\), and it is onto by definition.  If two representatives
\(m\) at stage \(n\) and \(a\) at stage \(r\) give the same fraction, then
after multiplying their equality by a sufficiently large power of \(2\),
their images agree at a common later stage.  Proposition~\ref{prop:direct-limit-elements}
therefore makes the induced map one-to-one.  This example turns the abstract
identification of forward representatives into the familiar rule for
equivalent fractions.
\end{example}

\begin{corollary}[Directed unions]
If \(M_i\subseteq M_j\) for \(i\leq j\) and the transition maps are
inclusions, then
\[
\varinjlim_iM_i\cong\bigcup_iM_i.
\]
\end{corollary}
\begin{proof}
The inclusions into the union are compatible.  The induced map is onto, and
the zero criterion in Proposition~\ref{prop:direct-limit-elements} makes it
one-to-one.
\end{proof}

\begin{theorem}[Exactness of directed direct limits]\label{thm:direct-limit-exact}
Suppose that
\[
A_i\xrightarrow{u_i}B_i\xrightarrow{v_i}C_i
\]
is exact for every \(i\), and all transition maps commute with these maps.
Then
\[
\varinjlim_iA_i\longrightarrow\varinjlim_iB_i
\longrightarrow\varinjlim_iC_i
\]
is exact.  Thus directed direct limits preserve injections and surjections
when these occur stagewise.
\end{theorem}
\begin{proof}
The composite is zero at every stage.  Conversely, suppose that
\(\iota_i(b)\) maps to zero.  The zero criterion gives \(j\geq i\) with
\[
v_j(\phi_i^j(b))=0.
\]
At stage \(j\), exactness supplies \(a\in A_j\) with
\(u_j(a)=\phi_i^j(b)\).  Therefore \(\iota_i(b)\) is the image of
\(\iota_j(a)\).  Applying this to sequences beginning or ending in zero gives
the final statement.
\end{proof}

\begin{definition}
An \emph{inverse system} of \(R\)-modules over the directed set \(I\) is a
family \((N_i)_{i\in I}\) and maps
\[
\psi_i^j:N_j\longrightarrow N_i\qquad(i\leq j)
\]
such that \(\psi_i^i=1_{N_i}\) and
\[
\psi_i^j\psi_j^k=\psi_i^k
\qquad(i\leq j\leq k).
\]
For \(i\leq j\leq k\), the corresponding reversed triangle commutes:
\[
\begin{tikzcd}[column sep=large,row sep=large]
& N_j \arrow[dr,"\psi_i^j"] & \\
N_k \arrow[ur,"\psi_j^k"] \arrow[rr,"\psi_i^k"'] && N_i
\end{tikzcd}
\]
Thus information at a later stage restricts consistently to every earlier
stage.  Unlike a direct limit, an inverse limit will retain only collections
of data that agree under all these restrictions.
\end{definition}

\begin{definition}
An \emph{inverse limit} of an inverse system is an \(R\)-module
\[
\varprojlim_iN_i,
\]
together with maps \(\pi_i:\varprojlim_iN_i\to N_i\) satisfying
\[
\psi_i^j\pi_j=\pi_i
\qquad(i\leq j),
\]
with the following universal property.  Whenever maps \(g_i:X\to N_i\)
satisfy \(\psi_i^jg_j=g_i\), there is a unique map
\[
u:X\longrightarrow\varprojlim_iN_i
\]
for which \(\pi_i u=g_i\) for every \(i\).
\end{definition}

\[
\begin{tikzcd}[column sep=large,row sep=large]
& X \arrow[d,dashed,"u"] \arrow[ddl,bend right=35,"g_i"']
  \arrow[ddr,bend left=35,"g_j"] & \\
& \varprojlim_iN_i \arrow[dl,"\pi_i"] \arrow[dr,"\pi_j"'] & \\
N_i && N_j \arrow[ll,"\psi_i^j"]
\end{tikzcd}
\]

\begin{proposition}[Uniqueness of inverse limits]
Two inverse limits of the same inverse system are uniquely isomorphic by an
isomorphism commuting with all projections.
\end{proposition}
\begin{proof}
Let \((N,\pi_i)\) and \((N',\pi_i')\) be inverse limits.  The universal
property of \(N'\), applied to the compatible family \((\pi_i)\), produces a
unique \(u:N\to N'\) with \(\pi_i'u=\pi_i\).  Similarly there is a unique
\(v:N'\to N\) with \(\pi_iv=\pi_i'\).  The maps \(vu\) and \(1_N\) have the
same composites with every \(\pi_i\); uniqueness gives \(vu=1_N\).  The
same argument gives \(uv=1_{N'}\).
\end{proof}

\begin{proposition}[Construction of inverse limits]
Every inverse system of \(R\)-modules has an inverse limit.  It is
\[
\varprojlim_{i\in I}N_i=
\left\{(n_i)\in\prod_{i\in I}N_i:
\psi_i^j(n_j)=n_i\text{ whenever }i\leq j\right\}.
\]
\end{proposition}
\begin{proof}
The displayed set is a submodule of the product: the compatibility equations
are preserved under addition and scalar multiplication.  Its coordinate
projections are compatible.  A compatible family \(g_i:X\to N_i\) gives the
map
\[
x\longmapsto(g_i(x))_{i\in I}
\]
to this submodule, because the assumed equations say precisely that this
tuple is compatible.  Its \(i\)-th coordinate is \(g_i(x)\), and a map to a
submodule of a product is determined by its coordinates.  This proves the
universal property.  Thus a direct limit is a quotient of a direct sum,
whereas an inverse limit is a compatible submodule of a direct product.
\end{proof}

\begin{proposition}[Maps induced on inverse limits]
Suppose \((M_i,\phi_i^j)\) and \((N_i,\psi_i^j)\) are inverse systems and
maps \(\rho_i:M_i\to N_i\) satisfy
\[
\psi_i^j\rho_j=\rho_i\phi_i^j
\qquad(i\leq j).
\]
Then there is a unique map
\[
\rho:\varprojlim_iM_i\longrightarrow\varprojlim_iN_i
\]
whose \(i\)-th coordinate is \(\rho_i\) applied to the \(i\)-th coordinate.
\end{proposition}
\begin{proof}
For \((m_i)_i\in\varprojlim_iM_i\), the compatibility calculation
\[
\psi_i^j\bigl(\rho_j(m_j)\bigr)
=\rho_i\bigl(\phi_i^j(m_j)\bigr)
=\rho_i(m_i)
\]
shows that \((\rho_i(m_i))_i\) is a compatible tuple.  The asserted formula
therefore defines \(\rho\); uniqueness follows from the coordinate
projections.
\end{proof}

\begin{example}[The \(2\)-adic and profinite integers]
The inverse limit of the reductions
\[
\mathbb Z/2^{n+1}\mathbb Z\longrightarrow\mathbb Z/2^n\mathbb Z
\]
is \(\mathbb Z_2\).  An element is a compatible sequence
\[
(a_n\bmod2^n)_{n\geq1};
\]
the ordinary integer \(5\), for instance, gives
\((1\bmod2,1\bmod4,5\bmod8,5\bmod16,\ldots)\).  Ordered by divisibility,
\[
\widehat{\mathbb Z}=\varprojlim_{n\geq1}\mathbb Z/n\mathbb Z.
\]
The former retains only powers of \(2\); the latter retains every modulus.
\(\mathbb Z_2\) also contains compatible sequences not coming from a
nonnegative ordinary integer in any fixed finite stage.  For example,
\[
(-1\bmod2,\,-1\bmod4,\,-1\bmod8,\,-1\bmod16,\ldots)
\]
is compatible and represents \(-1\), while the proposed first two entries
\[
(0\bmod2,\ 1\bmod4)
\]
cannot occur in a compatible sequence, because \(1\bmod4\) reduces to
\(1\bmod2\), not to \(0\bmod2\).  Compatibility is therefore a real
restriction, not merely a way of listing residues.
\end{example}

\begin{proposition}[Left exactness of inverse limits]
Inverse limits preserve kernels and hence are left exact, but need not
preserve surjectivity at the right end.
\end{proposition}
\begin{proof}
In the product description, a compatible family lies in the kernel of an
induced map exactly when every coordinate lies in the corresponding kernel.
The same compatibility equations remain, so the kernel is the inverse limit
of the kernels.  Compatible lifts of a compatible family need not exist,
which is why right exactness is not claimed.
\end{proof}

\begin{example}[Stagewise surjections need not lift compatibly]
For every \(n\geq1\), let \(X_n=\mathbb Z\) with identity transition maps,
let
\[
Y_n=\mathbb Z/2^n\mathbb Z
\]
with the reduction maps, and let \(q_n:X_n\to Y_n\) be the quotient map.
Each \(q_n\) is surjective, and the maps commute with the transitions.  Yet
\[
\varprojlim_nX_n\cong\mathbb Z,
\qquad
\varprojlim_nY_n\cong\mathbb Z_2,
\]
and the induced map \(\mathbb Z\to\mathbb Z_2\) is not surjective.  For
example, the compatible sequence
\[
(1\bmod2,\ 1\bmod4,\ 5\bmod8,\ 5\bmod16,\ 21\bmod32,\ldots)
\]
is not the image of an ordinary integer.  Its \(2\)-adic binary expansion has
infinitely many zeros and ones.  A nonnegative ordinary integer has only
finitely many nonzero binary digits, whereas a negative ordinary integer has a
\(2\)-adic binary expansion that is eventually all ones.  Thus compatible
stagewise lifts need not be selectable simultaneously.
\end{example}

\begin{remark}[Looking Ahead: the Mittag--Leffler condition]
The preceding proposition identifies the difficulty with inverse limits: one
can take compatible kernels coordinatewise, but compatible choices of
preimages need not exist.  For an inverse system \((M_i,\psi_i^j)\), the
\emph{Mittag--Leffler condition} says that for every fixed \(i\), there is a
\(j\geq i\) such that
\[
\operatorname{im}(\psi_i^k)=\operatorname{im}(\psi_i^j)
\qquad\text{for every }k\geq j.
\]
In other words, the images inside the fixed earlier module \(M_i\), as seen
from later and later stages, eventually stabilize.  Surjective transition
maps satisfy the condition because all these images are \(M_i\).  It also
holds when each \(M_i\) is finite as a set, since a descending family of
submodules of a fixed finite module stabilizes; for finite-dimensional vector
spaces, one may instead use the eventually constant nonincreasing dimensions
of the images.

Under standard hypotheses, a Mittag--Leffler condition on the kernel system
in an inverse system of short exact sequences restores surjectivity at the
right-hand end after taking inverse limits.  It is not needed for the present
Galois application: restriction maps between finite Galois stages are already
surjective.  The derived obstruction to exactness is usually denoted
\(\varprojlim{}^{1}\); its systematic treatment belongs to later homological
algebra.
\end{remark}

\section{Topological Preliminaries}

In finite Galois theory the group alone remembers all the relevant
information.  For an infinite extension, however, one must also remember when
automorphisms agree on larger and larger finite Galois subextensions.  Topology
is the concise language for this finite-stage approximation: in the Krull
topology, two automorphisms are close when they agree on a sufficiently large
finite Galois field.  We develop only the small amount of topology needed to
make this statement precise.

\subsection{Basic Point-Set Topology}

\begin{definition}
A \emph{topology} on a set \(X\) is a collection \(\tau\) of subsets of \(X\)
such that
\[
\varnothing,\ X\in\tau;
\]
an arbitrary union of members of \(\tau\) is in \(\tau\); and a finite
intersection of members of \(\tau\) is in \(\tau\).  The members of \(\tau\)
are the \emph{open sets}.  A set is \emph{closed} if its complement is open,
and a \emph{neighborhood} of \(x\in X\) is a set containing an open set that
contains \(x\).
\end{definition}

The discrete topology, in which every subset is open, is the important
example here: every finite Galois group will be given this topology.  At the
opposite extreme, the indiscrete topology has only \(\varnothing\) and \(X\)
open.  The usual topology on \(\mathbb R\) is familiar intuition, but no
analysis on \(\mathbb R\) is needed below.

\begin{definition}
A collection \(\mathcal B\) of open sets is a \emph{basis} if every open set
is a union of members of \(\mathcal B\).  A collection of neighborhoods of
\(x\) is a \emph{neighborhood basis at \(x\)} if every neighborhood of \(x\)
contains one of its members.  A map \(f:X\to Y\) is \emph{continuous} if
\(f^{-1}(U)\) is open whenever \(U\) is open.  A bijective continuous map
whose inverse is continuous is a \emph{homeomorphism}.
\end{definition}

For \(A\subseteq X\), its \emph{closure} \(\overline A\) is the intersection
of all closed subsets of \(X\) containing \(A\).  We say that \(A\) is
\emph{dense} when \(\overline A=X\).  For example, a nonempty subset of a
finite discrete space is closed, so it is dense exactly when it is the whole
space.

If \(Y\subseteq X\), the \emph{subspace topology} on \(Y\) declares
\[
U\subseteq Y\text{ open}\quad\Longleftrightarrow\quad
U=Y\cap V\text{ for some open }V\subseteq X.
\]
For a family of spaces, the product topology on
\[
\prod_{i\in I}X_i
\]
has as a basis the sets which restrict only finitely many coordinates to open
sets and leave every remaining coordinate unrestricted.  Thus a basic
neighborhood records finite information.  When every \(X_i\) is finite
discrete, it may simply prescribe the values of finitely many coordinates.
For example, in
\[
\prod_{n\geq1}\mathbb Z/2^n\mathbb Z,
\]
the condition
\[
a_1=1\bmod2,
\qquad
a_2=3\bmod4
\]
with all higher coordinates unrestricted is a basic open set of the full
product.  Compatibility is imposed only after one passes to the inverse-limit
subspace.
Inverse limits will be compatible subspaces of precisely such products.

\begin{definition}
A space is \emph{Hausdorff} if distinct points have disjoint neighborhoods.
It is \emph{compact} if every open cover has a finite subcover.  A space is
\emph{connected} if it cannot be written as the union of two disjoint nonempty
open sets, and \emph{totally disconnected} if every connected subset has at
most one point.
\end{definition}

Finite discrete spaces are Hausdorff, compact, and totally disconnected:
singletons separate points and are open, while an open cover of a finite set
has an evident finite subcover.

\begin{theorem}[Tychonoff's theorem {\rm(external)}]
An arbitrary product of compact spaces is compact in the product topology.
\end{theorem}

\begin{remark}[Foundational scope]
In standard set theory with the Axiom of Choice, Tychonoff's theorem gives
compactness for arbitrary products of compact spaces.  In fact, over ZF, this
full arbitrary-product form is equivalent to the Axiom of Choice.  The present
notes use it only as an external input for products of finite discrete spaces;
this is consistent with the earlier use of Zorn's Lemma in the construction of
algebraic closures.
\end{remark}

\subsection{Topological Groups}

\begin{definition}
A \emph{topological group} is a group \(G\) with a topology for which
\[
G\times G\longrightarrow G,\qquad (x,y)\longmapsto xy,
\]
and
\[
G\longrightarrow G,\qquad x\longmapsto x^{-1},
\]
are continuous.  A homomorphism of topological groups is \emph{continuous}
when it is continuous as a map of spaces; a \emph{topological-group
isomorphism} is a group isomorphism that is also a homeomorphism.
\end{definition}

\begin{proposition}
For \(g\in G\), left and right translation,
\[
L_g(x)=gx,\qquad R_g(x)=xg,
\]
are homeomorphisms of a topological group.  Consequently, the neighborhoods
of the identity determine all neighborhoods: if \(U\) is a neighborhood of
\(1\), then \(gU\) is a neighborhood of \(g\).
\end{proposition}
\begin{proof}
The map \(x\mapsto(g,x)\) is continuous, and left translation is its
composite with multiplication.  Similarly, \(x\mapsto(x,g)\) followed by
multiplication is right translation.  Their inverse translations are
\(L_{g^{-1}}\) and \(R_{g^{-1}}\), which are continuous by the same argument.
The final assertion follows by applying \(L_g\) to a neighborhood of \(1\).
\end{proof}

\begin{proposition}
If \(\phi:G\to H\) is a continuous homomorphism and \(H\) is discrete, then
\(\ker\phi\) is open.  If \(H\) is finite, this kernel has finite index.
Further, every subgroup with its subspace topology and every product of
topological groups with coordinatewise operations are topological groups.
\end{proposition}
\begin{proof}
Since \(\{1\}\) is open in a discrete group,
\[
\ker\phi=\phi^{-1}(\{1\})
\]
is open.  The First Isomorphism Theorem identifies \(G/\ker\phi\) with the
finite group \(\operatorname{im}\phi\).  The remaining assertions follow by
restricting the continuous group operations, respectively by checking
continuity coordinatewise on basic product opens.
\end{proof}

\subsection{Profinite Groups}

\begin{definition}
A \emph{profinite group} is an inverse limit of finite discrete groups, with
the subspace topology inherited from
\[
\prod_{i\in I}G_i.
\]
The topology says that two elements are close when they have the same image at
many finite stages.  The next proposition makes this intuition precise in the
particularly useful form of a neighborhood basis.
\end{definition}

One may think of a profinite group as an infinite symmetry group observed
through compatible finite stages.  No single stage sees the entire element,
but a compatible answer at every stage does.
The prototype is \(\widehat{\mathbb Z}\): an element gives a residue class
modulo every positive integer, with all reductions agreeing.  This viewpoint
is exactly suited to an infinite algebraic extension, whose individual
elements and automorphisms can be tested on finite subextensions.

\begin{proposition}
An inverse limit of finite discrete groups is a closed topological subgroup of
their product.  Consequently every profinite group is compact, Hausdorff, and
totally disconnected.
\end{proposition}
\begin{proof}
The product is a topological group and is compact by Tychonoff's theorem.  It
is Hausdorff: distinct points differ in one coordinate, and disjoint open
neighborhoods in that finite discrete factor pull back to disjoint product
neighborhoods.  It is totally disconnected for the same clopen-cylinder
reason: if \(x\ne y\), choose a coordinate where they differ; fixing that
coordinate to \(x_i\) gives a clopen set containing \(x\) but not \(y\), so
no connected subset can contain both points.  For \(i\leq j\), define
\[
\rho_i^j:\prod_kG_k\longrightarrow G_i\times G_i,
\qquad
(g_k)_k\longmapsto\bigl(\psi_i^j(g_j),g_i\bigr).
\]
This map is continuous.  The compatibility condition
\(\psi_i^j(g_j)=g_i\) is the inverse image under \(\rho_i^j\) of the
diagonal, which is closed because the finite discrete target is Hausdorff.
Intersecting these closed conditions over all \(i\leq j\) gives the inverse
limit.  A closed subspace of a compact space is compact, since an open cover
of it extends by its open complement to an open cover of the ambient space.
Hausdorffness and total disconnectedness pass to subspaces: separating open
sets and clopen cylinders may simply be intersected with the subspace.
\end{proof}

\begin{proposition}[The projection-kernel basis]\label{prop:profinite-open-subgroups}
Let
\[
G=\varprojlim_{i\in I}G_i
\]
be a profinite group, indexed by a directed set.  The subgroups
\[
\ker(\pi_i)\qquad(i\in I)
\]
form a neighborhood basis at the identity.  They are open normal subgroups,
and every open subgroup has finite index and is closed.
\end{proposition}
\begin{proof}
In the product topology, a basic neighborhood of the identity restricts only
finitely many coordinates, say those indexed by \(i_1,\ldots,i_r\).  Choose
\(k\) above all these indices.  Compatibility implies
\[
\ker(\pi_k)\subseteq
\bigcap_{\ell=1}^r\ker(\pi_{i_\ell}),
\]
so one projection kernel lies inside the given neighborhood.  Conversely,
each \(\ker(\pi_i)\) is open because it is the inverse image of the open
singleton \(\{1\}\subseteq G_i\).  It is normal, and its quotient is a
subgroup of the finite group \(G_i\).  An open subgroup contains one such
kernel, hence has finite index.  Its complement is a finite union of its open
cosets, so it is closed.
\end{proof}

\begin{example}[The topology of \(\mathbb Z_p\)]
For a prime \(p\), define
\[
\mathbb Z_p=\varprojlim_{n\geq1}\mathbb Z/p^n\mathbb Z
\]
using the usual reduction maps.  The projection
\[
\mathbb Z_p\longrightarrow\mathbb Z/p^n\mathbb Z
\]
has kernel customarily denoted \(p^n\mathbb Z_p\).  Thus
\[
p^n\mathbb Z_p\qquad(n\geq1)
\]
is a neighborhood basis at \(0\).  In particular, congruence modulo a high
power of \(p\) is the topological notion of being close in \(\mathbb Z_p\).
This is the additive prototype for the Krull topology below.
\end{example}

\section{Infinite Galois Theory}

Let \(E/F\) be algebraic Galois.  Each element of \(E\) lies in a finite
Galois subextension: take the splitting field in \(E\) of its separable
minimal polynomial.  If \(K/F\) and \(L/F\) are finite Galois subextensions,
their \emph{compositum} \(KL\) is the smallest subfield of \(E\) containing
both.  It is finite over \(F\), separable because it lies in the separable
extension \(E/F\), and normal because \(K\) and \(L\) are splitting fields
over \(F\) of finitely many separable polynomials, while \(KL\) is the
splitting field of the product of those polynomials.  Thus \(KL/F\) is again
finite Galois and contains both fields.  These subextensions are therefore
directed by compositum, and
\[
E=\bigcup_{K/F\ \mathrm{finite\ Galois}}K.
\]
For the running finite-field example, the finite stages inside
\(\mathbb F_q^{\mathrm{sep}}\) are the fields \(\mathbb F_{q^n}\), ordered
by divisibility of \(n\).  Restriction from a larger finite field to a smaller
one is the concrete model for every inverse-limit transition map used below.

\begin{lemma}[Extension to an algebraic closure]\label{lem:closure-auto-extension}
Let \(L/F\) be algebraic inside a fixed algebraic closure \(\overline F\).
Every \(F\)-embedding
\[
\sigma:L\longrightarrow\overline F
\]
extends to an element of \(\operatorname{Aut}_F(\overline F)\).
\end{lemma}
\begin{proof}
The embedding-extension theorem (Theorem~\ref{thm:embedding-extension}) extends \(\sigma\) to an \(F\)-embedding
\[
\widetilde\sigma:\overline F\longrightarrow\overline F.
\]
Its image is an algebraically closed subfield.  Since \(\overline F/F\) is
algebraic, \(\overline F/\widetilde\sigma(\overline F)\) is algebraic.  An
algebraically closed field has no proper algebraic extension: if \(B/A\) is
algebraic and \(A\) is algebraically closed, then the minimal polynomial over
\(A\) of any \(\beta\in B\) has a root in \(A\), so irreducibility forces it
to be linear.  Hence \(\beta\in A\).  Thus
\(\widetilde\sigma\) is onto, as required.
\end{proof}

\begin{lemma}[Extension through a normal algebraic extension]
If \(E/F\) is algebraic and normal and \(L/F\) is an algebraic subextension
of \(E/F\), every \(F\)-embedding \(L\to\overline F\) that maps \(L\) into
\(E\) extends to an \(F\)-automorphism of \(E\).  In particular, every
\(F\)-automorphism of a finite subextension \(L/F\) extends to an
\(F\)-automorphism of \(E\).
\end{lemma}
\begin{proof}
Extend the given embedding to
\(\widetilde\sigma\in\operatorname{Aut}_F(\overline F)\) by
Lemma~\ref{lem:closure-auto-extension}.  If \(x\in E\), then
\(\widetilde\sigma(x)\) is another root of the minimal polynomial of \(x\)
over \(F\).  Normality says that all such roots lie in \(E\), so
\(\widetilde\sigma(E)\subseteq E\).  Applying the same argument to
\(\widetilde\sigma^{-1}\) gives the reverse inclusion.  Hence
\(\widetilde\sigma|_E\) is the desired automorphism.
\end{proof}

\begin{theorem}[Infinite Galois groups as inverse limits]\label{thm:infinite-galois-limit}
For a Galois extension \(E/F\), restriction induces an isomorphism
\[
\operatorname{Gal}(E/F)\cong
\varprojlim_{K/F\ \mathrm{finite\ Galois}}\operatorname{Gal}(K/F).
\]
\end{theorem}
\begin{proof}
For \(K\subseteq L\), restriction
\(\operatorname{Gal}(L/F)\to\operatorname{Gal}(K/F)\) is well-defined and
surjective by extension through the finite normal extension \(L/F\), and the
composition law is immediate.

An automorphism of \(E\) gives compatible restrictions.  It is determined by
them because the finite Galois fields cover \(E\).  Conversely, for a
compatible family \((\sigma_K)_K\), define \(\sigma(x)=\sigma_K(x)\) at any
finite Galois stage containing \(x\).  A common larger stage proves this is
well-defined and respects addition and multiplication.  It fixes \(F\), and
the compatible inverse family defines its inverse.  Thus it is an
automorphism of \(E\).
\end{proof}

\begin{definition}
The \emph{Krull topology} on \(\operatorname{Gal}(E/F)\) has basic
neighborhoods of the identity
\[
\operatorname{Gal}(E/K),
\]
where \(K/F\) is finite Galois.  Their cosets are the other basic open sets.
\end{definition}

This topology is forced by the finite-stage viewpoint: two automorphisms
\(\sigma,\tau\) are close when they agree on a large finite Galois
subextension \(K/F\).  Equivalently,
\[
\tau^{-1}\sigma\in\operatorname{Gal}(E/K).
\]
Thus the displayed stabilizers are exactly the natural identity
neighborhoods, and their translates are all other basic neighborhoods.
They do form a basis: for finite Galois \(K_1/F\) and \(K_2/F\), the
directedness argument above gives a finite Galois compositum \(K_1K_2/F\),
and
\[
\operatorname{Gal}(E/K_1)\cap\operatorname{Gal}(E/K_2)
=\operatorname{Gal}(E/K_1K_2).
\]

Under Theorem~\ref{thm:infinite-galois-limit}, these subgroups are exactly the
kernels of finite-stage projections.  Thus the Krull topology is the
inverse-limit topology, the displayed isomorphism is an isomorphism of
topological groups, and \(\operatorname{Gal}(E/F)\) is profinite.
Concretely, in the finite-field case a basic neighborhood of the identity
consists of the automorphisms that fix a prescribed
\(\mathbb F_{q^n}\).  The topology therefore says that two automorphisms are
close when their Frobenius powers agree on a sufficiently large finite field.

The infinite correspondence follows the same finite-stage logic as the
finite theorem, with one new feature: closedness is what guarantees that a
subgroup is recovered from all of its finite-stage images.  The proof
first recovers an intermediate field from its stabilizer, then uses the
finite correspondence at every finite Galois stage, and finally translates
finiteness and normality into openness and normality of the subgroup.
In the finite case every subgroup is automatically closed.  In the infinite
case a nonclosed subgroup can omit automorphisms that its elements approximate
at every finite stage, so it cannot be recovered from its finite images.

\begin{theorem}[Infinite Galois correspondence]\label{thm:infinite-galois}
For a Galois extension \(E/F\), the assignments
\[
K\longmapsto\operatorname{Gal}(E/K),
\qquad
H\longmapsto E^H
\]
are inverse inclusion-reversing bijections between intermediate fields and
closed subgroups of \(\operatorname{Gal}(E/F)\).  Further,
\[
K/F\text{ finite}\quad\Longleftrightarrow\quad
\operatorname{Gal}(E/K)\text{ open},
\]
and
\[
K/F\text{ Galois}\quad\Longleftrightarrow\quad
\operatorname{Gal}(E/K)\text{ normal}.
\]
Consequently, finite Galois intermediate extensions correspond to open normal
subgroups.
\end{theorem}
\begin{proof}
For \(\alpha\in E\), choose a finite Galois subextension \(L/F\) containing
\(\alpha\).  Its point stabilizer contains the open subgroup
\(\operatorname{Gal}(E/L)\), hence is itself open.  It has finite index and
is closed by Proposition~\ref{prop:profinite-open-subgroups}.  Therefore, for
an intermediate \(K\),
\[
\operatorname{Gal}(E/K)=\bigcap_{\alpha\in K}\operatorname{Stab}(\alpha)
\]
is closed.  We have
\[
E^{\operatorname{Gal}(E/K)}=K. \tag{12.2}
\]
Indeed, if \(\alpha\notin K\), its minimal polynomial over \(K\) has degree
larger than one.  The extension \(E/K\) is normal and separable, so that
polynomial has in \(E\) a second, distinct root \(\beta\).  The map
\(K(\alpha)\to K(\beta)\) fixing \(K\) and taking \(\alpha\) to \(\beta\)
extends, by the preceding lemma applied over \(K\), to a \(K\)-automorphism of
\(E\).  It moves \(\alpha\), proving (12.2).

Let \(H\) be closed, and let \(H_L\) be its image in
\(\operatorname{Gal}(L/F)\) for a finite Galois \(L/F\).  The field
\(L^{H_L}\) lies in \(E^H\).  If \(\sigma\) fixes \(E^H\), then
\(\sigma|_L\) fixes \(L^{H_L}\), so the finite correspondence gives
\(\sigma|_L\in H_L\).  Thus the image of \(\sigma\) lies in the image of
\(H\) at every finite stage; equivalently, every basic open neighborhood of
\(\sigma\) meets \(H\).  Hence \(\sigma\) lies in the closure of \(H\), and so
in \(H\).  Therefore \(\operatorname{Gal}(E/E^H)=H\), completing the
correspondence.

If \(K/F\) is finite, then it is separable; by the
\hyperref[thm:finite-galois-closure]{finite Galois closure theorem} and
normality of \(E/F\), choose a finite Galois extension \(L/F\) in \(E\)
containing \(K\).
Then \(\operatorname{Gal}(E/L)\) is open and is contained in
\(\operatorname{Gal}(E/K)\).  Conversely, openness supplies a finite Galois
\(L/F\) with \(\operatorname{Gal}(E/L)\) contained in
\(\operatorname{Gal}(E/K)\).  Taking fixed fields in (12.2) gives
\(K\subseteq L\).
For the normality assertion, write \(G=\operatorname{Gal}(E/F)\) and
\(H=\operatorname{Gal}(E/K)\).  First observe that
\[
\sigma H\sigma^{-1}=\operatorname{Gal}(E/\sigma(K))
\qquad(\sigma\in G).
\]
Indeed, an element on the left fixes \(\sigma(K)\), and the converse follows
by conjugating back.  If \(K/F\) is Galois, each \(\sigma\in G\) preserves
\(K\), so this formula gives \(\sigma H\sigma^{-1}=H\).  Conversely, if
\(H\) is normal, the formula and (12.2) give
\[
\sigma(K)=E^{\operatorname{Gal}(E/\sigma(K))}=E^H=K
\qquad(\sigma\in G).
\]
To prove that \(K/F\) is normal, let \(\tau:K\to\overline F\) be an
\(F\)-embedding.  By Lemma~\ref{lem:closure-auto-extension}, it extends to
an automorphism \(\widetilde\tau\) of \(\overline F\) fixing \(F\).  Since
\(E/F\) is normal, \(\widetilde\tau(E)=E\), so its restriction belongs to \(G\).
The displayed stability of \(K\) therefore yields \(\tau(K)=K\).  Thus
\(K/F\) is normal; it is separable because it is an intermediate extension
of the separable extension \(E/F\).  Hence it is Galois.
\end{proof}

\begin{definition}
Inside a fixed algebraic closure, \(F^{\mathrm{sep}}\) is the
\emph{separable closure} of \(F\), and
\[
G_F=\operatorname{Gal}(F^{\mathrm{sep}}/F)
\]
is its \emph{absolute Galois group}.
\end{definition}

\begin{theorem}[The absolute Galois group of a finite field]
Let \(\mathbb F_q^{\mathrm{sep}}\) be a separable closure of \(\mathbb F_q\).
Then
\[
G_{\mathbb F_q}
\cong
\varprojlim_{n\geq1}\mathbb Z/n\mathbb Z
=\widehat{\mathbb Z}
\]
as topological groups.  Under this identification, Frobenius
\[
\operatorname{Fr}_q:x\longmapsto x^q
\]
corresponds to \((1\bmod n)_{n\geq1}\) and is a dense topological generator.
\end{theorem}
\begin{proof}
Every element of \(\mathbb F_q^{\mathrm{sep}}\) has finite degree over
\(\mathbb F_q\), and hence lies in some \(\mathbb F_{q^n}\).  The finite
extensions \(\mathbb F_{q^n}\) are directed by divisibility, so
Theorem~\ref{thm:infinite-galois-limit} gives
\[
G_{\mathbb F_q}
\cong
\varprojlim_{n\geq1}\operatorname{Gal}(\mathbb F_{q^n}/\mathbb F_q).
\]
By Proposition~\ref{prop:frobenius-order}, the \(n\)-th group is cyclic of
order \(n\), generated by \(\operatorname{Fr}_q\).  Identify it with
\(\mathbb Z/n\mathbb Z\) by sending the class of \(a\) to
\(\operatorname{Fr}_q^a\).  If \(n\mid m\), restriction from
\(\mathbb F_{q^m}\) to \(\mathbb F_{q^n}\) sends
\(\operatorname{Fr}_q^a\) to \(\operatorname{Fr}_q^a\), which is exactly
the reduction map
\[
\mathbb Z/m\mathbb Z\longrightarrow\mathbb Z/n\mathbb Z.
\]
The resulting inverse limit is \(\widehat{\mathbb Z}\).  Frobenius maps to
the compatible family \((1\bmod n)_n\).  Its cyclic subgroup maps onto every
finite quotient \(\mathbb Z/n\mathbb Z\), so it meets every basic open coset;
the projection-kernel basis in Proposition~\ref{prop:profinite-open-subgroups}
therefore shows that it is dense.
\end{proof}

Here dense generator is deliberately topological language.  Frobenius does
not generate \(\widehat{\mathbb Z}\) as an abstract cyclic group; rather, the
ordinary cyclic subgroup generated by Frobenius is dense in
\(\widehat{\mathbb Z}\).

Under this identification,
\[
\operatorname{Gal}(\mathbb F_q^{\mathrm{sep}}/\mathbb F_{q^n})
\]
corresponds to
\[
\ker\bigl(\widehat{\mathbb Z}\longrightarrow\mathbb Z/n\mathbb Z\bigr)
=n\widehat{\mathbb Z},
\]
the unique open subgroup of index \(n\); its uniqueness also follows from the
infinite Galois correspondence together with the uniqueness of the degree-
\(n\) extension \(\mathbb F_{q^n}/\mathbb F_q\).  Thus the finite fields,
Frobenius, inverse limits, and the open-subgroup part of the infinite Galois
correspondence meet in one explicit example.

\section{Further Topics in Algebra}

This final section is orientation rather than prerequisite material.  Its aim
is to show how the four viewpoints developed in these notes become the common
language of later algebra.  The universal properties of free objects,
quotients, tensor products, and limits encountered throughout these notes are
organized systematically by category theory, whose formal development belongs
to later study.

\subsection{Groups, Actions, and Linear Algebra}

At its most elementary, a representation of a group \(G\) is a homomorphism
\[
G\longrightarrow\operatorname{GL}(V).
\]
Thus representation theory combines group theory with linear algebra.
Actions from Part~I are actions on sets; representations are the especially
structured case in which the set is a vector space.  Invariant subspaces,
irreducible representations, characters, and induction and restriction
organize this richer form of symmetry.

Normal subgroups, quotients, semidirect products, and exact sequences lead to
the study of group extensions.  Group cohomology later organizes extension
and obstruction questions.  Groups also arise outside algebra: fundamental
groups, covering spaces, and deck transformations make group actions central
in algebraic topology.  Matrix groups, Lie groups, and Lie algebras are a
further meeting point of group theory and linear algebra.

\subsection{Rings, Local Structure, and Geometry}

Part~II is the entry point to commutative algebra.  Ideals, localization,
polynomial rings, prime and maximal ideals, and modules lead onward to
Noetherian rings, primary decomposition, integral dependence, Krull
dimension, completions, and local algebra.  Localization deserves special
emphasis: it is algebraic zooming-in near a prime ideal.

Affine algebraic geometry studies spaces encoded by commutative rings.
Informally, \(\operatorname{Spec}(R)\) has prime ideals as points;
localization describes what is visible near a point, residue fields record
its scalar data, and modules supply algebraic analogues of geometric linear
data.  Scheme theory systematically combines rings, localization, modules,
tensor products, and field extensions.

\subsection{Modules, Exactness, and Multilinear Algebra}

Exact sequences, projective and injective modules, flatness, Hom, tensor
products, and direct and inverse limits form the starting vocabulary of
homological algebra.  Projective and injective resolutions, \(\operatorname{Tor}\),
\(\operatorname{Ext}\), derived functors, and spectral sequences are later
tools built from this vocabulary.  The Mittag--Leffler remark and the symbol
\(\varprojlim{}^1\) are early signals of that theory.

Tensor products, dual spaces, exterior powers, and alternating forms also
lead to differential geometry: they organize differential forms, orientation,
volume forms, and tangent and cotangent spaces.  No manifold theory is needed
to recognize this continuation.

\subsection{Fields, Arithmetic, and Galois Symmetry}

Finite extensions of \(\mathbb Q\) lead to number fields.  Their rings of
algebraic integers, ideals in Dedekind domains, factorization of prime ideals,
ramification, and decomposition and inertia groups combine the themes of
rings, localization, field extensions, and Galois groups.  The present notes
also deliberately stop before the full solvability-by-radicals theorem;
Kummer theory is a natural next step when suitable roots of unity are present.

Absolute Galois groups act on algebraic objects, and Galois cohomology later
organizes these actions in field arithmetic, descent, and number theory.
Continuous homomorphisms
\[
G_F\longrightarrow\operatorname{GL}_n(R)
\]
for suitable topological coefficient rings are Galois representations, a
major meeting point of Galois theory, representation theory, and topology.
Algebraic geometry has a parallel profinite fundamental group: finite
\'etale covers play the role of finite covering spaces.  Differential Galois
theory and covering-space Galois theory offer further variants of the same
symmetry principle.

\subsection{A Synthesis}

\[
\begin{array}{rcl}
\text{Groups} &\longrightarrow& \text{symmetry and actions},\\
\text{Rings} &\longrightarrow& \text{arithmetic and local structure},\\
\text{Modules} &\longrightarrow& \text{linearized algebra and exactness},\\
\text{Fields} &\longrightarrow& \text{algebraic equations and Galois symmetry}.
\end{array}
\]
Later subjects arise by combining these viewpoints: representation theory is
groups plus linear algebra; homological algebra is modules plus exact
sequences and Hom/tensor; algebraic geometry is commutative rings plus
localization, modules, and field extensions; algebraic number theory is rings
plus fields and Galois theory; and arithmetic geometry joins algebraic
geometry, number theory, and Galois representations.  These are conceptual
slogans, not formal identities.  The recurring methods of the notes---quotient,
localization, universal property, action, kernel and image, exact sequence,
tensor construction, passage to finite stages, and symmetry---continue to
reappear throughout advanced mathematics.

\section*{Exercises}
\subsection*{Basic}
\begin{exercise}[Basic] Prove directly that the fixed set of field
automorphisms is a field.
\end{exercise}
\begin{exercise}[Basic] Verify Artin independence for the two embeddings of
\(\mathbb Q(\sqrt2)\) in \(\mathbb C\).
\end{exercise}
\begin{exercise}[Basic] Determine the correspondence for
\(\mathbb Q(\sqrt2,\sqrt3)/\mathbb Q\).
\end{exercise}
\begin{exercise}[Basic] Compute the direct limit of
\(\mathbb Z\xrightarrow{\times3}\mathbb Z\xrightarrow{\times3}\cdots\).
\end{exercise}
\begin{exercise}[Basic] Describe the compatible sequence in \(\mathbb Z_2\)
associated with each of \(-1,0,1,2\).
\end{exercise}
\begin{exercise}[Basic] Prove that every discrete space is Hausdorff, and
describe a basic open set in a product of finite discrete spaces.
\end{exercise}

\subsection*{Core}
\begin{exercise}[Core] Reprove Lemma~\ref{lem:artin-independence} using any
chosen pair of distinct embeddings.
\end{exercise}
\begin{exercise}[Core] List all subgroup-field pairs in the splitting field
of \(X^3-2\), including their degrees.
\end{exercise}
\begin{exercise}[Core] Work out the correspondence for
\(\mathbb Q(\zeta_5)/\mathbb Q\), including normality.
\end{exercise}
\begin{exercise}[Core] Prove the equality criterion in
Proposition~\ref{prop:direct-limit-elements} from the quotient construction.
\end{exercise}
\begin{exercise}[Core] Use the element criteria to prove exactness for a
specified directed system of short exact sequences.
\end{exercise}
\begin{exercise}[Core] Give an inverse system for which surjectivity at each
stage does not give surjectivity on inverse limits.
\end{exercise}
\begin{exercise}[Core] Show that projection kernels form a neighborhood basis
in a profinite group.
\end{exercise}
\begin{exercise}[Core] Prove that the kernel of a continuous homomorphism to a
discrete group is open, and use translations to describe its cosets as open
sets.
\end{exercise}
\begin{exercise}[Core] Verify directly the gluing step in
Theorem~\ref{thm:infinite-galois-limit}.
\end{exercise}
\begin{exercise}[Core] Give a direct alternative proof that
\(\operatorname{Gal}(E/K)\) is closed in the Krull topology.
\end{exercise}
\begin{exercise}[Core] Show that an open subgroup of a profinite group has
finite index.
\end{exercise}

\subsection*{Further}
\begin{exercise}[Further] Prove the root-of-unity radical proposition in full
for a single radical step.
\end{exercise}
\begin{exercise}[Further] Use the degree obstruction to show that a regular
heptagon is not straightedge-and-compass constructible.
\end{exercise}
\begin{exercise}[Further] Show that \(\widehat{\mathbb Z}\) maps onto
\(\mathbb Z/n\mathbb Z\) for every \(n>0\).
\end{exercise}
\begin{exercise}[Further] Prove surjectivity of restriction from an infinite
Galois group to a finite Galois stage.
\end{exercise}
\begin{exercise}[Further] Let \(E/F\) be Galois and let \(K/F\) be finite.
Choose a finite Galois extension \(L/F\) in \(E\) containing \(K\).  Use the
finite correspondence in \(L/F\) to decide when
\(\operatorname{Gal}(E/K)\) is normal in \(\operatorname{Gal}(E/F)\).
\end{exercise}

\subsection*{Looking Ahead}
\begin{exercise}[Looking Ahead] For \(n\mid m\), verify that the restriction
map
\[
\operatorname{Gal}(\mathbb F_{q^m}/\mathbb F_q)
\longrightarrow
\operatorname{Gal}(\mathbb F_{q^n}/\mathbb F_q)
\]
corresponds, under the Frobenius generators, to reduction
\[
\mathbb Z/m\mathbb Z\longrightarrow\mathbb Z/n\mathbb Z.
\]
\end{exercise}
\begin{exercise}[Looking Ahead] Explain why finite separable subextensions of
\(F^{\mathrm{sep}}/F\) correspond to open subgroups of \(G_F\).
\end{exercise}

\backmatter
\begin{thebibliography}{9}
\raggedright
\bibitem{DF} David S. Dummit and Richard M. Foote,
\emph{Abstract Algebra}, 3rd ed., Wiley, 2004.
\bibitem{MA} Tommy Chen, \href{https://github.com/tommychen99/mathematical-analysis/releases/tag/v1.0.1}{\emph{Lecture Notes on Mathematical Analysis}}, Version 1.0.1, August 23, 2026.
GitHub release.
\end{thebibliography}
\end{document}
